% book example for classicthesis.sty
\documentclass[12pt,a4paper,footinclude=true,headinclude=true]{scrbook} % KOMA-Script book
\usepackage[T1]{fontenc}                
\usepackage{lipsum}
\usepackage[linedheaders,parts,pdfspacing]{classicthesis} % ,manychapters
%\usepackage[osf]{libertine}

%行距
\renewcommand{\baselinestretch}{1.15}
\usepackage[top=2.25cm,bottom=1.5cm,hmargin=2cm]{geometry}

%packages
\usepackage{dsfont}
\usepackage{polynom}
\usepackage{empheq}
%\usepackage[latin1]{inputenc}
\usepackage[utf8]{inputenc}
%\usepackage[T1]{fontenc}
%\usepackage[ngerman]{babel}
\usepackage{subfiles}
\usepackage{graphicx}
\usepackage{systeme}
\usepackage{mathtools, stmaryrd}
\usepackage{tikz-cd}
\usepackage{extpfeil}
%\usetikzlibrary{decorations.markings}
\usepackage{CJK}
\usepackage{amsmath, amsthm, amssymb, amsfonts, mathrsfs, bm, fancyhdr, color, comment, graphicx, environ}
\usepackage{xcolor}
\usepackage{mdframed}
\usepackage[shortlabels]{enumitem}
\usepackage{indentfirst}
\usepackage{hyperref}
\usepackage{ifthen}
\usepackage{enumerate}

%custom commands: symbols
\newcommand{\notni}{\not\mathrel{\text{\reflectbox{$\in$}}}}
\newcommand{\trianglerightneq}{\mathrel{\ooalign{\raisebox{-0.5ex}{\reflectbox{\rotatebox{90}{$\nshortmid$}}}\cr$\triangleright$\cr}\mkern-3mu}}
\newcommand{\triangleleftneq}{\mathrel{\reflectbox{$\trianglerightneq$}}}
\renewcommand\qedsymbol{$\square$}
\DeclareMathOperator{\E}{E} % expectation operator
\newcommand\myeq[1]{\stackrel{\textnormal{#1}}{=}}

%custom commands: theorem
\newtheorem{innercustomthm}{Theorem}
\newenvironment{theorem}[2][]
{\renewcommand\theinnercustomthm{#2}\innercustomthm[#1]
\begin{mdframed}[backgroundcolor=yellow!20]\ }
{\end{mdframed}\endinnercustomthm}
\newtheorem*{innertheorem*}{Theorem}
\newenvironment{theorem*}
{\begin{innertheorem*}\begin{mdframed}[backgroundcolor=gray!20]\ }
{\end{mdframed}\end{innertheorem*}}
\newtheorem{innercustomprop}{Proposition}
\newenvironment{proposition}[2][]
{\renewcommand\theinnercustomprop{#2}\innercustomprop[#1] 
\begin{mdframed}[backgroundcolor=gray!20]\ }
{\end{mdframed}\endinnercustomprop}
\newtheorem{innercustomlemma}{Lemma}
\newenvironment{lemma}[2][]
{\renewcommand\theinnercustomlemma{#2}\innercustomlemma[#1]
\begin{mdframed}[backgroundcolor=gray!20]\ }
{\end{mdframed}\endinnercustomprop}
\newtheorem{innercustomcol}{Corollary}
\newenvironment{corollary}[2][]
{\renewcommand\theinnercustomcol{#2}\innercustomcol[#1]
\begin{mdframed}[backgroundcolor=gray!20]\ }
{\end{mdframed}\endinnercustomcol}
\newtheorem{innercustomalg}{Algorithm}
\newenvironment{algorithm}[2][]
{\renewcommand\theinnercustomalg{#2}\innercustomalg[#1]
\begin{mdframed}[backgroundcolor=gray!20]\ }
{\end{mdframed}\endinnercustomalg}

%custom command
\newenvironment{definition}[2][Definition]
    { \begin{mdframed}[backgroundcolor=gray!20] \textbf{#1 #2.}\ }
    {  \end{mdframed}}
\newenvironment{definition*}[2][Definition]
    { \begin{mdframed}[backgroundcolor=blue!20] \textbf{#1 #2.}\ }
    {  \end{mdframed}}
\newenvironment{remark}[2][Remark]
    { \begin{mdframed}[backgroundcolor=gray!20] \textbf{#1 #2.} }
    {  \end{mdframed}}
\newenvironment{remark*}[2][Remark]
    { \begin{mdframed}[backgroundcolor=gray!20] \textbf{#1 #2.}\ }
    {  \end{mdframed}}    
\newenvironment{nt}[2][Note]
    { \begin{mdframed}[backgroundcolor=gray!20] \textbf{#1 #2.} }
    {  \end{mdframed}}
\newenvironment{example}[2][Example]
    { \begin{mdframed}[backgroundcolor=gray!20] \textbf{#1 #2.} }
    {  \end{mdframed}}
\newenvironment{exe}[2][Exercise]
    { \begin{mdframed}[backgroundcolor=gray!20] \textbf{#1 #2.}\\ }
    {  \end{mdframed}}
\newenvironment{problem}[2][Problem]
    { \begin{mdframed}[backgroundcolor=gray!20] \textbf{#1 #2.} }
    {  \end{mdframed}}
\newenvironment{notation}[2][Notation]
    { \begin{mdframed}[backgroundcolor=gray!20] \textbf{#1 #2.} }
    {  \end{mdframed}}   
\newenvironment{motivation}[2][Motivation]
    { \begin{mdframed}[backgroundcolor=red!20] \textbf{#1 #2.} }
    {  \end{mdframed}}    
\newenvironment{proofthm}[2][Proof of Theorem]
    { \begin{mdframed}[backgroundcolor=gray!20] \textbf{#1 #2.} }
    {  \end{mdframed}}   
\newenvironment{summary}[2][]
    { \begin{mdframed}[backgroundcolor=gray!20] \textbf{#1 #2.} }
    {  \end{mdframed}}    

%Setup
\hypersetup{
    colorlinks=true,
    linkcolor=blue,
    filecolor=magenta,      
    urlcolor=blue,}
    
\makeatletter
\@addtoreset{chapter}{part}
\makeatletter    


\begin{document}
%	\pagestyle{scrheadings}
%	\manualmark
%	\markboth{\spacedlowsmallcaps{\contentsname}}{\spacedlowsmallcaps{\contentsname}}
	
	\tableofcontents 

	%\automark[section]{chapter}
	%\renewcommand{\chaptermark}[1]{\markboth{\spacedlowsmallcaps{#1}}{\spacedlowsmallcaps{#1}}}
	%\renewcommand{\sectionmark}[1]{\markright{\thesection\enspace\spacedlowsmallcaps{#1}}}

% use \cleardoublepage here to avoid problems with pdfbookmark
\cleardoublepage\part{Arithmetic Curves \& Algebraic Number Theory}

\chapter*{abstract}    
\begin{paragraph}{Introduction 1.}
    As a general rule in algebraic geometry and number theory you never know enough algebra. For this lecture course knowledge of basic ring theory and rudiments of field theory as taught in the Algebra course of your syllabus is indispensable. A good reference for this material is \textbf{S. Bosch, Algebra, Springer}. This introductory chapter summaries some algebraic background material.\\
    Moreover, we often have to appeal to the theory of modules and algebras which are direct generalisations of vector spaces and polynomial rings. Most of this is covered in the course GAGA A. Here a good standard reference is \textbf{M. Atiyah} and \textbf{I. McDonald, Introduction to Commutative Algebra, Addison-Wesley}. You can first safely admit these results and approach this course as a concrete and preliminary motivation for the abstract foundations of (commutative algebra).
\end{paragraph}    
\setcounter{chapter}{-1}     

\chapter{review of basic notions}
\section{Rings}
\begin{notation}{2}
    \begin{itemize}
        \item[(a)] We use these notations:
        \begin{itemize}[$\bullet$]
            \item $\mathbb{N}=\{0,1,2,\cdots\}=$ the monoid of natural numbers.
            \item $\mathbb{Z}=$ ring of integers.
            \item $\mathbb{Q}=$ field of rational numbers.
            \item $\mathbb{R}=$ field of real numbers.
            \item $\mathbb{C}=$ field of complex numbers.
            \item $\mathbb{F}_p=\mathbb{Z}/p\mathbb{Z}=$ field with $p$ elements, ($p$: prime).
        \end{itemize}
        \item[(b)] Notations for sets and mappings:
        \begin{itemize}[$\bullet$]
            \item The empty set is denoted by $\varnothing$.
            \item The cardinality of a set $S$ by $\# S$.
            \item A family of elements of a set $A$ indexed by a set $I$, written $\{a_i\}_{i\in I}$, is a function $I\to A,\ i\mapsto a_i$.
            \item Given an equivalence relation, $\bar{x}$ denotes the equivalence class of $x$.
            \item Further, $\xhookrightarrow{}$ and $\twoheadrightarrow$ denote the injective and surjective maps, respectively.
        \end{itemize}
        \item[(c)] Relations on sets:
        \begin{itemize}[$\bullet$]
            \item $X\subset Y$: $X$ is a subset (not necessarily proper) of $Y$.
            \item $X\simeq Y$: $X$ is isomorphic to $Y$.
            \item $X\cong Y$: $X$ and $Y$ are canonically isomorphic (i.e. there is a given or unique isomorphism).
        \end{itemize}
    \end{itemize}
\end{notation}
\newpage

\begin{definition}{3. (Ring)}
    \begin{itemize}
        \item[(a)] A \textbf{ring} $(A,+,\cdot)$ consists of an abelian group $(A,+)$ with neutral element $0$, and an associative multiplication $\cdot:A\times A\to A$ such that for all $a,b,c,\in A$,
        \begin{align*}
            a\cdot(b+c)=a\cdot b+a\cdot c\qquad\text{and}\qquad (b+c)\cdot a=b\cdot a+c\cdot a.
        \end{align*}
        In this course, a ring will always be assumed to be \textbf{commutative with 1}, i.e. multiplication $\cdot$ is commutative with neutral multiplicative element $1$.
        \item[(b)] A subgroup of a ring $A$ is a \textbf{subring} if it contains $1$ and is closed under $\cdot$.
        \item[(c)] A \textbf{unit} $u\in A\setminus\{0\}$ is an element such that there exists $u'\in A$ with $u\cdot u'=1$. \\ Units form a multiplicative subgroup of $A$ which we denote by $A^\times$. If $1\neq0$ and $A^\times=A\setminus\{0\}$, that is, all nonzero elements are units, then $A$ is called a \textbf{field} and usually denoted by $F$.
        \item[(d)] A \textbf{ring homomorphism} $\phi:A\to B$ is a group homomorphism which satisfies that $\phi(1_A)=1_B$ and $\phi(a\cdot b)=\phi(a) \cdot\phi(b)$. A \textbf{ring isomorphism} is a bijective ring homomorphism, (its inverse is necessarily a ring homomorphism).
    \end{itemize}
\end{definition}

\begin{definition}{4. (Ideal)}
    \begin{itemize}
        \item[(a)] A subset $\mathcal{I}\subset A$ is an \textbf{ideal} if it is an abelian subgroup of $A$ and for any $r\in A,\ a\in\mathcal{I}$, we have $r\cdot a\in\mathcal{I}$. We denote by $\mathcal{I}\triangleleft A$.
        \item[(b)] The ideal \textbf{generated} by the subset $\{a_i:i\in \Lambda\}$ of $A$ is the set of finite linear combinations of generators $a_i$, i.e. 
        \begin{align*}
            \mathcal{I}=\{\sum_{\text{finite}}r_ia_i:r_i\in A, i\in\Lambda\}.
        \end{align*}
        This is indeed an ideal and we denote $\mathcal{I}\coloneqq (a_i: i\in\Lambda)$.
        \item[(c)] If $\Lambda$ is a finite set, then we say $\mathcal{I}$ is \textbf{finitely generated}. 
        \item[(d)] An ideal is generated by a single element is called a \textbf{principal ideal}.
        \item[(e)] The \textbf{null ideal} is the ideal $(0)=\{0\}$ generated by $0$.
    \end{itemize}
\end{definition}

\begin{example}{5}
    If $\phi:A\to B$ is a ring homomorphism, the \textbf{kernel} $\mathsf{ker}\phi\coloneqq\phi^{-1}(0)$ is an ideal of $A$.
\end{example}

\begin{definition}{6. (Constructions of Ideals)}
    \begin{itemize}
        \item[(a)] Let $\mathcal{I}$ and $\mathcal{J}$ be two ideals of $A$. Their \textbf{sum} is the ideal 
        \begin{equation*}
            \mathcal{I}+\mathcal{J}\coloneqq\{a+b:a\in\mathcal{I}, b\in\mathcal{J}\}.
        \end{equation*}
        Their \textbf{product} is 
        \begin{equation*}
            \mathcal{I}\cdot\mathcal{J}\coloneqq(a\cdot b:a\in\mathcal{I}, b\in\mathcal{J}).
        \end{equation*}
        \item[(b)] Let $\phi:A\to B$ be a ring homomorphism and let $\mathfrak{b}\triangleleft B$. Then its \textbf{contraction}, denoted by $\mathfrak{b}^c\coloneqq\phi^{-1}(\mathfrak{b})$ is an ideal of $A$.
        \item[(c)] If $\mathfrak{a}\triangleleft A$ is an ideal, then $\phi(\mathfrak{a})$ is, generally, merely a subring of $B$, but not an ideal. The \textbf{extension} of $\mathfrak{a}$ by $\phi$ is the ideal $\mathfrak{a}^e\coloneqq\phi(\mathfrak{a})\cdot B$ generated by $\phi(\mathfrak{a})$ in $B$. 
    \end{itemize}
\end{definition}

\begin{example}{7}
    \begin{itemize}
        \item[(a)] If $\mathcal{I}\coloneqq(a_1,\cdots,a_r)$ and $\mathcal{J}\coloneqq(b_1,\cdots,b_s)$ are two finitely generated ideals, then so are their sum and their product. In fact,
        \begin{align*}
            \mathcal{I}+\mathcal{J}=(a_1,\cdots,a_r,b_1,\cdots,b_s)\ \text{and}\ \mathcal{I}\cdot\mathcal{J}=
            (a_ib_j:1\leq i\leq r, 1\leq j\leq s).
        \end{align*}
        \item[(b)] If $\mathfrak{a}\coloneqq(a)\triangleleft A$ is principal, then $\mathfrak{a}^e=\{\phi(a)\cdot b:b\in B\}$.\\
        For example, if $\mathfrak{a}\coloneqq(m)\subset\mathbb{Z} \xhookrightarrow{} \mathbb{Q}$, then $\mathfrak{a}^e=\mathfrak{a}\cdot\mathbb{Q}=\mathbb{Q}$ if $m\neq0$.
    \end{itemize}
\end{example}

\begin{definition}{8. (Classes of Rings)}
    \begin{itemize}
        \item[(a)] A \textbf{zero divisor} $a\in A$ divides $0$, i.e. there exists $b\in A\setminus\{0\}$ such that $a\cdot b=0$. An element $a\in A$ is \textbf{nilpotent} if $a^n=0$ for some $n\in\mathbb{N}$. In particular, a nilotent $a\in A$ is a zero divisor if $a\neq0$.
        \item[(b)] A nontrivial ring $A$ is an \textbf{integral domain (ID)}/ \textbf{reduced} if $A$ has no zero divisors/ no nilpotents other than $0$.
        \item[(c)] For an ID A, a nonunit $a\in A$ is \textbf{irreducible} if $a\neq0$ and $a=b\cdot c$ for $b,c\in A$ implies that either $b$ or $c$ is a unit. Further, a nonunit $a$ is called \textbf{prime} if $a\neq0$ and $a|(b\cdot c)$ implies either $a|b$ or $a|c$. Moreover, "prime" implies "irreducible", but the converse is flase in general.
        \item[(d)] An ID $A$ is a \textbf{unique factorisation domain (UFD)} or is \textbf{factorial} if every nonunit $0\neq a\in A$ admits a decomposition $a\coloneqq a_1\cdots a_r$ into primes unique up to their orders and units. In a UFD, irreducibles are prime.
        \item[(e)] An ID $A$ is a \textbf{principal ideal domain (PID)} if it has only principal ideals.
        \item[(f)] A \textbf{Euclidean domain (ED)} is an ID $A$ with a function $\nu:A\setminus\{0\}\to\mathbb{N}$ such that for any $f,g\in A,\ g\neq0$, there are elements $q,r\in A$ with $f=qg+r$ and either $r=0$ or $\nu(r)<\nu(g)$. We can such $r$ the \textbf{remainder} of $f$ by $g$ and write $r\equiv f\ (\mathrm{mod}\ g)$.
    \end{itemize}
\end{definition}

\begin{remark}{9}
    We have the following implications:
    \begin{align*}
        \textbf{Field}\ \Longrightarrow\ \textbf{ED}\ \Longrightarrow\ \textbf{PID}\ \Longrightarrow\ \textbf{UFD}\ \Longrightarrow\ \textbf{ID}.
    \end{align*}
    There are counterexamples to every converse implication.
\end{remark}

\begin{example}{10}
    Here are important examples of EDs:
    \begin{itemize}
        \item[(a)] The integers $\mathbb{Z}$ with $\nu(a)\coloneqq|a|$, the \textbf{absolute value} of $a$.
        \item[(b)] The \textbf{Gaussian integers} $\mathbb{Z}[i] \coloneqq\{a+ib:a,b\in\mathbb{Z}\}\subset\mathbb{C}$ with $\nu(a+ib)\coloneqq a^2+b^2$, the \textbf{norm} of $a+ib$. (cf. \textbf{Prop. 88}).
        \item[(c)] The polynomial ring $F[x]$ in one variable over the field $F$ with $\nu(f)\coloneqq\mathrm{deg}f$, the \textbf{degree} of the polynomial $f$.
    \end{itemize}
\end{example}

\begin{definition}{11. (Multivariable Polynomials)}
    \begin{itemize}
        \item A \textbf{monomial} in the variables $x_1,\cdots,x_n$ is a formal expression\\ \centerline{$x^{\alpha}\coloneqq x_1^{\alpha_1}\cdots x_n^{\alpha_n},\qquad \alpha=(\alpha_1,\cdots,\alpha_n)\in \mathbb{N}^n$.}\\
        The set of monomials $\mathrm{Mon}(x_1,\cdots,x_n)\coloneqq \{x^\alpha:\alpha\in\mathbb{N}^n\}$ forms a semi-group with neutral element $1=x_1^0\cdots x_n^0$ via 
        \begin{align*}
            x_1^{\alpha_1}\cdots x_n^{\alpha_n}\cdot x_1^{\beta_1}\cdots x_n^{\beta_n}=x_1^{\alpha_1+\beta_1}\cdots x_n^{\alpha_n+\beta_n}.
        \end{align*}
        \item The \textbf{degree} of $x^\alpha$ is defined to be the sum $|\alpha|=\alpha_1+ \cdots+\alpha_n$. A \textbf{polynomial} over $A$ in the variables $x_1,\cdots,x_n$ is a formal finite linear combination of monomials $x^\alpha$ over $A$: 
        \begin{align*}
            f\coloneqq\sum_{\alpha=(\alpha_1,\cdots,\alpha_n)} a_\alpha x_1^{\alpha_1}\cdots x_n^{\alpha_n},
        \end{align*}
        where $a_\alpha\in A,\ a_\alpha=0$ for all but finitely many $\alpha$.
        \item The \textbf{support} of $f$ is the set $\mathrm{supp}f\coloneqq \{x^\alpha:a_\alpha\neq0\}$. For $f\neq0$, the \textbf{degree} $\mathrm{deg}f$ of $f$ is the maximum of all $|\alpha|$ such that $x^\alpha\in\mathrm{supp}f$; by convention, we put $\mathrm{deg}\ 0 \coloneqq-1$ and \\ \centerline{$A[x_1,\cdots,x_n]\coloneqq$ set of polynomials in $n$ variables.}
    \end{itemize}
\end{definition}

\begin{definition}{12. (Ring Structure on the Polynomials)}
    A \textbf{term} is a polynomial of the form $ax^\alpha$. The ring structure on $A[x_1,\cdots,x_n]$ is defined as follows:
    \begin{align*}
        \sum a_\alpha x^\alpha+\sum b_\alpha x^\alpha
        &\coloneqq\sum(a_\alpha+b_\alpha)x^\alpha,\\
        (\sum a_\alpha x^\alpha)\cdot(\sum b_\alpha x^\alpha)
        &\coloneqq\sum_\gamma(\sum_{\alpha+\beta=\gamma}a_\alpha b_\beta)x^\gamma.
    \end{align*}
    If we identify $1\in A$ with $1=x^0$, then $A\subset A[x_1,\cdots,x_n]$ as a subring (with the obvious definitions). The elements of $A\subset A[x_1,\cdots,x_n]$ will be called the \textbf{constant polynomials}, and $A$ is the \textbf{base ring}.
\end{definition}

\begin{proposition}{13}
    \begin{itemize}
        \item[(a)] $(A[x_1,\cdots,x_{n-1}])[x_n]\cong A[x_1,\cdots,x_n]$.
        \item[(b)] If $A$ is an ID, then so is $A[x_1,\cdots,x_n]$.
        \item[(c)] (\textbf{Gau\ss\ Theorem}) If $A$ is a UFD, then so is $A[x_1,\cdots,x_n]$.
    \end{itemize}
\end{proposition}

An important construction method for rings is taking the quotient with respect to an ideal:
\begin{definition}{14. (Quotient Ring)}
    Let $\mathcal{I}\triangleleft A$ be an ideal. Then $a\sim b\ \xLeftrightarrow{\text{Def.}}\ a-b\in \mathcal{I}$ defines an equivalence class $\overline{a}=\overline{a+b},\ (b\in\mathcal{I})$. \\ The \textbf{quotient ring} $A/\mathcal{I}$ is the ring structure on the closed space $\{\overline{a}:a\in A\}$ defined by the following operations:
    \begin{align*}
        \overline{a}+\overline{b} = (a+A)+(b+A) &\coloneqq (a+b)+A =\overline{a+b}, \\ \overline{a}\cdot\overline{b} = (a+A)\cdot(b+A) &\coloneqq (a\cdot b)+A =\overline{a\cdot b}.
    \end{align*}
    In particular, the \textbf{projection} or \textbf{quotient map} $\pi_{\mathcal{I}}:A\to A/\mathcal{I},\ a\mapsto\bar{a}=a+A$ becomes a surjective ring homomorphism.
\end{definition}

\begin{example}{15}
    In $\mathbb{Q}[x,y]/(x^2,xy)$ we have $\bar{x}^2=0$ and $\bar{x}\cdot\bar{y}=0$. For instance, 
    \begin{equation*}
        \overline{5x^2y+3x^4+y+1}\cdot\overline{2x+xy+7y^3}=
        \overline{y+1}\cdot\overline{2x+7y^3}=2\overline{x}+7\overline{y}^3 +7\overline{y}^4.
    \end{equation*}
\end{example}
It is an important problem to understand the structure of ideals of a given ring as the set of ideals of a ring $A$ can be regarded as a measure for the complexity of $A$.

\begin{lemma}{16. (Comparison Lemma)}
    \begin{itemize}
        \item[(a)] A ring homomorphism $\phi:A\to B$ induces a ring isomorphism $A/\mathrm{ker}\phi\cong\mathrm{im} \phi$. In particular, if $\phi$ is surjective, then $B\cong A/\mathfrak{a}$ for some ideal $\mathfrak{a}\triangleleft A$.
        \item[(b)] The map $\mathfrak{b}\mapsto\mathfrak{b}^e =\pi_{\mathfrak{a}}(\mathfrak{b})\cdot(A/\mathfrak{a})$ induces a bijection 
        \begin{align*}
            \{\mathfrak{b}\triangleleft A:\mathfrak{b}\supset\mathfrak{a}\}\ &\longleftrightarrow\ \{\mathfrak{d}\triangleleft A/\mathfrak{a}\}, \\
            \pi_{\mathfrak{a}}^{-1}(\mathfrak{d})=\mathfrak{d}^c &\longmapsfrom\  \mathfrak{d}\triangleleft A/\mathfrak{a}.
        \end{align*}
    \end{itemize}
\end{lemma}
\newpage

\begin{theorem}{17. (Chinese Remainder Theorem)}
    Let $\mathfrak{a}_1,\cdots,\mathfrak{a}_r$ be ideals of $A$ which are pairwise coprimes, i.e. $\mathfrak{a}_i+\mathfrak{a}_j=A,\ \forall i\neq j$. Then $\mathfrak{a}\coloneqq\bigcap_j\mathfrak{a}_j=\prod_j\mathfrak{a}_j$ and the natural map 
    \begin{align*}
        A\longrightarrow\prod_j(A/\mathfrak{a}_j),\qquad a\longmapsto (a+\mathfrak{a}_1,\cdots,a+\mathfrak{a}_r)
    \end{align*}
    induces a ring isomorphism $A/\mathfrak{a}\cong\prod_j(A/ \mathfrak{a}_j)$.
\end{theorem}

The following types of ideals are extremely important in algebraic geometry and algebraic number theory:
\begin{definition}{18. (Prime and Maximal Ideals)}
    \begin{itemize}
        \item A proper ideal $\mathcal{I}\triangleleftneq A$ is \textbf{prime} if $a\cdot b\in\mathcal{I}$ implies either $a\in\mathcal{I}$ or $b\in\mathcal{I}$.
        \item A proper ideal $\mathcal{I}\triangleleftneq A$ is \textbf{maximal} if $\mathcal{I}\subset\mathcal{J}\subsetneq A$ implies $\mathcal{I}=\mathcal{J}$.
    \end{itemize}
\end{definition}

\begin{lemma}{19}
    \begin{itemize}
        \item[(a)] $\mathcal{I}\triangleleft A$ is prime $\Leftrightarrow\ A/\mathcal{I}$ is an ID. 
        \item[(b)] $\mathcal{I}\triangleleft A$ is maximal $\Leftrightarrow\ A/\mathcal{I}$ is a field.
        \item[(c)] In particular, a maximal ideal is prime.
    \end{itemize}
\end{lemma}

\begin{example}{20}
    In $\mathbb{Z}$ an ideal $\mathcal{I}=(m)$ is prime iff $m=0$, or $m$ is prime. Moreover, if $m$ is prime, then $\mathcal{I}$ is maximal. 
\end{example}

\begin{remark}{21}
    \begin{itemize}
        \item[(a)] For a field $K$ consider the ring homomorphism $\phi:\mathbb{Z}\to K,\ 1\mapsto1_K$. Then $\mathrm{ker}\phi =(m)$ for some $m\in\mathbb{Z}$. This integer is called the \textbf{characteristic} of $K$, and is denoted by $\mathrm{char}K\coloneqq m$. It follows that the characteristic of a field must be prime, for $\mathbb{Z}/\mathrm{ker}\phi\cong\mathrm{im}\phi$ is an ID. For example, $\mathrm{char}\ \mathbb{Q}=0$ and $\mathrm{char}\  \mathbb{Z}_p=p$.
        \item[(b)] If $\phi:A\to B$ is a ring homomorphism and $\mathfrak{q}\triangleleft B$ prime, then so is its contraction $\mathfrak{q}^c=\phi^{-1}(\mathfrak{q})$. On the other hand, contractions of maximal ideals are not necessarily maximal. (Consider the inclusion $\mathbb{Z}\xhookrightarrow{} \mathbb{Q}$ and $(0)\subset \mathbb{Q}$).
    \end{itemize}
\end{remark}

\begin{lemma}{22}
    \begin{itemize}
        \item[(a)] Let $\mathfrak{a},\ \mathfrak{b},\ \mathfrak{p}\triangleleft A$ be ideals with $\mathfrak{p}$ prime. If $\mathfrak{a}\nsubseteq\mathfrak{p}$ and $\mathfrak{a}\cdot \mathfrak{b}\subset\mathfrak{p}$, then $\mathfrak{b}\subset\mathfrak{p}$. 
        \item[(b)] Let $\mathfrak{a}_1,\cdots,\mathfrak{a}_n,\mathfrak{p}\triangleleft A$ be ideals with $\mathfrak{p}$ prime. If $\bigcap_{i=1}^n\mathfrak{a}_i\subset\mathfrak{p}$ [respectively, $=$], then $\mathfrak{a}_j\subset\mathfrak{p}$ [respectively, $=$] for some $j$.
        \item[(c)] (\textbf{Prime Avoidance}) Let $\mathfrak{p}_1,\cdots,\mathfrak{p}_n,\mathfrak{a}\triangleleft A$ be ideals with $\mathfrak{p}_i$ prime. If $\mathfrak{a}\subset\bigcup_{i=1}^n\mathfrak{p}_i$, then $\mathfrak{a}\subset\mathfrak{p}_j$ for some $j$.
    \end{itemize}
\end{lemma}

The structure of prime and maximal ideals of a ring is very important. From this point of view, the following class of rings is particularly simple.
\begin{definition}{23. (Local Ring)}
    A ring $A$ is \textbf{local} if it has exactly one maximal ideal $\mathfrak{m}$. The local ring is denoted by $(A,\mathfrak{m})$. The associated field $A/\mathfrak{m}$ is called the \textbf{residue field}.
\end{definition}

\begin{proposition}{24}
    The following properties on a ring $A$ are equivalent:
    \begin{itemize}
        \item[(a)] $(A,\mathfrak{m})$ is a local ring.
        \item[(b)] All the nonunits of $A$ form an ideal $\mathfrak{m}$.
        \item[(c)] There exists an ideal $\mathfrak{m}\neq(1)$ such that every $x\in A\setminus\mathfrak{m}$ is a unit in $A$.
        \item[(d)] There exists a maximal ideal $\mathfrak{m} \triangleleft A$ such that $1+\mathfrak{m}=\{1+m:m\in\mathfrak{m}\}\subset A^{\times}$.
    \end{itemize}
\end{proposition}

\begin{example}{25}
    \begin{itemize}
        \item[(a)] Every field is a local ring with maximal ideal $\mathfrak{m}=(0)$.
        \item[(b)] Let $\mathfrak{p}\triangleleft k[x_1,\cdots,x_n]$ be a prime ideal. Then 
        \begin{align*}
            A=k[x_1,\cdots,x_n]_{\mathfrak{p}}\coloneqq
            \{\tfrac{f}{g}:f,g\in k[x_1,\ \cdots,x_n],g\notin\mathfrak{p}\}
            \subset k(x_1,\cdots,x_n)
        \end{align*}
        is a local ring with ring structure induced by the field $k(x_1,\cdots,x_n)$. The maximal ideal is 
        \begin{align*}
            \mathfrak{m}=\mathfrak{p}\cdot k[x_1,\cdots,x_n]_{\mathfrak{p}}=\{\tfrac{f}{g}:f,g\notin \mathfrak{p}\}.
        \end{align*}
        This is just the extension with respect to the injection
        \begin{align*}
            k[x_1,\cdots,x_n]\xhookrightarrow{~~~~} A,\qquad f\longmapsto\tfrac{f}{1}.
        \end{align*}
        \item[(c)] (\textbf{Localisation}) Let $A$ be an ID, then there is a field $K\coloneqq\mathrm{Quot}A$ with canonical injection $A\xhookrightarrow{} K$, the so called \textbf{quotient field}. As a set, $K=\{\frac{a}{b}:a,b\in A,b\neq0\}$ with field operations given by the usual rules for fractions. A \textbf{multiplicative subset} $S$ of $A$ is a set which is closed under multiplication and which does not contain $0$. By $S^{-1}A$ we shall denote the set of quotients $\frac{a}{s}\in S^{-1}A$ with $s\in S$. This is again a ring with canonical injection $A\xhookrightarrow{} S^{-1}A$.
    \end{itemize}
\end{example}

\begin{example}{26}
    A typical example is $A_{\mathfrak{p}}\coloneqq S^{-1}A$, where $S\coloneqq A\setminus\mathfrak{p}$ and $\mathfrak{p}\triangleleft A$ prime ideal. Further, $A_{\mathfrak{p}}$ is a local ring with maximal ideal $\mathfrak{m}=\mathfrak{p}^e=\mathfrak{p}A_{\mathfrak{p}}$, which contains precisely of the elements $\frac{a}{s}$ with $a\in\mathfrak{p}$. Note that $\mathfrak{m}\cap A=\mathfrak{p}$.
\end{example}

\begin{proposition}{27}
    We consider the localisation $S^{-1}A$ of an ID. We consider extension and contraction with respect to the natural map $A\to S^{-1}A$. Then 
    \begin{align*}
        \mathfrak{b}^{ce} &=\mathfrak{b},\qquad \forall\ \mathfrak{b} \triangleleft S^{-1}A,\\
        \mathfrak{p}^{ec} &=\mathfrak{p},\qquad \forall\ \mathfrak{p} \triangleleft A\ \text{prime ideal with}\ \mathfrak{p}\cap S=\varnothing.
    \end{align*}
    In particular, the prime ideals of $S^{-1}A$ are precisely the extensions of the prime ideals of $A$ which have trivial intersection with $S$.
\end{proposition}

\begin{definition}{28. (Graded Ring)}
    \begin{itemize}
        \item A \textbf{graded ring} $A$ is a ring with a decomposition $A\coloneqq\bigoplus_{\nu\in\mathbb{N}}A_\nu$ as a direct sum of abelian groups such that $A_\nu\cdot A_\mu \subset A_{\nu+\mu},\ \forall\nu,\mu\in\mathbb{N}$.
        \item The $A_\nu$ are called the \textbf{homogeneous components} of degree $\nu$ of $A$.
        \item Elements of $A_\nu$ are said to be \textbf{homogeneous}.
        \item A \textbf{graded} $k$-\textbf{algebra} is a $k$-algebra which is a graded ring for which $A_0=k$ and $A_\nu$ is a $k$-vector space for all $n\in\mathbb{N}$.
    \end{itemize}
\end{definition}

\begin{example}{29}
    \begin{itemize}
        \item[(a)] Every ring $A$ carries a trivial graded ring structure by setting $A_0=A$ and $A_\nu=0$ for $\nu\geq1$.
        \item[(b)] If $\mathfrak{a}\triangleleft A$ is an ideal, then 
        \begin{align*}
            \mathrm{Gr}_\mathfrak{a}(A)\ =\ \bigoplus_{\nu\geq0}\mathrm{Gr}_\mathfrak{a}(A)_\nu
            \ \coloneqq\ \bigoplus_{\nu\geq0}\mathfrak{a}^\nu/\mathfrak{a}^{\nu+1}
        \end{align*}
        is a graded ring (where by convention, $\mathfrak{a}^0\coloneqq(1)=A$). Multiplication is given as follows: If $\overline{a}\in\mathfrak{a}^\nu/\mathfrak{a}^{\nu+1}$ and $\overline{b}\in\mathfrak{a}^\mu/\mathfrak{a}^{\mu+1}$, then $\overline{a}\cdot\bar{b}\coloneqq a\cdot b+\mathfrak{a}^{\nu\mu+1}\in\mathfrak{a}^{\nu+\mu}/\mathfrak{a}^{\nu\mu+1}$ (this is indeed well-defined). See also \textbf{Def. 304}.
        \item[(c)] Let $A=k[x_1,\cdots,x_n]$ and $w=(w_1,\cdots,w_n)$ be an $n$-tuple of positive integers. We let $\mathrm{wdeg}(x^\alpha)\coloneqq\sum w_i\alpha_i$, which is additive in the $\alpha_i$. Let $A_\nu$ be the $k$-vector space generated by monomials of $\mathrm{wdeg}=\nu$. Then $A=\bigoplus_{\nu\geq0}A_\nu$ is a graded $k$-algebra. The elements of $A_\nu$ are called \textbf{quasihomogeneous} or \textbf{weighted homogeneous} of degree $\nu$. 
    \end{itemize}
\end{example}
    
\newpage    

\section{Modules}
\begin{definition}{30. (Modules)}
    Let $A$ be a ring.
    \begin{itemize}
        \item An $A$-module is a set $M$ together with two operations
        \begin{align*}
            +:M\times M\to M\qquad\text{and}\qquad \cdot:A\times M\to M
        \end{align*}
        such that $(M,+)$ is an abelian group and for all $a,b\in A,\ m,n\in M$, we have
        \begin{align*}
            &\bullet\ (a+b)\cdot m=a\cdot m+b\cdot m
            &&\bullet\ a\cdot(m+n)=a\cdot m+a\cdot n \\
            &\bullet\ (ab)\cdot m=a\cdot(b\cdot m) &&\bullet\ 1\cdot m=m
        \end{align*}
        \item Let $M,\ N$ be $A$-modules. A map $\phi:M\to N$ is called an $A$-\textbf{module homomorphism} or is $A$-\textbf{liear} if for all $a\in A,\ m,n\in M$, 
        \begin{align*}
            \phi(a\cdot m)=a\cdot\phi(m)\qquad\text{and}\qquad
            \phi(m+n)=\phi(m)+\phi(n).
        \end{align*}
        In particular, $\phi$ is called an $A$-\textbf{module endomorphism} if $M=N$.
        \item $\mathrm{Hom}_A(M,N)$ [respectively, $\mathrm{End}_A(M)$] denotes the set of all $A$-module homomorphisms $M\to N$ [respectively, endomorphisms $M\to M$].
        \item $A$-modules together with $A$-linear maps form a category. Isomorphisms in this category are bijective $A$-linear maps $\phi:M\to N,\ \phi^{-1}$ being automatically $A$-linear. If such an isomorphism $M\to N$ exists, then $M$ and $N$ are called \textbf{isomorphic} and denote $M\simeq N$.
    \end{itemize}
\end{definition}

\begin{definition}{31. (Submodules)}
    \begin{itemize}
        \item Let $M$ be an $A$-module. An abelian subgroup $N$ of $M$ is an $A$-\textbf{submodule} of $M$ if it is closed under scalar multiplication, i.e. for all $a\in A,\ n\in N$ , we have $a\cdot n\in N$.
        \item Let $N\subset M$ be an $A$-submodule. Then $m\sim m'\xLeftrightarrow{\text{Def.}} m-m'\in N$ defines an equivalence relation on $M$ with equivalence class $\overline{m} =\overline{m+n},\ (n\in N)$. The \textbf{quotient module} $M/N$ is an $A$-module structure on the coset space $\{\overline{m}:m\in M\}$ defined by the following operations:
        \begin{align*}
            \overline{m}+\overline{m'}\coloneqq\overline{m+m'}\qquad \text{and} \qquad \overline{m}\cdot\overline{m'}\coloneqq\overline{m\cdot m'}.
        \end{align*}
        (cf. \textbf{Def. 14} of quotient ring)
    \end{itemize}
\end{definition}
\newpage
\begin{example}{32}
    \begin{itemize}
        \item[(a)] Every ring is a module over itself with scalar multiplication induced by the ring product. An ideal $\mathcal{I}\triangleleft A$ is an $A$-module with scalar multiplication induced by the ring product. In fact, $\mathcal{I}$ is an $A$-submodule.
        \item[(b)] Scalar multiplication $a\cdot\overline{b}\coloneqq \overline{a\cdot b}$ turns $A/\mathcal{I}$ into an $A$-module.
        \item[(c)] $A^n=\{(a_1,\cdots,a_n):a_i\in A\}$ is an $A$-module in the obvious way.
        \item[(d)] $A[x_1,\cdots,x_n]$ is an $A$-module with $a\cdot \sum a_\alpha x^\alpha\coloneqq\sum(a\cdot a_\alpha)x^\alpha$.
    \end{itemize}
\end{example}

\begin{definition}{33. (Graded Modules)}
    A module $M$ over a ring $A\coloneqq\bigoplus_{\nu\in\mathbb{N}}A_\nu$ is a \textbf{graded} $A$-\textbf{module} if there is a direct sum decomposition $M\coloneqq\bigoplus_{\mu\in\mathbb{Z}}M_\nu$ into abelian groups $M_\mu$ with $A_\nu M_\mu\subset M_{\mu+\nu}$.
\end{definition}

\begin{example}{34}
    Let $A\coloneqq\bigoplus_{\nu\in\mathbb{N}}A_\nu$ be a graded $k$-algebra and consider the free $A$-module $A^m=\bigoplus_{i=1}^m Ae_1$. Let $\mathrm{deg}\ e_i \coloneqq\nu_i\in\mathbb{Z}$ and $M_\mu$ be the $k$-vector space generated by $\{ae_i:a\in A_{\nu-\nu_i}\}$. Then $A^m=\bigoplus_{\nu\in\mathbb{Z}}M_\nu$ is a graded $A$-module.
\end{example}

\begin{lemma}{35. (Graded Submodules)}
    Let $M\coloneqq\bigoplus_{\nu\in\mathbb{Z}}M_\nu$ be a graded $A$-module and $N\subset M$ be an $A$-submodule. Then T.F.A.E.:
    \begin{itemize}
        \item[(a)] $N$ is graded by the induced grading, i.e. $N\coloneqq\bigoplus_{\nu\in\mathbb{Z}}N_\nu$ with $N_\nu=M_\nu\cap N$.
        \item[(b)] $N$ is generated by homogeneous elements.
        \item[(c)] Let $m\coloneqq\sum m_\nu,\ (m_\nu\in M)$. Then $m\in N.\ \Leftrightarrow\ m_\nu\in N,\ \forall\ \nu$. 
    \end{itemize}
    A submodule satisfying one and thus all of these conditions is called a \textbf{graded} or \textbf{homogeneous submodule} of $M$. Moreover, if $M=A$, where $A$ is a graded ring, we call $N$ a \textbf{graded} or \textbf{homogeneous ideal}.
\end{lemma}

\begin{remark}{36}
    Let $M\coloneqq\bigoplus_{\nu\in\mathbb{Z}}M_\nu$ be a graded $A$-module and $N\subset M$ be a graded submodule. Then the quotient $M/N$ is a graded module with $(M/N)_\nu=M_\nu/N_\nu\cong(M_\nu+N)/N$. Moreover, if $M=A$ is a graded ring and $N=\mathfrak{a}$ is a homogeneous ideal, then $A/\mathfrak{a}$ is a graded ring.
\end{remark}

\begin{definition}{37. (Homogeneous Degree)}
    Let $A\coloneqq\bigoplus_{\nu\geq0}A_\nu$ be a graded ring and $M\coloneqq\bigoplus_{\nu\in\mathbb{Z}}M_\nu,\ N\coloneqq \bigoplus_{\nu\in\mathbb{Z}}N_\nu$ be two graded $A$-modules. An $A$-linear map\\ $\phi:M\to N$ is \textbf{graded} or \textbf{homogeneous} of \textbf{degree} $d$ if $\phi(M_\nu)\subset N_{\nu+d}$. We call $\phi$ simply \textbf{homogeneous} if it is homogeneous of degree $0$.
\end{definition}

\begin{example}{38}
    Under the assumptions of the previous definition, multiplication with $f\in A_d$ defines a graded morphism $\phi:M\to M,\ m\mapsto f\cdot m$ of degree $d$.
\end{example}

\begin{remark}{39}
    If $\phi:M\to N$ is a homogeneous morphism, then $\mathrm{ker}\phi,\ \mathrm{im}\phi$ and $\mathrm{coker}\phi$ are graded $A$-submodules.
\end{remark}
\vspace{5mm}
Let $\phi:M\to N$ be an $A$-linear map. We know that its kernel $\mathrm{ker}\phi= \{m\in M:\phi(m)=0\}$ and its image $\mathrm{im}\phi=\{n\in N:n=\phi(m)\ \text{for some}\ m\in M\}$ are $A$-submodules of $M$ and $N$, respectively.
\begin{definition}{40. (Exact Sequences)}
    \begin{itemize}
        \item A sequence of $A$-modules with $A$-linear maps
        \begin{equation*}
            L\ \xrightarrow{~~\alpha~~}\ M\ \xrightarrow{~~\beta~~}\ N 
        \end{equation*}
        is called a \textbf{complex} if $\mathrm{im}\alpha\subset\mathrm{ker}\beta$, i.e. $\beta\circ\alpha=0_L$; and is said to be \textbf{exact} at $M$ if $\mathrm{im}\alpha=\mathrm{ker}\beta$.
        \item A sequence 
        \begin{equation*}
            \cdots\xrightarrow{~~~~~~}\ M_{k-1}\ \xrightarrow{~\alpha_{k-1}~}\
            M_{k}\ \xrightarrow{~~\alpha_{k}~~}\ M_{k+1}\ \xrightarrow{~~~~~~}\cdots
        \end{equation*}
        is \textbf{exact} [respectively, \textbf{complex}] if it is exact [respectively, complex] at every $M_k$.
        \item An exact sequence of the form
        \begin{equation*}
            0\ \xrightarrow{~~~~~~}\ L\ \xrightarrow{~~\alpha~~~}\ M\ \xrightarrow{~~\beta~~~}\ N\ \xrightarrow{~~~~~~}\ 0
        \end{equation*}
        is called a \textbf{short exact sequence} of $A$-modules.
    \end{itemize}
\end{definition}

\begin{remark}{41}
    \begin{itemize}
        \item[(a)] $0\xrightarrow{~~~} M\xrightarrow{~\alpha~} N$ is exact  $\Longleftrightarrow\ \mathrm{ker}\alpha=\{0\}\ \Longleftrightarrow\ \alpha$ is injective.
        \item[(b)] $M\xrightarrow{~~~} N\xrightarrow{~\beta~} 0$ is exact  $\Longleftrightarrow\ \mathrm{im}\beta=N\ \Longleftrightarrow\ \beta$ is surjective.
        \item[(c)] $0\xrightarrow{~~~} M\xrightarrow{~\phi~} N\xrightarrow{~~~} 0$ is exact at $M$ and $N\ \Longleftrightarrow\ \phi$ is an isomorphism.
    \end{itemize}
\end{remark}
\newpage
\begin{example}{42}
    If $\phi:M\to N$ is an $A$-module homomorphism, then the natural inclusions $\iota:\mathrm{ker}\phi\to M$ and $\kappa:\mathrm{im}\phi\to N$ induce the following short exact sequences:
    \begin{equation*}
        0\ \xrightarrow{~~~~~~}\ \mathrm{ker}\phi\ \xrightarrow{~~\iota~~~}\ M\ \xrightarrow{~~\phi~~~}\ \mathrm{im}\phi\ \xrightarrow{~~~~~~}\ 0
    \end{equation*}
    and
    \begin{equation*}
        0\ \xrightarrow{~~~~~~}\ \mathrm{im}\phi\ \xrightarrow{~~\kappa~~~}\ N\ \xrightarrow{~~~~~~}\ \mathrm{coker}\phi\ \xrightarrow{~~~~~~}\ 0.
    \end{equation*}
\end{example}

\begin{proposition}{43. (Split Exact)}
    Let $0\to L\xrightarrow{\alpha} M\xrightarrow{\beta} N\to0$ be a short exact sequence. Then T.F.A.E.:
    \begin{itemize}
        \item[(a)] There exists an isomorphism $\phi:L\oplus N\to M$ such that
        \begin{align*}
            \alpha(l)=\phi(i\oplus0)\qquad\text{and}\qquad \beta\circ\phi(l\oplus n)=n.
        \end{align*}
        \item[(b)] There exists a section of $\beta$, i.e. a map $\kappa:N\to M$ such that $\beta\circ\kappa=1_N$.
        \item[(c)] There exists a retraction of $\alpha$, i.e. a map $\iota:M\to N$ such that $\iota\circ\alpha=1_M$.
    \end{itemize}
    A short exact sequence which admits a section is called \textbf{split exact}.
\end{proposition}

\begin{definition}{44. (Localisation of Modules)}
    If $M$ is an $A$-module contained in some field $L\supset K=\mathrm{Quot}A$, then $S^{-1}M$ denotes the set $\{\frac{m}{s}:m\in M,\ s\in S\}$. This becomes an $S^{-1}A$-module in the obvious way and is called the \textbf{localisation} of $M$. For example, let $\mathfrak{b}$ be an ideal of $S^{-1}A\subset K=L$, then $S^{-1}(\mathfrak{b}\cap A)=\mathfrak{b}$ as $S^{-1}A$-modules.
\end{definition}

\begin{theorem}{45. (Nakayama's Lemma)}
    Let $(A,\mathfrak{m})$ be a local ring, $\mathfrak{a}\subset A$ be a proper ideal and $\mathfrak{a}\cdot M\coloneqq\{\sum a_im_i:a_i\in \mathfrak{a},\ m_i\in M\}$. If $N$ is an $A$-submodule of $M$ such that $N+\mathfrak{a}M=M$, then $N=M$. In particular, $\mathfrak{a}M =M$ implies $M=0$.
\end{theorem}

\begin{corollary}{46}
    Let $(A,\mathfrak{m})$ be a local ring and $M$ be a finitely generated $A$-module. If $m_1,\cdots,m_n\in M$ induce generators of the $k=A/\mathfrak{m}$-vector space $M/\mathfrak{m}M$, then they also generate $M$.
\end{corollary}

\begin{theorem}{47. (Modules over PID)}
    Let $A$ be a PID and $M$ a finitely generated $A$-module, then $M$ is a direct sum of cyclic modules, that is, modules of the form $A/(a)$. In particular, $M$ is free if it is torsion-free.
\end{theorem}

\newpage

\section{Algebras}
\begin{definition}{48. (Algebras)}
    \begin{itemize}
        \item An $A$-\textbf{algebra} is a ring $B$ which contains the ring $A$ as a subring. In particular, $B$ is an $A$-module by restricting multiplication to $A$.
        \item Let $B$ and $C$ be two $A$-algebras. A ring homomorphism $\phi:B\to C$ is called an $A$-\textbf{algebra homomorphism} if $\phi(a)=a,\ \forall a\in A$, i.e. $\phi|_A=1_A$. In particular, an $A$-algebra homomorphism is a ring homomorphism which is also $A$-linear.
    \end{itemize}
\end{definition}

\begin{remark}{49}
    \begin{itemize}
        \item[(a)] $A$-algebras together with $A$-algebra homomorphisms form a category, where isomorphisms are given by bijective $A$-algebra homomorphisms.
        \item[(b)] A ring homomorphism $\phi:A\to B$ turns $B$ into a $\phi(A)$-algebra which by abuse of language we also call an $A$-algebra. 
        \item[(c)] If $B$ and $C$ are $A$-algebras via the ring homomorphisms $\phi:A\to B$ and $\psi:A\to C$. Then the ring homomorphism $\alpha:B\to C$ is an $A$-algebra homomorphism iff $\alpha\circ\phi=\psi$.
        \item[(d)] In general, $\mathrm{Hom}_{A-\text{Alg.}}(B,C) \subsetneq\mathrm{Hom}_{\text{Ring}}(B,C)$ for $A$-algebras $B$ and $C$.
    \end{itemize}
\end{remark}

\begin{example}{50}
    Considering $A$ as the constants in the polynomial ring $A[x_1,\cdots,x_n]$, the latter becomes an $A$-algebra in a natural way. Moreover, the natural ring isomorphism $A[x_1,\cdots,x_{n-1}][x_n]\cong A[x_1,\cdots,x_n]$ is actually an $A$-algebra homomorphism.\\ On the other hand, $\phi:\mathbb{C}[x]\to \mathbb{C}[x],\ \sum\lambda_ix^i\mapsto\sum\bar{\lambda}_ix^i$ is clearly a ring homomorphism, but not a $\mathbb{C}$-algebra homomorphism.
\end{example}

\begin{proposition}{51}
    Let $B$ be an $A$-algebra and $\alpha_1,\cdots,\alpha_n\in B$. Then there exists a unique $A$-algebra homomorphism $\phi:A[x_1,\cdots,x_n]\to B$ determined by $\phi(x_i)=\alpha_i$, the so-called \textbf{substitution homomorphism}.
\end{proposition}

\begin{definition}{52. (Finitely Generated Algebras)}
    An $A$-algebra $B$ is finitely generated if there exists a surjective $A$-algebra homomorphism $\phi:A[x_1,\cdots,x_n]\twoheadrightarrow B$.\\ (cf. \textbf{Prop.51})
\end{definition}

\begin{remark}{53}
    If $B$ is a finitely generated $A$-algebra, then $B\cong A[x_1,\cdots,x_n]/\mathrm{ker}\phi$ as $A$-algebras. Setting $b_i\coloneqq\phi(x_i)\in B,\ (1\leq i\leq n)$, we have $\mathrm{im}\phi=\{\sum a_\alpha b^\alpha:a_\alpha\in A\}$, where $b^\alpha=b_1^{\alpha_1}\cdots b_n^{\alpha_n}$. We also write $B=A[b_1,\cdots,b_n]$ and call it the $A$-\textbf{algebra generated} by $b_1,\cdots,b_n\in B$.
\end{remark}

\section{Finiteness}
\begin{definition}{54. (Finite Modules \& Finite Type)}
    \begin{itemize}
        \item It is common to call an $A$-module $M$ \textbf{finite} over $A$ if $M$ is a finitely generated $A$-module. 
        \item An $A$-algebra $B$ is of \textbf{finite type} over $A$ if $B$ is a finitely generated $A$-algebra.
    \end{itemize}
    Now, we have the following:
    \begin{align*}
        &\bullet\ M\ \text{finite over}\ A. &&\Longleftrightarrow\ \phi:A^n\twoheadrightarrow M\ \text{$A$-linear}.\\
        & &&\Longleftrightarrow\ M=\{\sum a_i\phi(e_i):a_i\in A\}=(\phi(e_1),\cdots,\phi(e_n)).\\
        &\bullet\ B\ \text{finite type over}\ A. &&\Longleftrightarrow\ \phi:A[x_1,\cdots,x_n]\twoheadrightarrow B\ \text{$A$-algebra homomorphism}.\\
        & &&\Longleftrightarrow\ B=\{\sum a_\alpha\phi b^\alpha:a_\alpha\in A\}=A[b_1,\cdots,b_n],\ (b_i\coloneqq\phi(e_i)).
        \end{align*}
\end{definition}

\begin{remark}{55}
    An $A$-algebra $B$ which is finite as an $A$-module is also of finite type, but the converse is false in general. It holds in special situations such as $A\subset B$ is an integral ring extensions. (cf. \textbf{Rmk.} after \textbf{Def. 94})
\end{remark}

\begin{definition}{56. (Noetherian Rings \& Modules)}
    Let $M$ be an $A$-module. Then $M$ is called \textbf{Noetherian} if every $A$-submodule of $M$ is finitely generated. In particular, $A$ is called a \textbf{Noetherian ring} if it is Noetherian as an $A$-module over itself.
\end{definition}

\begin{theorem}{57. (Hilbert's Basissatz)}
    If $A$ is Noetherian, then so is $A[x]$. [AM, 7.5]
\end{theorem}

\begin{proposition}{58}
    \begin{itemize}
        \item[(a)] Submodules and quotient modules of Noetherian modules are Noetherian. [AM, 6.3]
        \item[(b)] Let $N\subset M$ be $A$-modules, then $M$ is Noetherian $\Longleftrightarrow$ $N$ and $M/N$ are Noetherian.
        \item[(c)] Let $M$ be an $A$-module, then T.F.A.E.:
        \begin{itemize}
            \item[(i)] $M$ is Noetherian.
            \item[(ii)] Every ascending chain of $A$-submodules $M_1\subset M_2\subset\cdots\subset M_k\subset\cdots$ becomes stationary.
            \item[(iii)] Every non-empty set of $A$-submodules of $M$ has a maximal element with respect to inclusion. [AM, 6.1]
        \end{itemize}
        \item[(d)] Let $A$ be a Noetherian ring and $M$ be a finitely generated $A$-module, then $M$ is a Noetherian $A$-module. [AM, 6.5]
    \end{itemize}
\end{proposition}

\newpage

\section{Tensor Products}
\begin{definition}{59. (Tensor Products)}
    Let $T$ be the free $A$-module generated by the pairs $(m,n)$, where $m\in M,\ n\in N$, and $U$ be the $A$-submodule of $T$ generated by the elements:
    \begin{align*}
        &\bullet\ (am,n)-a(m,n), &&\bullet\ (m+m',n)-(m,n)-(m',n),\\
        &\bullet\ (m,an)-a(m,n), &&\bullet\ (m,n+n')-(m,n)-(m,n')
    \end{align*}
    with $a\in A,\ m,m'\in M$ and $n,n'\in N$.\\
    The \textbf{tensor product} of $M$ and $N$ is the $A$-module $M\otimes_AN\coloneqq T/U$. The equivalence class of $(m,n)$ in $T/U$ will be denoted by $m\otimes n$; the set $\{m\otimes n:m\in M,\ n\in N\}$ is thus a generating system of $M\otimes_AN$. Hence, we may define an $A$-linear map $\iota:M\times N\to M\otimes_AN,\ (m,n)\mapsto m\otimes n$.
\end{definition}

\begin{theorem}{60. (Universal Property)}
    Let $M,\ N$ and $L$ be $A$-modules. If $\alpha:M\times N\to L$ is an $A$-bilinear map, then there exisets a unique map $\overline{\alpha}:M\otimes_AN\to L$ such that $\alpha=\overline{\alpha}\circ \iota$. In particular, for any other pair $(\widetilde{T},\tilde{\iota})$ consisting of an $A$-module $\widetilde{T}$ and an $A$-linear map $\tilde{\iota}:M\times N\to\widetilde{T}$ there exists a unique isomorphism $\psi:M\otimes_AN\to\widetilde{T}$ such that $\tilde{\iota}=\psi\circ\iota$.
\end{theorem}

\begin{corollary}{61}
    There is a canonical isomorphism of $A$-modules
    \begin{equation*}
        \mathscr{B}_A(M,N;P)\ \cong\ \mathrm{Hom}_A(M\otimes_AN,P).
    \end{equation*}
\end{corollary}

\begin{proposition}{62}
    Let $M,\ M',\ N,\ N'$ and $P$ be $A$-modules. We have the canonical isomorphisms:
    \begin{itemize}
        \item[(a)] $M\otimes_AN\cong N\otimes_AM,\ m\otimes n\mapsto n\otimes m$.
        \item[(b)] $(M\otimes_AN)\otimes_AP\cong M\otimes_A(N\otimes_AP),\ (m\otimes n)\otimes p\mapsto m\otimes(n\otimes p)$.
        \item[(c)] $A\otimes_AM\cong M,\ a\otimes m\mapsto am$.
        \item[(d)] $(M\oplus N)\otimes_AP\cong(M\otimes_AP)\oplus (N\otimes_AP),\ (m,n)\otimes p\mapsto(m\otimes p,\ n\otimes p)$.
    \end{itemize}
    Further, $A$-linear maps $\phi:M\to M'$ and $\psi:N\to N'$ define the $A$-linear map
    \begin{equation*}
        \phi\otimes\psi:M\otimes_AN\longrightarrow M'\otimes_AN',\ m\otimes n\longmapsto \phi(m)\otimes\psi(n).
    \end{equation*}
\end{proposition}

\begin{example}{63}
    \begin{itemize}
        \item[(a)] If $M=A^r$ and $N=A^s$ are two free $A$-modules of rand $r$ and $s$ with basis $\{e_i\}_{i=1}^r$ and $\{f_j\}_{j=1}^s$, respectively. Then $M\otimes_AN\cong A^{rs}$ is again a free $A$-module of rank $rs$ with basis $\{e_i\otimes f_j:1\leq i\leq r,\ 1\leq j\leq s\}$.
        \newpage
        \item[(b)] The subscript $A$ at $\otimes$ indicates which scalars can be shuffled around. For instance, consider the abelian groups (=$\mathbb{Z}$-modules) $\mathbb{Z}/m\mathbb{Z}$ and $\mathbb{Z}/n\mathbb{Z}$. Then $\mathbb{Z}\otimes_{\mathbb{Z}}(\mathbb{Z}/m\mathbb{Z})\cong \mathbb{Z}/m\mathbb{Z}$, etc. But $(\mathbb{Z}/m\mathbb{Z})\otimes_{\mathbb{Z}}(\mathbb{Z}/n\mathbb{Z})=0$ if $m$ and $n$ are coprime:\\
        Indeed, $am+bn=1$ for suitable integers $a,b\in\mathbb{Z}$ so that 
        \begin{equation*}
            \overline{x}\otimes\overline{y}=1\cdot\overline{x}\otimes\overline{y}=(am\cdot\overline{x})\otimes\overline{y}+\overline{x}\otimes(bn\cdot\overline{y})=0\otimes\overline{y}+\overline{x}\otimes0=0.
        \end{equation*}
        \item[(c)] Tensoring of maps does not preserve injectivity. Indeed, let $A=\mathbb{Z}$ and consider the map $\alpha:\mathbb{Z}\to\mathbb{Z},\ m\mapsto2m$. Then $\alpha\otimes1_{\mathbb{Z}/2\mathbb{Z}}$ is not injective since
        \begin{equation*}
            \alpha\otimes1_{\mathbb{Z}/2\mathbb{Z}}(m\otimes\overline{b})=2a\otimes\overline{b}=a\otimes(2\overline{b})=a\otimes0=0,
        \end{equation*}
        that is, $\alpha\otimes1_{\mathbb{Z}/2\mathbb{Z}}=0$. However, surjectivity is always preserved.
    \end{itemize}
\end{example}

\begin{definition}{64. (Tensor Product of Algebras)}
    \begin{itemize}
        \item Let $B$ and $C$ be two $A$-algebras via ring homomorphisms $f:A\to B$ and $g:A\to C$. Then $(b\otimes c)\cdot(b'\otimes c')\coloneqq bb'\otimes cc'$ defines a ring structure on the $A$-module $D\coloneqq B\otimes_AC$.
        \item Moreover, the ring homomorphism $A\to D,\ a\mapsto f(a)\otimes g(a)$ turns $D$ into an $A$-algebra with $A$-algebra homomorphisms
        \begin{equation*}
            \beta:B\to D,\ b\mapsto b\otimes1\qquad\text{and}\qquad \gamma:C\to D,\ c\mapsto1\otimes c.
        \end{equation*}
        \item The algebra $D$ has the following universal property: For any $A$-algebra $E$ with $A$-algebra homomorphisms 
        \begin{equation*}
            \beta':B\to E\qquad\text{and}\qquad \gamma':C\to E\qquad \text{satisfying}\qquad \beta'\circ f=\gamma'\circ g,
        \end{equation*}
        there exists a unique $A$-algebra homomorphism $\lambda:D\to E$ such that $\beta'=\lambda\circ\beta$ and $\gamma'=\lambda\circ \gamma$.
    \end{itemize}
    \begin{center}
        \begin{tikzcd}
            E & & \\
            & B \otimes_A C \arrow[ul, dashed, "\lambda"]
            & C \arrow[ull, bend right, "\gamma'"] \arrow[l, "\gamma"]\\
            & B \arrow[uul, bend left, "\beta'"] \arrow[u, "\beta"] 
            & A \arrow[l, "f"] \arrow[u, "g"]
        \end{tikzcd}
    \end{center}
\end{definition}
\newpage
\begin{example}{65}
    \begin{itemize}
        \item[(a)] (\textbf{Extension of Scalars}): If $B$ is an $A$-algebra and $M$ is an $A$-module, then $B\otimes_AM$ becomes a $B$-module via $b(b'\otimes m)\coloneqq(bb')\otimes m,\ (b,b'\in B,\ m\in M)$. If $M$ is a free $A$-module with basis $\{e_i\}_{i=1}^n$, then $B\otimes_AM$ is a free $B$-module with basis $\{1\otimes e_i\}_{i=1}^n$.\\
        For example, for a field extension $K\subset L$, we have a canonical isomorphism:
        \begin{equation*}
            L\otimes_KK[x_1,\cdots,x_n]\longrightarrow L[x_1,\cdots,x_n],\ \alpha\otimes f\longmapsto\alpha f.
        \end{equation*}
        \item[(b)] Let $B\coloneqq A[x_1,\cdots,x_n]/\mathfrak{a}$ and $C\coloneqq A[y_1,\cdots,y_m]/\mathfrak{b}$, then \\ \centerline{$B\otimes_AC=A[x_1,\cdots,x_n,y_1,\cdots,y_m]/(\mathfrak{a}+\mathfrak{b})$,}\\ where we consider $\mathfrak{a}$ and $\mathfrak{b}$ as ideals of $A[x_1,\cdots,x_n,y_1,\cdots,y_m]$ via extension.
        \item[(c)] These allow the computation of various tensor products. For instance, consider the field extension $\mathbb{Q}\subset\mathbb{Q}[\alpha]\subset\mathbb{C}$ with minimal polynomial $f$ of degree $r$, (cf. \textbf{Example 67}). Since $\mathbb{Q}[\alpha]$ is separable, $f$ splits completely into simple zeros, whence
        \begin{equation*}
            \mathbb{C}\otimes_{\mathbb{Q}}\mathbb{Q}[\alpha]\cong
            \mathbb{C}\otimes_{\mathbb{Q}}(\mathbb{Q}[x]/(f))\cong
            (\mathbb{C}[x]/(f))\cdot\mathbb{C}[x]\cong
            \prod_{i=1}^n\mathbb{C}[x]/(x-\alpha_i)
        \end{equation*}
        as $\mathbb{Q}$-algebras by \textbf{CRT}. As a $\mathbb{C}$-vector space, this is isomorphic to $\mathbb{C}^n$.
    \end{itemize}
\end{example}

\newpage

\section{Field Extensions}
\begin{definition}{66. (Field Extension)}
    A \textbf{field extension} is an embedding $K\xrightarrow{} L$ of two fields $K$ and $L$. In particular, we may view $L$ as a $K$-algebra and thus as a $K$-vector space. It is customary to write $L/K$ for the field extension and $[L:K]$ for $\mathrm{dim}_KL$, the \textbf{degree} of the field extension. We say a field extension is 
    \begin{itemize}
        \item[(i)] \textbf{finite}, if $\mathrm{dim}_KL=[L:K]<\infty$.
        \item[(ii)] \textbf{algebraic}, if for any $\alpha\in L$, there exists $f\in K[x]$ such that $f(\alpha)=0$.
    \end{itemize}
\end{definition}

\begin{example}{67}
    If $\alpha\in L$, then $\mathcal{I}\coloneqq\{f\in K[x]:f(\alpha)=0\}\triangleleft K[x]$ is an ideal. Unless $\mathcal{I}=(0)$, or $\mathcal{I}=(f)$ for an irreducible monic polynomial $f$, the so-called \textbf{minimal polynomial} of $\alpha$, since $K[x]$ is a PID. Hence $K(\alpha)$ (the subfield of $L$ generated by $K$ and $\alpha$) is isomorphic to $K[x]/(f)$. The field extension $K\subset K(\alpha)$ is called \textbf{simple} and is of degree $[K(\alpha):K]=\mathrm{deg}f$.
\end{example}

\begin{proposition}{68}
    A finite field extension is algebraic. But the converse is false!
\end{proposition}
\begin{proof}
    Consider the algebraic extension $\mathbb{Q}\subset\overline{\mathbb{Q}}= \{\alpha\in\mathbb{C}:\alpha\ \text{algebraic over}\ \mathbb{Q}\}\subset \mathbb{C}$.\\ But $\mathbb{Q}(\sqrt[n]{3})\subset\overline{\mathbb{Q}}$ has minimal polynomial $x^n-3$ since it is irreducible by \textbf{Einstein's Criterion}. It follows that $\mathrm{dim}_{\mathbb{Q}}
    \mathbb{Q}(\sqrt[n]{3})=n$. In particular, $\mathrm{dim}_{\mathbb{Q}} \overline{\mathbb{Q}}=\infty$.
\end{proof}

\begin{definition}{69. (Algebraically Closed)}
    A field $K$ is called \textbf{algebraically closed} if any polynomial $f\in K[x]$ splits into linear factors. For instance, $\mathbb{C}$ is algebraically closed (by \textbf{Fundamental Theorem of Algebra}).
\end{definition}

\begin{theorem}{70. (Algebraic Closure)}
    For any field $K$, there exists an algebraic field extension $K\subset L$ such that $L$ is algebraically closed. This field $L$ is determined up to $K$-isomorphism in 
    \begin{align*}
        \mathrm{Aut}_K(L)\coloneqq\{\sigma\in\mathrm{Hom}(L,L) :\sigma|_K=1_K\}.
    \end{align*}
    We therefore speak of the \textbf{algebraic closure} and denote it by $\overline{K}$.
\end{theorem}

\begin{example}{71}
    Let $K\subset L$ be an algebraic field extension with $L$ algebraically closed. We can realise the algebraic closure of $K$ by 
    \begin{align*}
        \overline{K}=\{\alpha\in L:\alpha\ \text{algebraic over}\ K\}.
    \end{align*}
    This is indeed a field and defines an algebraic extension $K\subset\overline{K}$.\\ For instance, $\overline{\mathbb{Q}}=\{\alpha\in \mathbb{C}:\alpha\ \text{algebraic over}\ \mathbb{Q}\}$ yields the algebraic closure of $\mathbb{Q}$.
\end{example}
\newpage
\begin{proposition}{72. (Normal Extension)}
    Let $K\subset L$ be an algebraic field extension. T.F.A.E.:
    \begin{itemize}
        \item[(a)] $L$ is a splitting field for a family of polynomials in $K[x]$, i.e. every polynomial from that family splits completely into linear factors in $L[x]$.
        \item[(b)] Every irreducible polynomial in $K[x]$ which has a zero in $L$ splits completely into linear factors in $L[x]$.
        \item[(c)] Every $K$-homomorphism $L\to\overline{K}$ restricts to an automorphism of $L$.
    \end{itemize}
    Any algebraic field extension which satisfies one and hence all of these conditions is called a \textbf{normal extension}.
\end{proposition}

An important example of normal extensions comes from finite fields:
\begin{theorem}{73}
    Let $\mathbb{F}$ be a finite field. Then $p=\mathrm{char}\mathbb{F}>0$ and $\mathbb{F}_p\subset\mathbb{F}$ is a normal extension of $\mathbb{F}_p$. In fact, it is the splitting field of $f=x^q-x$, where $q=\#\mathbb{F} =p^n,\ n\coloneqq[\mathbb{F}:\mathbb{F}_p]$. The elements of $\mathbb{F}$ are therefore just the zeros of $f$. Moreover, $\mathbb{F}$ is the unique finite field of $q=p^n$ elements (up to isomorphism), and any finite field is isomorphic to one of this type.
\end{theorem}

\begin{definition}{74. (Separable Polynomial)}
    A polynomial $f\in K[x]$ is \textbf{separable} if it has only simple roots. An irreducible polynomial $f\in K[x]$ is \textbf{separable} if either
    \begin{itemize}
        \item[(i)] $\mathrm{char}K=0$, or
        \item[(ii)] $\mathrm{char}K=p$ and $f$ is of the form $f(x)=g(x^{p^r}) ,\ \text{where}\ g\in K[x]\ \text{irreducible}$.
    \end{itemize}
    Every zero of $f$ is a $p^r$-th square root of a zero of $g$ has multiplicity $p^r$.
\end{definition}

\begin{definition}{75. (Separable Extension)}
    For an algebraic extension $L/K$, $\alpha\in L$ is \textbf{separable} if it is the zero of a separable polynomial. $L/K$ is called \textbf{separable} if every $\alpha\in L$ is separable. In particular, a field $K$ is called \textbf{perfect} if any algebraic extension of $K$ is separable.
\end{definition}

\begin{example}{76}
    The following types of fields are perfect: \\ 
    \centerline{(a) fields of characteristic 0; (b) finite fields; (c) algebraically closed fields.}
\end{example}

\begin{theorem}{77}
    Let $K\subset L$ be a finite extension. Then T.F.A.E.:
    \begin{itemize}
        \item[(a)] $L/K$ is separable.
        \item[(b)] There are separable elements $\alpha_1,\cdots,\alpha_r\in L$ such that $L=K(\alpha_1,\cdots,\alpha_r)$.
        \item[(c)] $\#\mathrm{Hom}_K(L,\overline{K})=[L:K]$, where $\mathrm{Hom}_K(L,\overline{K})\coloneqq$ field morphisms which restrict to the identity on $K$. 
    \end{itemize}
\end{theorem}

\begin{definition}{78. (Trace \& Norm)}
    Let $K\subset L$ be a finite field extension. The \textbf{trace} $\mathrm{Tr}_{L/K}(\alpha)$ and the \textbf{norm} $\mathrm{N}_{L/K}(\alpha)$ of $\alpha\in L$ is the trace and determinant of the $K$-linear map\\ $\mathsf{T}_\alpha:L\to L,\ x\mapsto\alpha\cdot x$, respectively.
\end{definition}

\begin{proposition}{79}
    Let $L/K$ be a finite separable field extension, $H\coloneqq \mathrm{Hom}_K (L,\overline{K})$ and $\alpha\in L$ with minimal polynomial $f_\alpha$. Then \\ \centerline{(a) $f_\alpha(\alpha)=\prod_{\sigma\in H}(t-\sigma(\alpha))$. 
    (b) $\mathrm{Tr}_{L/K}(\alpha)=\sum_{\sigma\in H}\sigma(\alpha)$. 
    (c) $\mathrm{N}_{L/K}(\alpha)=\prod_{\sigma\in H}\sigma(\alpha)$.}
\end{proposition}

\begin{remark}{80}
    If $K\subset L$ is a finite extension, then the minimal polynomial of $\alpha\in L$ is given by the minimal polynomial of the $K$-linear map $\mathrm{T}_\alpha$.
\end{remark}

\begin{theorem}{81. (Primitive Element Theorem)}
    If $K\subset L$ is a finite separable field extension, then $L=K(\alpha)$, that is, $K\subset L$ is a simple extension. Such an $\alpha\in L$ is called a \textbf{primitive element}.
\end{theorem}

\begin{definition}{82. (Galois Group)}
    A field extension $K\subset L$ is \textbf{Galois} if it is normal and separable. In this case, the \textbf{Galois group} $\mathrm{Gal}(L/K) \subset\mathrm{Aut}_KL$ of the field extension $K\subset L$ is the group of $K$-automorphisms of $L$.
\end{definition}

\begin{theorem}{83. (Main Theorem of Galois Theory)}
    Let $K\subset L$ be a finite Galois extension with Galois group $G\coloneqq \mathrm{Gal}(L/K)$. Then we have a correspondence between subgoups of $H\subset G$ and intermediate field extensions $K\subset E\subset L$ given by: 
    \begin{align*}
        H\mapsto E\coloneqq\{x\in L:\sigma(x)=x,\ \forall\sigma\in H\} \ \text{and}\ E\mapsto H\coloneqq\mathrm{Gal}(L/E).
    \end{align*}
\end{theorem}

\begin{definition}{84. (Transcendental Extension)}
    Consider a field extension $K\subset L$. Elements $\beta_1,\cdots,\beta_n\in L$ are \textbf{algebraically independent} over $K$ if the natural surjection \begin{align*}
        K[x_1,\cdots,x_n]\xtwoheadrightarrow{~~} K[\beta_1,\cdots,\beta_n],\ x_i\longmapsto\beta_i
    \end{align*}
    defines an injection $K[x_1,\cdots,x_n]\xhookrightarrow{} L\to0$. Equivalently, any polynomial relation of the form $f(\beta_1,\cdots,\beta_n)=0$ where $f\in K[x_1,\cdots,x_n]$, implies $f=0$. \\ A family $\mathfrak{B}\coloneqq\{\beta_i\}_{i\in I}$ is algebraically independent if the previous definition applies to any finite subset of $\mathfrak{B}$. It is called a \textbf{transcendence base} if in addition $K(\mathfrak{B})\subset L$ is algebraic. If $L=K(\mathfrak{B})$ for some transcendence base, then the field extension $K\subset L$ is called \textbf{purely transcendental}.
\end{definition}

\begin{proposition}{85. (Transcendental Degree)}
    Any field extension $K\subset L$ admits a (potentially empty) transcendence base $\mathfrak{B}$, and therefore can be factorised into a purely transcendental and an algebraic field extension $K\subset K(\mathfrak{B}) \subset L$. Furthermore, any two transcendental bases have the same cardinality. The latter is called he \textbf{transcendental degree} $\mathrm{trdeg}_KL$.
\end{proposition}

\chapter{Integral Ring Extensions \& Rings of Integers}

\section{The Gaussian Integers}
\begin{motivation}{86} 
    Consider the basic question, "When is a prime number the sum of two squares?"\ (e.g. the sequence $2=1^2+1^2,\ 5=1^2+2^2,\ 13=2^2+3^2$, etc.)
\end{motivation}

\begin{theorem}{87}
    For all prime numbers $p>2$, one has 
    \begin{align*}
        p=a^2+b^2,\ (a,b\in\mathbb{Z}).\qquad \Longleftrightarrow \qquad p\equiv1\ (\mathrm{mod}\ 4).
    \end{align*}
\end{theorem}

To prove \text{Thm. 87.} it is beneficial, as often in mathematics, to pass to a bigger domain of definition, and to consider the \textbf{Gaussian integers} $\mathbb{Z}[i]\coloneqq\{ a+bi:a,b\in\mathbb{Z}\}$, (cf. \textbf{Example 10}). 
The \textbf{norm} is the function $\mathrm{N}:\mathbb{Z}[i]\to\mathbb{N},\ a+bi\mapsto a^2+b^2=|a+bi|^2$. It's multiplicative, i.e. $\mathrm{N}(z\cdot w)= \mathrm{N}(z)\cdot\mathrm{N}(w)$.

\begin{proposition}{88}
    The ring $\mathbb{Z}[i]$ is Euclidean and in particular a UFD. Its group of units is given by $\mathbb{Z}[i]^\ast=\{\pm1,\pm i\}.$
\end{proposition}
\begin{proof}
    We prove the following two facts:\\
    \textsc{claim 1.}
    The restricted norm $\mathrm{N}\big|_{\mathbb{Z}[i]\setminus\{0\}}: \mathbb{Z}[i]\setminus\{0\}\to\mathbb{N}$ is a degree function.\\ Let $\alpha,\beta\in\mathbb{Z}[i],\ (\beta\neq0)$. If $\alpha=\gamma\cdot\beta$ for some $\gamma\in\mathbb{Z}[i]$, we are done. Otherwise, there are $\gamma,\rho\in\mathbb{Z}[i]$ with $\alpha=\gamma\cdot\beta+\rho,\ \rho\neq0,\ \mathrm{N}(\rho)<\mathrm{N}(\beta)$. Since $\mathbb{Z}[i]$ forms a Lattice in $\mathbb{C}$, we can find $\gamma\in\mathbb{Z}[i]$ with $|\frac{\alpha}{\beta}-\gamma|<\frac{\sqrt{2}}{2}<1$, (two points in a mesh of a lattice have a distance of at most half of diagonal). $\Rightarrow |\beta| |\frac{\alpha}{\beta}-\gamma|=|\alpha-\gamma\beta|<|\beta|$. Now put $\rho\coloneqq\alpha-\gamma\beta$.\\
    \textsc{claim 2.}
    $\alpha\in\mathbb{Z}[i]$ is a unit $\Leftrightarrow \mathrm{N}(\alpha) =1$.\\ "$\Rightarrow$"\ $\alpha\cdot\beta=1.\ \Rightarrow \mathrm{N}(\alpha)\cdot \mathrm{N}(\beta)=1.\ \Rightarrow \mathrm{N}(\alpha)$ is a unit in $\mathbb{Z}.\ \Rightarrow \mathrm{N}(\alpha)=1$. \\ "$\Leftarrow$"\ Write $1\coloneqq\beta\cdot \alpha+\rho$. If $\rho\neq0$, we may assume that $\mathrm{N}(\rho)<\mathrm{N}(\alpha)=1.\ \Rightarrow p=0$ and $\alpha\in\mathbb{Z}[i]$. Hence $\mathrm{N}(\alpha)=\mathrm{N}(a+bi)=a^2+b^2=1.\ \Rightarrow (a,b)=(\pm1,0)$ or $(0,\pm i)$.
\end{proof}
\begin{proof}
    (i) Wir zeigen, dass $\mathbb{Z}[i]$ euklidisch ist bzgl. der Funktion $\mathbb{Z}[i]\to\mathbb{N}\cup\{0\},\ \alpha\mapsto|\alpha|^2$. Sind $\alpha,\beta\in\mathbb{Z}[i],\ \beta\neq0$, so ist die Existenz von Gau\ss schen Zahlen $\gamma,\rho$ nachzuweisen mit\\ \centerline{$\alpha=\gamma\beta+\rho\qquad$ und$\qquad |\rho|^2<|\beta|^2$.}\\
    Es gen\"{u}gt offenbar, ein $\gamma\in\mathbb{Z}[i]$ zu finden mit $|\frac{\alpha}{\beta}-\gamma|<1$. Die Gau\ss schen Zahlen bilden nur ein \textbf{Gitter} (Lattice) in der komplexen Zahlenebene $\mathbb{C}$ (Punkte mit ganzzahligen Koordinaten bzgl. der Basis $1$ und $i$). Die komplexe Zahl $\frac{\alpha}{\beta}$ liegt in einer Masche des Gitters und hat vom n\"{a}chsten Gitterpunkt einen Abstand, der nicht gr\"{o}\ss er ist, als die halbe L\"{a}nge $\frac{\sqrt{2}}{2}$ der Diagonalen der Masche. Daher gibt es ein $\gamma\in\mathbb{Z}[i]$ mit $|\frac{\alpha}{\beta}-\gamma|<\frac{\sqrt{2}}{2}<1$. [N, I.1.2]
\end{proof}

\begin{proofthm}{87} 
    "$\Rightarrow$"\ $\because 2\nmid p.\ \therefore p$ must be the sum of squares of an odd and an even integer. $\Rightarrow p\coloneqq(2k+1)^2+(2l)^2=4k^2+4k+1+4l^2 \equiv1\ (\mathrm{mod}\ 4)$.\\
    "$\Leftarrow$"\ Suppose $p\equiv1\ (\mathrm{mod}\ 4)$ is a prime in $\mathbb{Z}$ but not in $\mathbb{Z}[i]$. Write $p\coloneqq1+4n$, then\\ $-1\equiv [(2n)!]^2\ (\mathrm{mod}\ p)$. Indeed, \textbf{Wilson's Thm.} implies that $-1\equiv(p-1)!\ (\mathrm{mod}\ p)$, whence 
    \begin{align*}
        -1 &\ \equiv (p-1)!\ \equiv (p-1)\cdot(p-2)\cdots (p-2n)\cdot(2n)\cdot(2n-1)\cdots1 \\ &\ \equiv\  (-1)^{2n}\cdot(2n)!\cdot(2n)!\ 
        \equiv\ [(2n)!]^2\qquad (\mathrm{mod}\ p).
    \end{align*}
    (Note: $(p-1)(p-2)\cdots(p-2n)=\sum_{i=1}^{2n}a_ip^i+(-1)^{2n}(2n)!$, for coefficients $a_i\in\mathbb{Z}$). \\ Setting $X\coloneqq [(2n)!]^2$, it follows that $p|x^2+1=(x+i)(x-i)$. But $p\nmid x\pm i$ since $\frac{x}{p}\pm\frac{i}{p}\notin\mathbb{Z}[i]$, hence $p$ is not prime in $\mathbb{Z}[i]$. 
    $\because\mathbb{Z}[i]$ is a UFD. $\therefore p\coloneqq\alpha\cdot \beta$ for some nonunits $\alpha,\beta\in\mathbb{Z}[i]$.\\ Moreover, $\mathrm{N}(\alpha),\mathrm{N}(\beta)>1$. Hence, $\mathrm{N}(p)= p^2=\mathrm{N}(\alpha)\cdot\mathrm{N}(\beta)$, which implies\\
    \centerline{$\mathrm{N}(\alpha)=\mathrm{N}(\beta)=p=a^2+b^2$,\ if $\alpha\coloneqq a+bi$ in $\mathbb{Z}[i]$.} \qed
\end{proofthm}

\begin{proposition}{89}
    Up to multiplying with units, the primes $\pi$ of $\mathbb{Z}[i]$ are:
    \begin{itemize}
        \item[(i)] $\pi=1+i$;
        \item[(ii)] $\pi=a+bi$ with $a^2+b^2=p,\ p\equiv1\ (\mathrm{mod}\ 4)$ and $a>|b|>0$;
        \item[(iii)] $\pi=p$, where $p\equiv3\ (\mathrm{mod}\ 4)$,
    \end{itemize}
    where $p$ denotes a prime number of $\mathbb{Z}$.
\end{proposition}
\begin{proof}
    The first two cases imply $\mathrm{N}(\pi)=p=\mathrm{N}(\alpha)\cdot \mathrm{N}(\beta)$ for $\pi=\alpha\cdot\beta$ and $p\in\mathbb{Z}$ prime. So either $\mathrm{N}(\alpha)$ or $\mathrm{N}(\beta)$ is a unit, whence $\pi$ is a prime. If $p\equiv3\ (\mathrm{mod}\ 4)$ and $\pi=p=\alpha\cdot\beta,\ (\alpha, \beta\in\mathbb{Z}[i]).\ \Rightarrow \mathrm{N}(\pi)=p^2=\mathrm{N}(\alpha)\cdot \mathrm{N}(\beta)$, which gives $p\equiv1\ (\mathrm{mod}\ 4)$, unless $\mathrm{N}(\alpha)$ or $\mathrm{N}(\beta)$ is a unit. \\
    Conversely, suppose that $\pi\in\mathbb{Z}[i]$ is prime. Then $\mathrm{N}(\pi)=\pi\cdot\overline{\pi}= p_1\cdot\cdots\cdot p_r,\ (p_i\in \mathbb{Z}$: prime). Hence $\pi|p$ for some $p=p_i$, and therefore $\mathrm{N}(\pi)|\mathrm{N}(p)=p^2$, so either $\mathrm{N}(\pi)=p$ or $\mathrm{N}(\pi)=p^2$.\ If $\mathrm{N}(\pi)=p$, then $\pi=a+bi$ with $a^2+b^2=p$ so that $p$ is of type (i) or (ii). If $\mathrm{N}(\pi)=p^2$, then writing $p\coloneqq\alpha\cdot\pi\ (\alpha\in\mathbb{Z}[i])$ gives $\mathrm{N}(\alpha) \cdot\mathrm{N}(\pi)=p^2$ and thus $\alpha$ must be a unit. But then $p$ is prime in $\mathbb{Z}[i]$  and they cannot be of the form $p=2=(1+i)(1-i)$ or $p\equiv1\ (\mathrm{mod}\ 4)$, whence $\pi$ is of type (iii). 
\end{proof}
\begin{proof}
    Die Zahlen unter (i) und (ii) sind prim, weil as einer Zerlegung $\pi=\alpha\cdot\beta$ in $\mathbb{Z}[i]$ die Gleichung \\
    \centerline{$p= \mathrm{N}(\pi)=\mathrm{N}(\alpha)\cdot\mathrm{N}(\beta)$}
    mit einer Primzahl $p$ folgt, so dass entweder $\mathrm{N}(\alpha)=1$ oder $\mathrm{N}(\beta)=1$, also entweder $\alpha$ oder $\beta$ eine Einheit (unit) ist. Die Zahlen $\pi=p,\ (p\equiv3\ \mathrm{mod}\ 4)$, sind prim in $\mathbb{Z}[i]$, weil eine Zerlegung $p=\alpha\cdot\beta$ in Nicht-Einheiten $\alpha,\beta$ zur Folge h\"{a}tte, dass $p^2=\mathrm{N}(\alpha)\cdot\mathrm{N}(\beta)$ ist, d.h. $p=\mathrm{N}(\alpha)= \mathrm{N}(a+bi)=a^2+b^2$, woraus sich nach \textbf{Satz. 87} $p\equiv1\ \mathrm{mod}\ 4$ erg\"{a}be.\\
    Nach dieser Feststellung haben wir zu zeigen, dass ein bliebiges Primelement $\pi$ von $\mathbb{Z}[i]$ assoziiert ist zu einem der Genannten. Zun\"{a}chst folgt aus \\ \centerline{$\mathrm{N}(\pi)=\pi\cdot\overline{\pi}= p_1\cdot\cdots\cdot p_r,\qquad p_i:$ Primzahl in $\mathbb{Z}$,} dass $\pi|p$ f\"{u}r ein $p=p_i$, also $\mathrm{N}(\pi)|\mathrm{N}(p)=p^2$, d.h. entweder $\mathrm{\pi}=p$ oder $\mathrm{\pi}=p^2$. Im Falle $\mathrm{\pi}=p$ ist $\pi=a+bi$ mit $a^2+b^2=p$, d.h. $\pi$ ist vom Typ (ii) oder, wenn $p=2$ ist, assoziiert zu $1+i$. Ist aber $\mathrm{\pi}=p^2$, so ist $\pi$ zu $p$ assoziiert, weil $\frac{p}{\pi}$ wegen  $\mathrm{\frac{p}{\pi}} =1$ eine Einheit ist. Es muss \"{u}berdies $p\equiv3\ \mathrm{mod}\ 4$ gelten, weil sonst $p=2$ oder $p\equiv1\ \mathrm{mod}\ 4$ und nach Satz. 87, $p=a^2+b^2=(a+bi)(a-bi)$ nicht prim w\"{a}re. Damit ist alles gezeigt. [N, I.1.4]
\end{proof}

\begin{nt}{89*}
    $\mathbb{Z}[i]=\mathbb{Z}\cdot1\oplus\mathbb{Z}\cdot i$ is a free $\mathbb{Z}$-module in $\mathbb{Q}(i)=\{ a+bi:a,b\in\mathbb{Q}\}$ and plays the same role as $\mathbb{Z}=\mathbb{Z}\cdot1$ in $\mathbb{Q}$. For a basis independent characterisation we have the following proposition.
\end{nt}

\begin{proposition}{90}
    $\mathbb{Z}[i]$ consists precisely of those elements in $\mathbb{Q}(i)$ which satisfy a monic polynomial equation $x^2+ax+b=0,\ (a,b\in\mathbb{Z})$.
\end{proposition}
\begin{proof}
    Note that $c+di\in\mathbb{Q}(i)$ is a root of the polynomial $x^2+\alpha x+\beta \in\mathbb{Q}[x]$ with 
    \begin{align*}
        \alpha\coloneqq-2c\qquad\text{and}\qquad \beta\coloneqq c^2+d^2.
    \end{align*}
    Hence $\alpha,\beta\in\mathbb{Z}$ if $c,d\in\mathbb{Z}$.\\
    Conversely, if $\alpha,\beta\in\mathbb{Z}$, then $2c,2d\in\mathbb{Z}$. (Note: $(2d)^2=4\beta-(2c)^2$ and a rational number whose square is in $\mathbb{N}$ must be an integer). It follows that $(2c)^2+(2d)^2=4\beta\equiv0\ (\mathrm{mod}\ 4)$. Since all squares must be either $0$ or $1\ (\mathrm{mod}\ 4)$, we must have $(2c)^2\equiv(2d)^2\equiv0\ (\mathrm{mod}\ 4)$, so that $c^2,d^2\in\mathbb{Z}$, whence $c,d\in\mathbb{Z}$. 
\end{proof}

\begin{remark}{90*}
    This result sets out the goal of this section: Given a finite field extension $\mathbb{Q}\subset K$, try to find a natural ring extension $\mathbb{Z}\subset A$ characterised in the vein of \textbf{Prop. 81}, and analyse its $\mathbb{Z}$-module structure.
\end{remark}

\section{Integral Ring Extensions}
\begin{definition}{91. (Integral Extension)}
    Let $A\subset B$ be a ring extension. We call $b\in B$ \textbf{integral} over $A$ if there is a monic polynomial $f\in A[x]$ such that $f(b)=0$. If every $b\in B$ is integral over $A$, then $B$ is said to be \textbf{integral} over $A$. In analogy with algebraic field extensions we also say that $A\subset B$ is an \textbf{integral extension}.
\end{definition}

\begin{example}{92}
    \begin{itemize}
        \item $\tau=\frac{1+\sqrt{5}}{2}\in\mathbb{Z}[\tau]$ (the golden ratio) is integral over $\mathbb{Z}$, where $\mathbb{Z}[\tau]\subset\mathbb{Q}$ is the ring generated by $\mathbb{Z}$ and $\tau$. Indeed, $\tau^2-\tau-1=0$.
        \item $\mathbb{Z}\subset\mathbb{Z}[\frac{1}{2}]= \{\sum_{i=0}^n \frac{m_i}{2^i}:m_i\in\mathbb{Z},\ n\in\mathbb{N}\}$ is not an integral extension. \\ Indeed, consider $x=\frac{p}{2}\in\mathbb{Z}[\frac{1}{2}]$ with $p>2$ prime. If we had a polynomial relation
        \begin{equation*}
            (\tfrac{p}{2})^n+c_{n-1}(\tfrac{p}{2})^{n-1}+\cdots+c_0=0,\qquad (c_i\in \mathbb{Z}),
        \end{equation*}
        then multiplying with $2^n$ shows that $p^n=-2(c_{n-1}(\frac{p}{2})^{n-1}+\cdots+c_02^{n-1})$, hence $2\big|p,\ (\to\gets)$.
        \item For an integral domain $A$ we have the natural ring extension $A\subset K=\mathrm{Quot}(A)$.\\ If $A$ is a UFD, then $x\in K$ is integral over $A$ iff $x\in A$. [Exercise]
    \end{itemize}
\end{example}

\begin{proposition}{93}
    Let $A\subset B$ be a ring extension, and let $b\in B$. Then T.F.A.E.:
    \begin{itemize}
        \item[(a)] b is integral over $A$.
        \item[(b)] The subring $A[b]\subset B$ is a finite $A$-module. [AM, 5.1]
    \end{itemize}
\end{proposition}
\begin{proof}
    "$(a)\Rightarrow(b)$"\  If $b$ satisfies a monic relation of the form $b^n=- \sum_{i=0}^{n-1}a_ib^i,\ (a_i\in A)$, then $A[b]$ is generated by $1,\ b,\cdots,\ b^{n-1}$. \\
    "$(b)\Rightarrow(a)$"\ Consider $b$ as a map $\mu_b:A[b]\to A[b],\ a\mapsto b\cdot a$. Since $A[b]$ is a finite $A$-module,\\ $\mu_b$ satisfies a monic relation $\mu_b^n+a_{n-1}\mu_b^{n-1}+\cdots+a_0=0,\ (a_i\in A)$ in $\mathrm{End}(M)$, by Cayley-Hamilton [AM, 2.4]. Evaluate at $1$ gives (a).
\end{proof}

\begin{definition}{94}
    A ring extension $A\subset B$ is \textbf{finite} if $B$ is a finite $A$-module.
\end{definition}
\begin{remark}{94*}
    It follows from \textbf{Prop. 93} that a finite ring extension is integral. In fact, any finite ring extension $A\subset B$ is of the form $B\coloneqq A[b_1,\cdots,b_n],\ (b_i:$ integral over $A$). Hence, an $A$-module $B$ with $A\subset B$ is of 
    \begin{align*}
        \textbf{ finite type}\ +\ \textbf{integral}\ \Longleftrightarrow\ \textbf{finite}. 
    \end{align*}
\end{remark}

\begin{proposition}{95}
    If $A\subset B\subset C$ are finite [integral] ring extensions, then $A\subset C$ is a finite [integral] ring extension. [AM, 5.4] 
\end{proposition}
\begin{proof}
    (a) By \textbf{Prop. 93} and \textbf{Rmk.}, $B$ is generated as an $A$-nodule by $b_1,\cdots,b_n\in B$, and $C$ is generated as a $B$-module by $c_1,\cdots,c_m\in C$. Hence $C$ is generated as an $A$-module by $b_ic_j$'s. \\ (b) Let $c\in C$ satisfy the relation 
    \begin{equation*}
        c^n+b_{n-1}c^{n-1}+\cdots+b_1c+b_0=0,\qquad (b_i\in B).
    \end{equation*}
    Then $c$ is also integral over $A[b_0,\cdots,b_{n-1}]$, whence $A[b_0,\cdots,b_{n-1},c]$ is finite over $A[b_0,\cdots,b_{n-1}]$. Since each $b_i$ is integral over $A$, $A[b_0,\cdots,b_{n-1}]$ is finite over $A$, hence $A[b_0,\cdots,b_{n-1},c]$ is finite over $A$. It follows that $c$ is integral over $A$.
\end{proof}
\newpage
\begin{lemma}{96}
    Let $A\subset B$ be an integral extension.
    \begin{itemize}
        \item[(a)] If $\mathfrak{b}\triangleleft B$, then $\frac{A}{\mathfrak{b}\cap A}\subset \frac{B}{\mathfrak{b}}$ is an integral extension.
        \item[(b)] If $B$ is an integral domain and $S$ is a multiplicatively closed subset of $A$, then $S^{-1}A\subset S^{-1}B$ is an integral ring extension.
        \item[(c)] $A[x]\subset B[x]$ is an integral ring extension. [AM, 5.6]
    \end{itemize}
\end{lemma}
\begin{proof}
    (a) Consider the map $\phi:A\xhookrightarrow{} B\twoheadrightarrow B/\mathfrak{b}$, then $\mathrm{ker}\phi=A\cap\mathfrak{b}\coloneqq \mathfrak{a}.\ \Rightarrow\overline{\phi}:A/\mathfrak{a} \xhookrightarrow{} B/\mathfrak{b}$. Let $\overline{x}\coloneqq x+\mathfrak{b}\in B/\mathfrak{b}, (x\in B).\ \because B$ is integral over $A.\ \therefore x$ satisfies 
    \begin{equation*}
        x^n+a_{n-1}x^{n-1}+\cdots+a_1x+a_0=0\ \text{in}\ B,\qquad (a_i\in A).
    \end{equation*}
    $\Rightarrow \overline{x}^n+\overline{a_{n-1}}\overline{x}^{n-1}+\cdots+ \overline{a_0}=\overline{0}$ in $B/\mathfrak{b}$. Now, we have $A/\mathfrak{a} \xhookrightarrow{\overline{\phi}} B/\mathfrak{b},\ a_i+\mathfrak{a}\mapsto a_i+\mathfrak{b}=\overline{a_i}$,\\ may regard $\overline{a_i}\in A/\mathfrak{a}.\ \Rightarrow \overline{x}\in B/\mathfrak{b}$ is integral over $A/\mathfrak{a}$, and hence $B/\mathfrak{b}$ is integral over $A/\mathfrak{a}$. \\ 
    (b) Let $\frac{x}{s}\in S^{-1}B,\ (x\in B,s\in S).\ \because B$ is integral over $A.\ \therefore x$ satisfies 
    \begin{equation*}
        x^n+a_{n-1}x^{n-1}+\cdots+a_1x+a_0=0\ \text{in}\ B,\qquad (a_i\in A).
    \end{equation*}
    Dividing $s^n$, we have 
    \begin{equation*}
        (\tfrac{x}{s})^n+\tfrac{a_{n-1}}{s}\cdot(\tfrac{x}{s})^{n-1}+\cdots +\tfrac{a_1}{s^{n-1}}\cdot\tfrac{x}{s}+\tfrac{a_0}{s^n}=0=\tfrac{0}{s^n}\ \text{in}\ S^{-1}B,\qquad (\tfrac{a_i}{s^{n-i}}\in S^{-1}A).
    \end{equation*}
    $\Rightarrow \frac{x}{s}\in S^{-1}B$ is integral over $S^{-1}A$, and hence $S^{-1}B$ is integral over $S^{-1}A$.\\
    (c) Let $f\coloneqq b_nx^n+\cdots+b_1x+b_0\in B[x],\ (b_i\in B).\ \because b_i$ is integral over $A.\ \therefore b_i$ is also integral over $A[x]$.  Moreover, $\because f\in A[b_0,\cdots,b_n][x]\cong A[x][b_0,\cdots,b_n].\ \therefore$ By \text{Prop. 84.}, $A[x][b_0,\cdots,b_n]$ is a finite $A[x]$-module. $\Rightarrow f$ is integral over $A[x]$.  
\end{proof}

\begin{proposition}{97. (Integral Closure)}
    The set $\overline{A}\coloneqq\{ b\in B:\ b\ \text{integral over}\ A\} \subset B$ is a subring of $B$. In particular, the sum and the product of two integral elements is again integral, and $\overline{\overline{A}}=\overline{A}$. We call $\overline{A}$ the \textbf{integral closure} of $A$ in $B$.\\
    Moreover, if $\overline{A}=A$, then $A$ is called \textbf{integrally closed} in $B$. [AM, 5.3, 5.5]
\end{proposition}
\begin{proof}
    (i) If $x,y\in\overline{A}$ , then $A[x,y]$ is finite over $A$, whence $x+y$ and $x\cdot y$ are integral over $A$ and thus in $\overline{A}$. \\
    (ii) Since $A\subset\overline{A}$ and $\overline{A}\subset \overline{\overline{A}}$ are integral ring extensions, it follows from \textbf{Prop. 95} that $\overline{A}\subset \overline{\overline{A}}\subset \overline{A}$, i.e. $\overline{\overline{A}}=\overline{A}$.
\end{proof}

\begin{example}{98}
    If $\overline{A}$ is the integral closure of $A$ in $B$, then $S^{-1}\overline{A}$ is the integral closure of $S^{-1}A$ in $S^{-1}B$. [Exercise][AM, 5.12]
\end{example}

\newpage

\section{Ring of Integers}

Now we generalise the ring extension $\mathbb{Z}[i]\subset\mathbb{Q}(i)$ as follows:

\begin{definition}{99. (Ring of Integers)}
    An \textbf{algebraic number field} $L$ is a finite extension field of $\mathbb{Q}$, i.e. $[L:\mathbb{Q}]<\infty$. In particular, $\mathrm{char}L=0$ so that $\mathbb{Z}\subset L$. The integral closure of $\mathbb{Z}$ in $L$ defines the \textbf{ring of integers} $\mathcal{O}_L$.
\end{definition}

\begin{lemma}{100}
    Let $K\coloneqq\mathrm{Quot}(A)$ be the field of fractions of an integral domain $A$, and let $L/K$ be a field extension. If $x\in L$ is algebraic over $K$, then there exists $d\in A$ such that $d\cdot x$ is integral over $A$.
\end{lemma}
\begin{proof}
    $\because x\in L$ is algebraic over $K.\ \therefore x$ satisfies $x^n+ c_{n-1}x^{n-1}+\cdots+c_1x+c_0=0,\ (c_i\in K)$. Let $d$ be the common denominator for the $c_i$'s so that $d\cdot c_i\in A,\ \forall i$. Now, multiply the equation by $d^n$, we have\\  \centerline{$d^nx^n+ c_{n-1}d^nx^{n-1}+ \cdots +c_1d^nx+c_0d^n=0.$}\\ Rewrite as \\ \centerline{$(d\cdot x)^n+(c_{n-1}d)(d\cdot x)^{n-1}+ \cdots+(c_1d^{n-1})(d\cdot x)+c_0d^n=0$.} \\ $\because c_i\cdot d^{n-i}\in A,\ \forall i.\ \therefore d\cdot x$ is integral over $A$.
\end{proof}

\begin{corollary}{101}
    Let $A$ be an integral domain, $K\coloneqq\mathrm{Quot}(A)$ and $L/K$ an algebraic field extension. If $\overline{A}$ is the integral closure of $A$ in $L$, then $L=\mathrm{Quot}(\overline{A})$. In particular, $L=\mathrm{Quot}(\mathcal{O}_L)$ for any algebraic number field $L$.
\end{corollary}
\begin{proof}
    If $x\in L.\ \because L/K$ is algebraic. $\therefore$ By \textbf{Prop. 100}, $\exists\ d\in A$ such that $b\coloneqq d\cdot x\in A\subset\overline{A}.\\  \Rightarrow x=\frac{b}{d}\in\mathrm{Quot}(\overline{A})$.  
\end{proof}

\begin{definition}{102. (Integrally Closed)}
    An integral domain $A$ is called \textbf{integrally closed} or \textbf{normal} if $A$ is integrally closed in its quotient field.
\end{definition}

\begin{example}{103}
    \begin{itemize}
        \item Every UFD is normal by \textbf{Example. 92}.
        \item The ring of integers $\mathcal{O}_K$ of an algebraic number field is integrally closed.\\ (cf. \textbf{Prop. 95}).
        \item $A\coloneqq k[x,y]/(y^2-x^3)$ is an integral domain with $\mathrm{Quot}(A)\cong k(t)$. \\
        Indeed, the map $k(t)\to\mathrm{Quot}(A),\ (\frac{f}{g})(t)\mapsto (\frac{f}{g})(\frac{\overline{y}}{\overline{x}})$ is an isomorphism. Since $k[t]$ is a UFD, it is integrally closed. However, $A$ is not. ($\tau=\frac{\overline{y}}{\overline{x}}\notin A$, but $\tau^2-\overline{x}=0$.)
    \end{itemize}
\end{example}

\newpage
\begin{proposition}{104}
    Let $A$ be an integral domain. Then T.F.A.E.:
    \begin{itemize}
        \item[(a)] $A$ is normal.
        \item[(b)] $A_\mathfrak{p}$ is normal for each prime ideal $\mathfrak{p} \triangleleft A$.
        \item[(c)] $A_\mathfrak{m}$ is normal for each maximum ideal $\mathfrak{m}\triangleleft A$.
    \end{itemize}
    In view of the second item, we also say that normality is a \textbf{local property}. [Exercise][AM, 5.13]
\end{proposition}

\begin{proposition}{105}
    Let $A$ be integrally closed, $K\coloneqq\mathrm{Quot}(A)$ and $L/K$ a finite field extension. Then an element of $L$ is integral over $A$ iff its minimum polynomial over $K$ has coefficients in $A$.
\end{proposition}
\begin{proof}
    "$\Leftarrow$"\ Obvious. "$\Rightarrow$"\ Suppose $u\in L$ is integral over $A$, then $u$ satisfies \\  \centerline{$u^n+a_{n-1}u^{n-1}+\cdots+a_1u+a_0=0,\ (a_i\in A)$.}\\ Let $f$ be the minimum polynomial of $u$ over $K$. For any root $v$ of $f(x)=0$, the fields $K[u]$ and $K[v]$ are both stem fields for $f$, and so there exists a $K$-isomorphism $\sigma:K[u]\to K[v],\ \sigma u\mapsto v$; On applying $\sigma$ to the above equation, we obtain the equation \\ \centerline{$v^n+a_{n-1}v^{n-1}+\cdots+a_1v+a_0=0$,} \\ which shows that $v$ is integral over $A$. Hence all the roots of $f$ are integral over $A$, and it follows that the coefficients of $f$ are all integral over $A$. Thus, they lie in $K$, but $A$ is integrally closed, so they lie in $A$.
\end{proof}

\begin{example}{106}
    Consider the simple field extension $\mathbb{Q}\subset\mathbb{Q}(\tau), \ \tau=\frac{1+\sqrt{5}}{2}$. The associated $\mathbb{Q}$-linear map $\mathrm{T}_\tau(x)=\tau\cdot x$ (cf. \textbf{Def. 78}) maps $1\mapsto \tau$ and $\tau\mapsto1+\tau$. By \textbf{Rmk. 80}, the coefficients of the minimal polynomial of $\tau$ are given by the trace and the norm of $\tau$ which is $1$ and $-1$, respextively. Hence $\mathbb{Z}\subset\mathbb{Z}[\tau]$ is integral. On the other hand, $\tau'\coloneqq\frac{1+\sqrt{3}}{2}$ is not integral (its norm is $-\frac{1}{2}$).
\end{example}
\newpage
\section{Discriminants}
\begin{definition}{107. (Discriminants)}
    Let $V$ be a finite dimensional $K$-vector space and $\psi:V\times V\to K$ a bilinear symmetric form. Fixing a basis $\{e_1,\cdots,e_n\}$ represents $\psi$ as the matrix $[\psi(e_i,e_j)]$. The \textbf{discriminant} of this basis is then defined by 
    \begin{align*}
        D(e_1,\cdots,e_n)\coloneqq\mathrm{det}[\psi(e_i,e_j)].
    \end{align*}
    If $f_i\coloneqq\mathsf{T}e_i$ obtained by the transformation $\mathsf{T}:V\to V,\ \forall1\leq i\leq n$, then 
    \begin{align*}
        D(f_1,\cdots,f_n)=(\mathrm{det}\mathsf{T})^2\cdot D(e_1,\cdots,e_n).
    \end{align*}
    For instance, let $L/K$ be finite. We consider the bilinear symmectric form
    \begin{align*}
        \psi:L\times L\to K,\ (\alpha,\beta)\mapsto \mathrm{Tr}_{L/K}(\alpha\beta).
    \end{align*}
    More generally, this definition applies to a finite ring extension $A\subset B$ with $B\simeq A^m$.\\ The \textbf{discriminant of elements} $\beta_1,\cdots,\beta_m\in B$ is given by 
    \begin{align*}
        D(\beta_1,\cdots,\beta_m)\coloneqq \mathrm{det(Tr}_{B/A}(\beta_i\beta_j)).
    \end{align*}
    Since two bases are related by a transformation $\mathsf{T}$ with $\mathrm{det}\mathsf{T}\in A^\ast$, this defines a nontrivial ideal $\mathrm{disc}(B/A)$ in $A$, called the \textbf{discriminant} of $B$ over $A$. In particular, $\mathrm{disc}(B/\mathbb{Z})\in\mathbb{Z}$ is a well-defined integer, (since $1$ is the only square of a unit in $\mathbb{Z}$).
\end{definition}

\begin{proposition}{108}
    Let $A\subset B$ be integral domains such that $B\simeq A^m$ and $\mathrm{disc}(B/A)\neq0$.\\ Then elements $b_1,\cdots,b_m\in B$ form a basis as an $A$-module iff as ideals in $A$, we have 
    \begin{align*}
        D(b_1,\cdots,b_m)= \mathrm{disc}(B/A).
    \end{align*}
\end{proposition}
\begin{proof}
    Let $\{e_1,\cdots,e_m\}$ be an $A$-basis for $B$. For each $i$, we write $b_i\coloneqq\sum a_{ij}e_j,\ (a_{ij}\in A)$. Then 
    \begin{align*}
        D(b_1,\cdots,b_m)= \mathrm{det}(a_{ij})^2\cdot D(e_1,\cdots,e_m).
    \end{align*}
    Note that $\mathrm{det}(a_{ij})$ is a unit $\Leftrightarrow \{b_1,\cdots,b_n\}$ is an $A$-basis for $B$.
\end{proof}

\begin{remark}{109}
    If $A=\mathbb{Z}$ and $B\simeq\mathbb{Z}^m$, then $b_1,\cdots,b_m\in B$ generates a submodule $N$ of finite index $(B:N)$ in $B$ iff $D(b_1,\cdots,b_m)\neq0$. In this case
    \begin{align*}
        D(b_1,\cdots,b_m)\ =\ (B:N)^2\cdot\mathrm{disc}(B/ \mathbb{Z})
    \end{align*}
\end{remark}
\begin{proof}
    It is well-known from the theory of finite $\mathbb{Z}$-modules that $\mathrm{det}(a_{ij})=(B:N)$, where $b_i=\sum a_{ij}e_j$ and $\mathrm{disc}(B/\mathbb{Z})=\mathrm{det}(\mathrm{Tr}_{B/\mathbb{Z}}(e_ie_j))$. (cf. proof of \textbf{Prop. 108}.)
\end{proof}
\newpage
\begin{proposition}{110}
    Let $L/K$ be a finite separable field extension of degree $[L:K]=m$, and $\sigma_1,\cdots,\sigma_m\in\mathrm{Hom}_K(L,\overline{K})$ be the distinct $K$-homomorphisms of $L$ into the algebraic closure of $K$. Then, for any $K$-basis $\{\beta_1,\cdots,\beta_m\}$ of $L$, 
    \begin{align*}
        D(\beta_1,\cdots,\beta_m) =\mathrm{det}(\sigma_i(\beta_j))^2\neq0.
    \end{align*}
    In particular, $\psi:L\times L\to K$ from \textbf{Def. 107} is non-degenerate.
\end{proposition}
\begin{proof}
    By direct calculation, we have 
    \begin{align*}
        D(\beta_1,\cdots,\beta_m) &\myeq{def.} \mathrm{det} (\mathrm{Tr}_{B/A}(\beta_i\beta_j)) 
        =\mathrm{det}(\sum_k\sigma_k(\beta_i\beta_j))
        =\mathrm{det}(\sum_k\sigma_k(\beta_i)\sigma_k(\beta_j)) \\
        & =\mathrm{det}(\sigma_k(\beta_i))\cdot\mathrm{det}(\sigma_k(\beta_j)) =\mathrm{det}(\sigma_k(\beta_i))^2.
    \end{align*}
    Suppose $\mathrm{det}(\sigma_i(\beta_j))=0$. Then there exist $c_1,\cdots,c_m\in\overline{K}$ such that $\sum_ic_i\sigma_i(\beta_j) =0,\ \forall j$.
\end{proof}

\begin{corollary}{111}
    Let $A$ be an integral domain, $K\coloneqq\mathrm{Quot}(A)$ and $L/K$ a finite separable extension of degree $m$. If the integral closure $B$ of $A$ in $L$ satisfies $B\simeq A^m$, then $\mathrm{disc}(B/A)\neq0$.
\end{corollary}
\begin{proof}
    Let $\{b_1,\cdots,b_n\}$ be an $A$-basis for $B$, then it generates $L$ as a $K$-vector space.\\ Indeed, if $x\in L$ is algebraic over $K$, by \textbf{Lemma. 100}, there exists $d\in A$ such that $d\cdot x\in B$,\\ i.e. $d\cdot x\coloneqq\sum a_ib_i,\ (a_i\in A).\ \Rightarrow x=\sum \frac{a_i}{d}\cdot b_i\in K-\mathrm{span}\{b_1,\cdots,b_n\}$.
\end{proof}

\begin{proposition}{112}
    Let $A$ be integrally closed, $K\coloneqq\mathrm{Quot}(A)$, $L/K$ a finite separable extension of degree $m$ and $B\coloneqq\overline{A}$ the integral closure of $A$ in $L$. Then there exist $A$-submodules $M$ and $M'$ of $L$ such that $M,M'\simeq A^m$ and $M\subset B\subset M'$.\\ In particular, $B$ is a finitely generated $A$-module if $A$ is Noetherian, and it is free of rank $m$ if $A$ is a PID.
\end{proposition}
\begin{proof}
    Let $\{\beta_1,\cdots,\beta_m\}$ be a $K$-basis for $L$. By \textbf{Lemma. 100}, there exists $d\in A$ such that $d\cdot\beta_i\in B,\ \forall i$. Now, we may assume that $\beta_i\in B,\ \forall i.\ \because \psi:L\times L\to K$ is non-degenerate by \textbf{Prop. 110}.  $\therefore\exists\ \beta'_1,\cdots\beta'_m$ such that $\psi(\beta_i \cdot\beta'_j)=\delta_{ij}.\ \Rightarrow \{\beta'_1,\cdots,\beta'_m\}$ is also a $K$-basis for $L$. 
    \textsc{claim.} $A\beta_1+\cdots+A\beta_m\subset B\subset A\beta'_1+\cdots+ A\beta'_m$. \\ The first inclusion is obvious; for the second, let $\beta\in B$, we have $\beta\coloneqq\sum b_j\beta'_j,\ (b_j\in K)$. Now, $\because \beta, \beta_i\in B.\ \therefore \beta_i\cdot\beta\in B.\ \Rightarrow \mathrm{Tr}_{L/K}(\beta_i\cdot\beta)=\sum_{k=1}^m\sigma_k(\beta_i\cdot\beta)\in A$. Moreover, $A$ is integrally closed, it follows that  $\mathrm{Tr}_{L/K}(\beta_i\cdot\beta)=\sum_jb_j\mathrm{Tr}(\beta_i\cdot\beta'j)=b_i\in A$.\\
    If $A$ is Noetherian, then $M'$ is a Noetherian $A$-module, and so $B$ is finitely generated as an $A$-module. If $A$ is a PID, then $B$ is free of rank $\leq m$ since it is contained in a free $A$-module of rank $m$, and it has rank $\geq m$ because it contains a free $A$-module of rank $m$.
\end{proof}

\begin{remark}{113. (Integral Basis)}
    The \textbf{Prop. 112} applies in particular to the ring of integers $\mathcal{O}_L$ of a number field $L$. A $\mathbb{Z}$-basis of $\mathcal{O}_L$ and thus a $K$-basis of $L$ is called an \textbf{integral basis} of $L$.
\end{remark}

\begin{corollary}{114}
    The ring of integers $\mathcal{O}_L$ in a number field $L$ is the largest subring of $L$ which is finitely generated as a $\mathbb{Z}$-module.
\end{corollary}
\begin{proof}
    We have just seen that $\mathcal{O}_L$ is a finitely generated $\mathbb{Z}$-module. Let $B$ be another subring of $L$ that is finitely generated as a $\mathbb{Z}$-module; then every element of $B$ is integral over $\mathbb{Z}$, and so $B\subset\mathcal{O}_L$.
\end{proof}

\section{Exercises}
\begin{exe}{6. (Integral Closure of Localisations)}
    Let $A\subset B$ be a ring extension and $S$ be a multiplicative subset of $A$. Then $S^{-1}\overline{A}$ is the integral closure of $S^{-1}A$ in $S^{-1}B$.
\end{exe}
\begin{proof}
    We claim that $\overline{S^{-1}A}=S^{-1}\overline{A}$. "$\supset$"\ is clear by \textbf{96(b)} that $S^{-1}\overline{A}$ is integral over $S^{-1}A$.\\
    "$\subset$"\ Suppose $\frac{b}{s}\in S^{-1}B,\ (b\in B,\ s\in S)$ is integral over $S^{-1}A$, then $\frac{b}{s}$ satisfies
    \begin{align*}
        (\tfrac{b}{s})^n+(\tfrac{a_1}{s_1})(\tfrac{b}{s})^{n-1}+\cdots+\tfrac{a_n}{s_n}=0\ \text{in}\ S^{-1}B,\qquad (a_i\in A,\ s_i\in S).
    \end{align*}
    Take $t\coloneqq s_1\cdots s_n$ and multiply the above equation by $(st)^n$, then
    \begin{align*}
        (bt)^n+c_1(bt)^{n-1}+\cdots+c_n=0\ \text{in}\ B,\qquad (c_i\coloneqq\tfrac{a_i}{s_i}(st)^i\in A).
    \end{align*}
    Thus, $bt\in B$ is integral over $A$, i.e. $bt\in\overline{A}$. It follows that $\frac{b}{s}=\frac{bt}{st}\in S^{-1}\overline{A}$.
\end{proof}

\begin{exe}{7. (Characterisation of Integral Elements inside Fields)}
    Let $A\subset K$ be a ring extension with $K$ a field. Then $\alpha\in K$ is integral over $A$ iff there exists a finitely generated $A$-module $M\neq0$ inside $K$ with $\alpha M\subset M$.
\end{exe}

\begin{exe}{8. (Normality is a Local Property)}
    T.F.A.E. for an integral domain $A$:
    \begin{itemize}
        \item[(i)] $A$ is normal.
        \item[(ii)] $A_\mathfrak{p}$ is normal for each prime ideal $\mathfrak{p}\triangleleft A$.
        \item[(iii)] $A_\mathfrak{m}$ is normal for each maximal ideal $\mathfrak{m}\triangleleft A$.
    \end{itemize}
\end{exe}
\begin{proof}
    We use the fact in [AM, 3.9] that an injective (so is a surjective) $A$-linear map is a local property. Let $K=\mathrm{Quot}A$ and $f:A\xhookrightarrow{} \overline{A}$ be the identity map. Then, by \textbf{Ex. 6}, we have 
    \begin{align*}
        A\ \text{is normal}\ &\Leftrightarrow\ f\ \text{is surjective}\ \Leftrightarrow\ f_\mathfrak{p}\ \text{(resp. $f_\mathfrak{m}$) is surjective}\\ 
        &\Leftrightarrow\ A_\mathfrak{p}\ \text{(resp. $A_\mathfrak{m}$) is normal},
    \end{align*}
    for each prime ideal $\mathfrak{p}$ (resp. maximal ideal $\mathfrak{m}$) of $A$.
\end{proof}

\begin{exe}{9. (Normal Rings in Number Theory)}
    Let $d\neq0,1$ be a square-free integer, (i.e. no square divides $d$ in $\mathbb{Z}$). Then the integral closure of $\mathbb{Z}$ in $\mathbb{Q}(\sqrt{d})=\{a+b\sqrt{d}:a,b\in\mathbb{Q}\}\subset\mathbb{C}$ is giver by
    \begin{align*}
        \overline{\mathbb{Z}}=\mathcal{O}_{\mathbb{Q}(\sqrt{d})}= \{a+b\sqrt{d}:a,b\in\mathbb{Q},\ -2a\in\mathbb{Z},\ a^2-db^2\in\mathbb{Z}\}.
    \end{align*}
\end{exe}

\chapter{Dedekind Domains}

\section{Discrete Valuation and Dedekind Rings}

\begin{motivation}{115}
    Our next goal is to understand the ring of integers $\mathcal{O}_K$ of an algebraic number field as a ring. In particular, we want to understand which of the good properties of $\mathbb{Z}$ survive. This is a nontrivial question as the subsequent proposition and remark show.
\end{motivation}

\begin{proposition}{116}
    Every nonzero nonunit element of a Noetherian integral domain can be written as a product of irreducible elements (but not necessarily a unique way as for a UFD).
\end{proposition}
\begin{proof}
    Recall that if $A$ is an integral domain, then $(a)\subset(b)\ \Leftrightarrow\ b|a$. And the equality holds iff $a\coloneqq b\cdot u$, for some $u\in A^\times$. We assume to the contrary that there exists $a\in A$ which cannot be written as a product so that the set of ideals\\
    \centerline{$\Sigma\coloneqq\{(a):a\in A\ \text{cannot be written as a prodct}\}\neq \varnothing$.} \\ $\because A$ is Noetherian. $\therefore\exists$ maximal element $(a)\in\Sigma$. In particular, $a$ cannot be irreducible and so $a=b\cdot c$ for some nonunits $b,c\in A$. Hence $(a)\subsetneq(b)$, so that by maximality of $(a)$, $b$ admits a decomposition into irreducibles. Interchanging the roles of $b$ and $c$ yields a decomposition of $c$. Thus $a$ is a product of irreducible elements, a contradiction. 
\end{proof}

\begin{remark}{117}
    The \textbf{Prop. 116} fails for instance for the ring $\mathcal{O}_{\overline{\mathbb{Q}}}$ of algebraic integers in the algebraic closure of $\mathbb{Q}$ in $\mathbb{C}$. Indeed, an algebraic integer $\alpha\in \mathcal{O}_{\overline{\mathbb{Q}}}$ cannot be irreducible for $\sqrt{\alpha}$ is also in $\mathcal{O}_{\overline{\mathbb{Q}}}$ and $\alpha=\sqrt{\alpha}\cdot \sqrt{\alpha}$. In particular, $\mathcal{O}_{\overline{\mathbb{Q}}}$ is neither a UFD nor Noetherian.
\end{remark}

However, we will see that $\mathcal{O}_K$ is a Dedekind ring and as such it still retains certain factorisation properties.

\begin{proposition}{118. (Discrete Valuation Ring)}
    Let $A$ be a PID. Then T.F.A.E.:
    \begin{itemize}
        \item[(a)] $A$ has exactly one nonzero prime ideal.
        \item[(b)] Up to multiplying with units, $A$ has exactly one prime element.
        \item[(c)] $A$ is local and not a field.
    \end{itemize}
    A PID satisfying one and thus all of these conditions is called a \textbf{discrete valuation ring (DVR)}.
\end{proposition}

\begin{example}{119}
    The localisation $\mathbb{Z}_{(p)}\coloneqq\{\frac{m}{n}\in\mathbb{Q}:p\nmid n\}$ is a DVR with $(p)^e$ as its unique nonzero prime ideal. Up to multiplying with units, $p$ is the unique prime element. Any other element in $\mathbb{Z}_{(p)}$ is of the form $p^k\cdot\frac{m}{n}$ with $p\nmid m,n$. The group of units is given by $\mathbb{Z}^\times_{(p)}=\{\frac{m}{n}\in\mathbb{Q}:p\nmid m,n\}$, that is, the element of $\mathbb{Z}_{(p)}$ with vanishing exponent $k$.
\end{example}

\begin{proposition}{120}
    An integral domain $A$ is a DVR iff $A$ is \\ \centerline{(i) Noetherian, (ii) integrally closed, and (iii) has exactly one nonzero prime ideal.}
\end{proposition}
\begin{proof}
    "$\Rightarrow$"\ is obvious. "$\Leftarrow$"\ Suppose $A$ is an integral domain satisfying (i)-(iii). We must show that every ideal in $A$ is principal. Note: (iii) implies that $A$ is a local ring.\\
    \textbf{Step 1.} The nonzero prime ideal is principal.\\ \textsc{proof.}
    Choose a nonunit element $0\neq c\in A$ and consider the $A$-module $M\coloneqq A/(c)$. For each $0\neq m\in M$, note that 
    \begin{align*}
        \mathrm{Ann}(m)\coloneqq\{a\in A:am=0\}\triangleleftneq A.
    \end{align*}
    $\because A$ is Noetherian. $\therefore$ We may choose an $m\in M$ so that $\mathrm{Ann}(m)$ is maximal among these ideals. Write $m\coloneqq b+(c)$ and $\mathfrak{p}\coloneqq\mathrm{Ann} (b+(c))$. Note that $c\in\mathfrak{p}$, thus
    \begin{align*}
        \mathfrak{p}=\{a\in A:\ c|ab\ \}\neq0.
    \end{align*}
    \textbf{Claim 1.} $\mathfrak{p}$ is prime. If not, there exist elements $x,y\in A\setminus\mathfrak{p}$ such that $xy\in \mathfrak{p}$. Then $yb+(c)\neq0+(c)$ in $M$ since $y\notin\mathfrak{p}$. Consider $\mathrm{Ann}(yb+(c))$, which obviously contains $\mathfrak{p}$ and $x$, but this contradicts the maximality of $\mathfrak{p}$ among ideals of the form $\mathrm{Ann}(m)$. Hence $\mathfrak{p}$ is prime, thus maximal by (iii).\\
    \textbf{Claim 2.} $\frac{b}{c}\notin A$. Otherwise $b=c\cdot\frac{b}{c}\in(c)$, and $m=0$ in $M$.\\
    \textbf{Claim 3.} $\frac{c}{b}\in A$ and $\mathfrak{p}=(\frac{c}{b})$. By definition, $\mathfrak{p}b\subset(c)$, and so $\mathfrak{p}\cdot \frac{b}{c}\subset A$, thus an ideal in $A$.\\ If $\mathfrak{p} \cdot\frac{b}{c}\subset\mathfrak{p}$, then $\frac{b}{c}$ is integral over $A$ by \textbf{Prop. 93}. Thus, by (ii), $\frac{b}{c}\in A$, which contradicts to the previous \textsc{Claim}. Thus, by (iii), $\mathfrak{p}\cdot\frac{b}{c}=A.\ \because A$ is local.  $\therefore A\setminus\mathfrak{p}=A^\times.\ \Rightarrow\mathfrak{p}= (\frac{c}{b})$.\\
    \textbf{Step 2.} Every proper ideal $\mathfrak{a}\triangleleftneq A$ is principal.\\ \textsc{proof.}
    Put $\pi\coloneqq\frac{c}{b}$ so that $\mathfrak{p}=(\pi)$. Consider the sequence
    \begin{align*}
        \mathfrak{a}\subset\mathfrak{a}\pi^{-1} \subset\mathfrak{a}\pi^{-2}\subset\cdots.
    \end{align*}
    If $\mathfrak{a}\pi^{-r}=\mathfrak{a}\pi^{-r-1}$ for some $r$, then $\pi^{-1}(\mathfrak{a}\pi^{-r})=\mathfrak{a}\pi^{-r}$, it follows that $\pi^{-1}$ is integral over $A$ by \textbf{Prop. 93}. Again by (ii), $\pi^{-1}\in A$, but this contradicts to \textbf{Claim. 2}. Therefore the sequence is strictly increasing, and it can't be contained in $A$, since $A$ is Noetherian.\\
    Let $m\in\mathbb{N}$ be the smallest integer such that $\mathfrak{a}\pi^{-m}\subset A$ but $\mathfrak{a}\pi^{-m-1} \nsubseteq A$. Then  $\mathfrak{a}\pi^{-m}\nsubseteq\mathfrak{p}$, (otherwise  $\mathfrak{a}\pi^{-m-1}\subset A$). $\because A$ is local. $\therefore\mathfrak{a}\pi^{-m}=A.\ \Rightarrow \mathfrak{a} =(\pi^m)$ is principal.
\end{proof}

\begin{definition}{121. (Dedekind Domain)}
    A \textbf{Dedekind domain} is an integral domain $A$, not equal to a field such that $A$ is \\ \centerline{(i) Noetherian, (ii) integrally closed, and (iii)  every nonzero prime ideal is maximal.}\\
    It follows that a local integral domain is Dedekind iff it is a DVR.
\end{definition}

\begin{example}{122}
    The algebraic integers $\mathcal{O}_K$ of a number field form a Dedekind domain. (cf. \textbf{Cor. 135.})
\end{example}

\begin{proposition}{123}
    Let $A$ be a Dedekind domain and $S$ be a multiplicative subset of $A$. Then $S^{-1}A$ is either a Dedekind domain or a field.
\end{proposition}
\begin{proof}
    Note: the prime ideals of $S^{-1}A$ correspond to those of $A$ which have empty intersection with $S$. Hence, localisation preserves condition (iii). It also preserves conditions (i) and (ii) as noted in \textbf{Chapter 1}.
\end{proof}

\begin{proposition}{124}
    For any nonzero prime ideal $\mathfrak{p}\neq0$ in a Dedekind domain $A$, the localisation $A_\mathfrak{p}$ is a DVR.
\end{proposition}
\begin{proof}
    Note: $A_\mathfrak{p}$ is local. Moreover, $A_\mathfrak{p}$ is integrally closed and Noetherian by \textbf{Prop. 114} and [AM, 7.4], respectively.
\end{proof}

\begin{theorem}{125}
    Let $A$ be a Dedekind domain. Every proper nonzero ideal $\mathfrak{a}
    \triangleleft A$ can be written in the form 
    \begin{align*}
        \mathfrak{a}\coloneqq 
        \mathfrak{p}_1^{r_1}\cdots\mathfrak{p}_n^{r_n}
    \end{align*}
    with uniquely determined distinct prime ideals $\mathfrak{p}_i\triangleleft A$ and $r_i>0$. Furthermore, \\ \centerline{$r_i>0\qquad\Leftrightarrow\qquad \mathfrak{a}A_{\mathfrak{p}_i}\subsetneq A_{\mathfrak{p}_i}\qquad\Leftrightarrow \qquad \mathfrak{a}\subset\mathfrak{p}_i$.}
\end{theorem}

The proof will require the following lemmas:

\begin{lemma}{126}
    Let $A$ be a Noetherian ring. Then every ideal $\mathfrak{a}\triangleleft A$ contains a product of nonzero prime ideals.
\end{lemma}
\begin{proof}
    Suppose not, and choose a maximal counterexample $\mathfrak{a} \triangleleft A$, (since $A$ is Noetherian). Then $\mathfrak{a}$ itself cannot be prime, and so there exist $x,y\in A\setminus \mathfrak{a}$ such that $xy\in\mathfrak{a}$. Now, the ideals $\mathfrak{a}+(x)$ and $\mathfrak{a}+(y)$ strictly contain $\mathfrak{a}$, but their product is contained in $\mathfrak{a}$. By maximality of $\mathfrak{a}$, each of $\mathfrak{a}+(x)$ and $\mathfrak{a}+(y)$ contains a product of prime ideals, thus, so does $\mathfrak{a}.\ (\to\gets)$
\end{proof}

\begin{lemma}{127}
    Let $A$ be a ring, and $\mathfrak{a},\ \mathfrak{b}$ be relatively prime ideals in $A$. Then for any $m,m\in\mathbb{N}$, the ideals $\mathfrak{a}^m$ and $\mathfrak{b}^n$ are relatively prime.
\end{lemma}
\begin{proof}
    If $\mathfrak{a}^m$ and $\mathfrak{b}^n$ are not relatively prime, then they are both contained in some prime (even maximal) ideal $\mathfrak{p}$, i.e. $\mathfrak{a}^m+\mathfrak{b}^n\subset
    \mathfrak{p}$. Now, both $\mathfrak{a}^m,\ \mathfrak{b}^n\subset
    \mathfrak{p}$, it follows that both $\mathfrak{a}, \mathfrak{b} \subset\mathfrak{p}$. But this contradicts that $\mathfrak{a}$ and $\mathfrak{b}$ are relatively prime, (i.e. $\mathfrak{a}+ \mathfrak{b} =A$).
\end{proof}
\newpage
\begin{lemma}{128}
    Let $\mathfrak{p}$ be a maximal ideal of a ring $A$, and $\mathfrak{q}\coloneqq \mathfrak{p}A_\mathfrak{p}$ be the maximal ideal of the localisation $A_\mathfrak{p}$. Then the map 
    \begin{align*}
        A/\mathfrak{p}^m\longrightarrow A_\mathfrak{p}/\mathfrak{q}^m,\ a+\mathfrak{p}^m\longmapsto a+\mathfrak{q}^m
    \end{align*}
    is an isomorphism.
\end{lemma}
\begin{proof}
    \textsc{\textbf{Claim 1.}} The map is 1-1. It suffices to show that $\mathfrak{q}^m\cap A=\mathfrak{p}^m$. \\ \textsc{proof.}
    $\because\mathfrak{q}^m=S^{-1}\mathfrak{p}^m,\ \text{where}\ S\coloneqq A\setminus\mathfrak{p}.\ \therefore$ We have to show that $\mathfrak{p}^m=(S^{-1}\mathfrak{p}^m)\cap A$. \\ "$\subset$"\ is clear since $a=\frac{a}{1}\in(S^{-1}\mathfrak{p}^m)\cap A$ for any $a\in\mathfrak{p}^m$. "$\supset$"\ Let $a\coloneqq\frac{b}{s}\in(S^{-1}\mathfrak{p}^m)\cap A,\ \text{where}\\ b\in \mathfrak{p}^m,\ s\in S,\ a\in A.\ \because sa=b\in\mathfrak{p}^m.\ \therefore \overline{sa}=\overline{0}$ in $A/\mathfrak{p}^m$. Note: The only maximal ideal containing $\mathfrak{p}^m$ is $\mathfrak{p}$ (since $\mathfrak{m}\supset\mathfrak{p}^m\ \text{implies}\ \mathfrak{m}\supset \mathfrak{p}$), so the only maximal ideal in $A/\mathfrak{p}^m$ is $\mathfrak{p}/\mathfrak{p}^m$; in fact, $A/\mathfrak{p}^m$ is a local ring. $\because s\in S=A\setminus\mathfrak{p}.\ \therefore \overline{s}=s+\mathfrak{p}^m\notin\mathfrak{p}/\mathfrak{p}^m.\ \Rightarrow \overline{s}\in(A/\mathfrak{p}^m)^\times.\\ \Rightarrow \overline{a}=\overline{s}^{-1}\cdot\overline{sa}=\overline{s}^{-1}\cdot \overline{0}=\overline{0}$ in $A/\mathfrak{p}^m.\ \Rightarrow a\in\mathfrak{p}^m$.\\
    \textsc{\textbf{Claim 2.}} The map is surjective. \\ \textsc{proof.}
    Let $\overline{\frac{a}{s}}\in A_\mathfrak{p}/\mathfrak{q}^m$, where $a\in A,\ s\in S.\ \because s\notin\mathfrak{p}$ and $\mathfrak{p}$ is maximal. $\therefore (s)+\mathfrak{p}=A$, i.e. $(s)$ and $\mathfrak{p}$ are relatively prime. By \textbf{Lemma. 127}, $(s)$ and $\mathfrak{p}^m$ are also relatively prime. $\Rightarrow\exists\ b\in A,\ q\in\mathfrak{p}^m$ such that $bs+q=1.\ \Rightarrow 
    1+\mathfrak{p}^m=bs+\mathfrak{p}^m\mapsto\overline{bs}=\overline{1}$ in $A_\mathfrak{p}/\mathfrak{q}^m.\\ \Rightarrow b+\mathfrak{p}^m\mapsto \overline{b}=\frac{1}{\overline{s}}$. Thus, $ba+\mathfrak{p}^m\mapsto \overline{ba}=\overline{b}\cdot\overline{a}=\frac{1}{\overline{s}}\cdot\overline{a}= \overline{\frac{a}{s}}$.
\end{proof}

\begin{proofthm}{125}
    We now prove that a nonzero ideal $0\neq \mathfrak{a}\triangleleft A$ can be factorised into a product of prime ideals. By \textbf{Lemma 126}, $\mathfrak{a}$ contains a product of nonzero prime ideals $\mathfrak{b}\coloneqq\mathfrak{p}_1^{s_1}\cdots\mathfrak{p}_n^{s_n}$.
    We may suppose that the $\mathfrak{p}_i$'s are distinct. Then, 
    \begin{align*}
        A/\mathfrak{b}\ \simeq\ A/\mathfrak{p}_1^{s_1}\times\cdots\times
        A/\mathfrak{p}_n^{s_n}\ \simeq\  A_{\mathfrak{p}_1}/\mathfrak{q}_1^{s_1}\times\cdots\times
        A_{\mathfrak{p}_n}/\mathfrak{q}_1^{s_n},
    \end{align*}
    where $\mathfrak{q}_i\coloneqq\mathfrak{p}_iA_{\mathfrak{p}_i}$ is the maximal ideal of $A_{\mathfrak{p}_i}$. Here, the 1st isomorphism gives by the \textbf{CRT} and \textbf{Lemma 127} and the 2nd gives by \textbf{Lemma 128}. Under this isomorphism, there is a correspondence
    \begin{align*}
        \mathfrak{a}/\mathfrak{b}\ \longleftrightarrow\  \mathfrak{q}_1^{r_1}/\mathfrak{q}_1^{s_1} \times\cdots\times\mathfrak{q}_n^{r_n}/\mathfrak{q}_1^{s_n}
    \end{align*}
    for some $r_i\leq s_i$ (recall that each $A_{\mathfrak{p}_i}$ is a DVR by \textbf{Prop. 124}). Since this ideal is also the image of $\mathfrak{p}_1^{r_1}\cdots\mathfrak{p}_n^{r_n}$ under the isomorphism, we see that 
    \begin{align*}
        \mathfrak{a}= \mathfrak{p}_1^{r_1}\cdots\mathfrak{p}_n^{r_n}\ \text{in}\ A/\mathfrak{b}.
    \end{align*}
    Both of these ideals contain $\mathfrak{b}$, which follows from the \textbf{Correspondence Thm.} that
    \begin{align*}
        \mathfrak{a}= \mathfrak{p}_1^{r_1}\cdots\mathfrak{p}_n^{r_n}\ \text{in}\ A.
    \end{align*}
    Moreover, this factorisation is unique since each $r_i$ is uniquely determined by the condition\\ \centerline{$\mathfrak{a} A_{\mathfrak{p}_i}=\mathfrak{q}_i^{r_i},\ \text{where}\  \mathfrak{q}_i \triangleleft A_{\mathfrak{p}_i}$ maximal ideal.} \qed
\end{proofthm}

\begin{corollary}{129. (Local Criterion for Equality)}
    Let $\mathfrak{a},\mathfrak{b}$ be ideals in $A$. Then 
    \begin{align*}
        \mathfrak{a}\subset\mathfrak{b}\qquad \Longleftrightarrow\qquad
    \mathfrak{a}A_\mathfrak{p}\subset\mathfrak{b}A_\mathfrak{p},\ \forall\ 0\neq\mathfrak{p}\triangleleft A\ \text{prime ideals}.
    \end{align*}
    In particular, $\mathfrak{a}=\mathfrak{b}$ iff $\ \mathfrak{a}A_\mathfrak{p}=
    \mathfrak{b}A_\mathfrak{p}$.
\end{corollary}
\begin{proof}
    "$\Rightarrow$"\ is obvious. "$\Leftarrow$"\ Factor $\mathfrak{a}$ and $\mathfrak{b}$ into 
    \begin{align*}
        \mathfrak{a}\coloneqq\mathfrak{p}_1^{r_1}\cdots\mathfrak{p}_n^{r_n},\qquad \mathfrak{b}\coloneqq\mathfrak{p}_1^{s_1} \cdots\mathfrak{p}_n^{s_n},\qquad (r_i,s_i\geq0).
    \end{align*}
    Then
    \begin{align*}
        \mathfrak{a}A_\mathfrak{p}\subset\mathfrak{b}A_\mathfrak{p}\qquad\Longleftrightarrow\qquad r_i\geq s_i\ \text{for all}\ i,
    \end{align*}
    (recall that $A_{\mathfrak{p}_i}$ is a DVR) and $r_i\geq s_i$ implies that $\mathfrak{a}\subset\mathfrak{b}$.
\end{proof}

\begin{corollary}{130}
    Let $A$ be an integral domain with only finitely many prime ideals. Then $A$ is a Dedekind domain iff it is a PID.
\end{corollary}
\begin{proof}
    Assume $A$ is a Dedekind domain. By \textbf{Thm. 125}, it suffices to show that the prime ideals are principal. Let $\mathfrak{p}_1,\cdots, \mathfrak{p}_m$ be all prime ideals in $A$. We consider the element $x\in\mathfrak{p}_1\setminus\mathfrak{p}_1^2$. By \textbf{CRT}, there is an element $x\in A$ such that 
    \begin{align*}
        x\equiv x_1\ (\mathrm{mod}\ \mathfrak{p}_1^2)\qquad \text{and}\qquad x\equiv1\ (\mathrm{mod}\ \mathfrak{p}_i),\ \forall i\neq1.
    \end{align*}
    Now, the ideals $\mathfrak{p}_1$ and $(x)$ generate the same ideals in $A_{\mathfrak{p}_i}$ for all $i$. Hence, by \textbf{Cor. 129}, we obtain  $\mathfrak{p}_1=(x)$ in $A$.
\end{proof}

\begin{corollary}{131}
    Let $0\neq\mathfrak{a}\subset\mathfrak{b}$ be two ideals in a Dedekind domain $A$, then $\mathfrak{b}\coloneqq\mathfrak{a}+(a)$, for some $a\in A$. In particular, if $0\neq a\in\mathfrak{b}$, then there exists $b\in\mathfrak{b}$ such that $\mathfrak{b}= (a,b)$.
\end{corollary}
\begin{proof}
    Let $\mathfrak{a}\coloneqq\mathfrak{p}_1^{r_1}\cdots\mathfrak{p}_n^{r_n},\ \mathfrak{b}\coloneqq\mathfrak{p}_1^{s_1} \cdots\mathfrak{p}_n^{s_n},\ (r_i,s_i\geq0).\ \because\mathfrak{a} \subset\mathfrak{b}.\ \therefore s_i\leq r_i,\ \forall i$. Now, for each $1\leq i\leq m,\ \exists\ x_i\in A$ such that $x_i\in \mathfrak{p}_i^{s_i}$ but $x_i\notin\mathfrak{p}_i^{s_i+1}$. By \textbf{CRT}, $\exists\ a\in A$ such that $a\equiv x_i\ (\mathrm{mod}\mathfrak{p}_i^{r_i}),\ \forall i$. Now one sees that $\mathfrak{b}+(a)=\mathfrak{a}$ by looking at the ideals they generated in $A_{\mathfrak{p}},$ for all prime ideal $\mathfrak{p}\triangleleft A$. Apply the first part to $0\neq(a)\subset\mathfrak{b}$, we get the last sentence.
\end{proof}

\begin{corollary}{132}
    Let $\mathfrak{a}$ be a nonzero ideal in a Dedekind domain $A$. Then there exists a nonzero ideal $\mathfrak{a}^\ast\triangleleft A$ such that $\mathfrak{a}\cdot \mathfrak{a}^\ast$ is principal. Moreover, $\mathfrak{a}^\ast$ can be chosen to be either relatively prime to any particular ideal $\mathfrak{b}\triangleleft A$, or it can be chosen so that $\mathfrak{a}\cdot\mathfrak{a}^\ast=(a)$ with $a$ any particular element of $\mathfrak{a}$ (but not both).
\end{corollary}
\begin{proof}
    Let $0\neq a\in\mathfrak{a}$, then $(a)\subset\mathfrak{a}$. By \textbf{Thm. 125}, we have 
    \begin{align*}
        (a)\coloneqq\mathfrak{p}_1^{r_1}\cdots\mathfrak{p}_n^{r_n},\ \text{and}\ \mathfrak{a}\coloneqq\mathfrak{p}_1^{s_1} \cdots\mathfrak{p}_n^{s_n},\qquad (s_i\leq r_i).
    \end{align*}
    If $\mathfrak{a}^\ast\coloneqq\mathfrak{p}_1^{r_1-s_1}\cdots\mathfrak{p}_n^{r_n-s_n}$, then $\mathfrak{a}\mathfrak{a}^\ast=(a)$.\\
    We now show that $\mathfrak{a}^\ast$ can be chosen to be prime to $\mathfrak{b}$. We have $\mathfrak{a}\supset\mathfrak{a}\mathfrak{b}$, so by \textbf{Cor. 131}, $\exists\ a\in\mathfrak{a}$ such that $\mathfrak{a}=\mathfrak{a}\mathfrak{b}+(a)$. As $\mathfrak{a}\subset (a)$, we have $(a)=\mathfrak{a}\cdot\mathfrak{a}^\ast$ for some ideal $\mathfrak{a}^\ast\triangleleft A$ (by the above argument); now, $\because \mathfrak{a}\mathfrak{b}+\mathfrak{a}\mathfrak{a}^\ast =\mathfrak{a}.\ \therefore \mathfrak{b}+\mathfrak{a}^\ast=A$. (Otherwise $\mathfrak{b}+\mathfrak{a}^\ast\subset\mathfrak{p}$ for some prime ideal $\mathfrak{p}\triangleleft A.\ \Rightarrow \mathfrak{a}\mathfrak{b}+\mathfrak{a}\mathfrak{a}^\ast=\mathfrak{a} (\mathfrak{b}+\mathfrak{a}^\ast)\subset\mathfrak{a}\mathfrak{p}\neq \mathfrak{a}$.)
\end{proof}

\begin{lemma}{133}
    Let $A\subset B$ be an integral ring extension of integral domains. Then\\ \centerline{$A$ is a field\qquad $\Longleftrightarrow\qquad B$ is a field.\qquad [AM, 5.7]}
\end{lemma}
\begin{proof}
    "$\Rightarrow$"\ Suppose $A$ is a field. Let $0\neq x\in B$, we claim: $x^{-1}$ exists and $x^{-1}\in B$.\\
    $\because x\in B$ and $B$ is integral over $A.\ \therefore x$ satisfies  $x^n+a_1x^{n-1}+ \cdots+a_n=0$ in $B$, for some $a_i\in A$, where $n\in\mathbb{N}$ may be chosen to be the smallest possibly degree. Note that if $a_n=0,\ \because B$ is an ID. $\therefore$ Either $x=0$ or $x^{n-1}+a_1x^{n-2}+\cdots+a_{n-1}= 0$. But $x\in B$ is chosen to be nonzero and $n\in\mathbb{N}$ is chosen to be the smallest. $(\to\gets)$. Thus, $a_n\neq0$ in $A.\ \because A$ is a field. $\therefore a_n^{-1}$ exists in $A$. Hence $x^{-1}\coloneqq-a_n^{-1}(x^{n-1}+a_1x^{n-2}+\cdots+a_{n-1})$ exists in $B$.\\
    "$\Leftarrow$"\ Suppose $B$ is a field. Let $0\neq x\in A\subset B$, then $\exists\ x^{-1}\in B$. We claim: $x^{-1}\in A$. \\
    $\because B$ is integral over $A.\ \therefore x^{-1}$ satisfies $x^{-m}+a_1x^{-m+1}+ \cdots+a_{m-1}x^{-1}+a_m=0$ in $B\ \text{where}\ a_i\in A$. Multiplying $x^{m-1}$, we have $x^{-1}+a_1+a_2x+ \cdots+a_mx^{m-1}=0$ in $B$. Hence, we obtain  $x^{-1}=-(a_1+a_2x+\cdots+a_mx^{m-1})\in A$.  
\end{proof}

\begin{theorem}{134. (Special Case of Akizuki-Krull)}
    Let $A$ be a Dedekind domain with $K\coloneqq\mathrm{Quot}(A)$ and $B$ be the integral closure of $A$ in a finite separable extension $L$ of $K$. Then $B$ is a Dedekind domain. [AM, 5.17]
\end{theorem}
\begin{proof}
    We check the 3 conditions of \textbf{Def. 121}.
    \begin{itemize}
        \item \textbf{($B$ is Noetherian).} $\because A$ is Noetherian. $\therefore$ By \textbf{Prop. 112}, $B$ is contained in a finitely generated $A$-module. It follows that every ideal in $B$ is finitely generated as an $A$-module (being a submodule of a Noetherian $A$-module), and a fortiori as an ideal (i.e. $B$-module). Thus, $B$ is Noetherian.
        \item \textbf{($B$ is integrally closed).} $\overline{B}=\overline{\overline{A}} =\overline{A}=B$.
        \item \textbf{(Every nonzero prime ideal $0\neq\mathfrak{q}\triangleleft B$ is maximal).}\\ Let $0\neq b\in\mathfrak{q},\ \because B$ is integral over $A.\ \therefore b$ satisfies $b^n+a_1b^{n-1}+\cdots+a_n=0$, for some $a_i\in A$, where $n\in\mathbb{N}$ may be chosen to be the smallest possibly degree. Then $a_n\neq0$, which comes from the proof in \textbf{Lemma 133}. $\because a_n\in(b)\cap A.\ \therefore \mathfrak{p} \coloneqq\mathfrak{q}\cap A\neq(0)$. Moreover, $\mathfrak{p}$ is a prime ideal of $A. \because A$ is local. $\therefore\mathfrak{p}$ must be maximal. $\Rightarrow A/\mathfrak{p}$ is a field. Note that the inclusion $A\xhookrightarrow{} B$ extends to an integral extension
        \begin{align*}
            A/\mathfrak{p}\xhookrightarrow{~~~~} B/\mathfrak{q},\qquad a+\mathfrak{p} \longmapsto a+\mathfrak{q}
        \end{align*}
        of IDs. $\because A/\mathfrak{p}$ is a field, so is $B/\mathfrak{q}$ by \textbf{Lemma 133}, and hence $\mathfrak{q} \triangleleft B$ is maximal.
    \end{itemize}
    Therefore, $B$ is a Dedekind domain.
\end{proof}

\begin{corollary}{135}
    The ring of integers $\mathcal{O}_K$ of an algebraic number field $K$ is a Dedekind domain. [AM, 9.5]
\end{corollary}
\newpage

\section{The Ideal Class Group}
\begin{definition}{136. (Fractional Ideal)}
    Let $A$ be a Noetherian integral domain. 
    \begin{itemize}
        \item A \textbf{fractional ideal} is a finitely generated $A$-module inside $K=\mathrm{Quot}(A)$.
        \item The fractional ideal $(b)=bA=\{b\cdot a:a\in A\}$, where $b\neq0$, is called \textbf{principal}.
        \item Fractional ideals inside $A$ (i.e. the usual ideals) are called \textbf{integral}.
        \item The product of two fractional ideals $\mathfrak{a},\ \mathfrak{b}$ of $A$ is defined by:
        \begin{align*}
            \mathfrak{a}\cdot\mathfrak{b}\coloneqq\{\sum_{\text{finite}}a_ib_i: a_i\in\mathfrak{a},\ b_i\in\mathfrak{b}\}.
        \end{align*}
    \end{itemize}
\end{definition}

\begin{remark}{137}
    By multiplying the generators of a fractional ideal $\mathfrak{a}$ with a common denominator, we can always find $d\in A$ such that $d\cdot\mathfrak{a}\subset A$. In fact, the $A$-submodules of $K$ which satisfy this condition are precisely the fractional ideals: $d\cdot\mathfrak{a}$ is finitely generated for $A$ Noetherian, and the map $\mathfrak{a}\to A,\ a\mapsto d\cdot a$ induces an $A$-linear isomorphism.
\end{remark}

\begin{example}{138}
    Let $A$ be a DVR with maximal ideal $\mathfrak{m}=(\pi)$ and field of fractions $K=\mathrm{Quot}(A)$. Every element $b\in K^\times$ can be written as $b\coloneqq u\pi^m$, where $m\in\mathbb{Z},\ u\in A^\times$. If $\mathfrak{a}$ is a fractional ideal, we therefore have $\pi^n\mathfrak{a}\subset A$ for some $n\in\mathbb{N}$, and consequently $\pi^n\mathfrak{a} =(\pi^m)$ for some $m\in\mathbb{N}$. It follows that $\mathfrak{a}=(\pi^{m-n})$ and the fractional ideals of $A$ form a group isomorphic to $\mathbb{Z}$.
\end{example}

\begin{theorem}{139}
    Let $A$ be a Dedekind domain. The set of fractional ideals $\mathsf{Id}(A)$ is the free abelian group generated by the set of prime ideals. The principal ideals form a subgroup $\mathsf{P}(A)$.
\end{theorem}
\begin{proof}
    Note that the composition law is well-defined and it is obviously commutative. For associativity, one may check that
    \begin{align*}
        (\mathfrak{a}\mathfrak{b})\mathfrak{c}\ =\ 
        \{\sum a_ib_ic_i:a_i\in\mathfrak{a},\ b_i\in\mathfrak{b},\ c_i\in\mathfrak{c}\}\ =\ \mathfrak{a}(\mathfrak{b}\mathfrak{c}).
    \end{align*}
    The ring $A$ plays the role of an identity element: $\mathfrak{a}A= \mathfrak{a}$. To show that $\mathsf{Id}(A)$ is a group, it remains to show the existence of an inverse. 
    Let $\mathfrak{a}$ be a nonzero integral ideal. By \textbf{Cor. 132}, there exists an ideal $\mathfrak{a}^\ast\triangleleft A$ and $a\in A$ such that $\mathfrak{a}\mathfrak{a}^\ast=(a).\ \because a\in A\subset K.\ \therefore \exists a^{-1}\in K.\ \Rightarrow \mathfrak{a}\cdot(a^{-1}\cdot \mathfrak{a}^\ast)=(1)=A$, thus $a^{-1}\mathfrak{a}^\ast$ is an inverse of $\mathfrak{a}$. Now, if $\mathfrak{a}$ is a fractional ideal, then $d\cdot\mathfrak{a}$ is an integral ideal of $A$ for some $d\in A$, and $d\cdot(d\mathfrak{a})^{-1}$ will be an inverse for $\mathfrak{a}$. \\
    It remains to show that the group $\mathsf{Id}(A)$ is freely generated by the prime ideals, i.e. each fractional ideal can be expressed in a unique way as a product of powers of prime ideals. \\
    Let $\mathfrak{a}$ be a fractional ideal, then $d\mathfrak{a}$ is an integral ideal for some $d\in A$. We write
    \begin{align*}
        d\mathfrak{a}\coloneqq\mathfrak{p}_1^{r_1}\cdots\mathfrak{p}_m^{r_m}
        \qquad\text{and}\qquad (d)\coloneqq \mathfrak{p}_1^{s_1}\cdots\mathfrak{p}_m^{s_m}.
    \end{align*}
    Thus $\mathfrak{a}=\mathfrak{p}_1^{r_1-s_1}\cdots\mathfrak{p}_m^{r_m-s_m}$. The uniqueness follows from the uniqueness of the factorisation for integral ideals.
\end{proof}

\begin{corollary}{140}
    Let $A$ be a Dedekind domain and $S\subset A$ be a multiplicative subset. Then the extension map $A\to S^{-1}A,\ \mathfrak{a} \mapsto\mathfrak{a}^e=S^{-1}\mathfrak{a}$ defines an isomorphism from the subfroup of $\mathsf{Id}(A)$ generated by prime ideals not meeting $S$  to the group $\mathsf{Id}(S^{-1}A)$.
\end{corollary}
\begin{proof}
    As we have recalled in \textbf{Prop. 27}, the prime ideals of $S^{-1}A$ (which generate $\mathsf{Id}(S^{-1}A)$) are precisely obtained by extending the prime ideals $\mathfrak{p}\triangleleft A$ with $\mathfrak{p}\cap S=\varnothing$ via $A\to S^{-1}A$. The assertion follows from $\mathsf{Id}(S^{-1}A)=\langle\  \mathfrak{q} \triangleleft S^{-1}A:\ \text{prime}\ \rangle$. 
\end{proof}

\begin{remark}{141}
    Since extension commutes with the product of ideals, $(\mathfrak{a}\cdot\mathfrak{b})^e=\mathfrak{a}^e\cdot\mathfrak{b}^e$, the decomposition $\mathfrak{a}\coloneqq\mathfrak{p}_1^{r_1}\cdots\mathfrak{p}_n^{r_n}$ with $\mathfrak{p}_i\cap S=\varnothing$ yields the decomposition $\mathfrak{a}^e=(\mathfrak{p}_1^e)^{r_1}\cdots (\mathfrak{p}_n^e)^{r_n}$ of $\mathfrak{a}^e$ in $S^{-1}A$.
\end{remark}

Note that a PID defines always a UFD, the converse being false: $k[x,y]$ is factorial for a field $k$, but $(x,y)$ cannot be principal, since $x$ and $y$ have no common divisor. However, we have the following:
\begin{proposition}{142}
    If a Dedekind domain is a UFD, then it is a PID.
\end{proposition}
\begin{proof}
    Suppose $A$ is a Dedekind UFD. It suffices to show that the nonzero prime ideals $\mathfrak{p}\triangleleft A$ are principal. It will contain a nonzero element, which is a product of irreducibles. Since $\mathfrak{p}$ is prime, it contains one of the irreducible factors $\pi$, also by \textbf{Cor. 132}, there exists an ideal $\mathfrak{p}^\ast\triangleleft A$ such that $\mathfrak{p}\mathfrak{p}^\ast=(\pi)$. 
    Again, from \textbf{Cor. 132}, there are ideals $\mathfrak{q},\ \mathfrak{q}^\ast \triangleleft A$ such that 
    \begin{align*}
        \mathfrak{p}\mathfrak{q}=(a),\qquad \mathfrak{q}+\mathfrak{p}^\ast=A;
        \qquad \mathfrak{q}\mathfrak{q}^\ast=(b),\qquad \mathfrak{q}^\ast+\mathfrak{p}=A,
    \end{align*}
    for some $a,\ b\in A$. 
    \begin{itemize}
        \item $\because(\pi b)=\mathfrak{p}\mathfrak{p}^\ast\mathfrak{q} \mathfrak{q}^\ast=(a)\mathfrak{p}^\ast\mathfrak{q}^\ast. \therefore a|\pi b.\ \Rightarrow \exists\ c\in A$ such that $\pi b=ac.\ \because A$ is UFD. $\therefore$ We have either $\pi|a$ or $\pi|c$.
        \item If $\pi|a$, then $\frac{a}{\pi}\in A.\ \Rightarrow (\frac{\pi}{a})\mathfrak{p}^\ast=\frac{\mathfrak{p}\mathfrak{q}}{\mathfrak{p}\mathfrak{p}^\ast}\mathfrak{q}=\mathfrak{q}$. Thus any prime ideal dividing $\mathfrak{p}^\ast$ will also divide $\mathfrak{q}$, which contradicts that $\mathfrak{q}$ and $\mathfrak{p}^\ast$ are relatively prime, i.e. $\mathfrak{q}+\mathfrak{p}^\ast =A.\ (\to\gets)$. Therefore, there is no such ideal $\mathfrak{q}\neq(0)$, so $\mathfrak{p}^\ast=A$.
        \item Similarly, if $\pi|c$, then $\frac{\pi}{c}\in A.\ \Rightarrow (\frac{\pi}{c})\mathfrak{p}=\mathfrak{q}^\ast.\ (\to\gets)$. 
    \end{itemize}
    Hence, $\mathfrak{p}=\mathfrak{p}A=\mathfrak{p}\mathfrak{p}^\ast=(\pi)$ is principal.
\end{proof}

\begin{definition}{143. (Ideal Class Group)}
    \begin{itemize}
        \item The \textbf{ideal class group} $\mathsf{Cl}(A)$ is the quotient $\mathsf{Id}(A)/\mathsf{P}(A)$.
        \item The \textbf{class number} of $A$ is the number $\# \mathsf{Cl}(A)$ if it is finite.
        \item If $A=\mathcal{O}_K$ for a number field $K$, we simply speak of the ideal class group, respectively, the class number of $K$.
    \end{itemize}
\end{definition}
\newpage
The ideal class group of an algebraic number field can be effectively computed by an algorithm. Unfortunately, we cannot prove this nor the following important result for the lack of time:
\begin{theorem}{144}
    The class number of a number field is finite.
\end{theorem}

We recall that any finitely generated $\mathbb{Z}$-module is the direct sum of its torsion submodule and a free $\mathbb{Z}$-module of rank $r$. We state this without proof:
\begin{theorem}{145}
    Let $A$ be a Dedekind domain.
    \begin{itemize}
        \item[(a)] Every finitely generated torsion-free $A$-module $M$ is isomorphic to a direct sum of fractional ideals.
        \item[(b)] Two finitely generated torsion-free $A$-modules $M\simeq\mathfrak{a}_1\oplus\cdots\oplus\mathfrak{a}_m$ and $N\simeq\mathfrak{b}_1\oplus\cdots\oplus\mathfrak{b}_n$ are isomorphic iff $m=n$ and $\prod\mathfrak{a}_i\equiv \prod\mathfrak{b}_i$ modulo principal ideals.
    \end{itemize}
\end{theorem}

In \textbf{Prop. 88} we determined the units of $\mathbb{Z}[i]=\mathcal{O}_{\mathbb{Q}(i)}$. Let us now consider a general ring of integers $\mathcal{O}_K$, ($K$: number field). We have $[K:\mathbb{Q}]=r+2s$ embeddings $K\xhookrightarrow{} \mathbb{C}$, where $r$ embeddings factorise via $\mathbb{R}$ and $2s$ which do not. For instance, if $K=\mathbb{Q}[\alpha]$, then $r$ is the number of real roots of the minimal polynomial $f_\alpha$, while $s$ is the number of complex-conjugated pairs of root which is why the number of embeddings into $\mathbb{C}$ is even: We can always apply complex conjugation.

\begin{theorem}{146. (Dirichlet's Unit Theorem)}
    The group of units $\mathcal{O}_K^\times$ in a number field $K$ is finitely generated. Considered as a $\mathbb{Z}$-module, the rank of its free summand is equal to $r+s-1$.
\end{theorem}

\begin{example}{147}
    For $\mathbb{Q}(i)$, we find $r=0$ and $s=1$. Hence $\mathbb{Z}[i]^\times$ is pure torsion. (cf. \textbf{Prop. 88})
\end{example}

\newpage

\section{Discrete Valuations}
\begin{definition}{148. (Discrete Valuation)}
    Let $K$ be a field. A \textbf{discrete valuation} on $K$ is a group homomorphism $\nu:K^\times\to\mathbb{Z}$ such that 
    \begin{align*}
        \nu(x+y)\geq\mathrm{min}\{\nu(x),\nu(y)\},\qquad\text{if}\ x+y\neq0.
    \end{align*}
    Its image is thus of the form $m\mathbb{Z}$. We say that $\nu$ is \textbf{normalised} if $m=1$.
\end{definition}

\begin{remark}{149}
    \begin{itemize}
        \item[(a)] It is sometimes convenient to extend the definition of $\nu$ to all of $K$ by setting $\nu(0)\coloneqq\infty$.
        \item[(b)] We can always pass to a normalised valuation by considering $x\mapsto m^{-1}\nu(x)$.
        \item[(c)] It follows from the definition that $\nu(x^{-1})=-\nu(x)$.
    \end{itemize}
\end{remark}

\begin{example}{150}
    \begin{itemize}
        \item[(a)] Let $A$ be a UFD, $K\coloneqq\mathrm{Quot}A$ and $p\in A$ a prime. Up to units, any element $x\in K^\times$ can be written as $x\coloneqq p^r\cdot\frac{a}{b}$, where $p\nmid a,b$, so that the integer is uniquely determined. Then $\nu_p(x)\coloneqq r$ defines a discrete valuation. Take, for example, $A=\mathbb{Z},\ K=\mathbb{Q}$ and $p\in\mathbb{Z}$ any prime number or $A=k[x],\ K=k(x)$ and $p=f\in k[x]$ any irreducible polynomial.
        \item[(b)] Consider the field $\mathcal{M}$ of meromorphic functions on a connected open subset $\mathcal{U}$ of $\mathbb{C}$. For each $P\in\mathcal{U}$ and $f\in\mathcal{M}^\times$, we define $\mathrm{ord}_P(f)$ to be $0,\ m$ or $-m$ according to whether or not $f(P)$ is a unit in $\mathbb{C}$, or has a zero or a pole of order $m$. Then $\mathrm{ord}_P:\mathcal{M}^\times\to\mathbb{Z}$ is a discrete valuation.
        \item[(c)] Let $A$ be a Dedekind domain and $\mathfrak{p}\triangleleft A$ be a prime ideal. For any $x\in K^\times$, where $K=\mathrm{Quot}A$, let $\mathfrak{p}^{\nu_\mathfrak{p}(x)}$ be the power of $\mathfrak{p}$ in the decomposition of $(x)$. By uniqueness, this defines a discrete valuation $\nu_\mathfrak{p}:K^\times\to\mathbb{Z}$. In particular, a DVR (= local Dedekind domain) gives rise to a natural discrete valuation on its quotient field.
    \end{itemize}
\end{example}
\newpage
Further, there is a converse to \textbf{Example 150(c)}:
\begin{proposition}{151}
    Let $\nu$ be a discrete valuation on $K$, then 
    \begin{align*}
        A\ \coloneqq\ \{a\in K:\nu(a)\geq0\}
    \end{align*}
    is a DVR with maximal ideal and group of units given by
    \begin{align*}
        \mathfrak{m}=\{a\in K:\nu(a)>0\}\qquad\text{and}\qquad A^\times=\{a\in K:\nu(a)=0\}.
    \end{align*}
    Moreover, $\mathfrak{m}$ is generated by any element $\pi$ with $\nu(\pi)=1$ and the ideals of $A$ are given by
    \begin{align*}
        \mathfrak{m}^n=(\pi^n)=\{x\in K:\nu(x)\geq n\},\qquad (n\geq0).
    \end{align*}
\end{proposition}
\begin{proof}
    By \textbf{Rmk. 149(b)}, we may pass $\nu$ to a normalised valuation $\nu:K^\times\twoheadrightarrow\mathbb{Z}$. Thus $\exists\ \pi\in K^\times$ such that $\nu(\pi)=1$. Moreover, with operations induced by $K$, $A$ is an ID.\\
    \textbf{Claim 1.} $A^\times=\{a\in K:\nu(a)=0\}$.\\ \textsc{proof.}
    "$\subset$"\ If $a\in  A^\times$, then $\exists\ b\in A$ such that $ab=1.\ \because \nu(1)=\nu(1\cdot1)=\nu(1)+\nu(1).\\ \therefore \nu(1)=0.\ \Rightarrow \nu(a)+\nu(b)=\nu(a\cdot b)=\nu(1)=0.\ \because$ Both $\nu(a)$ and $\nu(b)\geq0.\ \therefore \nu(a)=\nu(b)=0$. "$\supset$"\ Let $0\neq x\in K$ such that $\nu(x)=0$. Then $x\in A=\{a\in K^\times:\nu(a)\geq0\}.\\ \because 0=\nu(1)=\nu(x\cdot x^{-1})=\nu(x)+\nu(x^{-1}). \ \therefore \nu(x^{-1})=0.\ \Rightarrow x^{-1}\in A$. Hence $x\in A^\times$.\\
    \textbf{Remark.} $\forall0\neq x\in A=\{a\in K^\times:\nu(a)\geq0\}$, say $\nu(x)\coloneqq n\geq0$. Consider $\frac{x}{\pi^n}=x\cdot\frac{1}{\pi^n}.\\ \Rightarrow \nu(\frac{x}{\pi^n})=\nu(x)+\nu(\frac{1}{\pi^n})=n+(-n)=0$. By \textsc{claim 1} we have $\frac{x}{\pi^n}\in A$, it follows that $\exists\ u\in A^\times$ such that $x=u\cdot\pi^n$. Moreover, $\because \nu(\pi)=1.\ \therefore\pi\notin A^\times.\ \Rightarrow 0\neq(\pi) \triangleleftneq A$.\\ 
    \textbf{Claim 2.} For any $0\neq\mathfrak{a}\triangleleftneq A,\ \exists\ l\in\mathbb{N}$ such that $\mathfrak{a}=(\pi^l)$.\\ \textsc{proof.}
    Consider $\nu|_A:A\to\mathbb{Z}_{\geq0}$ with $\mathfrak{a}\twoheadrightarrow \nu(\mathfrak{a})$. We take $l\coloneqq\mathrm{min}\{\nu(a):a\in\mathfrak{a}\}$ and check that $(\pi^l)=\mathfrak{a}$. "$\subset$"\ By definition there is $a\in\mathfrak{a}$ such that $\nu(a)=l$. Note that from \textsc{remark}, $\nu(\frac{a}{\pi^l})=0.\ \Rightarrow \frac{a}{\pi^l}\in A^\times.\ \Rightarrow \exists\ u\in A^\times$ such that $a=u\cdot\pi^l$. Thus $\pi^l=u^{-1}a\in\mathfrak{a}$. "$\supset$"\ For any $0\neq a\in\mathfrak{a}$, from \textsc{remark} we write $a\coloneqq u\cdot\pi^{\nu(a)}$, where $\nu(a)\geq0$ (since $\nu(a)\geq l$). Note that $\because \mathfrak{a}\triangleleftneq A$ is proper.\\ $\therefore$ By \textsc{claim 1}, $\nu(a)>0$, further, from the definition of $l$, we have $\nu(a)\geq l$. Now consider $\frac{a}{\pi^l}\in K^\times$, then $\nu(\frac{a}{\pi^l})=\nu(a)-l\geq0.\ \Rightarrow \frac{a}{\pi^l}\in A.\ \Rightarrow \exists\ b\in A$ such that $a=b\cdot\pi^l\in(\pi^l)$. Thus $\mathfrak{a}\subset(\pi^l)$. To conclude, \textsc{claim 2} implies that $A$ is a PID.\\
    \textbf{Claim 3.} $A$ has only one nonzero prime ideal of $A$.\\
    \textsc{proof.} Let $0\neq\mathfrak{a}\triangleleftneq A$, from \textsc{claim 2}, we have $\mathfrak{a}\coloneqq(\pi^n)$ for some $n\in\mathbb{N}.\ \Rightarrow \mathfrak{a}=(\pi^n)=(\pi).\\ \because \nu(\pi)=1.\ \therefore$ \textsc{claim 1}, $(\pi)\triangleleftneq A$ is a maximal ideal. Now if $n\geq2,\ (\pi^2)$ is not a prime ideal (since $\pi\notin(\pi^2)$ but $\pi\cdot \pi\in(\pi^2)$). Thus, if $\mathfrak{a}=(\pi^n)$ is a prime ideal of $A$, we must have $n=1$.
\end{proof}

\begin{example}{152}
    If $K=\mathbb{Q}$ and $\nu_p$ is the valuation defined by the prime $p\in\mathbb{Z}$, we find $A=\mathbb{Z}_{(p)}$ and $\mathfrak{m}=(p)\mathbb{Z}_{(p)}$. 
\end{example}

\begin{proposition}{153}
    Let $a_1,\cdots,a_m$ be elements of a Dedekind domain $A$ and $\mathfrak{p}_1,\cdots ,\mathfrak{p}_m\triangleleft A$ be distinct prime ideals. Then, for any integer $n\in\mathbb{N}$, there is an $a\in A$ such that $\nu_{\mathfrak{p}_i}(a-a_i)>n$. 
\end{proposition}
\begin{proof}
    $\because \mathfrak{p}_i$ and $\mathfrak{p}_j$ are relatively prime $\forall i\neq j.\ \therefore$ By \textbf{Lemma 127}, the ideals $\mathfrak{p}_i^{n+1}$ are pairwise relatively prime. By \textbf{CRT}, $\exists\ x\in A$ such that 
    \begin{align*}
        a\equiv a_i\ (\mathrm{mod}\ \mathfrak{p}_i^{n+1}),\qquad (1\leq i\leq m).
    \end{align*}
    So $a-a_i\in\mathfrak{p}_i^{n+1}$, and hence $(a-a_i)\subset\mathfrak{p}_i^{n+1}$, so that $\mathfrak{p}_i$ appears in the decomposition of $(a-a_i)$ with power $\geq n+1$. Therefore $\nu_{\mathfrak{p}_i}(a-a_i)>n,\ \forall i$.
\end{proof}

\begin{remark}{154}
    As we are going to see in \textbf{\S 2.4}, $\nu(a)\gg n$ means morally a "{}small"\  $a$. In this way, we can regard \textbf{Prop. 153} as an approximation version of the \textbf{CRT}. (cf. \textbf{Theorem 257}).
\end{remark}
\newpage

\section{Ramification}
For the remainder of this section, let $A$ be a Dedekind domain, $K=\mathrm{Quot}A$, and $B$ be the integral closure of $A$ in a finite separable extension $L$ of $K$. Then $B$ is also Dedekind. The prototype of this setup will be $A=\mathbb{Z},\ K=\mathbb{Q}$ with $B=\mathbb{Z}[i]$ in $L=\mathbb{Q}(i)$.

\begin{definition}{155. (Ramification)}
    We first extend a prime ideal $\mathfrak{p}\triangleleft A$ to $B$, where it factorises into 
    \begin{equation*}
        \mathfrak{p}^e=\mathfrak{p}B\coloneqq\mathfrak{P}_1^{e_1}\cdots
        \mathfrak{P}_g^{e_g},\ (e_1\geq1).
    \end{equation*}
    \begin{itemize}
        \item If there is $i$ with $e_i>1$, we say that $\mathfrak{p}$ \textbf{ramifies} in $B$ (or $L$). The number $e_i$ is called the \textbf{ramification index} at $\mathfrak{P}_i$.
        \item We say that $\mathfrak{P}$ \textbf{divides} $\mathfrak{p}$ if $\mathfrak{P}$ occurs in the factorisation of $\mathfrak{p}$ in $B$. We then write $\mathfrak{P}_i|\mathfrak{p}$ and $e_i=e(\mathfrak{P}_i/\mathfrak{p})$ for the ramification index. 
        \item The \textbf{residue class degree} $f_i=f(\mathfrak{P}_i/\mathfrak{p})$ is defined to be $[B/\mathfrak{P}_i:A/\mathfrak{p}]$.
        \item A prime $\mathfrak{p}$ is said to \textbf{split} in $L$ if $e=f=1$ for all $i$; and is said to be \textbf{inert} in $L$ if $\mathfrak{p}B$ is a prime ideal in $B$, (i.e. $e=g=1$.)
    \end{itemize}
\end{definition} 

\begin{lemma}{157}
    A prime ideal $\mathfrak{P}\triangleleft B$ divides a prime ideal $\mathfrak{p}\triangleleft A$ iff $\mathfrak{p}=\mathfrak{P}\cap K$. In particular, $\mathfrak{P}$ divides at most one prime ideal.
\end{lemma}
\begin{proof}
    "$\Rightarrow$"\ Since $\mathfrak{p}B\subset\mathfrak{P}$, we have $\mathfrak{p}\subset\mathfrak{P}\cap K\subsetneq A$. Otherwise, if $\mathfrak{P}\cap K=A$, then $1\in\mathfrak{P}.\ (\to\gets)$. As $\mathfrak{p}$ is maximal, this implies that $\mathfrak{p}=\mathfrak{P} \cap K$.\\ "$\Leftarrow$"\ $\because\mathfrak{p}=\mathfrak{P}\cap K.\ \therefore\mathfrak{p}B\subset\mathfrak{P}$. Thus, by \textbf{Thm. 125}, $\mathfrak{P}$ occurs in the factorisation of $\mathfrak{p}B$.
\end{proof}

\begin{theorem}{158}
    Let $m\coloneqq[L:K]$ and let $\mathfrak{P}_1,\cdots,\mathfrak{P}_g \triangleleft B$ be the prime ideals dividing the prime ideal $\mathfrak{p}\triangleleft A$. Then $\sum_{i=1}^ge_if_i=m$.\\
    Moreover, if $L/K$ is Galois, then all the ramification indices and residue class degrees are equal, whence $efg=m$.
\end{theorem}
\begin{proof}
    We divide the proof into two parts.\\
    \textsc{claim 1.} To prove $\sum_{i=1}^ge_if_i=m$, we show that both sides equal $[B/\mathfrak{p}B:A/\mathfrak{p}]$.\\
    \textsc{(lhs):} Note that from \textbf{CRT}, $B/\mathfrak{p}B=B/ \prod\mathfrak{P}_i^{e_i}\simeq\prod(B/\mathfrak{P}_i^{e_i})$, thus it suffices to show that each $[B/\mathfrak{P}_i^{e_i}:A/\mathfrak{p}]= e_if_i$. From the definition of $f_i$, we know that $B/\mathfrak{P}_i$ is a field of degree $f_i$ over $A/\mathfrak{p}$. For each $r_i,\ \mathfrak{P}_i^{r_i}/\mathfrak{P}_i^{r_i+1}$ is a $B/\mathfrak{P}_i$-module (thus, vector space). Since there is no ideal $\mathfrak{b}\triangleleft B$ such that $\mathfrak{P}_i^{r_i+1}\subsetneq\mathfrak{b} \subsetneq\mathfrak{P}_i$,  it must have dimension $1$ as a $B/\mathfrak{P}_i$-vector space, and hence dimension $f_i$ as an $A/\mathfrak{p}$-vector space. Therefore each quotient in the chain 
    \begin{align*}
        B\supset\mathfrak{P}_i\supset\mathfrak{P}_i^2\supset\cdots\supset \mathfrak{P}_i^{e^i}
    \end{align*}
    has dimension $f_i$ over $A/\mathfrak{p}$. Thus $[B/\mathfrak{P}_i^{e_i}:A/\mathfrak{p}]= e_if_i$.\\
    \textsc{(rhs):} We first assume that $A$ is a PID, by \textbf{Rmk. 113} $B$ admits an integral basis which is a $K$-basis over for $L$, hence $B\cong A^m$. Consider the exact sequence $0\to A^m\xrightarrow{\sim} B\to0$, and apply the (right) exact functor $-\otimes_AA/\mathfrak{p}$, which yields that 
    \begin{equation*}
        0 \longrightarrow A^m\otimes_A(A/\mathfrak{p})\longrightarrow B\otimes_A(A/\mathfrak{p})\longrightarrow0.
    \end{equation*}
    Moreover, $A^m\otimes_A(A/\mathfrak{p})\cong(A\otimes_A(A/\mathfrak{p}))^m\cong(A/\mathfrak{p})^m$ as an $A$-module, thus as an $A/\mathfrak{p}$-module. Further, the short exact sequence $0\to\mathfrak{p}\hookrightarrow A\twoheadrightarrow A/\mathfrak{p}\to 0$ applying the functor $B\otimes_A-$ gives 
    \begin{equation*}
        B\otimes_A\mathfrak{p}\longrightarrow B\otimes_AA\cong B\xtwoheadrightarrow{~} B\otimes_AA/\mathfrak{p}\longrightarrow0.
    \end{equation*}
    This induces that $B/\mathfrak{p}B\cong B\otimes_AA/\mathfrak{p}$ (the image of $B\otimes_A\mathfrak{p}\to B,\ b\otimes a\mapsto a\cdot b$ is just $\mathfrak{p}B$), whence $B/\mathfrak{p}B\cong(A/\mathfrak{p})^m$ is free and thus $[B/\mathfrak{p}B:A/\mathfrak{p}]=m$.\\
    Now, if $A$ is not a PID, we let $S$ be a multiplicative subset of $A$ disjoint from $\mathfrak{p}$ such that $S^{-1}A$ is principal (the localisation of a Dedekind ring), (e.g. $S=A\setminus\mathfrak{p}$). Write $A'\coloneqq S^{-1}A$ and $B'\coloneqq S^{-1}B$, then, by \textbf{Rmk. 141}, we have $\mathfrak{p}B'=\prod(\mathfrak{P}_iB')^{e_i}$. It follows that $\sum_{i=1}^ge_if_i=[B'/\mathfrak{p}B':A'/\mathfrak{p}A']$ since $B'$ is a free $A'$-module; but $A'$ is also principal, thus $[B'/\mathfrak{p}B':A'/\mathfrak{p}A']=m$.\\
    \textsc{claim 2.} Now assume that $L/K$ is Galois. We note that
    \begin{itemize}
        \item $\sigma\in\mathrm{Gal}(L/K)$ maps $B$ isomorphically to itself, since an integral relation $x^n+\sum a_ix^i=0$ for $x\in L$ is mapped to an integral relation $\sigma(x)^n+\sum a_i\sigma(x)^i=0$, whence $\sigma(n)\in B$. 
        \item In particular, if $\mathfrak{P}\triangleleft B$ is a prime ideal, then so is $\sigma\mathfrak{P}\subset B$.
        \item If $\mathfrak{P}|\mathfrak{p}$, it follows from \textbf{Lemma 157} that $\mathfrak{p}=\sigma\mathfrak{P}\cap K$ (since $\sigma|_K=1_K$), so that $\sigma\mathfrak{P}|\mathfrak{p}$.
        \item Since $\sigma|_B$ is a ring isomorphism, we have $e(\sigma \mathfrak{P}/\mathfrak{p})=e(\mathfrak{P}/\mathfrak{p})$ and $f(\sigma\mathfrak{P}/\mathfrak{p})=f(\mathfrak{P}/\mathfrak{p})$
    \end{itemize}
    Thus, it remains to show that $\mathrm{Gal}(L/K)$ acts transitively on the prime ideals of $B$ dividing $\mathfrak{p}$.\\
    Suppose $\mathfrak{Q}$ also divides $\mathfrak{p}$ but is not conjugate to $\mathfrak{P}$, i.e. $\sigma\mathfrak{P}\neq\mathfrak{Q},\ \forall\sigma\in\mathrm{Gal}(L/K)$. By \textbf{CRT} we can choose $b\in B$ with $b\equiv0\ (\mathrm{mod}\ \mathfrak{Q})$ and $b\equiv1\ (\mathrm{mod}\ \sigma\mathfrak{Q}),\ \forall\sigma\in\mathrm{Gal}(L/K)$ so that $b\in\mathfrak{Q}$ but $b\notin\sigma\mathfrak{P}$, and thus $\sigma(b)\notin\mathfrak{P},\ \forall\sigma\in\mathrm{Gal}(L/K)$. By integrality of $b$, we have $a\coloneqq\mathrm{N}_{L/K}(b)=\prod\sigma (b)\in A$. Since $a=b\cdot\prod_{\sigma\neq\epsilon}\sigma(b)\in \mathfrak{Q}$, we have $a\in\mathfrak{Q}\cap A=\mathfrak{p}$. On the other hand, for each $\sigma\in\mathrm{Gal}(L/K),\ \sigma(b)\notin \mathfrak{P}$ and $a=\prod\sigma(b)\in\mathfrak{p}\subset\mathfrak{P}$ contradicts the primality of $\mathfrak{P}.\ (\to\gets)$ 
\end{proof}

\begin{remark}{159}
    This theorem has a very nice geometric interpretation in terms of ramified coverings of curves provided $f_i=1$ for all $i$, for instance, if the fields $A/\mathfrak{p}$ are algebraically closed as for $A=k[x],\ (k: \text{algebraically closed})$.
\end{remark}
\newpage
\begin{theorem}{160}
    Let $A$ be a Dedekind domain, $K\coloneqq\mathrm{Quot}A$, $L/K$ a finite field extension and $B$ be the integral closure of $A$ in $L$. Assume that $B\simeq A^m$ (e.g. if $A$: PID). Then a prime ideal $\mathfrak{p}\triangleleft A$ ramifies in $L$ iff $\mathfrak{p}|\mathrm{disc}(B/A)$. In particular, only finitely many prime ideals ramify.
\end{theorem}

Recall \textbf{Def. 107} for the definition of discriminant. We subdivide the proof into several lemmas:
\begin{lemma}{161}
    Let $A\subset B$ be a ring extension such that $B\cong A^m$ with finite $A$-basis $\{e_1,\cdots,e_m\}$. For any ideal $\mathfrak{a}\triangleleft A,\ \{\overline{e_1},\cdots,\overline{e_m}\}$ is an $A/\mathfrak{a}$-basis for $B/\mathfrak{a}B$ and \\
    \centerline{$D(\overline{e_1},\cdots,\overline{e_m})\ \equiv\  D(e_1,\cdots,e_m)\qquad(\mathrm{mod}\ \mathfrak{a})$.}
\end{lemma}
\begin{proof}
    As in the proof of \textbf{Thm. 158}, the isomorphism gives
    \begin{align*}
        A^m\longrightarrow B,\qquad (a_1,\cdots,a_m)\longmapsto\sum a_ie_i
    \end{align*}
    gives, when applying the right exact functor $-\otimes_A(A/\mathfrak{a})$, an isomorphism
    \begin{align*}
        (A/\mathfrak{a})^m\longrightarrow B/\mathfrak{a},\qquad (a_1,\cdots,a_m)\longmapsto\sum a_i\overline{e_i}
    \end{align*}
    which shows that $\{\overline{e_1},\cdots,\overline{e_m}\}$ is an $A/\mathfrak{a}$-basis for $B/\mathfrak{a}B$. Note that moding out $\mathfrak{a}B$ gives
    \begin{align*}
        \mathrm{det(Tr}_{(B/\mathfrak{a}B)/(A/\mathfrak{a})}(\overline{e_i}\cdot\overline{e_j}))=\overline{\mathrm{det(Tr}_{B/A}(e_i\cdot e_j))},
    \end{align*}
    whence the second result.
\end{proof}

\begin{lemma}{162}
    Let $A\subset B_1,\cdots,B_g$ be the ring extensions such that each $B_i$ is a free $A$-module of finite rank. Then\\
    \centerline{$\mathrm{disc}((\prod B_i)/A)\ =\ \prod\mathrm{disc} (B_i/A)$.}
\end{lemma}
\begin{proof}
    For each $1\leq i\leq g$, choose an $A$-basis $\beta_i=\{e_{ij}\}_{j=1}^{m_i}$ for $B_i\cong A^{m_i}$. Then the matrix $[\mathrm{Tr}_{(\prod B_i)/A}(e_{i_j}\cdot e_{i_k})]_{1\leq j,k\leq m_i}$ has the block form
    \begin{center}
        $\begin{pmatrix}
        \mathrm{Tr}_{B_1/A}(e_{1_j}e_{1_k}) & O & \cdots & O \\
        O & \ddots & \ddots  & \vdots \\
        \vdots & \ddots & \ddots & O \\
        O & \cdots & O & \mathrm{Tr}_{B_g/A}(e_{g_j}e_{g_k}) \\
        \end{pmatrix}$
    \end{center} 
    and the result follows from the standard properties of the determinant.
\end{proof}

\begin{lemma}{163}
    Let $k$ be a perfect field and $B$ be a $k$-algebra of finite dimension. Then\\  \centerline{$B$ is reduced $\Longleftrightarrow\ \mathrm{disc}(B/k)\neq0$.}
\end{lemma}
\textbf{Recall.} A ring is reduced if it has no nonzero nilpotents.
\begin{proof}
    "$\Leftarrow$"\ We suppose to the contrary. Let $0\neq\beta\in B$ be nilpotent and choose a $k$-basis $\{e_1=\beta,e_2,\cdots,e_m\}$ for $B$. Then $\beta e_i$ is nilpotent for all $i$, and so the $k$-linear map 
    \begin{align*}
        B\longrightarrow B,\qquad x\longmapsto\beta e_ix
    \end{align*}
    is nilpotent. Its matrix is also nilpotent, but a nilpotent matrix has trace $0$ \textemdash its minimum polynomial (and hence characteristic polynomial) is of the form $X^r$ \textemdash and so the first row of the matrix $[\mathrm{Tr}(e_ie_j)]$ vanishes. Therefore its determinant (the generator of $\mathrm{disc}(B/k)$) is $0.\ (\to\gets)$\\
    "$\Rightarrow$"\ Suppose $B$ is reduced. We first show that the intersection $\mathfrak{N}$ of the prime ideals of $B$ is $0$, i.e. the \textbf{nilradical} $\mathfrak{N}=\bigcap\{\mathfrak{p}\triangleleft B:\ \text{prime}\}=\{0\}$.  This, in fact, is true for any reduced Noetherian ring, cf. [AM, 7.15]. Now let $\mathfrak{p}\triangleleft B$ be a prime ideal, then the inclusion $k\subset B/\mathfrak{p}$ descends to an integral ring extension since $B$ is of finite type over $k$, thus $B/\mathfrak{p}$ is a field by \textbf{Lemma 133}. Therefore $\mathfrak{p}$ is maximal. Now, let $\mathfrak{p}_1,\cdots,\mathfrak{p}_r \triangleleft B$ be prime (thus, maximal) ideals, then they are pairwise relatively prime. By \textbf{CRT}, we have $B/(\bigcap\mathfrak{p}_i)\cong \prod(B/\mathfrak{p}_i)$. Note that
    \begin{align*}
        [B:k]\geq[B/(\bigcap\mathfrak{p}_i):k]=\sum[B/\mathfrak{p}_i:k]\geq r.
    \end{align*}
    Therefore $B$ has only finitely many prime ideals, say $\mathfrak{p}_1,\cdots, \mathfrak{p}_g$ where $g\leq[B:k]$ and $\bigcap_{i=1}^g\mathfrak{p}_i=0$ (since $B$ is reduced). Again by \textbf{CRT}, we have $B\cong\prod_{i=1}^g (B/\mathfrak{p}_i)$.\\
    For each $i$, since $B/\mathfrak{p}_i\supset k$ is a finite extension and $k$ is perfect, it is also separable extension. Hence $\mathrm{disc}((B/ \mathfrak{p}_i)/k)\neq0$ by \textbf{Prop. 110}. Applt \textbf{Lemma 162}, we deduce that $\mathrm{disc}(B/k)\neq0$. 
\end{proof}

\begin{proofthm}{160}
    The idea is that a prime ideal $\mathfrak{p}$ that ramifies in $B$ gives rise to a non-reduced ring $B/\mathfrak{p}^r$, and this can be detected via the discriminant. From \textbf{Lemma 161}, we see that
    \begin{align*}
        \mathrm{disc}(B/A)\ (\mathrm{mod}\ \mathfrak{p})\ =\  \mathrm{disc}((B/\mathfrak{p}B)/(A/\mathfrak{p})),
    \end{align*}
    and from \textbf{Lemma 163} that 
    \begin{align*}
        \mathrm{disc}((B/\mathfrak{p}B)/(A/\mathfrak{p}))=0\ \Longleftrightarrow\ B/\mathfrak{p}B\ \text{is not reduced}.
    \end{align*}
    Now let $\mathfrak{p}B\coloneqq\prod\mathfrak{P}_i^{e_i}$. Then $B/\mathfrak{p}B \simeq\prod(B/\mathfrak{P}_i^{e_i})$ and 
    \begin{align*}
        \prod(B/\mathfrak{P}_i^{e_i})\ \text{is reduced}\ \Longleftrightarrow\ B/\mathfrak{P}_i^{e_i}\ \text{is reduced},\ \forall i\ \Longleftrightarrow\ e_i=1,\ \forall i. 
    \end{align*}\qed
\end{proofthm}

Next we want to compute some explicit examples.
\begin{proposition}{164}
    Let $A$ be a Dedekind domain, $K\coloneqq\mathrm{Quot}A$, $L/K$ a finite separable extension and $B$ be the integral closure of $A$ in $L$. Furthermore, assume that $B=A[\alpha]$. Let $f=f_\alpha$ be the minimum polynomial of $\alpha$ over $K$ and $\mathfrak{p}\triangleleft A$ a prime ideal. Choose monic polynomials $g_1,\cdots,g_r\in A[x]$ that are distinct and irreducible modulo $\mathfrak{p}$, and such that $f\equiv \prod g_i^{e_i}\ (\mathrm{mod}\ \mathfrak{p})$. Then 
    \begin{align*}
        \mathfrak{p}B\ =\ \prod\ (\mathfrak{p},g_i(\alpha))^{e_i}
    \end{align*}
    is the factorisation of $\mathfrak{p}B$ into a product of distinct prime ideals. The residue field is 
    \begin{align*}
        B/(\mathfrak{p},g_i(\alpha))\ \cong\ (A/\mathfrak{p})[x]/(\overline{g}_i),
    \end{align*}
    hence the residue class degree is $f_i=\mathrm{deg}\ g_i$.
\end{proposition}
\begin{proof}
    Note that $f\in A[x]$ since $\alpha$ is integral over $A$, thus our assumption on $B$ gives the isomorphism 
    \begin{align*}
        A[x]/(f)\xrightarrow{~~\cong~~} A[\alpha]=B,\qquad x\longmapsto\alpha.
    \end{align*}
    Tensoring with the field $k=A/\mathfrak{p}$ yields the isomorphism
    \begin{align*}
        k[x]/(\bar{f})\xrightarrow{~~\cong~~} B/\mathfrak{p}B,\qquad
        x\longmapsto\alpha.
    \end{align*}
    Since $k[x]$ is a PID, the maximal ideals of $k[x]/(\overline{f})$ are given by $(\overline{g}_i),\ (1\leq i\leq r)$, which correspond to the ideals $(g_i(\alpha))+\mathfrak{p}B$ in $B/\mathfrak{p}B$, and these correspond to the ideals $\mathfrak{P}_i\coloneqq(\mathfrak{p},g_i(\alpha))$ in $B$. Thus $\{\mathfrak{P}_i\}_{i=1}^r$ is the complete set of prime ideals containing $\mathfrak{p}B$, and hence is the complete set of prime divisors of $\mathfrak{p}$ by \textbf{Thm. 125}.\\
    The decomposition $\mathfrak{p}B=\prod\mathfrak{P}_i^{r_i}$ means that $\mathfrak{p}B$ contains $\prod\mathfrak{P}_i^{r_i}$ but does not contain the product when any $r_i$ is replaced with a smaller value. Since $\mathfrak{p}B$ corresponds to $(\overline{f})$ in $k[x]$ and that $f\equiv \prod g_i^{e_i}\ (\mathrm{mod}\ \mathfrak{p})$, we have $\overline{f}=\prod \overline{g}_i^{e_i}$, and thus $\mathfrak{p}B=\prod\mathfrak{P}_i^{e_i}$.\\
    Now, observe the residue field
    \begin{align*}
        B/\mathfrak{P}_i\ \cong\ \frac{B/\mathfrak{p}B}{\mathfrak{P}_i/\mathfrak{p}B}\ \cong\ 
        \frac{B/\mathfrak{p}B}{(g_i(\alpha))+\mathfrak{p}B}\ \cong\ 
        \frac{k[x]/(\overline{f})}{(\overline{g}_i)}\ =\ k[x]/(\overline{g}_i).
    \end{align*}
    Note that the leading coefficient of $f$ and thus $g_i$ is a unit, hence we have
    \begin{align*}
        f_i=[B/\mathfrak{p}_i:k]=[k[x]/(\overline{g}_i):k]=\mathrm{deg}\ \overline{g}_i=\mathrm{deg}\ g_i.
    \end{align*}
\end{proof}

\begin{example}{165}
    For a square-free integer $m\neq1$ factorise $p\in\mathbb{Z}$ in $\mathbb{Q}[\sqrt{m}]$. \textbf{Prop. 164} allows us to characterise the following three possibilities:
    \begin{itemize}
        \item[(i)] $(p)^e=\mathfrak{P}^2$: $(p)$ ramifies and $e=2,\ f=1,\ g=1$. 
        \item[(ii)] $(p)^e=\mathfrak{P}$: $(p)$ stays prime and $e=1,\ f=2,\ g=1$.
        \item[(iii)] $(p)^e=\mathfrak{P}_1\mathfrak{P}_2$: $(p)$ splits and $e=1,\ f=1,\ g=2$. 
    \end{itemize}
\end{example}

\newpage
\section{Exercises}
\begin{exe}{1. (Counterexamples to Dedekind Rings)}
    Let $K$ be a field.
    \begin{itemize}
        \item[(a)] Is $K[x,y]$ a Dedekind domain?
        \item[(b)] Let $A=K[x^2,x^3]$ be the subring in $K[x]$. Then
        \begin{itemize}[$\bullet$]
            \item $K[x]$ is a finitely generated $K[x^2]$-module, hence a Noetherian $K[x^2]$-module. Conclude that $A$ is Noetherian.
            \item Every prime ideal in $A$ is maximal, but $A$ is not a Dedekind domain.
        \end{itemize}
    \end{itemize}
\end{exe}
\begin{proof}
    (a) $K[x,y]$ is not a Dedekind, for the prime ideal $(x)$ is not maximal (since $K[x,y]/(x)\cong K[y]$ is an ID but not a field). Moreover, $(x)$ is strictly contained in the maximal ideal $(x,y)$, (where $K[x,y]/(x,y)\cong K$ is a field).\\
    (b) Note that $K[x^2]$ is obviously a Noetherian ring and that the map
    \begin{align*}
        K[x^2]\times K[x^2]\to K[x],\ (f(x^2),g(x^2))\mapsto f(x^2)+xg(x^2),
    \end{align*}
    is $K[x^2]$-linear and surjective. Hence $K[x]$ is a finitely generated, thus Noetherian $K[x^2]$-module. Moreover, $A=K[x^2,x^3]\subset K[x]$ is a $K[x^2]$-submodule, whence $A$ is also a Noetherian $K[x^2]$-module. Since $K[x^2]\subset K[x^2,x^3]=A$, it is also a Noetherian $A$-module, i.e. $A$ is a Noetherian ring.\\
    Note that the rings
    \begin{align*}
        K[x^2,x^3]\ \cong\ K[x,y]/(y^2-x^3)
    \end{align*}
    are isomorphic (even as a $K$-algebra). Now, consider the $K$-algebra homomorphism
    \begin{align*}
        \phi:K[x,y]\to K[t],\qquad x\mapsto t^2\ \text{and}\ y\mapsto t^3.
    \end{align*}
    Then the image is the subring $K[t^2,t^3]\subset K[t]$, which is an integral domain. By Isomorphism Theorem, $\mathrm{ker}\phi$ is a prime ideal. Since $K[x,y]/\mathrm{ker}\phi\cong K[t^2,t^3]$ is not a field, we deduce from \textbf{Ch1 \#1} that $\mathrm{ker}\phi$ is principal, hence it is generated by an irreducible polynomial. Since $y^2-x^3$ is irreducible, it must generate $\mathrm{ker}\phi$, whence the assertion.\\
    The only prime ideals of $K[x,y]$ which survive by taking the quotients are the maximal ideals of the form $(p,g)$ (cf. \textbf{Ch1 \#1}). Thus, any prime ideal is maximal. However, $K[x,y]/(y^2-x^3)$ is not integrally closed (cf. \textbf{Example 103c}), so it follows that $A$ is not Dedekind.
\end{proof}

\begin{exe}{2. (Valuation Rings)}
    \newline Let $A\subset K$ be a ring extension with $K$ a field. We call $A$ a \textbf{valuation ring} of $K$ if for each $0\neq x\in K$, either $x\in A$ or $x^{-1}\in A$.
    \begin{itemize}
        \item[(a)] Let $\nu:K^\times\to\mathbb{Z}$ be a valuation. Then $\nu(x^{-1})=-\nu(x)$. Conclude that a DVR is a valuation ring.
        \item[(b)] Let $A$ be a valuation ring, then
        \begin{itemize}[$\bullet$]
            \item $K=\mathrm{Quot}A$.
            \item $A$ is a local ring.
            \item $A$ is integrally closed in $K$, i.e. any valuation ring is normal.
            \item If $B$ is a ring such that $A\subset B\subset K$, then $B$ is also a valuation ring. [AM, 5.18]
        \end{itemize}
    \end{itemize}
\end{exe}
\begin{proof}
    (a) Since $\nu:K^\times\to\mathbb{Z}$ is a group homomorphism, it follows that
    \begin{align*}
        \nu(x^{-1})+\nu(x)=\nu(x^{-1}\cdot x)=\nu(1)=0,
    \end{align*}
    thus $\nu(x^{-1})=-\nu(x)$. Since any $A$:\ DVR is of the form
    \begin{align*}
        A\coloneqq\{x\in K=\mathrm{Quot}A:\nu(x)\geq0\}
    \end{align*}
    for some valuation $\nu$ on $K$. Thus, either $\nu(x)\geq0$ or $\nu(x^{-1})\geq0$, i.e. either $x\in A$ or $x^{-1}\in A$, whence $A$ is a valuation ring.\\
    (b)(i) Since $K$ is a field containing $A$, we have $\mathrm{Quot}A\subset K$. Conversely, for $x\in K^\times$, if $x\notin A$, then $a=x^{-1}\in A$ since $A$ is a valuation ring, whence $x=\frac{1}{a}\in\mathrm{Quot}A$. Thus $K\subset\mathrm{Quot}A$.\\
    (ii) \textsc{claim.} $\mathfrak{m}=A\setminus A^\times$. As $A$ is a valuation ring, $x\in\mathfrak{m}\ \Leftrightarrow$ either $x=0$ or $x^{-1}\notin A$. \\
    Now, it suffices to show that $\mathfrak{m}$ is an ideal of $A$.
    \begin{itemize}
        \item \textbf{($A\mathfrak{m}\subset\mathfrak{m}$).} Let $a\in A,\ x\in\mathfrak{m}$. If $ax=0$, then clearly $ax\in\mathfrak{m}$. If $ax\neq0$, assume $(ax)^{-1}\in A$. Then $x^{-1}=(ax)^{-1}\cdot a\in A$, whence $x\in\mathfrak{m}$ is a unit, i.e. $\mathfrak{m}=B.\ (\to\gets)$. Thus $(ax)^{-1}\notin A$. From the previous claim, $ax\in\mathfrak{m}$.
        \item \textbf{(Closed under $+$).} Let $0\neq x,y\in\mathfrak{m}$, then $x^{-1},y^{-1}\notin A$ by the previous claim. Consider $xy^{-1}\in K=\mathrm{Quot}A$. Since $A$ is a valuation ring of $K$, either $xy^{-1}\in A$ or $yx^{-1}=(xy^{-1})^{-1}\in A$. If $xy^{-1}\in A$, then $x+y=(1+xy^{-1})\cdot y\in A\mathfrak{m}\subset\mathfrak{m}$. Similarly for the case $yx^{-1}\in A$. Thus, $x+y\in\mathfrak{m}$.
    \end{itemize}
    Hence $\mathfrak{m}$ is the unique maximal ideal of $A$, i.e. $A$ is a local ring.\\
    (iii) Suppose $0\neq x\in K$ is integral over $A$ and assume $x\notin A$. Since $A$ is a valuation ring, we have $x^{-1}\in A$. By assumption, $x$ satisfies
    \begin{align*}
        x^n+a_1x^{n-1}+\cdots+a_n=0,\qquad (a_i\in A).
    \end{align*}
    $\Rightarrow 1+\sum_{i=1}^n\frac{a_i}{x^i}=0.\ \Rightarrow 1=-\sum_{i=1}^n\frac{a_i}{x^i}$ in $A$. Thus $x=x\cdot1=-\sum_{i=1}^n\frac{a_i}{x^{i-1}}\in A.\ (\to\gets)$. Hence $x\in A$ and it follows that $A$ is integrally closed in $K$. \\
    (iv) Given $0\neq x\in K$. If $x^{-1}\notin B,\ \because A\subset B.\ \therefore x^{-1}\in A.\ \because A$ is a valuation ring. $\therefore x\in A\subset B$. If $x\notin B$, a similar argument shows that $x^{-1}\in A\subset B$. Hence $B$ is a valuation ring in $K$.
\end{proof}

\begin{exe}{3. (The Integral Closure of a Ring)}
    Let $A\subset K$ be a ring extension with $K$ a field. The integral closure $\overline{A}$ of $A$ in $K$ is the intersection of all valuation rings of $K$ containing $A$.
\end{exe}
\begin{proof}
    Let $L$ be the algebraic closure of $K$ and we consider the set 
    \begin{align*}
        \Sigma\coloneqq\{(A,f):A\subsetneq K\ \text{subring},\ f:A\to L\ \text{ring homomorphism}\}
    \end{align*}
    partially ordered by
    \begin{align*}
        (A_1,f_1)\leq(A_2,f_2)\qquad\text{if}\qquad A_1\subset A_2\qquad \text{and}\qquad f_2|_{A_1}=f_1.
    \end{align*}
    It has a maximal element $(B,g)$ by Zorn's Lemma and we use the fact from [AM, 5.21] that $B$ is a valuation ring in $K$. We prove that
    \begin{align*}
        \overline{A}=\bigcap\{B:\ A\subset B\subset K\ \text{valuation ring}\}.
    \end{align*}
    "$\subset$"\ Let $B$ be a valuation with $A\subset B\subset K$. Let $x\in\overline{A}$, then $x$ is integral over $A\subset B$. Since $B$ is a valuation ring, by \textbf{\#2(b.iii)} $A=\overline{A}$ in $K.\ \Rightarrow x\in B.\ \Rightarrow \overline{A}\subset B$. It follows that $\overline{A}\subset\bigcap B$.\\
    "$\supset$"\ Let $x\in K\setminus\overline{A}$, we claim: there is a valuation ring $A\subset B\subset K$ with $x\notin B$.\\
    $\because x\notin\overline{A}.\ \therefore x\notin A'\coloneqq A[x^{-1}]\subset K$. If $x\in A'$, then $x\coloneqq\sum_{i=0}^na_ix^{-i},\ (a_i\in A)$. Multiply by $x^n$, we have 
    \begin{align*}
        x^{n+1}-a_0x^n-\cdots-a_{n-1}x-a_n=0.
    \end{align*}
    Thus $x\in K$ is integral over $A$ and it follows that $x\in\overline{A}.\ (\to\gets)$ \\ $\because x^{-1}\in A'$ but $x\notin A'.\ \therefore x^{-1}$ is not a unit in $A'.\ \Rightarrow x^{-1}\in\mathfrak{m}'$ for some maximal ideal $\mathfrak{m}'\triangleleft A'$. Consider the map
    \begin{align*}
        f:A'\twoheadrightarrow A'/\mathfrak{m}'\coloneqq k'\xhookrightarrow{} k'.
    \end{align*}
    By the fact we noted above, there is a valuation ring $A\subset B\subset K$ with $(A,f|_A)\leq(B,g)$, i.e. $g|_A=f|_A.\ \because x^{-1}\in\mathfrak{m}'.\ \therefore f(x^{-1})=0$. If $x\in B,\ \because x^{-1}\in A'\subset B.\ \therefore x\cdot x^{-1}=1$ in $B$. Now, 
    \begin{align*}
        1=g(1)=g(x\cdot x^{-1})=g(x)g(x^{-1})=g(x)f(x^{-1})=g(x)\cdot0=0.\qquad(\to\gets)
    \end{align*}
    Thus, $x\notin B$. It follows that $\bigcap B\subset\overline{A}$.
\end{proof}

\chapter{Affine Schemes}

\section{Spectra of Rings}
\begin{motivation}{166}
    Consider $A=\mathbb{C}[x]$ as a ring of functions $\mathbb{C}\to\mathbb{C}$. Given a point $P\in\mathbb{C}$ we take the associate valuation on $\mathrm{ord}_P:\mathbb{C}(x)^\times\to\mathbb{Z}$, (cf. \textbf{Example 150}). The induced valuation ring is $\mathbb{C}[x]_{(x-P)}$. One should therefore think of a point $P$ in terms of the prime ideal $(x-P)$ it defines.
\end{motivation}

\begin{definition}{167. (Spectra)}
    The \textbf{(prime) spectrum} of a ring $A$ is defined by 
    \begin{align*}
        \mathrm{Spec}A\ \coloneqq\ \{\mathfrak{p}\triangleleft A:\ \text{prime ideal}\}.
    \end{align*}
    For each ideal $\mathfrak{a}\triangleleft A$, we define
    \begin{align*}
        \mathcal{V}(\mathfrak{a})\ \coloneqq\ \{\mathfrak{p}\in\mathrm{Spec}A:\mathfrak{a}\subset\mathfrak{p}\}
        \ \cong\ \mathrm{Spec}(A/\mathfrak{a}).
    \end{align*}
    Then 
    \begin{itemize}
        \item $\mathcal{V}(0)=\mathrm{Spec}A$ and $\mathcal{V}(A)=\varnothing$.
        \item If $\{\mathfrak{a}_i\}_{i\in I}$ is any family of ideals of $A$, then 
        \begin{align*}
            \mathcal{V}(\sum_{i\in I}\mathfrak{a}_i)=\bigcap_{i\in I}\mathfrak{a}_i.
        \end{align*}
        \item For any two ideals $\mathfrak{a},\ \mathfrak{b}\triangleleft A$, we have
        \begin{align*}
            \mathcal{V}(\mathfrak{a}\cap\mathfrak{b})=\mathcal{V}(\mathfrak{a}\cdot\mathfrak{b})=\mathcal{V}(\mathfrak{a})\cup\mathcal{V}(\mathfrak{b}).
        \end{align*}
    \end{itemize}
    In particular, the sets $\mathcal{V}(\mathfrak{a})$ define the closed sets of the so-called \textbf{Zariski topology} of $\mathrm{Spec}A$.
\end{definition}
\newpage
As usual in this course, we assume that $A$ is an integral domain with quotient field $K$, though everything remains true for arbitrary rings using a suitable notion of localisation.
\begin{example}{168}
    \begin{itemize}
        \item[(a)] $\mathrm{Spec}\ \mathbb{C}[x]=\{(0)\}\cup\{(x-P):P\in\mathbb{C}\}$ for the nontrivial prime ideals are maximal, (note: $\mathbb{C}[x]$ is a UFD). They correspond to the points in $\mathbb{C}$ and are called \textbf{closed points} since their closure is $\overline{\{(x-P)\}}=\{(x-P)\}$.\\
        Moreover, $(0)$ is called the \textbf{generic point} \textemdash its closure is all of $\mathrm{Spec}\ \mathbb{C}[x]$, i.e. $(0)$ is dense in $\mathrm{Spec}\ \mathbb{C}[x]$.\\
        The nonempty open sets of $\mathrm{Spec}\ \mathbb{C}[x]$ can be identified with the subsets of $\mathbb{C}$ having finite complement. In particular, the Zariski topology is co-finite and much coarser than the Euclidean topology.
        \item[(b)] Similarly, $\mathrm{Spec}\ \mathbb{Z}$ consists of the generic point $(0)$ and the closed points $(p)$, where $p\in\mathbb{Z}$ prime. As in the previous discussion, the nontrivial open sets are the subsets of $\mathbb{Z}$ with finite complement.
    \end{itemize}
\end{example}

\begin{remark}{169}
    We have $\overline{\{\mathfrak{p}\}}=\mathcal{V}(\mathfrak{p})$ for all $\mathfrak{p}\in\mathrm{Spec}A$. In particular, $\mathfrak{p}$ is a closed point iff $\mathfrak{p}$ is maximal. [Exercise]
\end{remark}

\begin{proposition}{170. (Radical Ideals)}
    For an ideal $\mathfrak{a}\triangleleft A$, we define its \textbf{radical} by 
    \begin{align*}
        \sqrt{\mathfrak{a}}\ \coloneqq\ \{a\in A:a^n\in\mathfrak{a}\ \text{for some}\ n\in\mathbb{N}_{\geq1}\}.
    \end{align*}
    This is an ideal which contains $\mathfrak{a}$. Let $X\subset\mathrm{Spec} A$. If we define $\mathcal{I}(X)\coloneqq\bigcap_{\mathfrak{p}\in X}\mathfrak{p}$, then 
    \begin{align*}
        \mathcal{V(I}(X))=\overline{X}\qquad\text{and}\qquad \mathcal{I(V}(\mathfrak{a}))=\sqrt{\mathfrak{a}}.
    \end{align*}
\end{proposition}
\begin{proof}
    We proof the second part. Note that $\mathcal{I(V}(\mathfrak{a}))=
    \bigcap_{\mathfrak{p}\in\mathcal{V}(\mathfrak{a})}\mathfrak{p}$.
    \\ "$\supset$"\ Let $a\in\sqrt{\mathfrak{a}}$, then $a^n\in\mathfrak{a}$ for some $n\in\mathbb{N}$. For any $\mathfrak{p}\in\mathcal{V}(\mathfrak{a}),\ \because a^n\in\mathfrak{a}\subset\mathfrak{p}.\ \therefore$ Either $a\in\mathfrak{p}$ or $a^{n-1}\in\mathfrak{p}$, it follows that $a\in\mathfrak{p}$. Hence, $a\in\mathcal{I(V}(\mathfrak{a}))$.\\
    "$\supset$"\ Assume $a\in\mathcal{I(V}(\mathfrak{a}))$ but $a\notin\sqrt{\mathfrak{a}}$. Indeed, if $a^n\notin\mathfrak{a},\ \forall n\in\mathbb{N}$, then $S\coloneqq\{f^i:i\in\mathbb{N}\}$ is a multiplicative system disjoint from $\mathfrak{a}.\ \because \mathrm{Spec}\ S^{-1}A\cong \{\mathfrak{p}\in\mathrm{Spec}A:\mathfrak{p}\cap S=\varnothing\}.\ \therefore \exists\ \mathfrak{p}\in\mathrm{Spec}A$ such that $\mathfrak{p}\supset\mathfrak{a}$ and $\mathfrak{p}\cap S=\varnothing$. Thus, $a\notin\bigcap_{\mathfrak{p}\supset\mathfrak{a}}\mathfrak{p}.\ (\to\gets)$
\end{proof}
\newpage
\begin{definition}{171. Structure Sheaf \& Affine Scheme}
    \begin{itemize}
        \item For any open set $\varnothing\neq U\subset\mathrm{Spec}A$, we let $\mathcal{O}_X(U)\coloneqq\bigcap_{\mathfrak{p}\in U}A_\mathfrak{p}\subset K$ and $\mathcal{O}_X(\varnothing)\coloneqq\{0\}$. 
        \item Further, for open sets $V\subset U$, the inclusion $\rho_{UV}:\mathcal{O}(U)\to\mathcal{O}(V)$ defines the so-called \textbf{restriction morphism} satisfying $\rho_{UU}=1_{\mathcal{O}(U)}$ and $\rho_{VW}\circ\rho_{UV}= \rho_{UW}$. We write $s|_V\coloneqq\rho_{UV}(s)$.
        \item A \textbf{presheaf of rings} is an assignment $U\mapsto\mathcal{O} (U)$ together with restriction morphisms. The elements $s\in\mathcal{O}_X(U)$ are called \textbf{sections}.
        \item In fact, $\mathcal{O}$ defines a \textbf{sheaf}, which is called the \textbf{structure sheaf}; the pair $(\mathrm{Spec}A,\mathcal{O})$ is called an \textbf{affine scheme}.
        \item Given $s\in\mathcal{O}_X(U)$ and a point $\mathfrak{p}\in U\subset X$, we define the "value" of $s(\mathfrak{p})$ by 
        \begin{align*}
            s(\mathfrak{p})\coloneqq\overline{s}\in A_\mathfrak{p}/\mathfrak{p}A_\mathfrak{p}=k(A_\mathfrak{p}).
        \end{align*}
        Of course there values $s(\mathfrak{p})$ live a priori in different (residue) fields which is why we call $s$ a section, not a function.
        For instance, the value of $f\in\mathbb{C}[x]$ at the generic point $(0)$ is $f$ "evaluated" at the unknown $x$. Namely, $f((0))=\overline{f}=f(x)\in\mathrm{Quot}(\mathbb{C}[x]/(0))$. Since $\mathfrak{p}=\{a\in A:a(\mathfrak{p})=0\}$, we can recover any point from the ideal of "functions" vanishing at it.
    \end{itemize}
\end{definition}

\begin{remark}{172}
    \begin{itemize}
        \item[(a)] The construction above makes sense for any ring. If $A$ is not an ID, this requires a more general localisation procedure, cf. [G, Ch6]. However, IDs are sufficient for our purposes.
        \item[(b)] We have $\mathcal{O}(\mathrm{Spec}A)=A$. [Exercise]
        \item[(c)] For $s\in\mathcal{O}_X(U)\subset K$, we consider its \textbf{zero locus} $\mathcal{V}(s)=\{\mathfrak{p}\in\mathrm{Spec}A:s(\mathfrak{p})=0\}$. If $s\in A$, then $\mathcal{V}(s)=\mathcal{V}((s)))=\{\mathfrak{p}\in\mathrm{Spec}A: \mathfrak{p}\supset(s)\}$. For example, $\mathcal{V}((s))=s^{-1}(0)$ in $\mathrm{Spec}\ \mathbb{C}[x]$ for $s\in\mathbb{C}[x]$.
    \end{itemize}
\end{remark}
\newpage
\begin{definition}{173. (Scheme Morphism)}
    \newline Let $X=\mathrm{Spec}A$ and $Y=\mathrm{Spec}B$. A \textbf{scheme morphism} $F:(Y,\mathcal{O}_Y)\to (X,\mathcal{O}_X)$ consists of a continuous map $f:\mathrm{Spec}B\to\mathrm{Spec}A$ and a family of ring homomorphisms $f_U^\ast:\mathcal{O}_X(U)\to\mathcal{O}_Y(f^{-1}(U))$, where $U\subset X$ open, such that 
    \begin{itemize}
        \item For any open sets $V\subset U$, we have $\rho_{f^{-1}(U)f^{-1}(V)} \circ f_U^\ast=f_V^\ast\circ\rho_{UV}$.
        \begin{center}
            \begin{tikzcd}
                \mathcal{O}_X(U)\rar["f_U^\ast"] \dar["\rho_{UV}"] 
                \drar[phantom, "\curvearrowright"] 
                & \mathcal{O}_Y(f^{-1}(U))\dar["\rho_{f^{-1}(U)f^{-1}(V)}"] \\
                \mathcal{O}_X(V)\rar[swap, "f_V^\ast"] 
                & \mathcal{O}_Y(f^{-1}(U))
            \end{tikzcd}
        \end{center}
        \item For any $\mathfrak{q}\in f^{-1}(U)$ and $s\in\mathcal{O}(U)$, we have
        \begin{align*}
            s(f(\mathfrak{q}))=0\qquad\Longrightarrow\qquad f_U^\ast(s)(\mathfrak{q})=0,
        \end{align*}
        that is, $f^{-1}(\mathcal{V}(s))\subset\mathcal{V}(f^\ast_U(s))$.
    \end{itemize}
    We often simply write $F:\mathrm{Spec}B\to\mathrm{Spec}A$.\\ In particular, a \textbf{scheme isomorphism} is a scheme morphism with a two-sided inverse.
\end{definition}

\begin{example}{174}
    A ring homomorphism $\phi:A\to B$ induces a scheme morphism\\ $F:(\mathrm{Spec}B,\ \mathcal{O}_Y)\to (\mathrm{Spec}A,\ \mathcal{O}_X)$ via
    \begin{align*}
        f(\mathfrak{q})\coloneqq\phi^{-1}(\mathfrak{q}) \qquad\text{and}\qquad f_U^\ast(s)(\mathfrak{q})\coloneqq\phi_{\mathfrak{q}}\circ s(\phi^{-1} (\mathfrak{q})).
    \end{align*}
    Here, $\phi_{\mathfrak{q}}$ is the naturally induced map $\phi_\mathfrak{q}:A_{\phi^{-1} (\mathfrak{q})}\to B_\mathfrak{q},\ \frac{a}{b}\mapsto\frac{\phi(a)}{\phi(b)}$.
\end{example}

\begin{example}{174*}
    We obtain a scheme morphism $F=(f,f^\ast):(Y,\mathcal{O}_Y)\to(X,\mathcal{O}_X)$, where $X=\mathrm{Spec}A$ and $Y=\mathrm{Spec}B$, from a ring morphism $\phi:A\to B$ as follows:
    \begin{itemize}
        \item For $f:Y\to X$ we define $f(\mathfrak{q})\coloneqq\phi^{-1}(\mathfrak{q})$. It is easy to see that $f^{-1}(\mathcal{V}(\mathfrak{a}))=\mathcal{V}(\phi(\mathfrak{a}))$, hence $f$ is continuous.
        \item To define $f^\ast$, we note that for every $\mathfrak{p}\in X$, there is a unique map
        \begin{center}
            \begin{tikzcd}
                A \arrow[r, "\phi"] \arrow[d] 
                & B \arrow[d]\\
                A_{f(\mathfrak{q})} \arrow[r, "\hat{\phi}_\mathfrak{q}"]
                & B_\mathfrak{q}\ ,
            \end{tikzcd}\qquad $f_{\mathfrak{q}}^\ast(f(\mathfrak{q})^e)=\mathfrak{q}^e$
        \end{center}
        by the universal property of localisation, (cf. AM, 3.1). 
        \item In fact, $\hat{\phi}_\mathfrak{q}(\frac{g}{h})= \frac{\phi(g)}{\phi(h)}$. Then define 
        \begin{align*}
            f(x)= \left\{\begin{array}{lc} 0 &\text{if}\ f^{-1}(U)=\varnothing \\ \hat{\phi}_\mathfrak{q}|_{\mathcal{O}_X(U)} &\text{if}\ \mathfrak{q}\in f^{-1}(U) \end{array}\right..
        \end{align*}
        This is compatible with restrictions by design. 
        \item Furthermore, let $s=\frac{g}{h}\in A_{f(\mathfrak{q})}$. If $s(f(\mathfrak{q}))=0$, then $\frac{g}{h}\in f(\mathfrak{q})A_{f(\mathfrak{q})}$, that is, $g\in f(\mathfrak{q})=\phi^{-1}(\mathfrak{q})$. It follows that $\phi(g)\in\mathfrak{q}$, i.e. $f_U^\ast(s)=\frac{\phi(g)}{\phi(h)} \in\mathfrak{q}A_\mathfrak{q}$, whence $f_U^\ast(s)(\mathfrak{q})=0$.
    \end{itemize}
\end{example}

\begin{theorem}{175}
    Every scheme morphism is of the form $(\mathrm{Spec}B,\mathcal{O}_Y)\to (\mathrm{Spec}A,\mathcal{O}_X)$.
\end{theorem}
\begin{proof}
    \textbf{Sketch}. Write $X=\mathrm{Spec}A,\ Y=\mathrm{Spec}B,\ K=\mathrm{Quot}A$ and $L=\mathrm{Quot}B$ and we denote by $\mathrm{Hom}(X,Y)$ the set of scheme morphisms $F=(f,f^\ast):Y\to X$. Now, consider the following map
    \begin{align*}
        \Phi:\mathrm{Hom}(X,Y)\longrightarrow\mathrm{Hom}(A,B),\ F\longmapsto\phi\coloneqq f_X^\ast:A\to B.
    \end{align*}
    Now that:
    \begin{itemize}
        \item \textbf{$\Phi$ is surjective.} This follows from an example in the lecture notes where we constructed a scheme morphism $F=(f,f^\ast): \mathrm{Spec}B\to\mathrm{Spec}A$ for every $\phi:A\to B$ such that $f_X^\ast=\phi$, (cf. \textbf{Example 174}).
        \item \textbf{$\Phi$ is injective.} This requires some sheaf theory. Namely, a sheaf morphism $\{f_U^\ast:\mathcal{O}_X(U)\to \mathcal{O}_{Y}(f^{-1}(U))\}$ gives rise to a ring morphism $f_{\mathfrak{q}}^\ast:A_{f(\mathfrak{q})}\to B_{\mathfrak{q}}$ such that $f_{\mathfrak{q}}^\ast(s)=f_U^\ast(s)$ whenever $s\in\mathcal{O}_X(U)$ and $f(\mathfrak{q})\in U$.\\
        The second condition in the definition of a scheme morphisms ensures that $f_{\mathfrak{q}}^\ast$ is a local morphism, that is, if $\mathfrak{p}^e$ and $\mathfrak{q}^e$ are the maximal ideals of $A_{f(\mathfrak{q})}$ and $B_\mathfrak{q}$, respectively, then $f_\mathfrak{q}^\ast(\mathfrak{p}^e) =\mathfrak{q}^e$. Indeed, if $s\in\mathfrak{a}\cap\mathcal{O}_X(U)$ for some $U\ni\mathfrak{a}$, then $s(f(\mathfrak{q}))=0.\\ \Rightarrow f_U^\ast(s)(\mathfrak{q})=f_\mathfrak{q}^\ast(s)(\mathfrak{q})=0$, i.e. $f_\mathfrak{q}^\ast(s)\in\mathfrak{q}$. It follows that $\mathfrak{p}^e \subset(f_\mathfrak{q}^\ast)^{-1}(\mathfrak{q}^e)$ and since $\mathfrak{p}$ is maximal in $A_\mathfrak{p}$, we have the equality.\\ Hence we obtain the commutative diagram
        \begin{center}
            \begin{tikzcd}
                A \arrow[r, "\phi"] \arrow[d] 
                & B \arrow[d]\\
                A_{f(\mathfrak{q})} \arrow[r, "f_{\mathfrak{q}}^\ast"]
                & B_\mathfrak{q}\ ,
            \end{tikzcd}\qquad $f_{\mathfrak{q}}^\ast(f(\mathfrak{q})^e)=\mathfrak{q}^e$
        \end{center}
        and by the universal property of  the localisation, we have $f_\mathfrak{q}^\ast(\frac{g}{h})=\hat{\phi}_\mathfrak{q}(\frac{g}{h})=\frac{\phi(g)}{\phi(h)}$. This shows that $f_\mathfrak{q}^\ast$ is determined by $\phi$ and thus $\{f_U^\ast\}$ is determined by $\phi$ since the $f_\mathfrak{q}^\ast$'s determine $\{f_U^\ast\}$, (this is the local native of schemes). \\
        Finally, consider that $f(\mathfrak{q})\in\mathrm{Spec}A$. This is determined by 
        \begin{align*}
            f(\mathfrak{q}) &=\{a\in A:a(f(\mathfrak{q}))=0\} \\
            &=\{a\in A:f_\mathfrak{q}^\ast(a)(\mathfrak{q})=0\} \\
            &=\{a\in A:f_X^\ast(a)(\mathfrak{q})=0\} \qquad\text{since}\ a\in\mathcal{O}_X(X)=A \\
            &=\{a\in A:\phi(a)\in\mathfrak{q}=0\} \\ &=\phi^{-1}(\mathfrak{q}).
        \end{align*}
        Hence, $\phi$ also determines $f$.
    \end{itemize}
\end{proof}

\begin{remark}{176}
    Everything we stated so far remains true for arbitrary commutative rings. The previous theorem then states that the category of affine schemes is equivalent to the category of commutative rings generalising the equivalence of categories between affine varieties over an algebraically closed field $k$ and the category of finitely generated $k$-algebras which are integral domains.\\ The additional difficulty of working with general rings which are not necessarily integral domains comes from nilpotent elements. For instance, compare $\mathrm{Spec}\ \mathbb{C}[x]/(x)$ (one closed point, "functions" are $a+b\overline{x},\ (a,b\in\mathbb{C})$, but $\overline{x}$ is a nilpotent function!) Though the geometric interpretation of nilpotent functions is difficult at first glance they prove to be extremely important from a geometric point of view as well.
\end{remark}

\begin{example}{177}
    \begin{itemize}
        \item \textbf{Fields.} For $X=\mathrm{Spec}\ k,\ (k:\text{field})$, consists of one (closed) point. Further, $\mathcal{O}_X(\varnothing)=0$ and $\mathcal{O}_X(\mathrm{Spec}\ k)=k$. Hence, the spectra of distinct fields are homeomorphic topological spaces, but not isomorphic as affine schemes.
        \item \textbf{DVRs.} For $X=\mathrm{Spec}A,\ (A:\ \text{DVR})$, consists of two points: The generic point $(0)$ and the closed point given by its maximal ideal $\mathfrak{m}$. Further, $\mathcal{O}((0))=\mathrm{Quot}A$ and $\mathcal{O}(\mathrm{Spec}A)=A$. We thus think of $X$ as an infinitesimal neighbourhood of $\mathfrak{m}$.
        \item \textbf{Dedekind Rings.} For $X=\mathrm{Spec}\ A,\ (A:\ \text{Dedekind ring})$, consists of the closed points given by the nontrivial prime ideals $\mathfrak{p}$ and the generic point $(0)$. The localisation $A_\mathfrak{p}$ gives a DVR with inclusion $A\xhookrightarrow{} A_\mathfrak{p}$, whence 
        \begin{align*}
            F:X_\mathfrak{p}\coloneqq \mathrm{Spec}A_\mathfrak{p}\to X.    
        \end{align*}
        Thinking of $A_\mathfrak{p}$ as "functions", then 
        \begin{align*}
            A_\mathfrak{p}\ =\ \{\tfrac{f}{g}:g(\mathfrak{p})\neq0\}
        \end{align*}
        is the \textbf{function germ} of $\mathrm{Spec}A$ at $\mathfrak{p}$. If $(0)\neq\mathfrak{q}\in X$, then $\mathfrak{q}+\mathfrak{p}=A$, so we find $g\in A$ with 
        \begin{align*}
            g\equiv1\ (\mathrm{mod}\ \mathfrak{p})\qquad\text{and}\qquad
            g\equiv0\ (\mathrm{mod}\ \mathfrak{q})
        \end{align*}
        by \textbf{CRT}. Hence $\frac{1}{g}\in A_\mathfrak{p}$, but $\frac{1}{g}$ is not defined on any open set containing $\mathfrak{q}$, thus $A_\mathfrak{p}$ is called the \textbf{local ring} of $\mathfrak{p}$.
        \item \textbf{Polynomial Rings.} $X\coloneqq\mathrm{Spec}A[x_1,\cdots,x_n]$ is called the \textbf{affine $n$-space} over $A$ and denoted by $\mathbb{A}_A^n$. A prominent example is $A=k$, ($k$: algebraically closed field). Here, maximal ideals correspond to points of $k^n$. (cf. \textbf{Cor. 196}).
    \end{itemize}
\end{example}

\begin{remark}{178}
    Recall that any DVR arises as the localisation of a Dedekind ring by \textbf{(121)}.
\end{remark}

\begin{nt}{179. (Geometric View on Ring Extensions)}
    We consider a Dedekind ring $A$ with quotient field $K$ and a finite separable extension $L/K$. Let $\overline{A}$ be the integral closure of $A$ in $L$. The integral extension of Dedekind rings $A\xhookrightarrow{} \overline{A}$ leads to a scheme morphism 
    \begin{align*}
        F=(f,f^\ast):\ X=\mathrm{Spec}\overline{A}\to Y=\mathrm{Spec}A
    \end{align*}
    with $f$ surjective, (cf. \textbf{Lying Over Thm. 190}).\\
    If $\mathfrak{p}\in\mathrm{Spec}A$, then $\mathfrak{p}^e\coloneqq \mathfrak{P}_1^{e_1}\cdots\mathfrak{P}_r^{e_r}$ with $\mathfrak{p}\coloneqq \mathfrak{P}_i\cap A=f(\mathfrak{P}_i)$. If the quotient fields $A/\mathfrak{p}$ are algebraically closed, then $f_i=1$. Hence $\sum e_i=[L:K]$, and $\mathfrak{p}$ ramifies iff the preimages $\mathfrak{p}^c$ under $f$ coincide. Topologically speaking, the map $f$ is a \textbf{ramified covering}.
\end{nt}

\begin{example}{180}
    This requires some basic knowledge of Riemann surfaces.\\
    Let $X$ and $Y$ be compact connected Riemann surfaces, and let $F:X\to Y$ be a nonconstant holomorphic mapping. Write $\mathcal{M}(X)$ and $\mathcal{M}(Y)$ to be the field of meromorphic functions on $X$ and $Y$, respectively. The map $f\mapsto F\circ f$ defines an inclusion $\mathcal{M}(Y)\xhookrightarrow{} \mathcal{M}(X)$ and thus a finite extension of fields of degree $m$. The map is $m$-to-$1$ except at a finite number of branch points.\\ Let $P\in Y$ and $\mathcal{O}_P$ be the set of meromorphic functions on $Y$ that are holomorphic at $P$. This is the DVR attached to the discrete valuation $\mathrm{ord}_P$, (cf. \textbf{Example 50}). Its maximal ideal $\mathfrak{q}$ is the set of meromorphic functions on $Y$ that vanish at $P$. \\ Let $B$ be the integral closure of $\mathcal{O}_P$ in $\mathcal{M}(X),\ F^{-1}(P)\coloneqq\{Q_i:1\leq i\leq g\}$, and $e_i$ be the number of sheets of $X$ over $Y$ that coincide at $Q_i$. Then 
    \begin{align*}
        \mathfrak{q}B=\prod_{i=1}^g\mathfrak{Q}_i^{e_i},\ \text{where}\ \mathfrak{Q}_i\coloneqq\{f\in B:f(Q_i)=0\}\triangleleft B\ \text{prime ideal}.
    \end{align*}
\end{example}
\newpage
\begin{remark}{181}
    In your Topology course you have seen unramified coverings, for instance the universal covering of a sufficiently good topological space. This has an algebraic counterpart.\\
    In the situation of the previous note, we assume additionally that $L/K$ is Galois with the Galois group $\mathrm{Gal}(L/K)$. Then any $\sigma\in\mathrm{Gal}(L/K)$ fixes $K$ and in particular $A$, and thus maps $\overline{A}$ to $\overline{A}$. (Apply $\sigma$ to an integral equation $f\in A[x]$ satisfied by any element $x\in\overline{A}$. Then $\sigma(f)\in A[x]$ will be an integral equation for $\sigma(x)$ over $A$). Hence $\sigma$ induces a scheme morphism $\mathrm{Spec}\overline{A}\to\mathrm{Spec} \overline{A}$, also denoted by $\sigma$, and we get a commutative diagram (with $F$ the scheme morphism from the previous note):
    \begin{center}
        \begin{tikzcd}
            \mathrm{Spec}\overline{A} \arrow[rd, "F"] \arrow[r, "\sigma"] 
            & \mathrm{Spec}\overline{A} \arrow[d, "F"] \\
            & \mathrm{Spec}A
        \end{tikzcd}
    \end{center}
    The transformation $\sigma$ is called a \textbf{Deck transformation} of the ramified covering $\mathrm{Spec}\widetilde{A}\to\mathrm{Spec}A$ corresponding to a field extension $\widetilde{K}/K$, (cf. \textbf{Cor. 277}). This scheme morphism is called the \textbf{universal covering}. One then defines the \textbf{fundamental group} of the scheme $\mathrm{Spec}A$ by $\pi_1(Y)\coloneqq\mathrm{Gal}(\widetilde{K}/K)$.
\end{remark}

\newpage

\section{Dimension}
As Figure suggests we should think of the spectrum of a Dedekind ring as a smooth, $1$-dimensional object. We set out to make these notions rigourous.
\begin{definition}{182. (Height \& Dimension)}
    Let $\mathfrak{p}\in\mathrm{Spec}A$. \\ Its \textbf{height} (or \textbf{codimension}) is defined by
    \begin{align*}
        \mathrm{height}\ \mathfrak{p}\coloneqq\sup\limits_{r\geq0} \{\mathfrak{p}_0\subsetneq\mathfrak{p}_1\subsetneq\cdots\subsetneq\mathfrak{p}_r=\mathfrak{p}:\mathfrak{p}_i\in\mathrm{Spec}A\}.
    \end{align*}
    The \textbf{(Krull) dimension} $\mathrm{dim}A$ of $A$ is 
    \begin{align*}
        \mathrm{dim}\ A\coloneqq\sup\limits_{\mathfrak{p}\in\mathrm{Spec}A}\mathrm{height}\ \mathfrak{p}.
    \end{align*}
\end{definition}

\begin{remark}{183}
    Since a maximal chain necessarily ends with a maximal ideal, it sufices to take the supremum over heights of maximal ideals:
    \begin{align*}
        \mathrm{dim}\ A=\sup\limits_{\mathfrak{m}\subset A\ \text{maximal}}\mathrm{height}\ \mathfrak{m}.
    \end{align*}
\end{remark}

\begin{example}{184}
    \begin{itemize}
        \item \textbf{Fields.} If $k$ is a field, then $\mathrm{dim}\ k=0$.
        \item \textbf{PIDs.} If $A$ is a PID, then $\mathrm{dim}\ A=1$.
        \item \textbf{Dedekind Rings.} Since any prime is maximal in a Dedekind ring, we have $\mathrm{dim}\ A=1$.
        \item \textbf{Local Rings.} By \textbf{Rmk. 183}, we have
        \begin{align*}
            \mathrm{dim}\ A_\mathfrak{p}=\mathrm{height}\ \mathfrak{p}A_\mathfrak{p}=\mathrm{height}\ \mathfrak{p}.
        \end{align*}
        \item \textbf{Noetherian Rings.} Finite dimension is not related to the Noetherian property. In [R, \S9.4(3)] an example of a Noetherian ring with $\mathrm{dim}\ A=\infty$ is discussed. Conversely, the ring $A\coloneqq k[x_1,x_2,\cdots]/\langle x_ix_j:1\leq i\leq j<\infty\rangle$ is not a Noetherian, but $(\overline{x_1},\cdots,\overline{x_n})$ is the only prime ideal, whence $\mathrm{dim}\ A=0$. 
    \end{itemize}
\end{example}

\begin{remark}{185}
    Since $\mathrm{dim}\ \mathbb{Z}=1$ but $\mathrm{dim}\ \mathbb{Q}=0$, we cannot conclude that $\mathrm{dim}\ A\leq\mathrm{dim}\ B$ if $A\subset B$. For rings the notion of dimension can thus be quite pathological.
\end{remark}

\begin{proposition}{186}
    The dimension is local in the sense that 
    \begin{align*}
        \mathrm{dim}\ A=\sup\limits_{\mathfrak{p}\in\mathrm{Spec}A}\mathrm{dim}\ A_\mathfrak{p}=\sup\limits_{\mathfrak{m}\subset A\ \text{max.}}\mathrm{dim}\ A_\mathfrak{m}.
    \end{align*}
\end{proposition}
\begin{proof}
    For any prime ideal $\mathfrak{p}\triangleleft A,\ A_\mathfrak{p}$ is local with maximal ideal $\mathfrak{p}A_\mathfrak{p}$, and from \textbf{Ex. 4(d)} we have 
    \begin{align*}
        \mathrm{Spec}A_\mathfrak{p}=\{\mathfrak{q}\in\mathrm{Spec}A:\mathfrak{q}\cap(A\setminus\mathfrak{p})=\varnothing\}=\{\mathfrak{q}\in\mathrm{Spec}A:\mathfrak{q}\subset\mathfrak{p}\}.
    \end{align*}
    Hence, it follows that
    \begin{align*}
        \mathrm{dim}\ A \myeq{\text{Def.}} \sup\limits_{\mathfrak{p}\in\mathrm{Spec}A}\mathrm{height}\ \mathfrak{p}
        \myeq{\text{(184)}} \sup\limits_{\mathfrak{p}\in\mathrm{Spec}A}\mathrm{dim}\ A_\mathfrak{p}
        \myeq{\text{(183)}} \sup\limits_{\mathfrak{m}\subset A\ \text{maximal}}\mathrm{dim}\ A_\mathfrak{m}.
    \end{align*}
\end{proof}

\begin{proposition}{187}
    Let $A$ be a ring of finite dimension in which all maximal chains of prime ideals have the same length, and let $\mathfrak{p}\in\mathrm{Spec}A$. Then $A/\mathfrak{p}$ is also a ring of finite dimension whose maximal chains of prime ideal have the same length, and 
    \begin{align*}
        \mathrm{dim}\ A=\mathrm{height}\ \mathfrak{p}+\mathrm{dim}\ A/\mathfrak{p}=\mathrm{dim}\ A_\mathfrak{p}+\mathrm{dim}\ A/\mathfrak{p}.
    \end{align*}
    In particular, $\mathrm{dim}\ A=\mathrm{dim}\ A_\mathfrak{p}$ if $\mathfrak{p}$ is maximal.
\end{proposition}
\begin{proof}
    From \textbf{Ex. 2(d)} we have $\mathrm{Spec}(A/\mathfrak{p})\cong\mathcal{V}(\mathfrak{p})$. In particular, any increasing chain $\overline{0}=\overline{\mathfrak{p}_0}\subsetneq\overline{\mathfrak{p}_1}\subsetneq\cdots\subsetneq\overline{\mathfrak{p}_r}=\overline{\mathfrak{q}}$ of prime ideals in $A/\mathfrak{p}$ corresponds to the chain $\mathfrak{p}=\mathfrak{p}_0\subsetneq\mathfrak{p}_1\subsetneq\cdots\subsetneq\mathfrak{p}_r=\mathfrak{q}$ of prime ideals in $A$. Completing $\mathfrak{q}$ into a maximal chain
    \begin{align*}
        0=\mathfrak{r}_0\subsetneq\mathfrak{r}_1\subsetneq\cdots\subsetneq\mathfrak{r}_s=\mathfrak{p}=\mathfrak{p}_0\subsetneq\mathfrak{p}_1\subsetneq\cdots\subsetneq\mathfrak{p}_r=\mathfrak{q}
    \end{align*}
    shows that $\mathrm{dim}\ A=s+r$ (since all maximal chains have equal length), where $\mathrm{height}\ \mathfrak{p}=s$ (otherwise we can still insert some $\mathfrak{r}_j$) and $\mathrm{dim}\ A/\mathfrak{p}=r$ (otherwise we can still insert some $\mathfrak{q}_i$; in particular, all maximal chains have same length in $a/\mathfrak{p}$). Finally, if $\mathfrak{p}$ is maximal, then $A/\mathfrak{p}$ is a field, whence $\mathrm{dim}\ A/\mathfrak{p}=0$ by \textbf{Example 184}.
\end{proof}

\begin{example}{188}
    The assumptions hold for instance for Dedekind rings or polynomial rings over a field, see \textbf{Prop. 197}.
\end{example}

We wish to show that for a field $k,\ \mathrm{dim}\ k[x_1,\cdots,x_n]=n$. For this we need the following going-up and -down theorems which generalise the corresponding statements for Dedekind respectively $1$-dimensional ring extensions to higher dimensional ring extensions.

\begin{proposition}{189. (Lying-Over)}
    Let $A\subset B$ be an integral ring extension.
    \begin{itemize}
        \item[(a)] Let $\mathfrak{q}\in\mathrm{Spec}B$, then $\mathfrak{p}\coloneqq\mathfrak{q}\cap A$ is maximal iff $\mathfrak{q}$ is maximal. [AM, 5.8]
        \item[(b)] Let $\mathfrak{q}\subset\mathfrak{q}'$ in $\mathrm{Spec}B$ such that $\mathfrak{q}\cap A=\mathfrak{q}'\cap A\coloneqq\mathfrak{p}$, then $\mathfrak{q}=\mathfrak{q}'$. [AM, 5.9]
        \item[(c)] Let $\mathfrak{p}\in\mathrm{Spec}A$, then there exists $\mathfrak{q}\in\mathrel{Spec}B$ such that $\mathfrak{p}=\mathfrak{q}\cap A$. [AM, 5.10]
    \end{itemize}
\end{proposition}
\begin{proof}
    (a) Note that from \textbf{Lemma 96(a)} $B/\mathfrak{q}$ is integral over $A/\mathfrak{p}.\ \because \mathfrak{p}$ and $\mathfrak{q}$ are both prime ideal of $A$ and $B$, respectively. $\therefore$ Both $A/\mathfrak{p}$ and $B/\mathfrak{q}$ are integral domains. Thus
    \begin{align*}
        \mathfrak{p}\ \text{is maximal}.\ \xLeftrightarrow{}\ A/\mathfrak{p}\ \text{is a field}.\ \xLeftrightarrow{\text{(133)}}\ B/\mathfrak{q}\ \text{is a field}.\ \xLeftrightarrow{}\ \mathfrak{q}\ \text{is maximal}.
    \end{align*}
    (b) By \textbf{Lemma 96(b)}, $B_\mathfrak{p}$ is integral over $A_\mathfrak{p}$. Let $\mathfrak{m}$ be the extension of $\mathfrak{p}$ in $A_\mathfrak{p}$ and $\mathfrak{n},\ \mathfrak{n}'$ be the extension of $\mathfrak{q}$ of $\mathfrak{q},\ \mathfrak{q}'$ in $B_\mathfrak{p}$, respectively. Then $\mathfrak{m}$ is the maximal ideal of $A_\mathfrak{p}$; $\mathfrak{n}\subset\mathfrak{n}'$; and $\mathfrak{n}^c=\mathfrak{n}'^c=\mathfrak{m}$. It follows from (a) that both $\mathfrak{n},\ \mathfrak{n}'$ are maximal, hence $\mathfrak{n}=\mathfrak{n}'$, hence $\mathfrak{q}=\mathfrak{q}'$ by [AM, 3.11iv].\\
    (c) Note that $B_\mathfrak{p}$ is integral over $A_\mathfrak{p}$, thus the diagram 
    \begin{center}
        \begin{tikzcd}
            A \arrow[r, hook, "~\text{integral}~"] \arrow[d, "\alpha"] 
            & B \arrow[d, "\beta"]\\
            A_{\mathfrak{p}} \arrow[r, hook, "~\text{integral}~"]
            & B_\mathfrak{p}
        \end{tikzcd}
    \end{center}
    is commutative. Let $\mathfrak{n}\triangleleft B_\mathfrak{p}$ be a maximal ideal, then $\mathfrak{m}\coloneqq\mathfrak{n}\cap A_\mathfrak{p}\triangleleft A_\mathfrak{p}$ is maximal by (a). Note: $A_\mathfrak{p}$ is local, hence $\mathfrak{m}$ is the unique maximal ideal. Now, take $\mathfrak{q}\coloneqq\beta^{-1}(\mathfrak{n})$, then $\mathfrak{q}\triangleleft B$ is prime and $\mathfrak{q}\cap A=\alpha^{-1}(\mathfrak{m})=\mathfrak{p}$.
    \begin{center}
        \begin{tikzcd}
            A \arrow[r, hook] \arrow[d, "\alpha"] 
            & B \arrow[d, "\beta"]
            & \mathfrak{p}=\alpha^{-1}(\mathfrak{m}) 
            \arrow[r, maps to] \arrow[d, maps to, "\alpha"] 
            & \mathfrak{q}\coloneqq\beta^{-1}(\mathfrak{n}) 
            \arrow[d, maps to,  "\beta"]\\
            A_{\mathfrak{p}} \arrow[r, hook]
            & B_\mathfrak{p}
            & \exists!\ \mathfrak{m} \arrow[r, maps to, "\text{(a)}"]
            & \mathfrak{n}
        \end{tikzcd}
    \end{center}    
\end{proof}

\begin{proposition}{190. (Incompatibility)}
    Let $A\subset B$ be an integral ring extension. Suppose $\mathfrak{q},\mathfrak{q}'\in\mathrm{Spec}B$ satisfy $\mathfrak{q}\cap A=\mathfrak{q}'\cap A$ but $\mathfrak{q}\neq\mathfrak{q}'$. Then $\mathfrak{q}\nsubseteq\mathfrak{q}'$ and $\mathfrak{q}'\nsubseteq\mathfrak{q}$. [AM, 5.9]
\end{proposition}

\begin{proposition}{191. (Going-Up)}
    Let $A\subset B$ be an integral ring extension. Suppose that  $\mathfrak{p}_1,\mathfrak{p}_2\in\mathrm{Spec}A$ and $\mathfrak{q}_1\in\mathrm{Spec}B$ such that $\mathfrak{q}_1\cap A=\mathfrak{p}_1$. Then there exists $\mathfrak{q}_2\in\mathrm{Spec}B$ with $\mathfrak{q}_1\subset\mathfrak{q}_2$ such that $\mathfrak{q}_2\cap A=\mathfrak{p}_2$. [AM, 5.11]
    \begin{center}
        \begin{tikzcd}
            \mathfrak{q}_1 \arrow[r, hook, dashed, "\subset", red] \arrow[d]
            & \color{red}{\mathfrak{q}_2} \arrow[r, hook, dashed, "\triangleleft", red] \arrow[d, dashed, red]
            & B \arrow[d, hookleftarrow, "\text{integral}"]\\
            \mathfrak{p}_1 \arrow[r, hook, "\subset"]
            & \mathfrak{p}_2 \arrow[r, hook, "\triangleleft"]
            & A \\
        \end{tikzcd}
    \end{center}
\end{proposition}
\begin{proof}
    Consider the map $\phi:A\xhookrightarrow{} B\twoheadrightarrow B/\mathfrak{q}_1$ with its kernel $\mathrm{ker}\phi=\mathfrak{q}_1\cap A=\mathfrak{p}_1$, then we have the inclusion $A/\mathfrak{p}_1\xhookrightarrow{} B/\mathfrak{q}_1$. Note that by \textbf{Example 96(a)}, $A/\mathfrak{p}_1\subset B/\mathfrak{q}_1$ is integral. From \textbf{Prop. 190}, we may assume that $\mathfrak{p}_1\subsetneq\mathfrak{p}_2$ and we consider
    \begin{center}
        \begin{tikzcd}
            \mathfrak{q}_1 \arrow[r, hook, "\subsetneq", red] \arrow[d, maps to]
            & \color{red}{\pi^{-1}(\mathfrak{P})} 
            \arrow[r, hook, "\triangleleft", red] \arrow[d, maps to, red]
            & B \arrow[ddd, bend left, two heads, "\pi"] \\
            \mathfrak{p}_1 \arrow[r, hook, "\subsetneq"] \arrow[d, maps to]
            & \mathfrak{p}_2 \arrow[r, hook, "\triangleleft"] \arrow[d, maps to]
            & A \arrow[u, hook, "\text{integral}"] \arrow[d, two heads]\\
            \overline{0} \arrow[r, hook, "\subsetneq"]
            & \overline{\mathfrak{p}_2} \arrow[r, hook, "\triangleleft"]
            & A/\mathfrak{p}_1 \\
            & \color{blue}{\mathfrak{P}} \arrow[r, hook, "\triangleleft", blue]
            \arrow[u, maps to, blue] \arrow[uuu, bend left, maps to, red]
            & B/\mathfrak{q}_1 \arrow[u, hookleftarrow, "\text{integral}"]
        \end{tikzcd}
    \end{center}
    where the existence of $\mathfrak{P}\in\mathrm{Spec}(B/\mathfrak{q}_1)$ satisfying $\mathfrak{P}\cap(A/\mathfrak{p}_1)=\overline{\mathfrak{p}_2}$ follows from \textbf{Lying-Over Thm}. We take $\mathfrak{q}_2\coloneqq\pi^{-1}(\mathfrak{P})\in\mathrm{Spec}B$, then clearly, $\mathfrak{q}_2\supsetneq\mathfrak{q}_1$. Moreover, one can show that $\mathfrak{P}=\overline{\mathfrak{q}_2}=\{y+\mathfrak{q}_1:y\in\mathfrak{q}_2\}$, and hence $\mathfrak{q}_2\cap A=\mathfrak{p}_2$.
\end{proof}

\begin{proposition}{192. (Going-Down)}
    Let $A\subset B$ be an integral ring extension with $A$ integrally closed and $B$ an integral domain. Suppose that  $\mathfrak{p}_1,\mathfrak{p}_2\in\mathrm{Spec}A$ and $\mathfrak{q}_2\in\mathrm{Spec}B$ such that $\mathfrak{q}_2\cap A=\mathfrak{p}_2$. Then there exists $\mathfrak{q}_1\in\mathrm{Spec}B$ with $\mathfrak{q}_1\subset\mathfrak{q}_2$ such that $\mathfrak{q}_1\cap A=\mathfrak{p}_1$. [AM, 5.16]
    \begin{center}
        \begin{tikzcd}
            \color{red}{\mathfrak{q}_1} \arrow[r, hook, dashed, "\subset", red] \arrow[d, dashed, red]
            & \mathfrak{q}_2 \arrow[r, hook, "\triangleleft"] \arrow[d]
            & B \arrow[d, hookleftarrow, "\text{integral}"]\\
            \mathfrak{p}_1 \arrow[r, hook, "\subset"]
            & \mathfrak{p}_2 \arrow[r, hook, "\triangleleft"]
            & A \\
        \end{tikzcd}
    \end{center}
\end{proposition}
\begin{proof}
    \textbf{(Sketch of Proof)}. 
    Consider the local ring $(B_{\mathfrak{q}_2},\ \mathfrak{q}_2B_{\mathfrak{q}_2})$. The idea is to find $\mathfrak{P}\in\mathrm{Spec}B_{\mathfrak{q}_2}$ such that
    \begin{align*}
        \mathfrak{P}\subset\mathfrak{q}_2B_{\mathfrak{q}_2}\qquad\text{and} \qquad \mathfrak{P}\cap A=\mathfrak{p}_1.    
    \end{align*}
    We may choose $\mathfrak{P}\coloneqq\mathfrak{p}_1B_{\mathfrak{q}_2}\subset\mathfrak{p_2}B_{\mathfrak{q}_2}$. Then we pick $\mathfrak{q}_1\coloneqq\mathfrak{P}\cap B$.
    \begin{center}
        \begin{tikzcd}
            &\color{blue}{\mathfrak{P}} \arrow[r, hook, "\subset", blue] \arrow[dd, blue] \arrow[dl, dashed, red] 
            & \mathfrak{q}_2B_{\mathfrak{q}_2} \arrow[r, hook, "\triangleleft"] \arrow[d]
            & B_{\mathfrak{q}_2} \arrow[d, hookleftarrow]\\
            \color{red}{\mathfrak{P}\cap A} \arrow[rr, hook, dashed, "~~~~~~~\subset", red] \arrow[dr, dashed, red] &
            & \mathfrak{q}_2 \arrow[r, hook, "\triangleleft"] \arrow[d]
            & B \arrow[d, hookleftarrow]\\
            &\mathfrak{p}_1 \arrow[r, hook, "\subset"]
            & \mathfrak{p}_2 \arrow[r, hook, "\triangleleft"]
            & A \\
        \end{tikzcd}
    \end{center}    
    Now, we obtain
    \begin{align*}
        \mathfrak{q}_1\cap A\myeq{\text{Def.}} (\mathfrak{P}\cap B)\cap A =\mathfrak{P}\cap A=\mathfrak{p}_1
    \end{align*}
    Thus, it remains to show: $\mathfrak{p}_1$ is a contracted prime ideal of the prime ideal $\mathfrak{P}\triangleleft B_{\mathfrak{q}_2}$. By [AM, 3.16], it is equivalent to show that $\mathfrak{p}_2^{ec}=\mathfrak{p}_2$.
\end{proof}

\begin{corollary}{193}
    For any integral ring extension $A\subset B$ we have $\mathrm{dim}\ A=\mathrm{dim}\ B$.
\end{corollary}
\begin{proof}
    "$\leq$"\ Let $\mathfrak{p}_0\subsetneq\mathfrak{p}_1\subsetneq\cdots\subsetneq\mathfrak{p}_n$ be a strict chain of prime ideals in $A$. Apply \textbf{Lying-Over} to $\mathfrak{p}_0$, there is a prime ideal $\mathfrak{q}_0\triangleleft B$ with $\mathfrak{q}_0\cap A=\mathfrak{p}_0$. Now, by applying \textbf{Going-Up} we obtain a strict chain of prime ideals in $B$.
    \begin{center}
        \begin{tikzcd}
            \color{blue}{\mathfrak{q}_0} \arrow[r, hook, "\subset", red] 
            & \color{red}{\mathfrak{q}_1} \arrow[rrr, hook, "\triangleleft", red] \arrow[d, "\text{Going-Up}", red] & &
            & B \arrow[d, hookleftarrow, "\text{integral}"]\\
            \mathfrak{p}_0 \arrow[r, hook, "\subset"]
            \arrow[u, leftarrow, "\text{Lying-Over}", blue]
            & \mathfrak{p}_1 \arrow[r, hook, "\subset"]
            & \cdots \arrow[r, hook, "\subset"]
            & \mathfrak{p}_n \arrow[r, hook, "\triangleleft"]
            & A \\
        \end{tikzcd}
    \end{center}
    "$\geq$"\ Let $\mathfrak{q}_0\subsetneq\mathfrak{q}_1\subsetneq\cdots\subsetneq\mathfrak{q}_n$ be a strict chain of prime ideals in $B$. Taking "contraction" under $A$ yields a chain of prime ideals in $A$ which is strict by "incompatibility".
\end{proof}

\newpage

\section{Noether Normalisation}
\begin{theorem}{194. (Noether Normalisation Lemma)}
    Let $k\subset B$ be a $k$-algebra of finite type. Then there exists algebraically independent elements $y_1,\cdots,y_l\in B$ such that $B$ is a finite $A=k[y_1,\cdots,y_l]$-module, that is, we have
    \begin{align*}
        k\subset A=k[y_1,\cdots,y_l]\subset B,
    \end{align*}
    where $A$ is a polynomial algebra over $k$ and $B$ a finite over $A$. Furthermore, if $k$ is infinite the variables $y_1,\cdots,y_l$ are obtained from a linear transformation of the $k$-generators $x_1,\cdots,x_l$ of $B$.  
\end{theorem}
\begin{theorem}{194*. (Generalization)}
    Let $A\subset B\subset C$ be rings. Suppose that $A$ is Noetherian, $C$ is finitely generated as an $A$-algebra and that $C$ is either \\
    \centerline{\textbf{(i)} finitely generated as a $B$-module or \textbf{(ii)} integral over $B$.}\\ Then $B$ is finitely generated as an $A$-algebra. [AM, 7.8]
\end{theorem}
\begin{remark}{194**}
    Since $C$ is of finite type over $A$ and $A\subset B$, so is $C$ over $B$. By \textbf{Rmk. 94*} the assumptions (i) and (ii) are equivalent.
\end{remark}
\begin{proofthm}{194*}
    We suppose $C$ satisfies the assumption (i), say $C\coloneqq By_1+\cdots+By_l, \ (y_j\in C).\ \because C$ is a finitely generated $A$-algebra. $\therefore C\coloneqq A[x_1,\cdots,x_m],\ (x_i\in C)$. For each $1\leq i\leq m$, we write
    \begin{align}
        x_i\coloneqq\sum_{j=1}^lb_{ij}y_j,\qquad (b_{ij}\in B).
    \end{align}
    Since $C$ is a ring, each $y_iy_j\in C$, so we may write
    \begin{align}
        y_iy_j\coloneqq\sum_{k=1}^lb_{ijk}y_k,\qquad (b_{ijk}\in B).
    \end{align}
    Now, let $B_0$ be the $A$-algebra generated by these $b_{ij}$ and $b_{ijk}$'s, then $B_0$ is finitely generated as an $A$-algebra such that $A\subset B_0\subset B\subset C$. \\
    \textbf{Claim.} $C$ is a finitely generated $B_0$-module. We first note that for each $x\in C=A[x_1,\cdots,x_m],\ x$ is a polynomial in $x_1,\cdots x_m$ with coefficients in $A$. From (1) and (2), $x$ is a linear combination of $y_1,\cdots,y_l$ with coefficients in $B_0$, that is, $x\in B_0-\mathrm{span}\{y_1,\cdots,y_l\}$. Thus $C=B_0y_1+\cdots+B_0y_l$ is finitely generated as a $B_0$-module.\\
    To complete the proof, since $A$ is Noetherian, so is its finitely generated algebra $B_0$, thus $C$ is a Noetherian $B_0$-module. $\because B$ is a $B_0$-submodule of $C.\ \therefore B$ is a finitely generated $B_0$-module. Moreover, $B_0$ is a finitely generated $A$-algebra, thus so is $B$.\qed
\end{proofthm}

\begin{theorem}{195. (Hilbert Weak Nullstellensatz)}
    Let $E$ be a finitely generated $k$-algera which is a field. Then $[E:k]<\infty$. [AM, 7.9]
\end{theorem}
\begin{proof}
    For the finitely generated $k$-algebra $E$, we write $E\coloneqq k[\alpha_1,\cdots,\alpha_n],\ (\alpha_i\in E)$.\\ \textbf{Note:} If each $\alpha_i$ is algebraic over $k$, then $E/k$ is algebraic, thus finite, done.\\ Now, suppose that $E/k$ is not an algebraic extension. After suitable renumbering indexes, we may assume that (i) $\{\alpha_i\}_{i=1}^r$ is algebraically $k$-independent and (ii) $\{\alpha_i\}_{i=r+1}^n$ is algebraic over $F\coloneqq k(\alpha_1,\cdots,\alpha_r).\ \because E/F$ is algebraic, thus finite extension. $\therefore E$ is a finitely generated $F$-module. Moreover, $E$ is a finitely generated $k$-algebra, it follows from \textbf{Thm. 194} that $F$ is a finitely generated $k$-algebra, say $F\coloneqq k[\beta_1,\cdots,\beta_s],\ (\beta_j\in F)$. From (ii), for each $1\leq j\leq s$, we have 
    \begin{align*}
        \beta_j\coloneqq\tfrac{f_j}{g_j}\in F\cong k(x_1,\cdots,x_r) =\mathrm{Quot}(k[x_1,\cdots,x_r]),
    \end{align*}
    where $f_j,g_j\in k[\alpha_1,\cdots,\alpha_r]\cong k[x_1,\cdots,x_r]$: UFD with $g_j\neq0$. Choose an irreducible polynomial $h\in k[\alpha_1,\cdots,\alpha_r]\cong k[x_1,\cdots,x_r]$ such that $\gcd{(h,g_j)}\sim1,\ \forall 1\leq j\leq s$, (e.g. $h\coloneqq g_1g_2\cdots g_s+1$). $\because F$ is a field and $h\neq0.\ \therefore h^{-1}\in F=k[\beta_1,\cdots,\beta_s]$. However, $h^{-1}$ cannot be written as a polynomial in variables $\beta_1,\cdots,\beta_s.\ (\to\gets)$. Thus $E$ must be algebraic over $k$.
\end{proof}

\begin{corollary}{196}
    If $k$ is an algebraically closed, then every maximal ideal in $k[x_1,\cdots,x_n]$ is of the form $(x_1-a_1,\cdots,x_n-a_n)$ for $a\coloneqq(a_1,\cdots,a_n)\in\mathbb{A}_k^n$.
\end{corollary}
\begin{proof}
    Since $(x_1-a_1,\cdots,x_n-a_n)$ is the kernel of the map
    \begin{align*}
        \mathrm{ev}_a:k[x_1,\cdots,x_n]\longrightarrow k,\ x_i\longmapsto a_i,
    \end{align*}
    it is a maximal ideal. Conversely, let $\mathfrak{m}\triangleleft k[x_1,\cdots,x_n]$ be a maximal ideal, then $E\coloneqq k[x_1,\cdots,x_n]/\mathfrak{m}$ is a field and a finitely generated $k$-algebra, whence $E/k$ is a finite extension by\textbf{Thm. 195}. Since $k$ is algebraically closed, we have an isomorphism $\iota:k\to k[x_1,\cdots,x_n]/\mathfrak{m}=E$. Choosing inverse images $a_1,\cdots,a_n\in k$ for $\overline{x_1},\cdots,\overline{x_n}$ shows that $(x_1-a_1,\cdots,x_n-a_n)$ is the kernel of the composition $\iota\circ\mathrm{ev}_a:k[x_1,\cdots,x_n]\to k[x_1,\cdots,x_n]/\mathfrak{m}$ and is thus equal to $\mathfrak{m}$.
\end{proof}

\begin{proposition}{197}
    Let $k$ be a field and $n\in\mathbb{N}$. Then all maximal chains of prime ideals in $k[x_1,\cdots,x_n]$ have length $n$. In particular, $\mathrm{dim}\ k[x_1,\cdots,x_n]=n$.
\end{proposition}
\begin{proof}
    We proof the both statements by induction on $n$. The cases for $n=0$ and $n=1$ are clear. Now let $n\geq2$ and $\mathfrak{p}_0\subsetneq\mathfrak{p}_1\subsetneq\cdots\subsetneq\mathfrak{p}_r$ be a strict chain of prime ideals in $k[x_1,\cdots,x_n]$. We wish to show that $r\leq n$ with equality iff the chain is maximal. Without loss of generality we may assume that (possibly inserting prime ideals)
    \begin{itemize}
        \item $\mathfrak{p}_0=(0)$.
        \item $\mathfrak{p}_1=(f)$, for $f\in k[x_1,\cdots,x_n]$ irreducible. ($k[x_1,\cdots,x_n]$ is factorial!)
        \item $\mathfrak{p}_r$ is maximal.
    \end{itemize}
    As in the proof of \textbf{Thm. 194*}, we may assume (by linearly transforming the coordinates $x_1,\cdots,x_n$ if necessary), that $f\in k[x_1,\cdots,x_{n-1}][x_n]$ is monic in $x_n$. In particular,
    \begin{align*}
        k[x_1,\cdots,x_{n-1}]\subset k[x_1,\cdots,x_{n-1}][x_n]/(f)
        =k[x_1,\cdots,x_n]/\mathfrak{p}_1
    \end{align*}
    is an integral ring extension. We push down the strict chain $\mathfrak{p}_0\subsetneq\mathfrak{p}_1\subsetneq\cdots\subsetneq\mathfrak{p}_r$ of $k[x_1,\cdots,x_n]$ to a chain $(0)= \mathfrak{q}_1\subsetneq\cdots\subsetneq\mathfrak{q}_r$ by first moding out $\mathfrak{p}_1$ and second by intersecting with $k[x_1,\cdots,x_{n-1}]$. This chain is strict by "incompatibility". Next, it is maximal if the chain $\mathfrak{p}_i$ is maximal. Indeed, $\mathfrak{q}_r$ is maximal since it is obtained by restricting a maximal ideal of an integral ring extension. Further, if we could insert some prime ideals $\mathfrak{q}$ into the chain $\mathfrak{q}_i$, this would be covered by a prime ideal $\mathfrak{p}\triangleleft k[x_1,\cdots,x_n]$ which would strictly extend the chain $\mathfrak{p}_i$ contradicting its maximality. $(\to\gets)$. (cf. \textbf{Exercise}).\\
    By induction hypothesis, $r-1=$ the length of the maximal chain $\mathfrak{q}_i$ and $\leq n-1$ with equality iff the chain is maximal. Hence, $m\leq n$ with equality iff the chain $\mathfrak{p}_i$ is maximal.
\end{proof}

\begin{corollary}{198}
    Any finitely generated $k$-algebra $A$ is of finite dimension. Furtermore, if $A$ is integral, then
    \begin{itemize}
        \item[(a)] all maximal chains of prime ideals in $A$ have length $\mathrm{dim}\ A$.
        \item[(b)] $\mathrm{dim}\ A=\mathrm{trdeg}_k\ \mathrm{Quot}(A)$.
    \end{itemize}
\end{corollary}
\begin{proof}
    (a) By \textbf{Thm. 194} we know that $A$ is a finite, thus an integral extension of $k[x_1,\cdots,x_n]$ for some $n$ so that $\mathrm{dim}\ A=\mathrm{dim}\ k[x_1,\cdots,x_n]=n$ from \textbf{Prop. 197}. Furthermore, if $A$ is integral, then $A\cong k[x_1,\cdots,x_r]/\mathfrak{p}$ for some $r$ and $\mathfrak{p}$ prime ideal. Since all maximal chains of prime ideals in $k[x_1,\cdots,x_r]$ have the same length, so do the maximal chains in the quotient above.\\
    (b) Next, the field extension $k(x_1,\cdots,x_n)\subset\mathrm{Quot}A$ is algebraic for the ring extension $k[x_1,\cdots,x_n]\subset A$ is integral. Indeed, if $a,b\in A$ are integral over $k[x_1,\cdots,x_n]$, then they are algebraic over $k(x_1,\cdots,x_n)$. It follows that $k(x_1,\cdots,x_n)\subset k(x_1,\cdots,x_n)(a,b)$ must be an algebraic (thus, finite) field extension. In particular, $\frac{a}{b}$ is algebraic over $k(x_1,\cdots,x_n)$ since algebraic elements form a field. Therefore, we have
    \begin{align*}
        \mathrm{trdeg}_k\mathrm{Quot}(A)=\mathrm{trdeg}_kk(x_1,\cdots,x_n)=n=\mathrm{dim}\ A.
    \end{align*}
\end{proof}

\newpage

\section{Smoothness}
\begin{definition}{199. (Tangent Spaces)}
    Let $k$ be a field. 
    \begin{itemize}
        \item Consider the \textbf{formal differential map} defined by
        \begin{align*}
            d:k[x_1,\cdots,x_n]\longrightarrow k[x_1,\cdots,x_n],\ f\longmapsto\sum_{i=1}^n\tfrac{\partial f}{\partial x_i}(0)x_i,
        \end{align*}
        where $\frac{\partial f}{\partial x_i}\in k[x_1,\cdots,x_n]$ is the formal derivative of the polynomial $f$ with respect to $x_i$. 
        \item We can regard $df$ as a linear function
        \begin{align*}
            df:k^n\longrightarrow k,\ a\coloneqq(a_1,\cdots,a_n)\mapsto df(a)=\sum_{i=1}^n\tfrac{\partial f}{\partial x_i}(0)a_i,
        \end{align*}
        that is, $df\in k^{n\vee}$.
        \item Next, let $\mathfrak{a}\triangleleft k[x_1,\cdots,x_n]$ be an ideal with $\mathfrak{a}\subset\mathfrak{p}\coloneqq(x_1,\cdots,x_n)$. Geometrically, this means that the closed point corresponding to $\mathfrak{p}$ lies in $\mathcal{V}(\mathfrak{a})$. We define the \textbf{tangent space} of $\mathfrak{a}$ at $\mathfrak{p}$ by 
        \begin{align*}
            T_\mathfrak{p}\mathfrak{a}\ \coloneqq\ \bigcap_{f\in\mathfrak{a}}\mathrm{ker}\ df\ \subset\ k^n.
        \end{align*}
        This is clearly a $k$-vector space with \textbf{annihilator} $N^\vee T_\mathfrak{p}\mathfrak{a}=\{df:f\in\mathfrak{a}\}\subset k^{n\vee}$.
    \end{itemize}
\end{definition}

\begin{remark}{200}
    \begin{itemize}
        \item[(a)] One immediately verifies the \textbf{derivation property} (Leibniz rule):
        \begin{align*}
            d(fg)\ =\ f(0)\cdot dg+g(0)\cdot df.
        \end{align*}
        \item[(b)] As the formal differential is a $k$-linear operator, it is enough to compute $df$ for a generating system of $\mathfrak{a}$.
    \end{itemize}
\end{remark}

\begin{example}{201}
    Let $f=x_1^2+(x_2-1)^2-1$ and consider $\mathfrak{a}=(f)\subset\mathbb{R}[x_1,x_2]$. We have $df=-2x_2$, whence $T\mathfrak{a}=\{(a_1,a_2)\in\mathbb{R}^2:a_2=0\}$. If we think of $\mathcal{V}(\mathfrak{a})$ in $\mathrm{Spec}\ k[x_1,\cdots,x_n]$ as
    \begin{align*}
        \mathcal{V}(f)=f^{-1}(0)=\ \text{the unit circle through}\ (0,0)\ \text{and}\ (0,2),
    \end{align*}
    (cf. \textbf{Rmk. 172}), we can see that $T\mathfrak{a}$ corresponds to the geometric tangent space $T_{(0,0)}\mathcal{V}(f)$ of $\mathcal{V}(f)$.
\end{example}

\begin{proposition}{202}
    The $k$-linear map defined by
    \begin{align*}
        \overline{\mathfrak{p}}\ \coloneqq\ \mathfrak{p}/\mathfrak{a}\  &\xrightarrow{~\phi~}\  T\mathfrak{a}^\vee\coloneqq\mathrm{Hom}_k(T\mathfrak{a},k),\\ \overline{f}\coloneqq f+\mathfrak{a} &\longmapsto df|_{T\mathfrak{a}},
    \end{align*}
    induces an isomorphism
    \begin{align*}
        T\mathfrak{a}^\vee\ \cong\ \overline{\mathfrak{p}}/(\overline{\mathfrak{p}})^2.
    \end{align*}
\end{proposition}
\begin{proof}
    We claim that $\phi$ is a well-defined surjective and $\mathrm{ker}\phi=(\overline{\mathfrak{p}})^2$.\\
    \begin{itemize}
        \item \textbf{(Well-defined).} If $\overline{f}=\overline{g},\ (f,g\in\mathfrak{p})$, then $f-g\in\mathfrak{a}$.\\ \textsc{claim.} $d(f-g)|_{T\mathfrak{a}}=0$. For any $a\in T\mathfrak{a}=\bigcap_{f\in\mathfrak{a}}\mathrm{ker}(df)\subset k^n.\ \because a\in\mathrm{ker}(d(f-g)).\\ \therefore d(f-g)(a)=0.\ \Rightarrow d(f-g)|_{T\mathfrak{a}}=0.\ \Rightarrow df|_{T\mathfrak{a}}=dg|_{T\mathfrak{a}}$, i.e. $\phi$ is well-defined. 
        \item \textbf{(Surjective).} Let $F\in\mathrm{Hom}_k(T\mathfrak{a},k)$, i.e. $F:T\mathfrak{a}\to k$ is $k$-linear. Note that $T\mathfrak{a}$ is a $k$-vector subspace of $k^n$, let $\{e_j\}_{j=1}^m$ be its $k$-basis where $m\leq n$. By suitable linear transformation we may assume the basis is standard. We consider $f\coloneqq\sum_{j=1}^mF(e_j)x_j\in\mathfrak{p}$, then
        \begin{align*}
            df=\sum \tfrac{\partial f}{\partial x_j}(0)x_j=\sum F(e_j)x_j=f.
        \end{align*}
        Now, we have
        \begin{align*}
            f|_{T\mathfrak{a}}(e_j)=f(e_j)=\sum_{k=1}^mF(e_k)\delta_{jk}=F(e_j),\ \forall 1\leq j\leq m.
        \end{align*}
        Thus, $\phi(\overline{f})=f|_{T\mathfrak{a}}=F$.
        \item \textbf{(Kernel).} "$\subset$"\ If $\overline{f}\in\mathrm{ker}\phi,\ (f\in\mathfrak{p})$, then $df|_{T\mathfrak{a}}=0$, i.e. $df\in N^\vee T\mathfrak{a}$ the annihilator. Thus there is $g\in\mathfrak{a}$ such that $df=dg$, it follows that $d(f-g)=0$, which means $f-g$ has no linear terms. Moreover, since $f,g\in\mathfrak{p}$, they have no constant terms. Hence $f-g$ is of the form
        \begin{align*}
            f-g\coloneqq\ \sum a_{ij}x_ix_j+\ \text{higher order terms}.
        \end{align*}
        Since $g\in\mathfrak{a}$, we conclude that
        \begin{align*}
            \overline{f}=\overline{f-g}=\sum a_{ij}\overline{x_i}\overline{x_j}+\ \text{higher order terms}\ \in(\overline{\mathfrak{p}})^2.
        \end{align*}
        "$\supset$"\ Conversely, note that $(\overline{\mathfrak{p}})^2$ is generated by $\{\overline{x_i}\overline{x_j}\}$ and that $(\overline{\mathfrak{p}})^2\subset\overline{\mathfrak{p}}$. Since
        \begin{align*}
            \phi(\overline{x_i}\overline{x_j})=d(x_ix_j)|_{T\mathfrak{a}}=\big(x_i(0)dx_j+x_j(0)dx_i\big)|_{T\mathfrak{a}}=0,
        \end{align*}
        we have $\overline{x_i}\overline{x_j}\in\mathrm{ker}\phi$, whence $(\overline{\mathfrak{p}})^2\subset\mathrm{ker}\phi$.
    \end{itemize}
    By Isomorphism Theorem we obatain the result.
\end{proof}

\begin{remark}{203}
    This suggests to define the tangent space at a closed point of an affine scheme as follows:
    \begin{itemize}
        \item For a maximal ideal $\mathfrak{p}\in X=\mathrm{Spec}A$, define the $k=A/\mathfrak{p}$-vector space $T_\mathfrak{p}X\coloneqq(\mathfrak{p}/\mathfrak{p}^2)^\vee$. The tangent space thus defined as a purely local object.
        \item Indeed, consider the local ring $(A_\mathfrak{p},\ \mathfrak{p}A_\mathfrak{p})$, (for this we assume again that $A$ is integral though that is not strictly necessary). 
        \item Since $S=A\setminus\mathfrak{p}$ maps to units in $k$, we have
        \begin{align*}
            \mathfrak{p}/\mathfrak{p}^2\ \cong\ S^{-1}(\mathfrak{p}/\mathfrak{p}^2)\ \cong\ S^{-1}\mathfrak{p}/S^{-1}\mathfrak{p}^2\ \cong\ \mathfrak{p}A_\mathfrak{p}/(\mathfrak{p}A_\mathfrak{p})^2,
        \end{align*}
        where the first isomorphism is as $k$-vector spaces, and the second $A$-linear isomorphism follows from exactness of the localisation functor (cf. [AM, 3.3]). In particular, $\mathfrak{p}/\mathfrak{p}^2\cong\mathfrak{p}A_\mathfrak{p}/(\mathfrak{p}A_\mathfrak{p})^2$ as $k$-vector spaces, so that $T_\mathfrak{p}X\cong(\mathfrak{p}A_\mathfrak{p}/(\mathfrak{p}A_\mathfrak{p})^2)^\vee$ can be solely computed from the local ring $A_\mathfrak{p}$.
    \end{itemize}
\end{remark}

\begin{proposition}{204}
    Let $(A,\mathfrak{m})$ be a local Noetherian ring with the residur field $k$. Then
    \begin{itemize}
        \item[(a)] The number $\mathrm{dim}_k(\mathfrak{m}/\mathfrak{m}^2)$ is the minimal number of generators for $\mathfrak{m}$.
        \item[(b)] $\mathrm{dim}\ A\leq\mathrm{dim}_k(\mathfrak{m}/\mathfrak{m}^2)<\infty$.
    \end{itemize}
\end{proposition}
\begin{proof}
    \textsc{Note.} Since $A$ is Noetherian, there is a finite minimum number $n$ of generators of $\mathfrak{m}$. Further, as $\mathfrak{m}$ is finitely generated, so is  $\mathfrak{m}/\mathfrak{m}^2$ over $A/\mathfrak{m}=k$, i.e. $\mathrm{dim}_k(\mathfrak{m}/\mathfrak{m}^2)<\infty$.\\
    (a) Suppose $\mathfrak{m}\coloneqq(a_1,\cdots,a_m)$, then $\overline{a_i}\coloneqq a_i+\mathfrak{m}^2$ generates $\mathfrak{m}/\mathfrak{m}^2$ as a $k$-vector space, whence $n\geq\mathrm{dim}_k(\mathfrak{m}/\mathfrak{m}^2)$. If $\{\overline{a_i}\}_{i=1}^n$ is $k$-dependent, by possibly relabeling $\{\overline{a_i}\}_{i=1}^{n-1}$ still generates $\mathfrak{m}/\mathfrak{m}^2$. By \textbf{Nakayama's Lemma} [AM, 2.8], $\{a_i\}_{i=1}^{n-1}$ generates $\mathfrak{m}.\ (\to\gets)$.\\ 
    (b) As $A$ is local, from \textbf{Example 184} we have
    \begin{align*}
        \mathrm{dim}\ A\ =\ \mathrm{dim}\ A_\mathfrak{m}\ =\ \mathrm{height}\ \mathfrak{m}.
    \end{align*}
    But $\mathrm{height}\ \mathfrak{m}\leq n$ by Exercise.
\end{proof}

\begin{definition}{205. (Regular Ring)}
    A local ring $(A,\mathfrak{m})$ is \textbf{regular} if $A$ is Noetherian and $\mathrm{dim}\ A=\mathrm{dim}_k(\mathfrak{m}/\mathfrak{m}^2)$.
\end{definition}
\newpage
\begin{example}{206}
    \begin{itemize}
        \item[(a)] Since a field $(k,\ 0)$ is zero dimensional, it is always regular.
        \item[(b)] For a DVR $(A,\mathfrak{m})$, we have $\mathfrak{m}/\mathfrak{m}^2\cong A/\mathfrak{m}=k$, whence it is regular.
        \item[(c)] The localisation $(k[x_1,\cdots,x_n]_\mathfrak{p},\ \mathfrak{p}k[x_1,\cdots,x_n])$ is regular for
        \begin{align*}
            \mathfrak{p}k[x_1,\cdots,x_n]/(\mathfrak{p}k[x_1,\cdots,x_n])^2\cong \{\sum a_ix_i:a_i\in k\}
        \end{align*}
        is $n$-dimensional. On the other hand, from \textbf{Prop. 187}, we have
        \begin{align*}
            n=\mathrm{dim}\ k[x_1,\cdots,x_n]=\mathrm{dim}\ k[x_1,\cdots,x_n]_\mathfrak{p}.
        \end{align*}
    \end{itemize}
\end{example}

\begin{proposition}{207}
    A regular ring is an integral domain.
\end{proposition}
\begin{proof}
    Let $(A,\mathfrak{m})$ be a regular local ring with residue field $k=A/\mathfrak{m}$. We prove by induction on $n=\mathrm{dim}_k(\mathfrak{m}/\mathfrak{m}^2)=\mathrm{dim}\ A$.\\
    For $n=0$, we have $\mathfrak{m}=\mathfrak{m}^2$. Thus, by \textbf{Nakayama}, $\mathfrak{m}=0$. Since $A^\times=A\setminus\mathfrak{m}=A\setminus\{0\}$, $A$ is a field and thus an integral domain.\\ Now, let $n\geq1$ and $\mathfrak{p}_i$ be the minimal primes over $(0)$, i.e. for any prime $(0)\subset\mathfrak{p}\subset\mathfrak{p}_i$, we have either $\mathfrak{p}=(0)$ or $\mathfrak{p}=\mathfrak{p}_i$. Since $\mathrm{dim}\ A\geq1$, it follows that either $\mathfrak{m}$ is minimal (so that $(0)$ is prime, thus $A\cong A/(0)$ is an integral domain) or $\mathfrak{m}$ strictly contains these $\mathfrak{p}_i$. Assuming this, we will now proceed in three steps:\\
    \textsc{step 1.} \textit{We can find an element $a\in\mathfrak{m}$ not contained in $\mathfrak{m}^2$ nor in any $\mathfrak{p}_i$.}\\
    If not, then $\mathfrak{m}$ is contained in $\mathfrak{m}^2\cup(\bigcup\mathfrak{p}_i)$. It follows from \textbf{(22c)} that either $\mathfrak{m}\subset\mathfrak{m}^2$ (thus $\mathfrak{m}=0$ by \textbf{Nakayama}) or $\mathfrak{m}\subset\mathfrak{p}_i$ for some $i$ (contradicting the assumption).\\
    \textsc{step 2.} \textit{The Noetherian local ring $(A/(a),\mathfrak{m}/(a))$ is regular of dimension $n-1$.}\\
    Indeed, since $a\notin\mathfrak{m}^2$, we can extend $\overline{a}\in\mathfrak{m}/\mathfrak{m}^2$ to a $k$-basis $\{\overline{a}=\overline{a}_1,\cdots,\overline{a}_n\}$ of $\mathfrak{m}/\mathfrak{m}^2$. Hence $\{\overline{a}_i\}_{i=2}^n$ generates the $k$-vector space $V\coloneqq\frac{\mathfrak{m}/(a)}{\mathfrak{m}^2/(a)}$, so $\mathrm{dim}_kV\leq n-1$. On the other hand, let 
    \begin{align*}
        \mathfrak{q}_0\subsetneq\mathfrak{q}_1\subsetneq\cdots\subsetneq\mathfrak{q}_n=\mathfrak{m}
    \end{align*}
    be a maximal chain of prime ideals in $A$. It must necessarily start with a minimal prime ideal $\mathfrak{q}_0=\mathfrak{p}_i$. Since $a\in\mathfrak{m}$, we can choose the prime ideals in such a way that $a\in\mathfrak{q}_1$, (cf. [Ch3, \#10]). Therefore,
    \begin{align*}
        \mathfrak{q}_0/(a)\subsetneq\mathfrak{q}_1/(a)\subsetneq\cdots\subsetneq\mathfrak{q}_n/(a)=\mathfrak{m}/(a)
    \end{align*}
    is a chain of prime ideals in $A/(a)$, whence $\mathrm{dim}\ A/(a)\geq n-1$. Since $\mathrm{dim}_kV\geq\mathrm{dim}\ A/(a)$ by previous result and $\mathrm{dim}\ A/(a)\geq n-1\geq\mathrm{dim}_kV$, we thus conclude that $\mathrm{dim}\ A/(a)=\mathrm{dim}_kV=n-1$, i.e. $A/(a)$ is regular.\\
    \textsc{step 3.} \textit{The induction step.} From induction hypothesis we conclude that $A/(a)$ is an integral domain. In particular, $(a)$ is a prime ideal. Hence we must have $\mathfrak{p}_i\subset(a)$ for some $i$, so that any $b\in\mathfrak{p}_i$ is of the form $b=ac,\ (c\in A)$, whence $b\in\mathfrak{m}\cdot\mathfrak{p}_i$. Now, the finitely generated $A$-module $\mathfrak{p}_i$ satisfies $\mathfrak{p}_i=\mathfrak{m}\cdot\mathfrak{p}_i$. Again, by \textbf{Nakayama} we have $\mathfrak{p}_i=(0)$, which implies that $(0)$ is a prime ideal, hence $A$ is an integral domain.
\end{proof}
\newpage
\begin{remark}{208}
    In fact, regular rings are UFDs, which is the \textbf{Auslander-Buchsbaum Theorem}, cf. [Commutative Algebra, Eisenbud, 19.19].
\end{remark}

\begin{summary}{Algebraic Summary. 209}
    We obtain the following notions of rings
    \begin{center}
        \begin{tikzcd}
            & &\textbf{Noether} &\\
            &\color{red}{\textbf{PID of}\ \mathrm{dim}A=1} \arrow[r, Rightarrow] &\color{red}{\textbf{Dedekind}} \arrow[u, Rightarrow] \arrow[dr, Rightarrow] &\\
            \color{red}{\textbf{DVR}} \arrow[ur, Rightarrow] \arrow[dr, Rightarrow] & & &\textbf{normal}\\
            &\textbf{regular} \arrow[r, Rightarrow] \arrow[d, Rightarrow] &\textbf{UFD} \arrow[ur, Rightarrow] &\\ 
            &\textbf{Noether local} & &\\
        \end{tikzcd}
    \end{center}
    where rings of dimension $1$ are marked in red.
\end{summary}

\begin{summary}{Geometric Summary. 210}
    Both dimension and regularity are local properties in the algebraic and geometric sense:
    \begin{itemize}
        \item From \textbf{Example 184}, Dedekind rings are $1$-dimensional. Moreover, the localisation of a Dedekind ring $A$ is a DVR. They are $1$-dimensional as PIDs in accordance to the local character of the dimension. They are also regular in view of \textbf{Example 206}. Hence the affine scheme $X=\mathrm{Spec}A$ represents a smooth curve.
        \item By \textbf{Prop. 187}, the polynomial ring $k[x_1,\cdots,x_n]$ is $n$-dimensional, and its localisation at any prime ideal is regular by \textbf{Example 206}. The resulting $n$-dimensional and smooth affine scheme $\mathrm{Spec}\ k[x_1,\cdots,x_n]=\mathbb{A}_k^n$ is called the \textbf{affine space} over $k$.
    \end{itemize}
\end{summary}

\begin{proposition}{211}
    Let $(A,\mathfrak{m})$ be a Noetherian local domain of dimension $1$. Then T.F.A.E.:\\ \centerline{(a) $A$ is a DVR.\qquad (b) $A$ is integrally closed.\qquad (c) $A$ is regular. [Exercise]}
\end{proposition}

\begin{example}{212}
    For a nonregular example, take $A$ to be the localisation of the ring
    \begin{align*}
        \mathbb{C}[x,y]/(y^2-x^3)
    \end{align*}
    at the maximal ideal $(\overline{x},\overline{y})$. Then $\frac{\overline{x}}{\overline{y}}\in\mathrm{Quot}A$ is integral over $A$, (consider $t^2-\overline{x}\in A[t]$), but not integral over $A$. Hence $A$ is \textbf{singular}, i.e. not regular.
\end{example}
\newpage
To highlight the interplay between algebra and geometry, let us again consider the ring of integers of a number field. This has the desirable property of being Dedekind. But already a ring as simple and natural as $\mathbb{Z}[\sqrt{5}]$ is not algebraically closed in $\mathbb{Q}(\sqrt{5})$. On the other hand, the curve as in \textbf{Example 212} is singular and thus not given by a Dedekind ring. We therefore need to enlarge our class of rings.

\begin{definition}{213. (Order)}
    Let $K$ be an algebraic number field with $n=[K:\mathbb{Q}]$. An \textbf{order} of $K$ is a subring $\mathcal{O}$ of $\mathcal{O}_K$ which contains an integral basis of length $n$. The ring $\mathcal{O}_K$ is obviously an order and is called the \textbf{maximal order} of $K$.
\end{definition}

\begin{proposition}{214}
    An order $\mathcal{O}$ of an algebraic number field $K$ is a $1$-dimensional Noetherian integral domain.
\end{proposition}
\begin{proof}
    Since $\mathcal{O}$ is a finitely generated $\mathbb{Z}$-module of rank $n=[K:\mathbb{Q}]$ and $\mathbb{Z}$ is Noetherian, it follows that every ideal $\mathfrak{a}\triangleleft\mathcal{O}$ is also a finitely generated $\mathbb{Z}$-module. In particular, it is a finitely generated $\mathcal{O}$-module so that $\mathcal{O}$ is Noetherian. Now, if $(0)\neq\mathfrak{p}\triangleleft\mathcal{O}$ is prime and $0\neq a\in\mathfrak{p}\cap\mathbb{Z}$, then $a\mathcal{O}\subset\mathfrak{p}\subset\mathcal{O}$. Thus, by rescaling the lattice base appropriately, we obtain a $\mathbb{Z}$-basis of $\mathfrak{p}$ which is therefore of rank $n$. Hence $\mathcal{O}/\mathfrak{p}$ is finite and an integral domain, thus a field (since the map $f_{\overline{a}}:\mathcal{O}/\mathfrak{p}\to\mathcal{O}/\mathfrak{p},\ \overline{b}\mapsto\overline{a}\cdot\overline{b}$ is bijective). Consequently, $\mathfrak{p}\triangleleft\mathcal{O}$ is maximal. As $\mathcal{O}$ is an integral domain, it follows that $\mathrm{dim}\ \mathcal{O}=1$.
\end{proof}

\begin{proposition}{215. (Akizki-Krull)}
    Let $A$ be a $1$-dimensional Noetherian domain, $K=\mathrm{Quot}A$, and $L/K$ be a finite extension. If $B=\overline{A}$ is the integral closure of $A$ in $L$, then $B$ is a Dedekind domain. [Exercise]
\end{proposition}
\begin{proof}
    We need to show the following three conditions:\\
    \textsc{claim 1.} \textit{$B$ is integrally closed.} Clearly, since it is the integral closure of $A$ in $L$.\\
    \textsc{claim 2.} \textit{Every nontrivial prime ideal is maximal.} Since $A\subset B$ is an integral extension of integral domain, we have $\mathrm{dim}\ B=\mathrm{dim}\ A=1$. In particular, if $(0)\neq\mathfrak{q}\triangleleft B$, then $\mathfrak{q}$ must be maximal, otherwise we have the strict chain $(0)\subsetneq\mathfrak{p}\subsetneq\mathfrak{n}$ in $B$ so that $\mathrm{dim}\ B>1.\ (\to\gets)$.\\
    \textsc{claim 3.} \textit{$B$ is Noetherian.} Since $\mathrm{Quot}B=L$, we can pick a $K$-basis $\{b_i\}_{i=1}^n$ of $L$ which is contained in $B$. Then $A[b_1,\cdots,b_n]$ is Noetherian (as $A$ is), and is $1$-dimensional (since $A\subset A[b_1,\cdots,b_n]$ is an integral extension). Replacing $A$ by $A[b_1,\cdots,b_n]$ if necessary, we may assume that $L=K$. In particular, it follows that $\mathfrak{b}\cap A\neq(0)$ for any ideal $\mathfrak{b}\triangleleft B$, (if $x=\frac{r}{s}\in\mathfrak{b}$ with $r,s\in\ A$, then $sx\in\mathfrak{b}\cap A$). Now, it remains to show that $\mathfrak{b}$ is a finitely generated $B$-module. Let $0\neq b\in\mathfrak{b}\cap A$. Since $B/(b)B$ is a finitely generated $A$-module, it is Noetherian. But then $\mathfrak{b}/(b)B$ is a Noetherian $A$-module as an $A$-submodule, and a fortiori it is Noetherian, and thus a finitely generated $B$-module. Since $(b)B$ is a finitely generated $B$-module, so is $\mathfrak{b}$.
\end{proof}

\begin{corollary}{216. (Normalisation)}
    The integral closure $\mathcal{O}$ of an order $o$ inside its ring of fractions is a Dedekind ring. Such $\mathcal{O}$ is called the \textbf{normalisation} of $o$.
\end{corollary}

\begin{nt}{217}
    Let $o$ be a $1$-dimensional Noetherian domain. Singular points of $\mathrm{Spec}\ o$ are by definition primes $\mathfrak{p}$, where the localisation $o_\mathfrak{p}$ is not regular. Passing to the normalisation yields an integral extension $o\xhookrightarrow{} \mathcal{O}$ and thus a surjective morphism of curves $\mathrm{Spec}\ \mathcal{O}\to\mathrm{Spec}\ o$. \\
    Since $\mathcal{O}$ is Dedekind, we may think of this morphism as a \textit{resolution of singularities}. The different points $\mathfrak{P}_i\in\mathrm{Spec}\ \mathcal{O}$ which are mapped to a singular point $\mathfrak{p}\in\mathrm{Spec}\ o$ are given by the decomposition
    \begin{align*}
        \mathfrak{p}\mathcal{O}\ \coloneqq\ \mathfrak{P}_1^{e_1}\cdots\mathfrak{P}_r^{e_r}.
    \end{align*}
\end{nt}

\newpage

\section{Exercises}
\begin{exe}{1. (Prime ideals avoiding multiplicative subsets)}
    Let $S$ be a multiplicative subset of an integral domain $A$ and $\mathfrak{a}\triangleleft A$ an ideal with $\mathfrak{a}\cap S=\varnothing$. Then $\mathfrak{a}$ is contained in a prime ideal $\mathfrak{p}\triangleleft A$ with $\mathfrak{p}\cap A=\varnothing$.
\end{exe}
\begin{proof}
    We consider the set 
    \begin{align*}
        \Sigma\coloneqq\{\mathfrak{q}\triangleleft A:\mathfrak{q}\supset\mathfrak{a}\ \text{and}\ \mathfrak{q}\cap S=\varnothing\}
    \end{align*}
    partially ordered by set inclusion. Then $\Sigma$ is nonempty since it contains $\mathfrak{a}$, and it is bounded above by $A$. Thus the \textbf{Zorn's Lemma} (cf. [G,2.16-17]) gives the existence of such an ideal $\mathfrak{p}\in\Sigma$ which is maximal. 
    \textbf{Claim.} $\mathfrak{p}\triangleleft A$ is prime. If not, then there exist $a,b\notin\mathfrak{p}$ with $ab\in\mathfrak{p}$. By maximality of $\mathfrak{p}$, the ideals $\mathfrak{p}+(a)$ and $\mathfrak{p}+(b)$ must intersect $S$. We let $p+ra$ and $q+sb$ be the elements in the intersections, respectively, where $p,q\in\mathfrak{p},\ r,s\in A$. Since $S$ is multiplicative, we have
    \begin{align*}
        p'+rsab\coloneqq(pq+p\cdot sb+q\cdot ra)+rsab=(p+ra)(q+sb)\in S.
    \end{align*}
    Now that $p'=pq+(sb)p+(ra)q\in\mathfrak{p}$ and $(rs)ab\in\mathfrak{p}$ by assumption, we have $p'+rsab\in\mathfrak{p}$, i.e. $\mathfrak{p}\cap S\neq\varnothing.\ (\to\gets)$. Hence $\mathfrak{p}\triangleleft A$ must be prime.
\end{proof}

\begin{exe}{2. (The Zariski Topology)}
    Let $A$ be a general ring. Then
    \begin{itemize}
        \item[(a)] $\mathcal{V}((0))=\mathrm{Spec}A$ and $\mathcal{V}(A)=\varnothing$.
        \item[(a*)] If the ideals $\mathfrak{a}\subset\mathfrak{b}$, then $\mathcal{V}(\mathfrak{a})\supset\mathcal{V}(\mathfrak{b})$.
        \item[(b)] If $\{\mathfrak{a}_i\}_{i\in I}$ is any family of ideals of $A$, then $\mathcal{V}(\sum_{i\in I}\mathfrak{a}_i)=\bigcap_{i\in I}\mathcal{V}(\mathfrak{a}_i)$.
        \item[(c)] $\mathcal{V}(\mathfrak{a}\cap\mathfrak{b})=\mathcal{V}(\mathfrak{a}\cdot\mathfrak{b})=\mathcal{V}(\mathfrak{a})\cup\mathcal{V}(\mathfrak{b})$ for any two ideals $\mathfrak{a},\mathfrak{b}\triangleleft A$.
        \item[(d)] $\mathrm{Spec}(A/\mathfrak{a})\cong\mathcal{V}(\mathfrak{a})$ and $\mathrm{Spec}(S^{-1}A)=\{\mathfrak{p}\in\mathrm{Spec}A:\mathfrak{p}\cap S=\varnothing\}$. (For the last assertion assume that $A$ is an integral domain).
    \end{itemize}
\end{exe}
\begin{proof}
    (a) $\mathcal{V}((0))=\{\mathfrak{p}\in\mathrm{Spec}A:\mathfrak{p}\supset(0)\}=\mathrm{Spec}A$ and $\mathcal{V}(A)=\{\mathfrak{p}\in\mathrm{Spec}A:\mathfrak{p}\supset A\}=\varnothing$.\\
    (a*) Clearly, if $q\in\mathcal{V}(\mathfrak{b})$, then $\mathfrak{q}\supset\mathfrak{b}\supset\mathfrak{a}$, whence $\mathfrak{q}\in\mathcal{V}(\mathfrak{a})$.\\
    (b) We have
    \begin{align*}
        \mathfrak{p}\in\mathcal{V}(\sum_{i\in I}\mathfrak{a}_i)\ \Longleftrightarrow\  \mathfrak{p}\supset\mathfrak{a}_i,\ \forall i\in I\ \Longleftrightarrow\  \mathfrak{p}\in\bigcap_{i\in I}\mathcal{V}(\mathfrak{a}_i). 
    \end{align*}
    (c) First, we note that $\mathfrak{a}\cdot\mathfrak{b}\subset\mathfrak{a}\cap\mathfrak{b}$, by (a*), we have $\mathcal{V}(\mathfrak{a})\cup\mathcal{V}(\mathfrak{b})\subset\mathcal{V}(\mathfrak{a}\cdot\mathfrak{b})\subset\mathcal{V}(\mathfrak{a}\cap\mathfrak{b})$. Moreover, since $\mathfrak{p}\triangleleft A$ is a prime, $\mathfrak{p}\supset\mathfrak{a}\cap\mathfrak{b}$ implies either $\mathfrak{p}\supset\mathfrak{a}$ or $\mathfrak{p}\supset\mathfrak{b}$ by \textbf{Lemma 22}. Hence, we have
    \begin{align*}
        \mathcal{V}(\mathfrak{a})\cup\mathcal{V}(\mathfrak{b})\subset\mathcal{V}(\mathfrak{a}\cdot\mathfrak{b})\subset\mathcal{V}(\mathfrak{a}\cap\mathfrak{b})\subset\mathcal{V}(\mathfrak{a})\cup\mathcal{V}(\mathfrak{b}).
    \end{align*}
    (d) The first part $\mathrm{Spec}(A/\mathfrak{a})\cong\{\mathfrak{p}\in\mathrm{Spec}A:\mathfrak{p}\supset\mathfrak{a}\}=\mathcal{V}(\mathfrak{a})$ follows from the \textbf{Comparison Lemma 16}. The second part follows from \textbf{Prop. 27} under the natural map $\mathrm{Spec}(S^{-1}A)\to\mathrm{Spec}A,\ \mathfrak{q}\mapsto\mathfrak{q}^c$ induced by the inclusion $A\xhookrightarrow{} S^{-1}A$.
\end{proof}

\begin{exe}{3. (Closed Points)}
    Let $X\subset\mathrm{Spec}A$. Show that $\mathcal{V(I}(X))=\overline{X}$ and conclude that $\overline{\{\mathfrak{p}\}}=\mathcal{V}(\mathfrak{p})$ for any $\mathfrak{p}\in\mathrm{Spec}A$. In particular, $\mathfrak{p}$ is a closed point iff $\mathfrak{p}$ is maximal. 
\end{exe}
\begin{proof}
    (i)"$\supset$"\ If $\mathfrak{p}\in X$, then $\mathfrak{p}\supset\bigcap_{\mathfrak{q}\in X}\mathfrak{q}=\mathcal{I}(X)$, whence $\mathfrak{p}\in\mathcal{V(I}(X))$. Hence $X\subset\mathcal{V(I}(X))$. Moreover, since $\mathcal{V(I}(X))$ is closed, we have $\overline{X}\subset\mathcal{V(I}(X))$.\\
    "$\subset$"\ Let $Y\subset\mathrm{Spec}A$ be any closed set containing $X$, then $Y=\mathcal{V}(\mathfrak{a})\supset X$ for some $\mathfrak{a}\triangleleft A$. If $\mathfrak{p}\in X\subset Y$, then $\mathcal{I}(X)=\bigcap_{\mathfrak{q}\in X}\mathfrak{q}\supset\bigcap_{\mathfrak{q}\in Y}\mathfrak{q}=\mathcal{I}(Y)\supset\mathfrak{a}$. By \textbf{Ex. 2(a*)}, we have $\mathcal{V(I}(X))\subset\mathcal{V}(\mathfrak{a})=Y$. Thus, by definition of "closure", $Y=\overline{X}$.\\
    (ii) Since $\mathcal{I}(\{\mathfrak{p}\})=\mathfrak{p}$, from the previous part, we have $\mathcal{V}(\mathfrak{p})=\mathcal{V(I}(\{\mathfrak{p}\})=\overline{\{\mathfrak{p}\}}$.
\end{proof}

\begin{exe}{4. (Basic Open Sets for the Zariski Topology)}
    For each $a\in A$, let $X_a$ denote the complement of $\mathcal{V}(a)\coloneqq\mathcal{V}((a))$ in $X=\mathrm{Spec}A$. In particular, $X_a$ is open and defines the so-called \textbf{basic open set}. Show that 
    \begin{itemize}
        \item[(a)] $\{X_a\}_{a\in A}$ forms a basis of open sets for the Zariski topology, i.e. any open set is a union of open sets of the form $X_a$.
        \item[(b)] $X_a\cap X_b=X_{ab}$.
        \item[(c)] $X_a=\varnothing.\ \Longleftrightarrow\ a$ is nilpotent.
        \item[(d)] $X_a=\mathrm{Spec}A.\ \Longleftrightarrow\ a$ is a unit.
        \item[(e)] $X_a=X_b.\ \Longleftrightarrow\ \sqrt{(a)}=\sqrt{(b)}$.
        \item[(f)] $X$ is quasi-compact, i.e. every open covering of $X$ has a finite subcovering.
    \end{itemize}
    \textbf{Remark.} In addition to "quasi-compact", "compact" spaces are usually required to be Hausdorff.
\end{exe}
\begin{proof}
    \textsc{note.} By definition, $X_a=X\setminus\mathcal{V}(a)=\{\mathfrak{p}\in\mathrm{Spec}A:f\notin\mathfrak{p}\}$.\\ (a) It follows from $\mathcal{V}(\mathfrak{a})=\bigcap_{a\in\mathfrak{a}}\mathcal{V}(a)$ by taking complements.\\ (b) $\because \mathcal{V}(a)\cup\mathcal{V}(b)\myeq{\text{Ex.2(c)}} \mathcal{V}(ab).\ \therefore X_a\cap X_b=X\setminus [\mathcal{V}(a)\cup\mathcal{V}(b)]=X\setminus\mathcal{V}(ab)=X_{ab}$.\\
    (c) \textsc{recall.} $\mathfrak{N}\coloneqq$ set of all nilpotents of $A$, is an ideal called the \textbf{nilradical} of $A$. Then,
    \begin{align*}
        X_a=\varnothing\ \Leftrightarrow\ \mathcal{V}(f)=X\ \Leftrightarrow\ f\in\mathfrak{p},\ \forall\mathfrak{p}\in X\ \Leftrightarrow\ f\in\bigcap_{\mathfrak{p}\in X}\mathfrak{p}\myeq{\text{[AM, 1.8]}} \mathfrak{N}.
    \end{align*}
    (d) We have
    \begin{align*}
        X_f=X\ \Leftrightarrow\ \mathcal{V}(f)=\varnothing\ \Leftrightarrow\ f\notin\mathfrak{p},\ \forall\mathfrak{p}\in X.\ \xLeftrightarrow{~\ast~}\ f\in A^\times.
    \end{align*}
    \textsc{proof of $\ast$.} "$\Rightarrow$"\ Assume $f\notin A^\times$. By [AM, 1.5], $\exists\mathfrak{m}\triangleleft A$ maximal with $f\in\mathfrak{m}.\ (\to\gets)$. \\ "$\Leftarrow$"\ Assume $\exists\mathfrak{p}\in X$ with $f\in\mathfrak{p}.\ \because f\in A^\times.\ \therefore\mathfrak{p}=A.\ (\to\gets)$.\\
    (e) We have 
    \begin{align*}
        X_a=X_b\ \Leftrightarrow\ \mathcal{V}(a)=\mathcal{V}(b)\ \xLeftrightarrow{~\ast~}\ \sqrt{(a)}=\sqrt{(b)}.
    \end{align*}
    \textsc{proof of $\ast$.} "$\Rightarrow$"\ Clearly, by \textbf{Prop. 170}, $\sqrt{(a)}=\mathcal{I(V}(a))=\mathcal{I(V}(b))=\sqrt{(b)}$.\\ "$\Leftarrow$"\ Assume that there is a prime $\mathfrak{p}\triangleleft A$ with $(a)\subset\mathfrak{p}\subset\sqrt{(a)}$. If $x\in\sqrt{(a)}$, then $x^n\in(a)\supset\mathfrak{p}$ for some $n\in\mathbb{N}_{\geq1}$, whence $x\in\mathfrak{p}$. It follows that $\mathfrak{p}=\sqrt{(a)}$. By [AM, Ex.1.15i], we have $\mathcal{V}(a)=\mathcal{V}(\sqrt{(a)})=\mathcal{V}(\sqrt{(b)})=\mathcal{V}(b)$.\\
    (f) \textsc{claim.} $X_a$ is quasi-compact, for any $a\in A$. Given an open cover $\{X_{a_i}\}_{i\in I}$ of $X_a$. For each $i,\ \mathfrak{a}_i\coloneqq(a_i)$, then 
    \begin{align*}
        X_a\subset\bigcup_iX_{a_i} &\Leftrightarrow\  \mathcal{V}(\sum_i\mathfrak{a}_i)=\bigcap_i\mathcal{V}(\mathfrak{a}_i)\subset\mathcal{V}(f)\ \xLeftrightarrow{~\text{(e)}~}\ \sqrt{(a)}\subset\sqrt{\sum_i\mathfrak{a}_i}\\
        &\xLeftrightarrow{~\text{Def.}~} \exists m\geq1:\ a^m\in\sum_i\mathfrak{a}_i\ \xLeftrightarrow{~\text{finite sum}~}\ \exists m\geq1:\ a^m\coloneqq\sum_{j\in J}r_ja_j,\ |J|<\infty.
    \end{align*}
    If there is a prime $\mathfrak{p}\in\mathrm{Spec}A$ with $a_j\in\mathfrak{p},\ \forall j\in J$, then $a^m\in\mathfrak{p}.\ \because \mathfrak{p}\triangleleft A$ is prime. $\therefore a\in\mathfrak{p}$. Equivalently, if $a\notin\mathfrak{p}$, then $a_j\notin\mathfrak{p}$ for some $j\in J$, whence $X_f\subset\bigcup_{j\in J}X_{a_i}$. Hence $X_a$ is quasi-compact. We take $a=1$, then $X=X_1$ is quasi-compact.
\end{proof}

\begin{exe}{5. (The Spectrum of $A_f$)}
    Let $f\in A$ and $S_f\coloneqq\{f^n\}_{n\geq0}$ the multiplicative system. We denote $A_f\coloneqq S_f^{-1}A$ and $X=\mathrm{Spec}A$. Show that 
    \begin{itemize}
        \item[(a)] $X_f$ is homeomorphic to $\mathrm{Spec}A_f$, that is, the basic open sets have the same topology as affine schemes.
        \item[(b)] $\mathcal{O}_X(X_f)=A_f$. Conclude that $\mathcal{O}_X(X)=A$.
    \end{itemize}
    \textbf{Remark.} In fact, $X_f$ is a basic example of a scheme for which one can show that it is not only homeomorphic, but isomorphic as a scheme to $\mathrm{Spec}A_f$.
\end{exe}
\begin{proof}
    (a) Consider the natural map $\phi:A\to A_f,\ a\mapsto \frac{a}{1}$. Then the associated map
    \begin{align*}
        \phi^\ast:\mathrm{Spec}A_f\longrightarrow\mathrm{Spec}A,\ \mathfrak{q}\longmapsto\mathfrak{q}^c    
    \end{align*}
    is injective and has image $X_f=\{\mathfrak{p}\in\mathrm{Spec}A:f\notin\mathfrak{p}\}$. Now, the map
    \begin{align*}
        \psi:X_f\longrightarrow\mathrm{Spec}A_f,\ \mathfrak{p}\longmapsto\mathfrak{p}^e=\mathfrak{p}A_f
    \end{align*}
    is the inverse to $\phi^\ast$. It remains to show that $\psi$ is continuous. Let $\mathcal{V}(\mathfrak{a})\subset\mathrm{Spec}A_f$ be a closed set, then
    \begin{align*}
        \psi^{-1}(\mathcal{V}(\mathfrak{a}))=\{\mathfrak{p}\in X_f:\mathfrak{p}\supset\mathfrak{a}^c\}=X_f\cap\mathcal{V}(\mathfrak{a}^c),
    \end{align*}
    i.e. $\psi^{-1}(\mathcal{V}(\mathfrak{a}))$ is closed in $X_f$, whence $\psi$ is continuous.\\ Indeed, if $\phi(\mathfrak{p})\in\mathcal{V}(\mathfrak{a})$ for some $\mathfrak{p}\notni f$, then $\mathfrak{a}\subset\mathfrak{p}^e$, whence $\mathfrak{a}^c\subset\mathfrak{p}^{ec}=\mathfrak{p}$. Conversely, if $\mathfrak{a}^c\subset\mathfrak{p}$ for some $\mathfrak{p}\notni f$, then $\psi(\mathfrak{p})=\mathfrak{p}^e\supset\mathfrak{a}^{ce}=\mathfrak{a}$.\\
    (b) "$\supset$"\ Note that for any $\mathfrak{p}\in X_f$, we have $\frac{a}{f^n}\in A_\mathfrak{p}$, whence $\frac{a}{f^n}\in\bigcap_{\mathfrak{p}\in X_f}A_\mathfrak{p}=\mathcal{O}_X(X_f)$.\\ "$\subset$"\ Suppose that $F\in\mathcal{O}_X(X_f)=\bigcap_{\mathfrak{p}\notni f}A_\mathfrak{p}$, and consider the ideal $\mathfrak{a}\coloneqq\{a\in A:aF\in A\}$. Let $\mathfrak{p}\in X_f$, then $F\coloneqq\frac{a}{s}\in A_\mathfrak{p},\ (a\in A,\ s\notin\mathfrak{p})$, whence $bF=a$, i.e. $b\in\mathfrak{a}$. It follows that if $\mathfrak{p}\notni f$, then $\mathfrak{p}\not\supset\mathfrak{a}$, which means $\mathcal{V}(\mathfrak{a})\subset\mathcal{V}(f)$. By \textbf{Ex. 2(e)}, we have $\sqrt{(f)}\subset\sqrt{\mathfrak{a}}.\ \Rightarrow f^n\in\mathfrak{a}$ for some $n$, hence $F\in A_f$. By taking $f=1$, we obtain $X=X_1$ and $A=A_1$, thus the result follows.
\end{proof}

\begin{exe}{6. (The Spectrum of $\mathbb{Z}[i]$)}
    Compute explicitly the continuous map $f:\mathrm{Spec}\ \mathbb{Z}[i]\to\mathrm{Spec}\ \mathbb{Z}$ from the scheme morphism $F=(f,f^\ast)$ induced by the inclusion $\mathbb{Z}\xhookrightarrow{} \mathbb{Z}[i]$.
\end{exe}

\begin{exe}{7. (The Associated Map with an Integral Extension)}
    Let $\iota:A\xhookrightarrow{} B$ be an integral ring extension (considering $\iota$ as an inclusion). Then the associated map
    \begin{align*}
        \iota^\ast:\mathrm{Spec}B\longrightarrow\mathrm{Spec}A,\ \mathfrak{p}\longmapsto\mathfrak{p}\cap A
    \end{align*}
    is a closed map, i.e. it maps closed sets to closed sets.
\end{exe}

\begin{exe}{8. (Scheme Morphisms)}
    Let $B$ be a $k$-algebra of finite type with $n=\mathrm{dim}\ B$. Then there exists a scheme morphism $F=(f,f^\ast)$ such that $f$ is a closed surjective map.
\end{exe}

\begin{exe}{9. (Inserting Prime Ideals into Chains)}
    Let $A\subset B$ be an integral extension where $A$ is a finitely generated $k$-algebra. Assume there are given prime ideals $\mathfrak{q}_0$ and $\mathfrak{q}_2$ of $B$ and let $\mathfrak{p}_0\coloneqq\mathfrak{q}_0\cap A$ and $\mathfrak{p}_2\coloneqq\mathfrak{q}_2\cap A$. If $\mathfrak{p}_1$ is a prime ideal of $A$ with $\mathfrak{p}_0\subsetneq\mathfrak{p}_1\subsetneq\mathfrak{p}_2$, then there exists a prime ideal $\mathfrak{q}_1$ in $B$ with $\mathfrak{q}_0\subsetneq\mathfrak{q}_1\subsetneq\mathfrak{q}_2$.
\end{exe}

\begin{exe}{10. (Modifying Chains of Prime Ideals)}
    Let $n\geq1$ and $\mathfrak{p}_0\subsetneq\mathfrak{p}_1\subsetneq\cdots\subsetneq\mathfrak{p}_n$ be a chain of prime ideals in a Noetherian ring $A$. If $a\in\mathfrak{p}_n$, then there exists a chain of prime ideals $\mathfrak{p}'_0\subsetneq\mathfrak{p}'_1\subsetneq\cdots\subsetneq\mathfrak{p}'_{n-1}\subsetneq\mathfrak{p}_n$ with $a\in\mathfrak{p}'_1$.
\end{exe}

\begin{exe}{11. (Codimension \& Generators)}
    Let $A$ be Noetherian and $a_1,\cdots,a_n\in A$. Then every minimal prime ideal $\mathfrak{p}$ over $(a_1,\cdots,a_n)$ satisfies $\mathrm{height}\ \mathfrak{p}\leq n$.
\end{exe}

\chapter{Completions \& Local-Global Principles}
\begin{motivation}{218}
    Consider a polynomial $f=\frac{f}{1}\in\mathbb{C}[x]_{(x-a)}$ given by 
    \begin{align*}
        f(x)=a_0+a_1(x-a)+\cdots+a_n(x-a)^n.
    \end{align*}
    At $\mathfrak{p}=(x-a)$ we assign to $f$ the "value" $a_o\equiv f\ \mathrm{mod}\ \mathfrak{p}$. Moding out higher powers $(x-a)^n$ yields the higher order terms of $f$ corresponding to the successive derivatives of $f$ at $\mathfrak{p}$. More generally, we can expand a rational function $\frac{f}{g}$ with $g\not\equiv0\ \mathrm{mod}\ \mathfrak{p}$ into a Laurent series $\sum_{i=m}^\infty a_i(x-a)^i,\ (m\in\mathbb{Z})$, (for the moment this notation is purely formal and stands for the sequence $(s_n)$ of partial sums $s_n\coloneqq\sum_{i=0}^na_i(x-a)^i$).\\
    Similarly, consider $f\in\mathbb{Z}_{(p)}=\{\frac{a}{b}\in\mathbb{Q}:p\nmid b\}$ for a prime $p\in\mathbb{Z}$. By division algorithm we have,  
    \begin{align*}
        f=pf_1+a_0,\ f_1=pf_2+a_1,\ \cdots,\ f_n=a_n,\qquad (0\leq a_i\leq p-1)
    \end{align*}
    gives the $p$-adic expansion $f=a_0+a_1p+\cdots+a_np^n$. More generally, we can expand rational numbers into Laurent series, leading eventually to $p$-\textbf{adic numbers}.\\
    These examples are obtained by a process called \textbf{completion}. Morally, this allows us to focus on local information of global data in a sense that we yet have to specify. An informal example is provided by the ring of holomorphic functions $\mathcal{O}(U)$ or the field of meromomorphic fnctions $\mathcal{M}(U)$ on an open subset $U\subset\mathbb{C}$. We should think of these as "global" objects. Locally, we can develop these functions into power or Laurent series. This is a powerful tool in the theory of functions which we want to make available in our algebraic context.
\end{motivation}

\section{$p$-Adic Integers}
In this section $p$ always denotes a positive prime number in $\mathbb{Z}$.
\begin{definition}{219. ($p$-Adic Integer)}
    \newline A $p$-\textbf{adic integer} is a formal series 
    \begin{align*}
        a_0+a_1p+a_2p^2+\cdots+,\ (0\leq a_i<p).
    \end{align*}
    Here, $\mathbb{Z}_p$ denotes the set of all $p$-adic integers.
\end{definition}

\begin{lemma}{220}
    For $f\in\mathbb{Z}$, the residue classes $f\ \mathrm{mod}\ p^n$ have a unique decomposition of the form
    \begin{align*}
        f\equiv a_0+a_1p+a_2p^2+\cdots+a_{n-1}p^{n-1}\ \mathrm{mod}\ p^n
    \end{align*}
    with each $0\leq a_i<p$.
\end{lemma}
\begin{proof}
    Induction on $n$. For $n=1$,  we have
    \begin{align*}
        f\equiv a_0\ \mathrm{mod}\ p\ \text{for a unique}\ 0\leq a_0<p.
    \end{align*}
    For $n>1$, assume that there are unique $a_i$'s with $0\leq a_i<p$ such that
    \begin{align*}
        f\equiv a_0+a_1p+\cdots+a_{n-1}p^{n-1}\qquad\mathrm{mod}\ p^n.
    \end{align*}
    Then $f=a_0+a_1p+\cdots+a_{n-1}p^{n-1}+bp^n$ for some $b\in\mathbb{Z}$, whence $\exists!\ 0\leq a_n<p$ such that $b=a_n+p\cdot c,\ (c\in\mathbb{Z})$. By moding $p^{n+1}$, we have 
    \begin{align*}
        f\equiv a_0+a_1p+\cdots+a_{n-1}p^{n-1}+a_np^n\qquad\mathrm{mod}\ p^{n+1}.
    \end{align*}
\end{proof}

\begin{definition}{221. ($p$-Adic Expansion)}
    For $f\in\mathbb{Z}$ the decomposition of the previous lemma gives rise to the sequence 
    \begin{align*}
        s_1=a_0,\ s_2=a_0+a_1p,\ s_3=a_0+a_1p+a_2p^2,\ \cdots.
    \end{align*}
    We succinctly write $\sum_{i=0}^\infty a_ip^i$ for the sequence $(s_i)$ and call this expression the $p$-\textbf{adic expansion} of $f$.\\
    Note that $p\nmid f$ implies that $\gcd{(f,p^n)}=1$ so that $a\cdot f+b\cdot p^n=1$ for some $a,b\in\mathbb{Z}$, and thus $a\cdot f\equiv1\ \mathrm{mod}\ p^n$, or equivalently, $f^{-1}\equiv a\ \mathrm{mod}\ p^n$. Hence we can assign a development to $f^{-1}$ as before. In particular, this yields an expansion for every element in $\mathbb{Z}_p$, (cf. Exercise).
\end{definition}

\begin{definition}{222. ($p$-Adic Number)}
    \newline A $p$-\textbf{adic number} is a formal Laurent series $\sum_{i=-m}^\infty a_ip^i$ with $m\in\mathbb{Z},\ 0\leq a_i<p$. We then write $\mathbb{Q}_p$ for the set of all $p$-adic numbers.
\end{definition}

\begin{remark}{223}
    For $f\in\mathbb{Q}$, we write $f=\frac{g}{h}\cdot p^{-m},\ (m,g,h\in\mathbb{Z})$ with $\gcd{(gh,p)}=1$. Then $\frac{g}{h}\in\mathbb{Z}_p$ has an expansion $\sum_{i=0}^\infty a_ip^i$ and we assign to $f$ the expansion $\sum_{i=-m}^\infty a_{i+m}p^i$.\\
    It follows that we have an injection $\mathbb{Q}\xhookrightarrow{} \mathbb{Q}_p$ which restricts to $\mathbb{Z}\xhookrightarrow{} \mathbb{Z}_p$, (this is not a morphism yet since we defined $\mathbb{Z}_p$ and $\mathbb{Q}$ as sets, not as rings or fields). In this way we obtained an algebraic counterpart for the development of a holomorphic or meromomorphic function into a power or Laurent series.
\end{remark}

\begin{example}{224}
    \begin{itemize}
        \item[(a)] $3$-adic expansion of $216$:
        \begin{align*}
            216=3^3\cdot8,\qquad 8=2+3\cdot2,
        \end{align*}
        so that $216=2\cdot3^3+2\cdot3^4$. This is sometimes also written as $0,0022$ \textemdash the higher the power, the smaller the contribution. Similarly, we find for the $5$-adic expansion
        \begin{align*}
            216=1+5\cdot43,\qquad 43=3+5\cdot8,\qquad 8=3+5,
        \end{align*}
        so we have $216=1,331$.
        \item[(b)] We have $-1=\sum_{i=0}^\infty(p-1)p^i$, for $-1\equiv p-1+p(p-1)\ \mathrm{mod}\ p^2$, etc.
        \item[(c)] We have
        \begin{align*}
            \tfrac{1}{p-1}\equiv1+p+\cdots+p^{n-1}\qquad \mathrm{mod}\ p^n
        \end{align*}
        for $1=(1+p+\cdots+p^{n-1})(p-1)+p^n$.
    \end{itemize}
\end{example}

The construction of $\mathbb{Z}_p$ provides an example for the \textit{completion} of a ring. We are now going to investigate this notion systematically.
\begin{definition}{225. (Null \& Cauchy Sequence)}
    \begin{itemize}
        \item Let $A$ temporarily denote an abelian topological group, i.e. $A$ carries a topology for which addition and taking inverses are continuous operations.
        \item Since the translation $T_b:A\to A,\ a\mapsto a+b$ for a fixed $b\in A$ defines a homeomorphism, the topology of $A$ is determined by the neighbourhoods of $0$. 
        \item In addition, we assume that there is a countable fundamental system $\{U_j\}_{j\in\mathbb{N}}$ of neighbourhoods of $0$, i.e. any open set containing $0$ eventually contains some $U_j$.
        \item A \textbf{null sequence} $(a_n)$ in $A$ is a sequence such that for all $j\in\mathbb{N}$ there exists $N(j)\in\mathbb{N}$ with $a_n\in U_j$ provided $n\geq N(j)$. 
        \item A \textbf{Cauchy sequence} $(a_n)$ in $A$ is a sequence such that for all $j\in\mathbb{N}$ there exists $N(j)\in\mathbb{N}$ with $a_n-a_m\in U_j$ provided $n,m\geq N(j)$. Moreover, two Cauchy sequences $(a_n)$ and $(b_n)$ are \textbf{equivalent}, write $(a_n)\sim(b_n)$, if $(a_n-b_n)$ is a null sequence.
    \end{itemize}
\end{definition}
\newpage
\begin{definition}{226. (Completions)}
    \begin{itemize}
        \item We define the \textbf{completion} of $A$ by 
        \begin{align*}
            \widehat{A}\coloneqq\ \text{Cauchy sequences in $A$}/\sim.
        \end{align*}
        \item To define a topology on $\widehat{A}$ we exhibit a fundamental system $\{\widehat{U}_j\}_{j\in\mathbb{N}}$ of neighbourhoods at $\widehat{0}$, the equivalence class of the constant sequence $(a_n)=(0)$. Namely, $U_j$ consists of equivalence classes $\overline{(a_n)}$ with $a_n\in U_j\subset A$ for all except finitely many $n\in\mathbb{N}$.
        \item Addition of Cauchy sequences makes $\widehat{A}$ a topological group. The natural map $\phi:A\to\widehat{A},\ a\mapsto\overline{(a_n)}=\overline{(a)}$, is a continuous group homomorphism. If $\phi$ is an isomorphism, we say that $A$ is \textbf{complete}.
        \item If $\phi:A\to B$ is a continuous group homomorphism, then $\phi$ maps Cauchy sequences to Cauchy sequences and induces a continuous ring homomorphism $\widehat{\phi}:\widehat{A}\to\widehat{B}$. The assignment $A\to\widehat{A},\ \phi\mapsto\widehat{\phi}$ is thus functorial.
    \end{itemize}
\end{definition}

\begin{remark}{227}
    If we start with a topological ring, then its completion is also a topological ring using the natural product of Cauchy sequences.
\end{remark}

From an algebraic point of view, a natural way to obtain a topological group is by choosing a \textit{filtration} for $A$.
\begin{definition}{228. (Filtration)}
    \newline Let $A$ be an abelian group. A \textbf{filtration}
    \begin{align*}
        A=\mathfrak{a}_0\supset\mathfrak{a}_1\supset\mathfrak{a}_2\supset\cdots
    \end{align*}
    is a descending chain of subgroups. It gives $A$ the structure of a topological group by taking $U_j=\mathfrak{a}_j$ as a system of open neighbourhoods of $0$. In particular, let $A$ be a ring and $\mathfrak{a}\triangleleft A$ an ideal. The $\mathfrak{a}$-\textbf{adic filtration} is given by $\mathfrak{a}_j\coloneqq\mathfrak{a}^j$. We denote the resulting completion by $\widehat{A}_\mathfrak{a}$.
\end{definition}
\newpage
\begin{nt}{228*}
    If $A$ is a topological group whose topology is determined by a filtration $\mathfrak{a}_0\supset\mathfrak{a}_1\supset\cdots$ and $(a_n)$ is a Cauchy sequence, then its image in $A/\mathfrak{a}_j$ is fixed, becomes eventually constant and is equal to $\xi_j\equiv a_N\ (\mathrm{mod}\ \mathfrak{a}_j)$, where $N$ is minimal among all $M$ such that $a_k-a_l\in\mathfrak{a}_j$ for all $k,l\geq M$.\\
    Note that $\xi_j\equiv a_M\ (\mathrm{mod}\ \mathfrak{a}_j)$ for all $M\geq N$. The resulting sequence is therefore \textit{coherent} in the sense that the canonical projection
    \begin{align*}
        \theta_{j+1}:A/\mathfrak{a}_{j+1}\xtwoheadrightarrow{} A/\mathfrak{a}_j,\qquad a+\mathfrak{a}_{j+1}\longmapsto a+\mathfrak{a}_j
    \end{align*}
    maps $\xi_{j+1}\mapsto\xi_j$. It follows that we can identify the completion $\widehat{A}$ with the space of coherent sequences. This motivates the following construction.
\end{nt}

\begin{definition}{229. (Inverse System \& Inverse Limit)}
    \newline Let $\{A_i\}_{i\in\mathbb{N}}$ be a family of abelian groups together with group homomorphisms $\theta_{ij}:A_j\to A_i$ such that
    \begin{itemize}
        \item $\theta_{ii}=1_{A_i}$ for all $i\in\mathbb{N}$.
        \item $\theta_{ik}=\theta_{ij}\circ\theta_{jk}:A_k\to A_i$ whenever $i\leq j\leq k$.
    \end{itemize}
    We also say that $(A_i,\ \theta_{ij})$ (or simply $A_i$ for short) defines an \textbf{inverse system}. The abelian group
    \begin{align*}
        \varprojlim_{i\in I}A_i\coloneqq\{a\in\prod_{i\in I}A_i:a_i=\theta_{ij}(a_j),\ \forall i\leq j\ \text{in}\ I\}
    \end{align*}
    is called the \textbf{inverse limit} of the family.
\end{definition}

\begin{remark}{230}
    The inverse limit comes with natural projections
    \begin{align*}
        \pi_j:\varprojlim A_i\xtwoheadrightarrow{} A_j \qquad\text{satisfying}\qquad \theta_{ij}\circ\pi_j=\pi_i.
    \end{align*}
    This even characterises the inverse limit via the following universal property: If $A$ is an abelian group together with morphisms $\psi_i:A\to A_i$ such that $\theta_{ij}\circ\psi_j=\psi_i$, then there exists a unique morphism 
    \begin{align*}
        \psi:A\to\varprojlim A_i\qquad\text{such that}\qquad \pi_j\circ\psi=\psi_j,\ \forall j\in I.
    \end{align*}
    \begin{center}
        \begin{tikzcd}
            & A \arrow[d, "\psi_j"] \\
            \varprojlim A_i \arrow[ru, leftarrow, dashed, "\exists!\ \psi", red]
            & A_j \arrow[l, leftarrow, "\pi_j"] \\
        \end{tikzcd}
    \end{center}
\end{remark}

\begin{nt}{231. (Completions as Inverse Limits)}
    \newline Let $A$ be an abelian group whose topology is defined by the filtration
    \begin{align*}
        A=\mathfrak{a}_0\supset\mathfrak{a}_1\supset\mathfrak{a}_2\supset\cdots.
    \end{align*}
    Let $\pi_{ij}:A/\mathfrak{a}_j\xtwoheadrightarrow{} A/\mathfrak{a}_i,\ (i\leq j)$, be the natural projection map. Since the completion $\widehat{A}$ with respect to this filtration equals the set of coherent sequences which is just the inverse limit of the system $\{A/\mathfrak{a}_i\}$, we thus obtain
    \begin{align*}
        \widehat{A}\ \cong\ \varprojlim A/\mathfrak{a}_i\ =\ \{a\in\prod_{i\in I}A/\mathfrak{a}_i:a_i=\pi_{ij}(a_j),\ \forall i\leq j\ \text{in}\ I\}.
    \end{align*}
\end{nt}

\begin{remark}{232}
    If $A$ is a ring and $\{\mathfrak{a}_j\}_{j\in\mathbb{N}}$ is a filtration by ideals, then restricting multiplication from the product ring $\prod_{j\in\mathbb{N}}A/\mathfrak{a}_j$ yields a ring structure in $\widehat{A}$.
\end{remark}

\begin{example}{233}
    \begin{itemize}
        \item[(a)] \textbf{$p$-Adic Integers.} Let $p$ be a prime number. Then we have
        \begin{align*}
            \mathbb{Z}_p\ \cong\ \widehat{\mathbb{Z}}_{(p)}\ \cong\ \varprojlim_{n\in\mathbb{N}}\mathbb{Z}/p^{n+1}\mathbb{Z}.
        \end{align*}
        \textsc{proof.} We obtain a bijection by associating to every $p$-adic integer $a\coloneqq\sum_{i=0}^\infty a_ip^i$ the sequence $(\overline{s_n})$ given by $\overline{s_n}\coloneqq a+(p^{n+1})\in\mathbb{Z}/p^{n+1}\mathbb{Z}$, (cf. \textbf{Lemma 220}). In particular, this induces a ring structure on $\mathbb{Z}_p$. Note that addition is done by "carrying", not "termwise". \\
        For instance, 
        \begin{align*}
            1+[p-1+ap]= \left\{\begin{array}{lc} (a+1)p &\text{if}\ a<p-1 \\ 
            p^2 &\text{if}\ a=p-1 \end{array}\right..
        \end{align*}
        \item[(b)] \textbf{Power Series.} Let $A=B[x_1,\cdots,x_n]$ be a polynomial ring over $B$ and take the maximal ideal $\mathfrak{m}=(x_1,\cdots,x_n)$. Then
        \begin{align*}
            \widehat{A}_\mathfrak{m}\ \cong\ \varprojlim_{j\in\mathbb{N}}A/\mathfrak{m}^{j+1}\ \cong\ B\llbracket x_1,\cdots,x_n\rrbracket.
        \end{align*}
        \textsc{proof.} Indeed, we obtain a bijection by sending as in the previous item 
        \begin{align*}
            B\llbracket x_1,\cdots,x_n\rrbracket &\longrightarrow\ \widehat{A}_\mathfrak{m},\\ f\coloneqq\sum_{\alpha\geq0}a_{\alpha}x^\alpha &\longmapsto\ (\overline{s_j}):\ \overline{s_j}\coloneqq  f+\mathfrak{m}^{j+1}\in A/\mathfrak{m}^{j+1}.
        \end{align*}
    \end{itemize}
\end{example}
\newpage
Now, we see that the inverse limit has good functorial properties:
\begin{definition}{234}
    \begin{itemize}
        \item An inverse system $\{A_i\}$ is called \textbf{surjective} if for any $i\leq j,\ \theta_{ij}:A_j\to A_i$ is surjective.
        \item An \textbf{exact sequence} of inverse systems
        \begin{align*}
            0\ \to\ \{A'_j\}\ \to\ \{A_j\}\ \to\ \{A''_j\}\ \to\ 0
        \end{align*}
        consists of a commutative diagram
        \begin{center}
        \begin{tikzcd}
            0\rar[] &A'_{j+1}\rar[] \dar["\theta'_{j,j+1}"]
            &A_{j+1}\rar[] \dar["\theta_{j,j+1}"]
            &A''_{j+1}\rar[] \dar["\theta''_{j,j+1}"] &0 \\
            0\rar[] &A'_j\rar[] &A_j\rar[] &A''_j\rar[] &0 \\
        \end{tikzcd}
    \end{center}
    of exact sequences where the maps $\theta'_{ij},\ \theta_{ij}$ and $\theta''_{ij}$ are provided by the inverse systems $\{A'_i\},\ \{A_i\}$ and $\{A''_i\}$, respectively.
    \end{itemize}
\end{definition}

The followings are the rather straightforward homological arguments involving the so-called "\textit{Snake Lemma}":
\begin{proposition}{235}
    If $0\to\{A'_j\}\to\{A_j\}\to\{A''_j\}\to0$ is an exact sequence of inverse systems, then taking inverse limit is a "left exact functor", i.e.
    \begin{align*}
        0\ \to\ \varprojlim A'_j\ \to\ \varprojlim A_j\ \to\ \varprojlim A''_j\
    \end{align*}
    is exact. Furthermore, if $\{A_j\}$ is a surjective system, then
    \begin{align*}
        0\ \to\ \varprojlim A'_j\ \to\ \varprojlim A_j\ \to\ \varprojlim A''_j\ \to\ 0
    \end{align*}
    is exact. [AM, 10.2]
\end{proposition}

\begin{corollary}{236}
    Let $0\to A'\to A\xtwoheadrightarrow{p} A''\to0$ be an exact sequence of groups. Let $\{\mathfrak{a}_j\}$ be a filtration of subgroups in $A$, and endow $A'$ and $A''$ with the \textbf{induced filtrations} given by $\{A'\cap\mathfrak{a}_j\}$ and $\{p(\mathfrak{a}_j)\}$, respectively. Then we have the exact sequence of groups
     \begin{align*}
        0\ \to\ \widehat{A}'\ \to\ \widehat{A}\ \to\ \widehat{A}''\ \to\ 0
    \end{align*}
    completed with respect to these filtrations. [Exercise]
\end{corollary}

\begin{corollary}{237}
    If we complete with respect to the induced filtrations, then $\widehat{\mathfrak{a}_j}$ is a subgroup of $\widehat{A}$ and $\widehat{A}/\widehat{\mathfrak{a}_j}\cong A/\mathfrak{a}_j$. In particular, completing $\widehat{A}$ with respect to $\{\widehat{\mathfrak{a}_j}\}$ yields $\widehat{\widehat{A}}\cong\widehat{A}$. [Exercise]
\end{corollary}
\newpage
Before we investigate completions further, we give an important application, see, for example, [Commutative Algebra, Eisenbud, 7.3] for a proof:
\begin{theorem}{238. (Hensel's Lemma Ver. I)}
    \newline Let $A\cong\widehat{A}_\mathfrak{m}$ be a ring that is complete with respect to the ideal $\mathfrak{m}$ and let $f\in A[x]$. If $a$ is an approximate root of $f$ in the sense that $f(a)\equiv0\ (\mathrm{mod}\ f'(a)^2\mathfrak{m})$, then there exists a root $b\in A$ of $f$ which is near $a$ in the sense that $b\equiv a\ (\mathrm{mod}\ f'(a)\mathfrak{m})$. If $f'(a)$ is not a zero-divisor in $A$, then $b$ is unique with this property. In particular, if $f'(a)$ is a unit, then $f(a)\equiv0\ (\mathrm{mod}\ \mathfrak{m})$ implies that there exists a unique $b\in A$ such that
    \begin{align*}
        f(b)=0\qquad\text{and}\qquad b\equiv a\ \mathrm{mod}\ \mathfrak{m}.
    \end{align*}
\end{theorem}

\begin{example}{239}
    \begin{itemize}
        \item[(a)] We use \textbf{Hensel's Lemma} to determine the $p$-adic integer $c$ which admits a square root in $\mathbb{Z}_p$. First write $c=p^nb$ for some $n\in\mathbb{N}$ and $b\in\mathbb{Z}_p^\times$. Hence $c$ is a square iff $n$ is even and $b$ admits a square root in $\mathbb{Z}_p$. \\
        \textsc{proof.} If $b$ is not a square in $\mathbb{Z}_p$, we are done. Otherwise, $\overline{b}=\overline{a}^2$ in $\mathbb{Z}_p/(p)\mathbb{Z}_p$ for a unit $a\in\mathbb{Z}_p^\times$. Since $\mathbb{Z}_p/(p)\mathbb{Z}_p\cong\mathbb{Z}/(p)$, (where $\widehat{(p)}\cong(p)\mathbb{Z}_p$ follows from \textbf{Lemma 296} below), we can effectively decide whether or not $\overline{b}$ is a square by computing its \textit{Legendre symbol} from elementary number theory.\\
        Then $f=x^2-b\in\mathbb{Z}_p[x]$ has a root. If $p=2$, this follows from elementary considerations, [Exercise]. Otherwise, $f'(a)=2a$ is a unit in $\mathbb{Z}_p$. Further, $f(a)\equiv0\ (\mathrm{mod}\ (p)\mathbb{Z}_p)$ by design of $a$. \textbf{Hensel's Lemma} therefore guarantees the existence of a unique root $b\equiv a\ (\mathrm{mod}\ (p)\mathbb{Z}_p)$.
        \item[(b)] If $f(t,x)\in k[t,x]$ is a polynomial in two variables, and $x=a$ is a simple root of $f(0,x)$, then there is a unique power series $x(t)\in k\llbracket t\rrbracket$ with $x(0)=a$ and $f(t,x(t))=0$ identically.\\
        \textsc{proof.} Indeed, consider $f(t,x)=f_t(x)$ as an element of $k\llbracket t\rrbracket[x]=\widehat{k[t]}_{(t)}[x]$. Since $f_t(a)\equiv f_0(a)\equiv0\ (\mathrm{mod}\ (t)\widehat{k[t]}_{(t)})$, there exists a unique $x(t)\in k\llbracket t\rrbracket$ with $f_t(x(t))=0$ and $x(t)\equiv a\ (\mathrm{mod}\ (t))$, i.e. $x(0)=a$.
    \end{itemize}
\end{example}

\begin{remark}{240}
    The condition that $f(0,a)$ is a simple root means that $\partial_xf(0,a)\neq0$. In this sense, \textbf{Hensel's Lemma} is the algebraic (and equally powerful) analogue of \textit{Implicit Function Theorem}.
\end{remark}
\newpage
\begin{example}{241}
    To illustrate the previous remark further and to emphasise the (yet to make precise) "local" nature of completions, we consider the inclusion
    \begin{align*}
        k[x]\xhookrightarrow{} A=k[x,y]/(y^2-x-1),\ x\mapsto x\ \mathrm{mod}\ (y^2-x-1).
    \end{align*}
    Taking spectra we obtain a ramified finite cover which is generically 2-1. In a neighbourhood without branching points (e.g. near $(x,y)=(0,-1)$), one should be able to invert this map and to find local sections of this covering \textemdash this is certainly trur if $k=\mathbb{R}$ or $\mathbb{C}$ using the \textit{Inverse Function Theorem}.\\
    In general, we face the problem that the map
    \begin{align*}
        y\ \mathrm{mod}\ (y^-x-1)\ \longmapsto\ \sqrt{x+1}
    \end{align*}
    is not defined at $\sqrt{x+1}\notin k[x]$. However, 
    \begin{align*}
        \sqrt{x+1}\ =\ 1+\tfrac{x}{2}-\tfrac{x^2}{8}+\cdots
    \end{align*}
    has a development in the completion $k\llbracket x\rrbracket$. Therefore, it follows from \textbf{Cor. 298} that the inverse map
    \begin{align*}
        \widehat{A}_{(\overline{x},\overline{y+1})}\cong k\llbracket x,y\rrbracket/(y^2-x-1)^e &\to k\llbracket x\rrbracket,\\
        x &\mapsto x;\\ y &\mapsto 1+\tfrac{x}{2}-\tfrac{x^2}{8}+\cdots
    \end{align*}
    is defined at the level of completions.
\end{example}
\newpage

\section{Nonarchimedean Valuations}
Completing fields is also extremely useful; in particular, this yields the $p$-adic numbers $\mathbb{Q}_p$. For this we will use valuations.
\begin{definition}{242. (Nonarchimedean Valuation)}
    \newline Let $K$ be a field. A \textbf{(multiplicative) valuation} on $K$ is a function
    \begin{align*}
        K\longrightarrow\mathbb{R},\ x\longmapsto|x|
    \end{align*}
    such that
    \begin{itemize}
        \item $|x|>0$ except for $|0|=0$.
        \item $|xy|=|x|\cdot|y|$.
        \item (Triangle Inequality): $|x+y|\leq|x|+|y|$.
    \end{itemize}
    Further, we call $|\cdot|$ a \textbf{nonarchimedean valuation} if $|x+y|\leq\mathrm{max}\{|x|,|y|\}$ holds. 
\end{definition}

\begin{remark}{243}
    \begin{itemize}
        \item[(a)] The first two axioms of a valuation imply that a valuation restricts to a multiplicative group hmomorphism $K^\times\to\mathbb{R}_{>0}$. Since $\mathbb{R}_{>0}$ is torsion-free, any root of unity in $K^\times$ maps to $1$. In particular, $|-1|=|1|$ and $|-x|=|x|$. Further, the third axiom generalises to
        \begin{align*}
            |\sum x_i|\ \leq\ \mathrm{max}\{|x_i|\}.
        \end{align*}
        \item[(b)] A valuation satisfies the usual properties of the absolute modulus on $\mathbb{R}$ or $\mathbb{C}$, so one can apply the standard estimates from analysis.
    \end{itemize}
\end{remark}

\begin{example}{244}
    \begin{itemize}
        \item[(a)] Any field admits the trivial valuation given by
        \begin{align*}
            |a|= \left\{\begin{array}{lc} 1 &\text{if}\ a\neq0 \\ 0 &\text{if}\ a=0 \end{array}\right..
        \end{align*}
        On a finite field $\mathbb{F}$ there is no other valuation as all $\mathbb{F}^\times$ consists of roots of unity.
        \item[(b)] For any number field $K$, there exists an embedding $\sigma:K\xhookrightarrow{} \mathbb{C}$. We get a valuation on $K$ by setting $|x|\coloneqq|\sigma(x)|$, where the latter valuation denotes the absolute modulus of a complex number.
        \item[(c)] For an (additive) discrete valuation $\nu:K^\times\to\mathbb{Z}$, let $e\in\mathbb{R}$ with $e>1$. Then
        \begin{align*}
            |a|\coloneqq \left\{\begin{array}{lc} e^{-\nu(a)} &\text{if}\ a\neq0 \\ 0 &\text{if}\ a=0 \end{array}\right..
        \end{align*}
        is a nonarchimedean valuation on $K$. A popular choice is the $p$-\textbf{adic valuation} on $\mathbb{Q}$ induced by $\nu_p$, where $p\in\mathbb{Z}$ a prime (cf. \textbf{Example 150}), together with $e=p$. Thus, if $a=a_0p^r$, then $|a|=p^{-r}$. It reflects our intuition that numbers which are divisible by high powers of $p$ should be thought of as small quantities.
        \item[(d)] One can generalise (c) to any number field $K$. Fix  $\mathfrak{p}\in\mathrm{Spec}\mathcal{O}_K$ and let $\mathbb{N}(\mathfrak{p})\coloneqq p^f$ (the \textbf{numerical norm} of $\mathfrak{p}$), where $(p)=\mathfrak{p}\cap\mathbb{Z}$ and $f$ the residue class degree (cf. \textbf{Prop. 155}). Using the $\mathfrak{p}$-adic discrete valuation $\nu_\mathfrak{p}$ from \textbf{(150a)} yields the multiplicative $\mathfrak{p}$-\textbf{adic valuation} $|x|_\mathfrak{p}\coloneqq\mathbb{N}(\mathfrak{p})^{-\nu_\mathfrak{p}(x)}$.
    \end{itemize}
\end{example}

\begin{proposition}{245}
    A valuation $|\cdot|$ is nonarchimedean iff it takes bounded values on the set  $\{m1:m\in\mathbb{Z}\}$ in $K$.
\end{proposition}
\begin{proof}
    "$\Rightarrow$"\ Let $m\in\mathbb{Z}_{>0}$, since $|\cdot|$ is nonarchimedean, we have
    \begin{align*}
        |m1|=|1+\cdots+1|\leq|1|=1.
    \end{align*}
    As we noted above, $|-1|=|1|$, thus $|-m1|=|m1|\leq1$.\\
    "$\Leftarrow$"\ Conversely, for $x,y\in K$ we note that
    \begin{align*}
        |x+y|^n &=|(x+y)^n|=|\sum_r\tbinom{n}{r}x^ry^{n-r}|\leq\sum_r|\tbinom{n}{r}||x|^r|y|^{n-r}\\ &\leq\sum_r|\tbinom{n}{r}|\cdot\mathrm{max}\{|x|^n,|y|^n\} =\sum_r|\tbinom{n}{r}|\cdot\mathrm{max}\{|x|,|y|\}^n.
    \end{align*}
    Suppose $|m1|$ is bounded by $N$ for all $m$. Since $\binom{n}{r}$ is an integer, we have
    \begin{align*}
        |x+y|^n\leq N(n+1)\cdot\mathrm{max}\{|x|,|y|\}^n.
    \end{align*}
    On taking $n$th roots we find that
    \begin{align*}
        |x+y|\leq N^{\frac{1}{n}}(n+1)^{\frac{1}{n}}\cdot\mathrm{max}\{|x|,|y|\}.
    \end{align*}
    Let $n\to\infty$ and take $\log$, we can see that $N^{\frac{1}{n}}(n+1)^{\frac{1}{n}}\to1$. Thus $|\cdot|$ is a nonarchimedean valuation.
\end{proof}

\begin{corollary}{246}
    If $\mathrm{char}K\neq0$, then $K$ has only nonarchimedean valuations.
\end{corollary}
\begin{proof}
    Suppose $\mathrm{char}K\coloneqq n$, then the set  $\{m1:m\in\mathbb{Z}\}\cong\mathbb{Z}/n\mathbb{Z}$ is finite, thus bounded. 
\end{proof}
\newpage
\begin{proposition}{247}
    Let $|\cdot|$ be a nontrivial nonarchimedean valuation and let $\log$ be the logarithm taken with respect to a real base $e>1$. Then the map 
    \begin{align*}
        \nu:K^\times\longrightarrow\mathbb{R},\ x\mapsto-\log|x|
    \end{align*}
    satisfies
    \begin{align*}
        \text{(i)}\ \nu(xy)=\nu(x)+\nu(y)\qquad \text{and\qquad (ii)}\ \nu(x+y)\geq\mathrm{min}\{\nu(x),\nu(y)\}.
    \end{align*}
    If $|\cdot|$ is \textbf{discrete}, i.e. $\nu(K^\times)$ is a discrete subgroup of $\mathbb{R}$, then $\nu$ is a multiple of some discrete valuation $K^\times\to\mathbb{Z}\subset\mathbb{R}$.
\end{proposition}
\begin{proof}
    The two properties (i) and (ii) are clear by the definition. For the last statement, note that $\nu(K^\times)$ is an (additive) subgroup of $\mathbb{R}$. Thus, if it is discrete, then it must be a lattice, i.e. $\nu(K^\times)\coloneqq m\mathbb{Z}$ for some $m\in\mathbb{R}_{>0}$. Hence
    \begin{align*}
        \mathrm{ord}:K^\times\xtwoheadrightarrow{} \mathbb{Z},\qquad x\longmapsto\tfrac{1}{m}\cdot\nu(x)
    \end{align*}
    is an additive discrete valuation.
\end{proof}

\begin{remark}{248}
    By \textbf{Thm. 270} and \textbf{Rmk. 271} we can extend the $p$-adic valuation $|\cdot|_p$ on $\mathbb{Q}$ to its algebraic closure $\overline{\mathbb{Q}}$. Since $p^{\frac{1}{n}}\in\overline{\mathbb{Q}}$ and $|p^{\frac{1}{n}}|_{\mathbb{Q}}=p^{-\frac{1}{n}}$ tends to $1$ as $n\to\infty$, a nonarchimedean valuation is not necessarily discrete.
\end{remark}

\begin{proposition}{249. (Local Parameter)}
    \newline Let $|\cdot|$ be a nonarchimedean valuation and $\nu=-\log|\cdot|$ as above. Then
    \begin{itemize}
        \item[(a)] $A\coloneqq\{a\in K:|a|\leq1\}=\{a\in K:\nu(a)\geq0\}$ is a valuation ring.
        \item[(b)] $A^\times=\{a\in K:|a|=1\}=\{a\in K:\nu(a)=0\}$ is its group of units.
        \item[(c)] $\mathfrak{m}\coloneqq\{a\in K:|a|<1\}=\{a\in K:\nu(a)>0\}$ is its unique maximal ideal.
    \end{itemize}
    Furthermore, $|\cdot|$ is discrete iff $A$ is a DVR. In this case, the maximal ideal $\mathfrak{m}=(\pi)$, where $\nu(K^\times)=\nu(\pi)\mathbb{Z}$ and we call such $\pi$ a \textbf{local (uniformising) parameter}.
\end{proposition}
\begin{proof}
    (a) Note that $|a|\cdot|a^{-1}|=|a\cdot a^{-1}|=1$, we have $|a^{-1}|=|a|^{-1}$, thus if $a\notin A$, then $|a|>1$, whence $|a^{-1}|=|a|^{-1}<1$, i.e. $a^{-1}\in A$. Hence $A$ is a valuation ring.\\
    (b) Since $|a|\cdot|a^{-1}|=|a\cdot a^{-1}|=1$, we have
    \begin{align*}
        a\in A^\times\ \Leftrightarrow\ a^{-1}\in A\ \Leftrightarrow\ |a|^{-1}\leq1\ \Leftrightarrow\ |a|=1.
    \end{align*}
    (c) Note that $A\setminus\mathfrak{m}=A^\times$, it suffices to show that $\mathfrak{m}$ is an ideal.\\
    If $x,y\in\mathfrak{m}$, then $|x+y|\leq\mathrm{max}\{|x|,|y|\}<1$, i.e. $x+y\in\mathfrak{m}$. If $a\in A$, then $|ax|=|a|\cdot|x|<1$, i.e. $ax\in\mathfrak{m}$. Hence $\mathfrak{m}$ is an ideal.\\
    "$\Rightarrow$"\ If $|\cdot|$ is discrete, then this follows from \textbf{Prop. 247} that $\nu(K^\times)=m\mathbb{Z},\ (m>0)$ and there exists $\pi\in K^\times$ with $\nu(\pi)=m>0$, thus $\pi\in\mathfrak{m}\subset A$. Now let $a\in A$, then $0\leq\nu(a)\coloneqq k\nu(\pi)$ for some $k\in\mathbb{R}_{\geq0}$. Now, since $0=\nu(a)-k\nu(\pi)=\nu(a\pi^{-k})$, we have $u\coloneqq a\pi^{-k}\in A^\times$, whence $a=u\cdot\pi^k\in(\pi^k)$. Thus every ideal is of the form $(\pi^k)$, so that $A$ is a PID and thus a DVR as being local.\\
    "$\Leftarrow$"\ Conversely, if $A$ is a DVR. Since $\nu(x^{-1})=-\nu(x)$, it suffices to show that: if $\nu(x)>0$, then $\nu(x)\in\nu(\pi)\mathbb{Z}$, where $\mathfrak{m}\coloneqq(\pi)$. Indeed, $\because \nu(x)>0.\ \therefore x\in A.\ \because A$ is a DVR. $\therefore x\coloneqq u\pi^k$ for some $u\in A^\times,\ k\in\mathbb{N}$. Thus,
    \begin{align*}
        \nu(x)=\nu(u)+k\nu(\pi)=k\nu(\pi)\in\nu(\pi)\mathbb{Z}.
    \end{align*}
\end{proof}

\begin{nt}{250. (Metrics)}
    \newline A valuation $|\cdot|$ defines a metric on $K$ with the distance function $d(a,b)=|a-b|$ and thus induces a metric topology on $K$. A countable local base at $0$ is given by 
    \begin{align*}
        U_n\ \coloneqq\ \{x\in K:|x|<\tfrac{1}{n}\}.
    \end{align*}
    For instance, $|a-b|$ is small with respect to the metric on $\mathbb{Q}$ defined by $|\cdot|_p$ if $p^n|a-b$ for large $n$.
    Now, let $K$ be a number field and fix $\mathfrak{p}\in\mathrm{Spec}\mathcal{O}_K$. The topology on $K$ defined by the $\mathfrak{p}$-adic valuation $|\cdot|_\mathfrak{p}$ from \textbf{(244)} is called the $\mathfrak{p}$-\textbf{adic topology} on $K$.
\end{nt}

\begin{proposition}{251. (Equivalent Valuations)}
    \newline Let $|\cdot|_1$ and $|\cdot|_2$ be two valuations on $K$ with $|\cdot|_1$ nontrivial. Then T.F.A.E.:
    \begin{itemize}
        \item[(i)] $|\cdot|_1$ and $|\cdot|_2$ define the same topology.
        \item[(ii)] $|\alpha|_1<1$ implies $|\alpha|_2<1$ for any $\alpha\in K$.
        \item[(iii)] $|\cdot|_2=|\cdot|_1^a$ for some $a>0$.
    \end{itemize}
    The valuations $|\cdot|_1$ and $|\cdot|_2$ are \textbf{equivalent} if they satisfy one and thus all of these conditions.
\end{proposition}
\begin{proof}
    "(i) $\Rightarrow$ (ii)"\ Since $|\alpha^n|=|\alpha|^n$, clearly $\alpha^n\to0$ iff $|\alpha|<1$. Thus (i) implies that
    \begin{align*}
        |\alpha|_1<1\ \Longleftrightarrow\ |\alpha|_2<1.
    \end{align*}
    "(ii) $\Rightarrow$ (iii)"\ Since $|\cdot|_1$ is nontrivial, there exists $y\in K$ with $|y|>1$. Consider\\ $a\coloneqq\log|y|_2/\log|y|_1$, so that $\log|y|_2=a\cdot\log|y|_1.\ \Rightarrow |y|_2=|y|_1^a$. Now, for any $0\neq x\in K$, choose $b\coloneqq\log|x|_1/\log|y|_1$, then $|x|_1=|y|_1^b$. \textsc{claim.} $|x|_2=|y|_2^b$.\\
    Suppose the rational $\frac{m}{n}>b,\ (n>0)$. Then $|x|_1=|y|_1^b<|y|_1^{m/n}.\ \Rightarrow |\frac{x^n}{y^m}|_1<1$. From assumption, this implies $|\frac{x^n}{y^m}|_2<1$, so $|x|_2<|y|_2^{m/n}$. Moreover, this is true for all rationals $\frac{m}{n}>b$, thus $|x|_2\leq|y|_2^b$. A similar argument with rationals $\frac{m}{n}<b$ shows that $|x|_2\geq|y|_2^b$. Thus we have the equality and finally
    \begin{align*}
        |x|_2=|y|_2^b=|y|_1^{ab}=|x|_1^a.
    \end{align*}
    "(iii) $\Rightarrow$ (i)"\ By assumption, a sequence converges with respect to  $|\cdot|_1$ iff it converges with respect to $|\cdot|_2$. Hence these valuations define the same topology.
\end{proof}
\newpage
Now, we give a complete list of the valuations on $\mathbb{Q}$ (up to equivalence).
\begin{theorem}{252. (Ostrowski)}
    Let $|\cdot|$ be a nontrivial valuation on $\mathbb{Q}$.
    \begin{itemize}
        \item[(a)] If $|\cdot|$ is archimedean, then $|\cdot|$ is equivalent to $|\cdot|_\infty$, the valuation defined by the restriction of the usual absolute modulus on $\mathbb{R}$ to $\mathbb{Q}$.
        \item[(b)] If $|\cdot|$ is nonarchimedean, then $|\cdot|$ is equivalent to $|\cdot|_p$, for unique prime $p$.
    \end{itemize}
\end{theorem}
\begin{proof}
    Let $m,n>1$ be integers. Then we can write
    \begin{align*}
        m\coloneqq a_0+a_1n+\cdots+a_rn^r,\qquad(0\leq a_i<n,\ n^r<m).
    \end{align*}
    Note that
    \begin{align*}
        |a_i|=|\underbrace{1+\ldots+1}_{a_i\text{-times}}|\leq a_i\cdot|1|=a_i<n.
    \end{align*}
    Pick $N\coloneqq\max\{1,|n|\}$, then we have
    \begin{align*}
        |m|\ \leq\ \sum_{i=0}^r|a_i||n|^i\ \leq\ \sum_{i=0}^rnN^r=(r+1)nN^r\ 
        \leq\ (1+\tfrac{\log m}{\log n})nN^{\frac{\log m}{\log n}}.
    \end{align*}
    Now, replace $m$ by $m^t (t\in\mathbb{N})$, we have
    \begin{align*}
        |m|\leq\ (1+\tfrac{t\log m}{\log n})^{\frac{1}{t}}n^{\frac{1}{t}}N^{\frac{\log m}{\log n}}\to1\qquad\text{as}\qquad t\to\infty.
    \end{align*}
    Hence, $|m|\leq N^\frac{\log m}{\log n}$. Now we distinguish cases.\\
    \textsc{case 1.} For all integers $n>1$ with $|n|>1$. In this case, $N=|n|$ and the above result yields: $|m|^{1/\log m}\leq|n|^{1/\log n}$. By symmetry, we must have the equality, and so for some $c>1$ we have $c=|m|^{\frac{1}{\log m}}=|m|^{\frac{1}{\log n}},\ \forall m,n\in\mathbb{N}_{>1}$. Hence, for any integers $n>1$, we have 
    \begin{align*}
        |n|=c^{\log n}=e^{\log c\log n}=n^{\log c}=|n|_\infty^a\qquad \text{where}\ a\coloneqq\log c.
    \end{align*}
    Since both $|\cdot|$ and $|\cdot|_\infty^a$ are homomorphisms $\mathbb{Q}^\times\to\mathbb{R}_{>0}$, the fact that they agree on a set of generators (the primes and $-1$) for the group $\mathbb{Q}^\times$ implies that they agree on all of $\mathbb{Q}^\times$, hence these two valuations are equivalent.\\
    \textsc{case 2.} For some integer $n>1$ with $|n|\leq1$. In this case we have $N=1$ and the above result implies $|m|\leq1$ for all integers $m$, whence the valuation is nonarchimedean by \textbf{(245)}. Let $(A,\mathfrak{m})$ be the associated local ring as in \textbf{(249)}, then $\mathbb{Z}\subset A$. Note that 
    \begin{itemize}
        \item $\mathfrak{m}\cap\mathbb{Z}$ is a prime ideal in $\mathbb{Z}$ since $\mathfrak{m}$ is prime.
        \item $\mathfrak{m}\cap\mathbb{Z}$ is nonzero otherwise the valuation would be trivial.
    \end{itemize}
    Hence $\mathfrak{m}\cap\mathbb{Z}=(p)$ for some prime $p\in\mathbb{Z}$. It follows that $|m|=1,\ \forall p\nmid m$, and so $|\frac{m}{n}\cdot p^r|=|p|^r$ if $m,n\in\mathbb{Z}$ with $p\nmid m,n$. Thus, if $a>0$ satisfies that $|p|=(\frac{1}{p})^a$, then $|x|=|x|_p^a$ for all $x=\frac{m}{n}\in\mathbb{Q}$.
\end{proof}

\begin{theorem}{253. (Product Formula)}
    \newline For $p\in\mathcal{P}\coloneqq\{\text{prime numbers}\}\cup\{\infty\}$, let $|\cdot|_p$ be the corresponding $p$-adic valuation on $\mathbb{Q}$. Then
    \begin{align*}
        \prod_{p\in\mathcal{P}}|x|_p=1,\qquad\text{for all}\ x\in\mathbb{Q}^\times.
    \end{align*}
\end{theorem}
\begin{proof}
    Note that the $p$-adic valuation $|\cdot|_p$ is defined by
    \begin{align*}
        \mathbb{Q}^\times\longrightarrow\mathbb{R},\ x\coloneqq \tfrac{a}{b}\cdot p^r\longmapsto p^{-r},\qquad(p\nmid a,b).
    \end{align*}
    Let $x=\frac{a}{b},\ (a,b\in\mathbb{Z})$. Then $|x|_p=1$ unless $p|a$ or $p|b$. Hence $|x|_\nu\neq1$ for all but finitely many $\nu$, so the product is actually finite. Now, consider the well-defined (multiplicative) group homomorphism 
    \begin{align*}
        \pi:\mathbb{Q}^\times\longrightarrow\mathbb{R}^\times,\ x\longmapsto\prod|x|_\nu
    \end{align*}
    We factorise $x\coloneqq p_1^{r_1}\cdots p_k^{r_k},\ (p_i\mathbb{Z},\ r\in\mathbb{Z})$ into primes and note that
    \begin{align*}
        |x|_p=|p_1^{r_1}\cdots p_k^{r_k}|_p=|p_1|_p^{r_1}\cdots|p_k|_p^{r_k}.
    \end{align*}
    Since $\pi(-1)=1$ (as $|-1|_p=1$ for any $p$), and from the definition of the $p$-adic valuation, we have
    \begin{align*}
        |p|_p=\tfrac{1}{p};\qquad |p|_q=1\ \text{if}\ q\neq p;\qquad |p|_\infty=p.
    \end{align*}
    Hence, $\pi(p_i)=1$ for all $i$, it follows that the product $\pi(x)=1$.
\end{proof}

Let $K$ be an algebraic number field. We can generalise the previous results as follows:
\begin{definition}{254. (Primes of a Number Field)}
    \newline An equivalence class of valuations on $K$ is called a \textbf{prime} or \textbf{place} of $K$. We usually denote the primes of $K$ by $\upsilon$.\\
    For instance, the primes of $\mathbb{Q}$ are the primes of $\mathbb{Z}$ and the valuation $|\cdot|_\infty$ induced by the real embedding $\sigma:\mathbb{Q}\xhookrightarrow{} \mathbb{C}$, i.e. $\sigma$ factorises via $\mathbb{R}$.
\end{definition}

\begin{theorem}{255}
    There exists exactly one prime of $K$, for each
    \begin{itemize}
        \item[(a)] prime ideal $\mathfrak{p}\subset\mathrm{Spec}\mathcal{O}_K:\ |x|_{\mathfrak{p}}=\mathbb{N}(\mathfrak{p})^{-\nu_{\mathbb{p}}(x)}$.
        \item[(b)] real embedding $K\xhookrightarrow{\sigma} \mathbb{R}\subset\mathbb{C}:\ |x|=|\sigma(x)|$.
        \item[(c)] complex pair of complex embedding $K\xhookrightarrow{\sigma} \mathbb{C}:\ |x|=|\sigma(x)|^2$.
    \end{itemize}
\end{theorem}

\begin{theorem}{256. (Product Formula)}
    For each prime $\upsilon$ of $K$, we let $|\cdot|_\upsilon$ be the valuation defined in \textbf{(255)}. Then
    \begin{align*}
        \prod_{\upsilon}|x|_\upsilon=1,\qquad\text{for all}\ x\in K^\times.
    \end{align*}
\end{theorem}

\begin{lemma}{257}
    If $|\cdot|_1,|\cdot|_2,\cdots,|\cdot|_n$ are nontrivial inequivalent valuations of a field $K$, there exists $a\in K$ such that $|a|_1>1$ and $|a|_i<1$ for all $i\neq1$.
\end{lemma}
\begin{proof}
    For $n=2$, since $|\cdot|_1$ and $|\cdot|_2$ are inequivalent, by \textbf{(251ii)}, there exist $b,c\in K$ such that
    \begin{align*}
        |b|_1<1,\qquad &|b|_2\geq1;\\ |c|_1\geq1,\qquad &|c|_2<1.\\
    \end{align*}
    Take $a\coloneqq\frac{c}{b}\in K$, then done. Now, we proceed by induction assuming that it is true for $n-1$ valuations. There exist $b,c\in K$ such that
    \begin{align*}
        |b|_1>1,\qquad &|b|_i<1,\ \forall\ 2\leq i\leq n-1;\\ |c|_1<1,\qquad &|c|_n>1.\\
    \end{align*}
    If $|b|_n\leq1$, then $a\coloneqq cb^r$ works for sufficiently large $r$ since \begin{align*}
        |a|=|c|\cdot|b|^r\longrightarrow
        \left\{\begin{array}{lc} |c| &\text{if}\ |b|=1 \\ 0 &\text{if}\ |b|<1 \\ \infty &\text{if}\ |b|>1
        \end{array}\right.
        \qquad\text{as}\ r\to\infty.
    \end{align*}
    If $|b|_n>1$, then $a\coloneqq\frac{cb^r}{1+b^r}$ works for sufficiently large $r$ since
    \begin{align*}
        |a|=|c|\cdot\tfrac{|b|^r}{|1+b|^r}\longrightarrow
        \left\{\begin{array}{lc} |c| &\text{if}\ |b|>1 \\ 0 &\text{if}\ |b|<1 \end{array}\right.
        \qquad\text{as}\ r\to\infty.
    \end{align*}
\end{proof}

\begin{lemma}{258}
    If $|\cdot|_1,|\cdot|_2,\cdots,|\cdot|_n$ are nontrivial inequivalent valuations of a field $K$ and $\epsilon>0$, we can find an $a\in K$ with 
    \begin{align*}
        |a-1|_1<\epsilon\qquad\text{and}\qquad |a|_i<\epsilon,\ \forall i\neq1.
    \end{align*}
\end{lemma}
\begin{proof}
    Choose $a\in K$ as in the previous lemma and set $a_r\coloneqq\frac{a^r}{1+a^r}$. Then 
    \begin{align*}
        |a_r-1|_1 &=\tfrac{1}{|1+a^r|_1}\leq\tfrac{1}{|a|_1^r-1}\longrightarrow0; \\
        |a_r|_i &=\tfrac{|a|_i^r}{|1+a|_i^r}\leq\tfrac{|a|_i^r}{1-|a|_i^r}\longrightarrow0,\qquad\text{for}\ i\geq2
    \end{align*}
    as $r\to\infty$.
\end{proof}
\newpage
\begin{theorem}{259. (Weak Approximation)}
    In the situation of \textbf{(257)}, let $a_1,\cdots,a_n\in K$. For any $\epsilon>0$, there is an element $a\in K$ such that $|a-a_i|_i<\epsilon$ for all $i$.
\end{theorem}
\begin{proof}
    Choose $b_i,\ (1\leq i\leq n)$ close to $1$ for $|\cdot|_i$ and close to $0$ for $|\cdot|_j,\ (j\neq i)$. Then $a\coloneqq\sum_{j=1}^na_jb_j$ works since $|a-a_i|_i\leq|a_i|_i|b_i-1|_i+\sum_{j\neq i}|a_j|_j|b_j|_j \longrightarrow0+0=0$.
\end{proof}

\begin{remark}{260}
    \begin{itemize}
        \item[(a)] This theorem can be regarded as a valuation form of the CRT.
        \item[(b)] In particular, it follows that there is no finite product formula: If 
        \begin{align*}
            \prod_{i=1}^n|a|_i^{r_i}=1,\ (r_i\in\mathbb{R}),\qquad\text{for all}\ a\in K^\times,
        \end{align*}
        then $r_i=0$ for all $i$.
    \end{itemize}
\end{remark}
\begin{proof}
    If $r_i\neq0$, we can choose $a$ such that $|a|_i$ is sufficiently large and $|a|_j$ sufficiently small for $j\neq i.\ (\to\gets)$
\end{proof}

\begin{nt}{261. (The Completion of Fields)}
    \newline An equivalence class $\upsilon$ of $|\cdot|$ turns $K$ into a topological field and determines a completion $K_\upsilon\coloneqq\widehat{K}$ which is unique up to homeomorphism, (cf. \textbf{Thm. 266}). Then every $x\in\widehat{K}$ is the limit of a Cauchy sequence $(x_n)$ in $K$. Since
    \begin{align*}
        |x_n-x_m|\ \geq\ ||x_n|-|x_m||,
    \end{align*}
    $(|x_n|)$ is a Cauchy sequence in $\mathbb{R}$ and thus has a limit. We can therefore extend $|\cdot|$ to a valuation on $\widehat{K}$ by setting $|x|\coloneqq\lim_n|x_n|$. In the same way we can extend the associated additive valuation $\nu=-\log|\cdot|$ to $\widehat{K}$.\\
    Now, let $|\cdot|$ be a discrete nonarchimedean valuation on $K$. Since $\nu(K^\times)=c\mathbb{Z}$ is closed in $\mathbb{R}$, we have $\nu(\widehat{K}^\times)=\nu(K^\times)$. Hence $|\cdot|$ extends to a  discrete valuation on $\widehat{K}$ which is normalised if $\nu$ is normalised. The associated DVR is 
    \begin{align*}
        \widehat{A}\ \coloneqq\ \{a\in\widehat{K}:|a|\leq1\}.
    \end{align*}
    It is the completion of the DVR $A=\{a\in K:|a|\leq1\}$ with respect to the $\mathfrak{m}$-adic topology defined by its maximal ideal 
    \begin{align*}
        \mathfrak{m}\ =\ \{a\in K:|a|<1\}\ =\ (\pi).
    \end{align*}
    In particular, $\widehat{A}$ is complete. The completion of $\mathfrak{m}$ is 
    \begin{align*}
        \widehat{\mathfrak{m}}\ =\ \{a\in\widehat{K}:|a|<1\}\ =\ (\pi)\widehat{A},
    \end{align*}
    and that $\widehat{A}/\widehat{\mathfrak{m}}^n\cong A/\mathfrak{m}^n$. [Exercise]
\end{nt}

\begin{proposition}{262}
    Let $S\subset A$ be a set of representatives of $A/\mathfrak{m}$ with $0\in S$, and let $\mathfrak{m}=(\pi)$. Then every element $0\neq x\in\widehat{K}$ has a unique representation as a convergent series
    \begin{align*}
        x\ =\ \pi^m(a_0+a_1\pi+a_2\pi^2+\cdots),\qquad (a_i\in S,\ a_0\neq0,\ m\in\mathbb{Z}).
    \end{align*}
\end{proposition}
\begin{proof}
    Since $|\cdot|$ is discrete and $\mathfrak{m}=(\pi)$, we may write $x\coloneqq u\cdot\pi^m,\ (u\in\widehat{A}^\times,\ m\in\mathbb{Z})$. Then $\nu(x)=-m\nu(\pi)$, so $\nu(x\pi^m)=0$. Since $u\in\widehat{A}^\times$ and $\widehat{A}/\widehat{\mathfrak{m}}\cong A/\mathfrak{m}$, there is a unique representative $0\neq a_0\in S$ for $u\ \mathrm{mod}\ \widehat{\mathfrak{m}}$. Hence $u=a_0+b_1\pi$ for some $b_1\in\widehat{A}$. Continuing this process we get
    \begin{align*}
        u=a_0+a_1\pi+\cdots+a_{n-1}\pi^{n-1}+b_n\pi^n,\qquad(b_n\in\widehat{A}).
    \end{align*}
    Then there is a unique representative $a_n\in S\cong A/\mathfrak{m}\cong\widehat{A}/\widehat{\mathfrak{m}}$ for $b_n\ \mathrm{mod}\ \widehat{\mathfrak{m}}$, whence $b_n=a_n+b_{n+1}\pi^{n+1}$ for some $b_{n+1}\in\widehat{A}$. Now, it remains to show that the series converges to $x$. But
    \begin{align*}
        |x-\pi^m\sum_{i=0}^na_i\pi^i|=|\pi|^m|b_{n+1}\pi^{n+1}|\leq|\pi|^{m+n+1}\longrightarrow0\qquad\text{as}\ n\to\infty
    \end{align*}
    for $|\cdot|=|\cdot|_{(\pi)}=|\cdot|_\mathfrak{m}$ the $\mathfrak{m}$-adic valuation. 
\end{proof}

\begin{example}{263. (The $p$-Adic Numbers)}
    \newline $\mathbb{Q}_p$ is precisely the completion of $\mathbb{Q}$ by the prime $|\cdot|_p$, ($p\in\mathbb{Z}$ prime).\\ A set of representatives of $\mathbb{Z}_p/(p)\mathbb{Z}_p\cong\mathbb{Z}/p\mathbb{Z}$ is given by $\{0,1,\cdots,p-1\}$. Hence \textbf{(262)} interprets a $p$-adic number in $\mathbb{Q}_p$ as a Laurent series $\sum_{i=-n}^\infty a_i\pi^i,\ (n\in\mathbb{N},\ 0\leq a_i<p)$, in complete analogy with holomorphic or meromorphic functions.
\end{example}

\begin{remark}{264}
    \begin{itemize}
        \item Primes of a number field $K$ which correspond to nonarchimedean valuations $|\cdot|_\mathfrak{p},\ (\mathfrak{p}\in\mathrm{Spec}\mathcal{O}_K)$, are called \textbf{finite}. 
        \item Primes which correspond to an archimedean prime are called \textbf{infinite}. This reflected in the notation $|\cdot|_\infty$ for the valuation on $\mathbb{Q}$ which is usually considered in your analysis course which yields the real numbers $\mathbb{R}$.
    \end{itemize}
\end{remark}

\begin{example}{265. (The Field of Laurent Series)}
    \newline Our initial motivation concerned an algebraic analogue of the meromorphic functions in one variable. \\
    Consider a \textbf{rational function field} $k(t)\coloneqq\mathrm{Quot}\ k[t]$ and the valuation induced by the prime $\mathfrak{p}\coloneqq(t-a)\triangleleft k[t]$. A set of representatives for the completion $k\llparenthesis t\rrparenthesis\coloneqq k(t)_{(t-a)}$ is given by
    \begin{align*}
        k\llbracket t\rrbracket/(t-a)^e\cong k[t]/(t-a)\cong k.
    \end{align*}
    In particular, any rational function $\frac{f(t)}{g(t)}$ has a development into a Laurent series
    \begin{align*}
        \sum_{i=-n}^\infty a_i(x-a)^i\ \text{in}\ k\llparenthesis t\rrparenthesis,\qquad (a_i\in k,\ n\in\mathbb{N}).
    \end{align*}
\end{example}

Let $A$ be a DVR with maximal ideal $\mathfrak{m}=(\pi)$ and residue field $k=A/\mathfrak{m}$, and we assume that $A$ is complete with respect to the $\mathfrak{m}$-adic topology. For $f\in A[x]$, let $\overline{f}\in k[x]$ denote the reduction of the coefficients of $f$ modulo $\mathfrak{m}$. The first version of \textbf{Hensel's Lemma} asserts that if $\overline{f}$ has a simple root, then we can split off a linear factor from $f$. This generalises the following:
\begin{theorem}{266. (Hensel's Lemma Ver. II)}
    \newline We use the previous notation and consider a monic polynomial $f\in A[x]$. If $\overline{f}$ factors as $\overline{f}=g_0h_0$ in $k[x]$ where $g_0$ and $h_0$ are monic and relatively prime, then there exist $g,h\in A[x]$ monic which are uniquely determined and such that $f=gh,\ \overline{g}=g_0$ and $\overline{h}=h_0$. Furthermore, we have $(g,h)=A[x]$. [Exercise]
\end{theorem}

\begin{remark}{267}
    The lemma holds in fact for any valuation ring or a field which is complete with respect to a nonarchimedean valuation, (cf. [N, II.4.6]).
\end{remark}

\begin{example}{268}
    \newline We claim that the $p$-adic numbers $\mathbb{Q}_p$ contain the $p$th roots of unity. Indeed, consider the polynomial $f\coloneqq x^{p-1}-1\in\mathbb{Z}_p[x]$. The residue field is $\mathbb{Z}_p/(p)\mathbb{Z}_p\cong\mathbb{Z}/p\mathbb{Z}$ and so splits into linear factors over $k=\mathbb{F}_p$. Repeated application of \textbf{Hensel's Lemma} allows us to split off a linear factor of $f$ at each stage so that $x^{p-1}-1=\prod_{i=1}^{p-1}(x-\xi_i)$.
\end{example}

\begin{corollary}{269}
    Let $K$ be complete with respect to a discrete nonarchimedean valuation and let $A$ be a DVR. Then for any irreducible polynomial $f\coloneqq a_0+a_1x+\cdots+a_nx^n\in K[x]$, we have
    \begin{align*}
        |f|\coloneqq\max\{|a_0|,\cdots,|a_n|\}=\max\{|a_0|,|a_n|\}.
    \end{align*}
    In particular, $a_n=1$ and $a_0\in A$ imply $f\in A[x]$.
\end{corollary}
\begin{proof}
    Since $K=\mathrm{Quot}A$ (where $A$ is a valuation ring), multiplying by a suitable element yields $x\cdot f\in A[x]$. Furthermore, dividing by a coefficient $a_j$ with $|a_j|\coloneqq\max\{|a_0|,\cdots,|a_n|\}$ yields
    \begin{align*}
        |\tfrac{x}{a_j}\cdot f|=1\qquad\text{and}\qquad\tfrac{x}{a_j}\cdot f\in A[x].
    \end{align*}
    Since rescaling of $f$ does not affect the assertion, we may straight away assume that $f\in A[x]$ and $|f|=1$. Now, let $a_r$ be the first coefficient such that $|a_r|=1$. Then 
    \begin{align*}
        f\equiv x^r(a_r+a_{r+1}x+\cdots+a_nx^{n-r})\qquad\mathrm{mod}\ \mathfrak{m}.
    \end{align*}
    If $0<r<n$, this contradicts the \textbf{Hensel's Lemma} since $f$ is irreducible. Thus we get the result $|f|=\max\{|a_0|,|a_n|\}$.
\end{proof}

\begin{theorem}{270}
    Let $K$ be complete with respect to a discrete nonarchimedean valuation $|\cdot|_K$ with $A$: DVR and let $L/K$ be a finite extension of degree $n$. Then $|\cdot|_K$ extends uniquely to a discrete valuation $|\cdot|_L$ on $L$ for which $L$ is complete and whose DVR is the integral closure of $A$ in $L$. Furthermore,
    \begin{align*}
        |x|_L=|\mathrm{N}_{L/K}(x)|_K^{1/n},\qquad\text{for all}\ x\in L.
    \end{align*}
\end{theorem}
\begin{proof}
    Let $B$ be the integral closure of $A$ in $L$ and $\mathfrak{p}$ be the maximal ideal of $A$. Then $B$ is a Dedekind domain and the valuations of $L$ extending $|\cdot|_\mathfrak{p}$ correspond to the ideals of $B$ lying over $\mathfrak{p}$.\\
    Suppose that there are distinct prime ideals $\mathfrak{P}_1$ and $\mathfrak{P}_2$ in $B$ dividing $\mathfrak{p}$, then there exists $b\in B$ such that $\mathfrak{P}_1\cap A[b]\neq\mathfrak{P}_2\cap A[b]$, (e.g. choose $b\in\mathfrak{P}_1$ but $b\notin\mathfrak{P}_2$). Let $f$ be the minimum polynomial of $b$ over $K$, so that $A[b]\cong A[x]/(f)$. Since $f\in A[x]$ is irreducible and $A$ is complete, \textbf{Hensel's Lemma} show that $\overline{f}$ (the image of $f$ in $k[x],\ (k=A/\mathfrak{p})$) must be a power of an irreducible polynomial. Then
    \begin{align*}
        A[b]/\mathfrak{p}A[b]\ \cong\ k[x]/(\overline{f})
    \end{align*}
    is a local ring, which contradicts the fact that $A[b]$ has two prime ideals containing $\mathfrak{p}$. Hence $|\cdot|_\mathfrak{p}$ extends uniquely to a valuation $|\cdot|_L$ on $L$.\\
    Clearly, $|\cdot|_\mathfrak{p}$ also extends to the Galois closure $\overline{L}$ of $L$. For each $\sigma\in\mathrm{Gal}(L/K)$, consider the map
    \begin{align*}
        L\xhookrightarrow{} \mathbb{C},\ b\longmapsto|\sigma(b)|.
    \end{align*}
    This is again a valuation of $L$, and so the uniqueness implies that $|b|_L=|\sigma(b)|$. Now,
    \begin{align*}
        |\mathrm{N}_{L/K}(b)|_K=|\prod\sigma(b)|_K=\prod|\sigma(b)|=\prod|b|_L=|b|_L^n
    \end{align*}
    implies the required formula.\\
    Finally, we have to show that $L$ is complete. Let $\{e_i\}_{i=1}^n$ be an $A$-basis for $B$ (as an $A$-module), and let $(\alpha(m))$ be a Cauchy sequence in $L$. We write $\alpha(m)\coloneqq\sum_{i=1}^na_i(m)e_i,\ (a_i(m)\in K)$, then for each $i,\ (a_i(m))$ is a Cauchy sequence in $K$. As $K$ is complete, say $a_i\coloneqq\lim_ma_i(m)$, then $\alpha\coloneqq\sum_{i=1}^na_ie_i=\lim_m(\alpha(m))$, whence $L$ is complete.
\end{proof}

\begin{remark}{271}
    \begin{itemize}
        \item[(a)] If $|\cdot|_K$ is nonarchimedean, then so is $|\cdot|_L$.
        \item[(b)] The statement holds for any valuation, whether it is discrete or not, (see [N, II.4.8]).
        \item[(c)] If $|\cdot|$ is a valuation on $K$, then $|\cdot|$ extends uniquely to any finite subextension of an algebraic separable extension $L/K$ and thus to a nonarcimedean valuation $|\cdot|_L$ on $L$. However, $|\cdot|_L$ need not be discrete, (cf. \textbf{Rmk. 248}).
    \end{itemize}
\end{remark}

\begin{corollary}{272}
    Let $K$ and $L$ be as in \textbf{Thm. 270} with DVRs $(A,\mathfrak{m})$ and $(B,\mathfrak{n})$, respectively; then $n=ef$, where $n=[L:K],\ e$ the ramification index of $\mathfrak{m}$, and $f=[B/\mathfrak{n}:A/\mathfrak{m}]$ the degree of the residue field extension.
\end{corollary}
\begin{proof}
    Consider $\mathfrak{p}\in\mathrm{Spec}A$ and the associated $\mathfrak{p}$-adic valuation $|\cdot|_\mathfrak{p}$. If $\mathfrak{a}$ is an ideal of $A$, the decomposition $\mathfrak{a}\coloneqq\mathfrak{p}_1^{r_1}\cdots\mathfrak{p}_g^{r_g}$ extends to $\mathfrak{a}^e\coloneqq(\mathfrak{p}_1^e)^{r_1}\cdots(\mathfrak{p}_g^e)^{r_g}$, so that every $\mathfrak{P}\in\mathrm{Spec}B$ which divides $\mathfrak{p}$ yields a $\mathfrak{P}$-adic valuation on $L$ which extends $|\cdot|_\mathfrak{p}$. By uniqueness of the extension \textbf{(270)}, we conclude that there exists precisely one $\mathfrak{P}$ dividing $\mathfrak{p}$, whence the result of the general formula $n=\sum e_if_i$ follows from \textbf{Thm. 158}.
\end{proof}

\begin{remark}{273}
    In this situation, let $\nu_L$ and $\nu_K$ denote the induced additive valuations. Since $\mathfrak{m}=(\pi)$ and $\mathfrak{n}=(\omega)$, we have $(\pi)B=\mathfrak{m}B=\mathfrak{n}^e=(\omega^e)$, we conclude that $\nu_K(\pi)=\nu_L(\pi)=e\nu_L(\omega)$. It follows that $e=(\nu_L(L^\times):\nu_K(K^\times))$.
\end{remark}

\begin{definition}{274. (Ramification)} 
    For the case in \textbf{Cor. 272}:
    \begin{itemize}
        \item If $e=n=[L:K]$, i.e. $f=1$ or $\mathfrak{m}B=\mathfrak{n}^n$, we say that $L$ is \textbf{totally ramified}.
        \item If $f=n=[L:K]$, i.e. $e=1$, we say that $L$ is \textbf{unramified} over $K$.
    \end{itemize}
\end{definition}

For the rest of this section we assume that:
\begin{itemize}
    \item $K$ is complete with respect to a discrete valuation $\nu$ with the associated DVR $(A,\mathfrak{m})$.
    \item $K$ and the residue field $A/\mathfrak{m}$ are perfect. In particular, any algebraic extension of these fields is separable.
\end{itemize}
These assumptions hold, for instance, for a $p$-adic number field which is sufficient for our purpose. Moreover, we have already seen that the valuation $|\cdot|$ of $K$ can extend uniquely to a (not necessarily discrete) valuation on $L$ if $L/K$ is a (possibly infinite) algebraic extension, (cf. \textbf{Rmk. 271}). In particular, we get the valuation ring $(B,\mathfrak{n})$:
\begin{align*}
    B=\{x\in L:\nu_L(x)\leq1\}\qquad\text{and}\qquad \mathfrak{n}=\{x\in B:\nu_L(x)<1\}
\end{align*}
and the residue field $\kappa(L)\coloneqq B/\mathfrak{n}$ of $L$.
\newpage
\begin{proposition}{275}
    Let $L/K$ be an algebraic extension and write $\ell=\kappa(L)$ the residue field of $L$. The assignment $K'\mapsto k'=\kappa(K)'$ sending an unramified extension $K'$ of $K$ contained in $L$ to its residue field $k'$ yields a 1-1 correspondence between the sets:
    \begin{align*}
        \{K'\subset L:\text{finite and unramified over}\ K\}\longleftrightarrow \{k'\subset\ell:\text{finite over}\ k\}.
    \end{align*}
    Moreover, this correspondence preserves the following properties:
    \begin{itemize}
        \item[(a)] \textbf{(Inclusion).} $K'\subset K''\ \Leftrightarrow\ \kappa(K')\subset\kappa(K'')$.
        \item[(b)] \textbf{(Normality).} $K'$ is Galois over $K\ \Leftrightarrow\ k'=\kappa(K')$ is Galois over $k$. In this case we have the canonical isomorphism $\mathrm{Gal}(K'/K)\cong\mathrm{Gal}(k'/k)$.
    \end{itemize}
\end{proposition}
\begin{proof}
    \textsc{claim 1.} The map is surjective. Let $k'\subset\ell$ be a finite extension of $k$, we may write $k'\coloneqq k(a),\ (a\in L)$ since $k$ is perfect. Let $f_0$ be the minimum polynomial of $a$ over $k$ and $f$ be any lifting of $f_0$ to $A[x]$, i.e. $\overline{f}=f_0$. As $a$ is a simple root of $f_0$, there exists a unique $\alpha\in L$ such that $f(\alpha)=0$ and $\alpha\equiv a\ \mathrm{mod}\ \mathfrak{n}$. Hence $K'\coloneqq K(\alpha)\subset L$ has residue field $k'$.\\
    \textsc{claim 2.} The map is injective. Suppose that $K'$ and $K''$ are unramified extensions of $K$ in $L$ with the same residue field $k'$. Then the compositum $K'\cdot K''$ (the smallest subfield of $L$ containing $K'$ and $K''$) is an unramified extension of $K$ with residue field $k'$. Hence
    \begin{align*}
        [K'\cdot K'':K]=[k':k]=[K'=K],
    \end{align*}
    whence $K''=K'$, i.e. the map is 1-1.\\
    (a) We have just seen that if $K'\leftrightarrow k'=\kappa(K')$, then there exists $\alpha'\in B'$ (where $(B',\mathfrak{n}')$ the valuation ring of $K'$) such that $K'=K(\alpha')$ and $k'=k(\overline{\alpha}')$ where $\overline{\alpha}'\equiv\alpha'\ \mathrm{mod}\ \mathfrak{n}'$. The assertion is then obvious. \\
    (b) We assume that $K'$ is Galois over $K$, then $\mathrm{Gal}(K'/K)$ preserves $A'$ (the valuation ring in $K'$) and its maximal ideal, thus we get a map $\mathrm{Gal}(K'/K)\to\mathrm{Aut}(k'/k)$. Write $k'=k(a)$ and let $g\in A[x]$ be such that $\overline{g}\in k[x]$ is the minimum polynomial of $a$. Let $\alpha\in A'$ be the unique root of $g$ such that $\overline{\alpha}=a$ (as above). Since $K'$ is Galois over $K$, $g$ splits in $A'[x]$ and this implies that $\overline{g}$ splits in $k[x]$, whence $k'$ is Galois over $k$. Now let $r=[k':k]=[K':K]$ and $\alpha_1,\cdots,\alpha_r$ be the roots of $g$. Then 
    \begin{align*}
        \{\alpha_1,\cdots,\alpha_r\}\ =\ \{\sigma(\alpha):\mathrm{Gal}(L/K)\}.
    \end{align*}
    As $\overline{g}$ is separable, the $\alpha_i$'s are distinct modulo $\mathfrak{n}$, which shows that the image of the map $\mathrm{Gal}(K'/K)\to\mathrm{Gal}(k'/k)$ has order $r$, and hence is an isomorphism. Conversely, suppose that $k'/k$ is Galois. Again we write $k'=k(a)$ and $\alpha\in A'$ lifts $a$. It follows from \textbf{Hensel's Lemma} that $A'$ contains the conjugates of $\alpha$, whence $K'$ is Galois over $K$.
\end{proof}
\newpage
\begin{example}{276}
    Let $K$ be a $p$-adic number field. As we will see in \textbf{Prop. 278}, its residue field $k$ is finite, say of order $q$. From the theory of finite fields, we find for each $n\in\mathbb{N}$ an extension field $k_n$ of $k$ of degree $n$ which is just the splitting field of $x^{q^n}-x$ (up to $k$-isomorphism). The Galois group is $\mathrm{Gal}(k_n/k)\cong\mathbb{Z}/n\mathbb{Z}$, which is generated by the Frobenius morphism $x\mapsto x^q$.\\ It follows that for all $n\in\mathbb{N}$, there exists an unramified extension field $K_n$ of $K$ of degree $n$ with the associated DVR $(B_n,\mathfrak{P}_n)$ which is unique up to $K$-isomorphism. The Galois group $\mathrm{Gal}(k_n/k)\cong\mathbb{Z}/n\mathbb{Z}$ is thus generated by the map $\sigma_n:K_n\to K_n$ determined by the property
    \begin{align*}
        \sigma_n(a)\equiv a^q\ \mathrm{mod}\ \mathfrak{P}_n,\qquad\text{for all}\ a\in B_n.
    \end{align*}
\end{example}

\begin{corollary}{277. (Largest Unramified Extension)}
    \newline Let $K_0\subset L$ be the compositum of all unramified extensions of $K$ in $L$. If $k$ is finite, then $K_0$ is obtained from $K$ by adjoining al roots of $1$ in $L$ of order prime to $\mathrm{char}\ k$. It carries a discrete valuation $|\cdot|_0$ with $\nu_0(K_0^\times)=\nu(K)$. Its residue field is the separable closure of $k$. The field $K_0$ is called the \textbf{largest unramified extension} of $K$ in $L$. \\ In particular, the largest unramified extension $K_{nr}$ ("non-ramifi\'{e}e") of $K$ is the largest unramified extension in its algebraic closure $\overline{K}$ and its residue field is $\overline{k}$, the algebraic closure of $k$.
\end{corollary}
\begin{proof}
    Let $k_0$ denote the residue field $\kappa(K_0)$ and $a\in\ell$ be separable over $k$. We need to show that $a\in k$. Let $\overline{g}\in k[x]$ be the minimal polynomial of $a$ such that $\overline{g}\equiv g\ \mathrm{mod}\ \mathfrak{m}$ for some $g\in A[x]$. By \textbf{Hensel}, there is a root $\alpha\in L$ with $\alpha\equiv a\ \mathrm{mod}\ \mathfrak{n}$, (note that the roots of $\overline{g}$ are simple). But then 
    \begin{align*}
        [K(\alpha):K]=\mathrm{deg}\ g=\mathrm{deg}\ \overline{g}=[k(a):k],
    \end{align*}
    i.e. $K(\alpha)$ is unramified over $K$. Hence we have $K(\alpha)\subset K_0,\ \alpha\in K_0$ and so $\overline{\alpha}=a\in k_0$, which follows that $k_0$ is a separable element of $k$. \\ Since every element of $K_0$ is contained in some finite extension field $K'$ with $K\subset K'\subset K_0$, we may assume that $K_0/K$ is finite. Then
    \begin{align*}
        [K_0:K]\geq ef=(\nu_{K_0}(K_0^\times):\nu_K(K^\times))[k_0:k]
        =(\nu_{K_0}(K_0^\times):\nu_K(K^\times))[K_0:K],
    \end{align*}
    whence $(\nu_{K_0}(K_0^\times):\nu_K(K^\times))=1$, i.e. $\nu_{K_0}(K_0^\times)=\nu_K(K^\times)$.\\
    If $k$ is finite with $p=\mathrm{char}\ k$, then $f_m\coloneqq x^m-1$ is separable if $p\nmid m$. Since $f_m(1)=0$ and $k_0$ is the separable closure of $k$, it follows that $k_0$ contains all the roots of $f_m$. By \textbf{Hensel's Lemma}, these lift to $K_0$. Since these roots of unity generate any finite extension of $k$ and thus any finite extension of $K$, they also generate $K_0$ since $K_0$ is exhausted by finite field extensions.\\ Finally, consider the algebraic closure $\overline{K}$ of $K$. If $\overline{f}\in k[x]$ with $f\in A[x]$ and $\overline{f}\equiv f\ \mathrm{mod}\ \mathfrak{m}$, then $f\in\overline{K}[x]$ and thus splits completely in $\overline{K}$, further, in $\kappa(\overline{K})[x]$. Since $\kappa(\overline{K})$ is algebraic over $k$, we have $\kappa(\overline{K})=k$.
\end{proof}
\newpage

\section{Local \& Global Fields}
\begin{proposition}{278}
    \newline Let $K$ be complete with respect to a nonarchimedean discrete valuation and $A$ be the valuation ring of $K$ with maximal ideal $\mathfrak{m}$. Then $A$ is compact iff its residue field $k=A/\mathfrak{m}$ is finite.
\end{proposition}
\begin{proof}
    Let $S$ be a set of representatives for $k=A/\mathfrak{m}$. \textsc{claim.} $A$ is compact $\Leftrightarrow\ S$ is finite. \\
    "$\Rightarrow$"\ Clearly, $\mathfrak{m}=\{x\in K:|x|<1\}$ is open in $K$. As $A=\bigsqcup_{s\in S}(s+\mathfrak{m})$ is the disjoint union of the open sets $s+\mathfrak{m}$ covering $A$. Thus $S$ is finite if $A$ is compact.\\
    "$\Leftarrow$"\ Recall that a metric space $X$ is compact iff it is complete and totally bounded (i.e. for any $r>0$, there is a finite covering of $X$ by open balls of radius $r$). But every element of $A$ can be written as
    \begin{align*}
        s_0+s_1\pi+s_2\pi^2+\cdots+s_n\pi^n+\cdots,\qquad(s_i\in S).
    \end{align*}
    For fixed $n\in\mathbb{N}$, since $S$ is finite, there are only finitely many sums
    \begin{align*}
        s_0+s_1\pi+s_2\pi^2+\cdots+s_n\pi^n,\qquad(s_i\in S),
    \end{align*}
    and that every element of $A$ is at distance at most $|\pi^{n+1}|$, which tends to $0$ as $n\to\infty$.\\ Thus $A$ is compact.
\end{proof}

\begin{definition}{279. (Local Field)}
    \newline A \textbf{local field} is a field $K$ with a valuation $|\cdot|$ which is locally compact (and hence complete) with respect to the induced topology.
\end{definition}

\begin{remark}{280}
    \begin{itemize}
        \item[(a)] A topological vector space, (e.g. a field $K$ with a valuation), is locally compact iff $0$ has a neighbourhood whose closure is compact.\\
        \textsc{proof.} Since a local base at $0$ is given by $U_n=\{x\in K:|x|<\frac{1}{n}\}$ by \textbf{(250)}, $K$ is locally compact iff $A$ is compact, where $A$ being the closure of its maximal ideal $\mathfrak{m}$.
        \item[(b)] It follows from the \textbf{Gelfand-Mazur Theorem} that a field which is locally compact with respect to an archimedean valuation is isomorphic to $\mathbb{R}$ or $\mathbb{C}$, (see [Algebra, Lang, Ch.XII]).
    \end{itemize}
\end{remark}

\begin{proposition}{281}
    The local fields which are locally compact with respect to a discrete nonarchimedean valuation are precisely the finite extensions of $\mathbb{Q}_p$ and $\mathbb{F}_p\llparenthesis t\rrparenthesis$.
\end{proposition}
\begin{proof}
    We first claim that: $L$ is a finite extension of $K\ \Leftrightarrow\ L$ is a  local field.\\ "$\Rightarrow$"\ Indeed, by \textbf{(270)}, $L$ is complete with respect to the discrete nonarchimedean valuation $|x|_L=|\mathrm{N}_{L/K}(x)|_K^{1/n}$, where $n=[L:K]$. Now, it remains to show that $\ell=\kappa(L)$ is finite. For this, it suffices to show that $\ell/\mathbb{F}_p$ is finite for $k=\kappa(K)=\mathbb{F}_p$ is finite. We let $\{\overline{x_i}\}_{i=1}^n$ be $k$-independent in $\ell$ with $\overline{x_i}\equiv x_i\ \mathrm{mod}\ \mathfrak{n},\ (x_i\in L)$. If 
    \begin{align*}
        a_1x_1+\cdots+a_nx_n=0,\qquad(a_i\in K)
    \end{align*}
    is a nontrivial relation, then dividing by the $a_i$ with the highest valuation gives a relation with coefficients in the DVR $(A,\mathfrak{m})$ of $K$. Hence, we get a nontrivial relation of $\overline{x_i}$ over $k.\ (\to\gets)$. Therefore, $\{\overline{x_i}\}_{i=1}^n$ is $k$-independent, whence $\ell$ is finite, which implies $L$ is local.\\
    "$\Leftarrow$"\ Since $L$ is local, $\ell$ is finite, thus of $\mathrm{char}\ \ell=p>0$. We distinguish into two cases:\\
    \textsc{case 1.} $\mathrm{char}L=0$. Then $\mathbb{Q}\subset L$ and if we restrict the additive valuation $\nu$ of $L$ to $\mathbb{Q}$, then we have $\nu|_{\mathbb{Q}}=\nu_p$, (since $p\in\mathfrak{n}$ for the DVR $(B,\mathfrak{n})$ where $\mathrm{char}B/\mathfrak{n}=p$, and thus $\nu(p)>0$). Hence the completion of $\mathbb{Q}$ with respect to $\nu_p$, namely $\mathbb{Q}_p$, is contained in $L$. This turns $L$ into a locally compact topological $\mathbb{Q}_p$-vector space which necessarily has finite $\mathbb{Q}_p$-dimension.\\
    \textsc{case 2.} $\mathrm{char}L>0$. Since $p\cdot1=0$ in $B$ (for $\mathrm{char}B/\mathfrak{n}=p$), we conclude that $\mathrm{char}L=p$. Now, we write $\ell\coloneqq\mathbb{F}_p(a)$ since $\mathbb{F}_p$ is perfect, and let $f\in\mathbb{F}_p[x]$ be the minimal polynomial of $a$ over $\mathbb{F}_p$. Since $\mathbb{F}_p$ is perfect, $\mathbb{F}_p(a)/\mathbb{F}_p$ is separable, whence $f$ splits completely in $\mathbb{F}_p(a)[x]$ and $\mathbb{F}_p(a)=\mathbb{F}_p(b)$ for any root of $f$. Considering $\mathbb{F}_p\subset B$ and thus $\mathbb{F}_p[x]\subset B[x]$, $f$ has a root $b\in B$ with $b\equiv a\ \mathrm{mod}\ \mathfrak{n}$. Hence $\mathbb{F}_p(b)=\mathbb{F}_p(a)=\ell\subset L$. But $\ell$ serves as a representative set for itself. Since any $x\in L$ can be developed into a Laurent series with coefficients in $\ell$, we conclude that $\mathbb{F}_p\llparenthesis t\rrparenthesis\subset L$ for a local parameter $t\in B$ with $(t)=\mathfrak{n}$. Moreover, $L/\mathbb{F}_p\llparenthesis t\rrparenthesis$ is finite (follows as above).
\end{proof}

\begin{definition}{282. (Global Field)}
    \newline A \textbf{global field} is a finite extension of $\mathbb{Q}$ (i.e. a number field) or $\mathbb{F}_p(t)$ (i.e. a \textbf{function field} over $\mathbb{F}_p$ in one variable).
\end{definition}

\begin{summary}{Local-Global Principal. 283}
    \newline $X=\mathrm{Spec}\mathcal{O}_K$ is a smooth curve for the ring of integers of a number field $K$, (cf. \textbf{Ch3}); $\mathcal{O}_K$ is the ring of global functions on $X$, and $K$ is the field of \textbf{rational functions}. This corresponds to the ring [resp. field] of global holomorphic [resp. meromorphic] functions of a complex curve. The local fields completed with respect to $|\cdot|_\mathfrak{p}$ correspond to Laurent series of meromorphic functions centered at that given point $\mathfrak{p}$.\\
    So a local field $K_\nu$ obtained by completion should be thought of as local information. Usually, it is easier to extract information/ solve problems locally than globally. Of course, passing from local information/ solutions to global information/ solutions is usually nontrivial (think for instance of the \textit{Mittag-Leffler Problem} to stay in the picture). Losely speaking, a \textbf{local-global principal} holds whenever the passage from local to global is unobstructed.
\end{summary}
\newpage
A famous instance is the \textbf{local-global principle} for the quadratic forms which we sketch next, see [F, Ch5] for details
\begin{definition}{284. (Isotropic Quadratic Forms)}
    \newline Let us consider quadratic forms over a field $K$ of $\mathrm{char}K\neq2$. A \textbf{quadratic form} is a function $Q:K^n\to K$ of the form $Q=\sum_{i,j=1}^na_{ij}x_ix_j$ with $a_{ij}=a_{ji}\in K$. It is an elementary fact that via a suitable coordinate change, any quadratic form can be transformed into 
    \begin{align*}
        Q(a_1,\cdots,a_n)\ =\ \sum_{i=1}^na_ix_i^2.
    \end{align*}
    We say that $Q$ is \textbf{isotropic} if $Q\neq0$ and there exists a nontrivial solution $x=(x_1,\cdots,x_n)$ such that $Q(a_1,\cdots,a_n)(x)=0$.
\end{definition}

Given a quadratic form $Q(a_1,\cdots,a_n),\ (a_i\in K)$, we now try to solve the quadratic equation $Q(a_1,\cdots,a_n)(x)=b$ in $x=(x_1,\cdots,x_n)$ for arbitrary $b\in K$.
\begin{lemma}{285}
    If $Q(a_1,\cdots,a_n)$ is isotropic, then for every $b\in K$, there exists $x=(x_1,\cdots,x_n)$ such that $Q(a_1,\cdots,a_n)(x)=b$.
\end{lemma}
\begin{proof}
    Let $x=(x_1,\cdots,x_n)\in K^n$ with $\sum_{i=1}^na_ix_i^2=0$, (WLOG may assume that $a_1x_1\neq0$). Let $b\in K$ and $t\in K$. Put 
    \begin{align*}
        y_i=\left\{\begin{array}{lc} x_i(1+t) &\text{if}\ i=1 \\ x_i(1-t) &\text{if}\ i\geq2 \end{array}\right..
    \end{align*}
    Then
    \begin{align*}
        \sum a_iy_i^2=(1+t^2)\sum a_ix_i^2+2ta_1x_1^2-2t\sum_{i=2}^na_ix_i^2=4ta_1x_1^2.
    \end{align*}
    By setting $t\coloneqq\frac{b}{4a_1x_1^2}$ yields the required result.
\end{proof}

\begin{theorem}{286. (Hasse's Local-Global Principle for Quadratic Forms)}
    \newline A quadratic form $Q$ over $\mathbb{Q}$ is isotropic iff $Q$ is isotropic when considered as a quadratic form over $\mathbb{R}$ and $\mathbb{Q}_p$.
\end{theorem}

We will not prove this theorem in full generality for lack of time and content ourselves with the special case $n=2$. In this case, $Q(a_1,a_2)=a_1x_1^2+a_2x_2^2\neq0$ iff $a_{1,2}\in K^\times$. Furthermore, $Q$ is isotropic iff $-\frac{a_1}{a_2}\in(K^\times)^2$, i.e. $-\frac{a_1}{a_2}$ is a square in $K^\times$. Hence we need to prove the following:
\begin{proposition}{287. (Hasse for $n=2$)}
    A nonzero rational number is a square in $\mathbb{Q}$ iff it is a square in $\mathbb{R}$ and $\mathbb{Q}_p$.
\end{proposition}
\begin{proof}
    Let $\frac{m}{n}\in\mathbb{Q}^\times$. then $\frac{m}{n}=\pm p_1^{r_1}\cdots p_g^{r_g}$ (with $r_i\in\mathbb{Z}$) is a square if and only if $\frac{m}{n}>0$ and $2|r_i$ for all $i$.
    If $\frac{m}{n}>0$, clearly $\frac{m}{n}$ is a square in $\mathbb{R}$; if $2|r_i$ for all $i$, then $\nu_p(\frac{m}{n})\equiv0\ \mathrm{mod}\ 2$, so $\frac{m}{n}=p^{2l}\cdot u, (u\in\mathbb{Z}_p^\times)$, whence $\frac{m}{n}$ is a square in $\mathbb{Q}_p$.
\end{proof}
\newpage

\section{Completions of Modules}
We already considered the completion of ideals. Next we consider the completion of modules in general.
\begin{definition}{288. ($\mathfrak{a}$-Filtration)}
    \newline Let $M$ be an $A$-module and $\mathfrak{a}\triangleleft A$ an ideal. 
    \begin{itemize}
        \item The $\mathfrak{a}$-\textbf{adic filtration} of $M$ is given by $M_n\coloneqq\mathfrak{a}^n\cdot M$. 
        \item The $\mathfrak{a}$-adic completion of $M$ is the $\widehat{A}$-module defined by 
        \begin{align*}
            \widehat{M}_\mathfrak{a}\ \coloneqq\ \varprojlim_{n\in\mathbb{N}}M/\mathfrak{a}^nM.
        \end{align*}
        \item An $A$-linear map $\phi:M\to N$ maps $M_n$ to $N_n$ and induces the map $\widehat{\phi}:\widehat{M}_\mathfrak{a}\to\widehat{N}_\mathfrak{a}$.
        \item A descending filtration $(M_n)$ is called an $\mathfrak{a}$-\textbf{filtration} if $\mathfrak{a}M_n\subset M_{n+1}$ for all $n$.
        \item If we have equality for all sufficiently large $n$, then the $\mathfrak{a}$-filtration is called \textbf{stable}. For instance, the $\mathfrak{a}$-adic filtration $(\mathfrak{a}^nM)$ is trivially $\mathfrak{a}$-stable.
    \end{itemize}
\end{definition}

\begin{lemma}{289}
    If $(M_n)$ and $(M'_n)$ are stable $\mathfrak{a}$-filtrations, then there exists an integer $k\in\mathbb{N}$ such that $M_{n+k}\subset M'_n\subset M_{n-k}$ for all $n\geq k$, i.e. both filtrations have \textbf{bounded difference}.\\ In particular, all stable $\mathfrak{a}$-filtrations induce the same topology on $M$, namely the $\mathfrak{a}$-topology. [AM, 10.6]
\end{lemma}
\begin{proof}
    By cutting off a sufficiently large number of submodules $M'_1,\cdots.M'_{n_0}$, we may assume that $M'_n=\mathfrak{a}^nM$. As $(M_n)$ is an $\mathfrak{a}$-filtration, the definition $\mathfrak{a}M_n\subset M_{n+1}$ for all $n$ implies that
    \begin{align*}
        M_n\supset\mathfrak{a}M_{n-1}\supset\cdots\supset\mathfrak{a}^nM_0=\mathfrak{a}^nM,\qquad\text{for all}\ n.
    \end{align*}
    Moreover, $(M_n)$ is stable, so there is a sufficiently large $k\in\mathbb{N}$ such that for all $n\geq0$ we have
    \begin{align*}
        M_{n+k}=\mathfrak{a}^nM_k\subset\mathfrak{a}^nM=M'_n\subset M_n\subset M_{n-k}.
    \end{align*}
\end{proof}

In the sequel, all completions will be taken with respect to some $\mathfrak{a}$-adic topology unless mentioned otherwise.\\
Let $A$ be a ring and $\mathfrak{a}\triangleleft A$ an ideal. Then 
\begin{align*}
    A^\ast\ \coloneqq\ \bigoplus_{n\geq0}\mathfrak{a}^n
\end{align*}
defines a graded ring. For instance, if $A$ is Noetherian so that $\mathfrak{a}=(a_1,\cdots,a_n)$ is finitely generated, then $A^\ast=A[a_1,\cdots,a_n]$. In particular, $A^\ast$ is Noetherian. Now, if $M$ is an $A$-module with an $\mathfrak{a}$-filtration $(M_n)$, then 
\begin{align*}
    M^\ast\ \coloneqq\ \bigoplus_{n\geq0}M_n
\end{align*}
is a graded $A^\ast$-module.

\begin{lemma}{290}
    Let $A$ be a Noetherian ring and $M$ be a finitely generated $A$-module with an $\mathfrak{a}$-filtration $(M_n)$. T.F.A.E.:
    \begin{itemize}
        \item[(a)] $M^\ast$ is a finitely generated $A^\ast$-module.
        \item[(b)] The filtration $(M_n)$ is stable.
    \end{itemize}
    In particular, any $\mathfrak{a}$-filtration of a finitely generated $A^\ast$-module $M$ with $A$ Noetherian induces the same topology on $M$. [AM, 10.8]
\end{lemma}
\begin{proof}
    Note that $M$ is a Noetherian $A$-module, so each $M_n$ is finitely generated, whence each $Q_n\coloneqq\bigoplus_{i=0}^nM_n\subset M^\ast$ is a finitely generated $A$-module (but not in general an $A^\ast$-module). However, for each $n\geq0$,
    \begin{align*}
        M_n^\ast\ \coloneqq\ Q_n\oplus(\bigoplus_{k\geq1}\mathfrak{a}^kM_n)
    \end{align*}
    is an $A^\ast$-submodule of $M^\ast$. Since $Q_n$ is finitely generated over $A$, and the finite set of $A$-generators of $M_n$ generates $\bigoplus_{k\geq1}\mathfrak{a}^kM_n$ over $A^\ast$, it follows that $M_n^\ast$ is finitely generated over $A^\ast$. Now, $(M_n^\ast)$ forms an ascending chain whose union is $M^\ast$. Since $A^\ast$ is Noetherian, we have
    \begin{align*}
        M^\ast\ \text{is finitely generated over}\ A^\ast. 
        &\Leftrightarrow\ \text{The chain $(M_n^\ast)$ stops}.\\
        &\Leftrightarrow\ M^\ast=M_{n_0}^\ast\ \text{for some}\ n_0.\\
        &\Leftrightarrow\ M_{n_0+r}=\mathfrak{a}^rM_{n_0}\ \text{for all}\ r\geq0.\\
        &\Leftrightarrow\ \text{The $\mathfrak{a}$-filtration $(M_n)$ is stable}.
    \end{align*}
\end{proof}

\begin{proposition}{291. (Artin-Rees Lemma)}
    Let $A$ be a Noetherian ring, $M$ a finitely generated $A$-module, and $(M_n)$ a stable $\mathfrak{a}$-filtration of $M$. If $N$ is an $A$-submodule of $M$, then $(N\cap M_n)$ is a stable $\mathfrak{a}$-filtration on $N$.
    In particular, taking $M_n=\mathfrak{a}^nM$, there exists an integer $k$ such that
    \begin{align*}
        N\cap\mathfrak{a}^nM=\mathfrak{a}^{n-k}(N\cap\mathfrak{a}^kM),
    \end{align*}
    for all $n\geq k$. [AM, 10.9-10]
\end{proposition}
\begin{proof}
    Note that for any $n\geq0$,
    \begin{align*}
        \mathfrak{a}(N\cap M_n)\subset\mathfrak{a}N\cap\mathfrak{a}M_n\subset N\cap M_{n+1},
    \end{align*}
    thus $(N\cap M_n)$ is an $\mathfrak{a}$-filtration. Hence it defines a graded $A^\ast$-module $N^\ast\coloneqq\bigoplus_{n\geq0}(N\cap M_n)$, which is a submodule of $M^\ast$ and thus finitely generated (for $M^\ast$ is Noetherian $A^\ast$-module by \textbf{Lemma 290}). Again, by \textbf{(290)}, $(N\cap M_n)$ is stable. For the last statement, we note that $(M_n)=(\mathfrak{a}^nM)$ is a stable $\mathfrak{a}$-filtration.
\end{proof}
\newpage
\begin{theorem}{292}
    Let $A$ be a Noetherian ring, $M$ a finitely generated $A$-module, and $N$ an $A$-submodule of $M$. Then the filtrations $(\mathfrak{a}^nN)$ and $(N\cap\mathfrak{a}^nM)$ have bounded difference. In particular, the $\mathfrak{a}$-topology of $N$ coincides with the topology induced by the $\mathfrak{a}$-adic topology of $M$. [AM, 10.11]
\end{theorem}
\begin{proof}
    Clearly, $(\mathfrak{a}^nN)$ is a stable $\mathfrak{a}$-filtration, thus the results follows from \textbf{(289)}. Coinciding topology follows from the fact that in a bounded difference we can always find an open neighbourhood of one topology which is contained in an open neighbourhood of the other topology.
\end{proof}

\begin{proposition}{293}
    Let
    \begin{align*}
        0\ \to\ M'\ \to\ M\ \to\ M''\ \to\ 0
    \end{align*}
    be an exact sequence of finitely generated modules over a Noetherian ring $A$. Then the sequence of $\mathfrak{a}$-adic completions
    \begin{align*}
        0\ \to\ \widehat{M}'_\mathfrak{a}\ \to\ \widehat{M}_\mathfrak{a}\ \to\ \widehat{M}''_\mathfrak{a}\ \to\ 0
    \end{align*}
    is exact. [AM, 10.12]
\end{proposition}
\begin{proof}
    We have seen from \textbf{(236)} that if $(M_n)$ is an $\mathfrak{a}$-filtration of $M$, then completing $M',\ M''$ by $(M'\cap M_n),\ (p(M_n))$, respectively, yields the exactness at the level of completions. Since $p:M\twoheadrightarrow M''$ is $A$-linear, from some sufficiently large $k$, we have
    \begin{align*}
        \mathfrak{a}p(M_n)\subset p(M_{n+1})=p(\mathfrak{a}M_n)\subset\mathfrak{a}p(M_n),\qquad\forall n\geq k.
    \end{align*}
    Thus, $(p(M_n))$ is stable and it follows from \textbf{(289)} that this induces the same $\mathfrak{a}$-topology, so that the completions with respect to these filtrations are isomorphic.
\end{proof}

The completion $\widehat{A}$ is a natural $A$-algebra via the completion map $A\to\widehat{A}$. In particular, given an $A$-module $M$, we can form the $\widehat{A}$-module $\widehat{A}\otimes_{\widehat{A}}M$. Since $\widehat{M}$ is an $\widehat{A}$-module by design, we get the sequence of $\widehat{A}$-modules
\begin{align*}
    \widehat{A}\otimes_AM\ \to\ \widehat{A}\otimes_A\widehat{M}\ \to\ \widehat{M}=\widehat{A}\otimes_{\widehat{A}}\widehat{M}.
\end{align*}
This behaves particularly well for $A$ Noetherian and finitely generated $M$.

\begin{proposition}{294}
    If $M$ is finitely generated, then $\widehat{A}\otimes_AM\to\widehat{M}$ is surjective. If, in addition, $A$ is Noetherian, then $\widehat{A}\otimes_AM\to\widehat{M}$ is an isomorphism. [AM, 10.13]
\end{proposition}
\begin{proof}
    Since $M$ is finitely generated, we have an exact sequence $0\to N\to F\to M\to0$ for a free $A$-module $F\cong A^n$. By the exactness of the completion functor \textbf{(236)} it follows that $\mathfrak{a}$-adic completion commutes with finite direct sums. Hence 
    \begin{align*}
        \widehat{A}\otimes_AF\cong(\widehat{A}\otimes_AA)^n\cong(\widehat{A})^n\cong\widehat{F}.
    \end{align*}
    \newpage This gives rise to the commutative diagram
    \begin{center}
        \begin{tikzcd}
            &\widehat{A}\otimes_AN\rar[] \dar["\gamma"]
            &\widehat{A}\otimes_AF\rar[] \dar["\beta"]
            &\widehat{A}\otimes_AM\rar[] \dar["\alpha"] &0 \\
            0\rar[] &\widehat{N}\rar[] &\widehat{F}\rar["\delta"] 
            &\widehat{M}\rar[] &0 \\
        \end{tikzcd}
    \end{center}
    in which the top row is exact (for $\widehat{A}\otimes_A-$ is a right exact functor) and the bottom row is exact by \textbf{(236)}. Hence $\delta$ is surjective. Since $\beta$ is an isomorphism, it follows that $\alpha$ is surjective, whence the first assertion.\\
    Now, if $A$ is Noetherian, then $N$ is also finitely generated so that $\gamma$ is also surjective, and that the bottom row is exact by \textbf{(293)}. A standard diagram chasing argument implies that $\alpha$ is injective, thus an isomorphism.
\end{proof}

\begin{proposition}{295}
    If $A$ is a Noetherian ring, then the functor $T_{\widehat{A}}\coloneqq\widehat{A}\otimes_A-$ is exact on the category of finitely generated $A$-modules. In particular, the $\mathfrak{a}$-adic completion $\widehat{A}$ of $A$ is a flat $A$-algebra. [AM, 10.14]
\end{proposition}
\begin{proof}
    Since $A$ is Noetherian, it follows from \textbf{(293-294)} that the functor $T_{\widehat{A}}=\widehat{A}\otimes_A\cdot\cong\widehat{\cdot}$ is an exact functor on the category of finitely generated $A$-modules.
\end{proof}

Next, we study the ring theoretic properties of $\widehat{A}$ in more detail.
\begin{lemma}{296}
    If $A$ is Noetherian with $\mathfrak{a}$-adic completion $\widehat{A}$. Then
    \begin{itemize}
        \item[(a)] $\widehat{\mathfrak{a}}\cong\widehat{A}\otimes_A\mathfrak{a}\cong\widehat{A}\mathfrak{a}=\mathfrak{a}^e$.
        \item[(b)] $\widehat{\mathfrak{a}^n}=(\widehat{\mathfrak{a}})^n$.
        \item[(c)] $\widehat{A}/\widehat{\mathfrak{a}}^n\cong A/\mathfrak{a}^n$ and $\widehat{\mathfrak{a}}^n/\widehat{\mathfrak{a}}^{n+1}\cong\mathfrak{a}^n/\mathfrak{a}^{n+1}$.
        \item[(d)] $\widehat{\mathfrak{a}}\subset\mathfrak{R}(\widehat{A})$, the Jacobson radical of $\widehat{A}$. [AM, 10.15]
    \end{itemize}
    \textbf{Recall.} The \textbf{Jacobson radical} of a ring $A$ is defined $\mathfrak{R}(A)\coloneqq\bigcap\{\mathfrak{m}\triangleleft A:\ \text{maximal}\}$. 
\end{lemma}
\begin{proof}
    (i) Since $A$ is Noetherian, $\mathfrak{a}$ is finitely generated so that we have $\widehat{A}\otimes_A\mathfrak{a}\cong\widehat{\mathfrak{a}}$ by \textbf{(294)}. Moreover, $\widehat{A}$ is flat by \textbf{(295)}, it follows that the injection $0\to\mathfrak{a}\xhookrightarrow{} A$ induces an injection
    \begin{align*}
        0\to\mathfrak{a}\otimes_A\widehat{A}\xhookrightarrow{} A\otimes_A\widehat{A}\cong\widehat{A},\ x\otimes a\mapsto x\cdot a,
    \end{align*}
    and thus $\mathfrak{a}\otimes_A\widehat{A}\cong\mathfrak{a}\cdot\widehat{A}=\mathfrak{a}^e$.\\ (ii) Applying (i) to $\mathfrak{a}^n$, we see that
    \begin{align*}
        \widehat{\mathfrak{a}}\cong\mathfrak{a}^n\cdot\widehat{A}=(\mathfrak{a}^n)^e=(\mathfrak{a}^e)^n\cong(\widehat{\mathfrak{a}})^n.
    \end{align*}
    (iii) We showed the first assertion in \textbf{(237)}. Further, completing the exact sequence $0\to\mathfrak{a}^{n+1}\to\mathfrak{a}^n\to\mathfrak{a}^n/\mathfrak{a}^{n+1}\to0$ gives the second isomorphism.\\
    (iv) For any $x\in\widehat{\mathfrak{a}}$, the sequence $a_n=\sum_{i=0}^nx^i$ is Cauchy in $A$ with respect to its $\mathfrak{a}$-topology and thus converges in $\widehat{A}$. For its limit we have $(1-x)\cdot\sum_{i=0}^\infty x^i=1$, i.e. $1-x$ is a unit. Since $\widehat{\mathfrak{a}}$ is an ideal, we can prove that $1-ax$ is a unit for all $a\in\widehat{A}$. Hence, if $x$ were not in some maximal ideal $\widehat{\mathfrak{m}}$, we have $(x)+\widehat{\mathfrak{m}}=(1)$, whence $\widehat{\mathfrak{m}}=\widehat{A}.\ (\to\gets)$. Therefore $x\in\mathfrak{R}(\widehat{A})$.
\end{proof}

\begin{corollary}{297}
    Let $(A,\mathfrak{m})$ be a Noetherian local ring. Then the $\mathfrak{m}$-adic completion $\widehat{A}$ of $A$ is a local ring with maximal ideal $\widehat{\mathfrak{m}}$. [AM, 10.16]
\end{corollary}
\begin{proof}
    By \textbf{(296c)} we have $\widehat{A}/\widehat{\mathfrak{m}}\cong A/\mathfrak{m}$, thus $\widehat{A}/\widehat{\mathfrak{m}}$ is a field and it follows that $\widehat{\mathfrak{m}}$ is maximal, so $\mathfrak{R}(\widehat{A})\subset\widehat{\mathfrak{m}}$. On the other hand, we have seen in \textbf{(296d)} that $\widehat{m}\subset\mathfrak{R}(\widehat{A})$, hence we have the equality and so $\widehat{\mathfrak{m}}$ is the unique maximal ideal. Therefore $(\widehat{A},\widehat{\mathfrak{m}})$ is local.
\end{proof}

\begin{corollary}{298}
    Let $A$ be a Noetherian ring and $\mathfrak{a}=(a_1,\cdots,a_n)\triangleleft A$ be an ideal. Then the $\mathfrak{a}$-adic completion is
    \begin{align*}
        \widehat{A}_\mathfrak{a}\ \cong\ A\llbracket x_1,\cdots,x_n\rrbracket/(x_1-a_1,\cdots,x_n-a_n).
    \end{align*}
    In particular, the completion of $A=K[x_1,\cdots,x_n]/\mathfrak{a}$ for some field $K$ and ideal $\mathfrak{a}\subsetneq K[x_1,\cdots,x_n]$ with respect to the maximal ideal $((x_1-a_1,\cdots,x_n-a_n)+\mathfrak{a})/\mathfrak{a}$ is
    \begin{align*}
        \widehat{A}_\mathfrak{p}\ \cong\ K\llbracket x_1,\cdots,x_n\rrbracket/\mathfrak{a}\cdot K\llbracket x_1,\cdots,x_n\rrbracket.
    \end{align*}
\end{corollary}

\begin{theorem}{299. (Cohen Structure Theorem)}
    \newline If $(A,\mathfrak{m})$ is a complete regular local ring of dimension $n$ containing some field, then $A\cong k\llbracket x_1,\cdots,x_n\rrbracket$, (the ring of formal power series over the residue field $k=A/\mathfrak{m}$ of $A$).
\end{theorem}
\begin{proof}
    See [Commutative Ring Theory, H.Matsumura, Thm. 29.8].
\end{proof}

\begin{theorem}{300. (Krull's Intersection Theorem)}
    \newline Let $A$ be a Noetherian ring, $M$ a finitely generated $A$-module and $\widehat{M}$ the $\mathfrak{a}$-completion of $M$. Then the kernel $E=\bigcap_{n=1}^\infty\mathfrak{a}^nM$ of the completion map $M\to\widehat{M}$ consists of those $x\in M$ annihilated by some element of $1+\mathfrak{a}$. [AM, 10.17]
\end{theorem}
\begin{proof}
    "$\supset$"\ If $(1-a)x=0$ for some $a\in\mathfrak{a}$, then 
    \begin{align*}
        x=ax=a^2x=\cdots\in\bigcap_{n=1}^\infty\mathfrak{a}^nM=E.
    \end{align*}
    "$\subset$"\ Since $E$ is the intersection of all neighbourhoods of $0\in M$, the subspace topology induced on $E$ is trivial, i.e. $E$ is the only neighbourhood of $0\in E$. By \textbf{Artin-Rees}, the induced topology on $E$ coincides with its $\mathfrak{a}$-topology. In particular, since $\mathfrak{a}E$ is a neighbourhood of $0$ in the $\mathfrak{a}$-topology, it follows that $\mathfrak{a}E=E$. Since $M$ is finitely generated and $A$ is Noetherian, $E$ is also finitely generated, whence by [AM, 2.5], $(1-a)E=0$ for some $a\in\mathfrak{a}$.
\end{proof}

\begin{remark}{301}
    \textbf{Krull's Thm.} may fail when $A$ is not Noetherian\\
    \textsc{proof.} Consider $\mathscr{C}^\infty(\mathbb{R})$ and let $\mathcal{P}$ be the maximal ideal of functions vanishing at the origin. By \textit{Taylor's Thm.}, $\mathcal{P}=(x)$ so that $E=\bigcap\mathcal{P}^k$ consists of functions whose derivative at the origin vanishes up to any order. On the other hand, $f\in\mathscr{C}^\infty(\mathbb{R})$ is annihilated by some element in $1+\mathcal{P}$ iff $f$ vanishes identically near $0$. However, $f=e^{-1/x^2}\in N$ does not vanish for $x>0$.
\end{remark}

\begin{corollary}{302}
    Let $A$ be a Noetherian integral domain and $\mathfrak{a}\triangleleftneq A$ an ideal. Then $\bigcap_{n=1}^\infty\mathfrak{a}^n=0$. In particular, the $\mathfrak{a}$-topology on $A$ is Hausdorff. [AM, 10.18]
\end{corollary}
\begin{proof}
    The ideal $1+\mathfrak{a}$ contains no nonzero zero-divisors as $A$ is an integral domain. The last assertion follows from \textbf{(A.5)}.
\end{proof}

\begin{corollary}{303}
    Let $A$ be a Noetherian ring, $\mathfrak{a}\triangleleft A$ an ideal containing in the Jacobson radical $\mathfrak{R}(A)$ and $M$ a finitely generated $A$-module. Then the $\mathfrak{a}$-topology of $M$ is Hausdorff, i.e. $\bigcap\mathfrak{a}^nM=0$. In particular, this applies to the situation of a Noetherian local ring $(A,\mathfrak{m})$ with the $\mathfrak{m}$-topology on $M$. [AM, 10.19-20]
\end{corollary}
\begin{proof}
    Since $\mathfrak{a}\subset\mathfrak{R}(A)$, every element of $1+\mathfrak{a}$ is a unit, (cf. [AM, 1.9]). Otherwise, $1+a,\ (a\in\mathfrak{a})$, is contained in some maximal ideal $\mathfrak{m}\triangleleft A$. Since $a\in\mathfrak{a}\subset\mathfrak{m}$, we have $1\in\mathfrak{m}.\ (\to\gets)$. But if $1+a$ is a unit, $x\mapsto(1+a)\cdot x$ has trivial kernel.
\end{proof}

Our final aim is to show that the $\mathfrak{a}$-adic completion of a Noetherian is again Noetherian.
\begin{definition}{304. (Associated Graded Ring)}
    \newline Let $A$ be a ring and $\mathfrak{a}\triangleleft A$ an ideal. We define the \textbf{associated graded ring} of $A$ by
    \begin{align*}
        \mathrm{Gr}_\mathfrak{a}(A)=\bigoplus_{n\geq0}\mathrm{Gr}_{\mathfrak{a},n}(A)\coloneqq\bigoplus_{n\geq0}\mathfrak{a}^n/\mathfrak{a}^{n+1},\qquad(\mathfrak{a}^0=A).
    \end{align*}
    If the underlying ideal $\mathfrak{a}$ is clear from the context, we simply write $\mathrm{Gr}(A)$. This is indeed a graded ring in the sense: if $x_n\in\mathfrak{a}^n$ and $x_m\in\mathfrak{a}^m$, then
    \begin{align*}
        (x_n+\mathfrak{a}^{n+1})\cdot(x_m+\mathfrak{a}^{m+1})\coloneqq x_nx_m+\mathfrak{a}^{n+m+1}.
    \end{align*}
    Similarly, if $M$ is an $A$-module with $\mathfrak{a}$-filtration $(M_n)$, we define
    \begin{align*}
        \mathrm{Gr}(M)=\bigoplus_{n\geq0}\mathrm{Gr}_n(M)\coloneqq\bigoplus_{n\geq0}M_n/M_{n+1}.
    \end{align*}
    This is a graded $\mathrm{Gr}(A)$-module. We write $\mathrm{Gr}_\mathfrak{a}(M)$ if $M_n=\mathfrak{a}^nM$ is the $\mathfrak{a}$-adic filtration.
\end{definition}

For example, if $A$ is Noetherian and $\mathfrak{a}=(x_1,\cdots,x_r)$, then 
\begin{align*}
    \mathrm{Gr}(A)=(A/\mathfrak{a})[\overline{x}_1,\cdots,\overline{x}_r],\qquad\text{(where $\overline{x}_i\coloneqq x_i+\mathfrak{a}^2)$}.
\end{align*}
\textsc{proof.} Indeed, $\mathrm{Gr}_{\mathfrak{a},n}(A)=\mathfrak{a}^n/\mathfrak{a}^{n+1}$ is finitely generated as an $A/\mathfrak{a}$-module by $(\overline{x}_1^n,\cdots,\overline{x}_r^n)$. Hence
\begin{align*}
    \mathrm{Gr}(A)=\bigoplus_{n\geq0}(\overline{x}_1^n,\cdots,\overline{x}_r^n)=(A/\mathfrak{a})[\overline{x}_1,\cdots,\overline{x}_r].
\end{align*}

\begin{proposition}{305}
    Let $A$ be a Noetherian ring and $\mathfrak{a}\triangleleft A$ an ideal. Then
    \begin{itemize}
        \item[(a)] $\mathrm{Gr}_\mathfrak{a}(A)$ is Noetherian.
        \item[(b)] $\mathrm{Gr}_\mathfrak{a}(A)$ and $\mathrm{Gr}_{\widehat{\mathfrak{a}}}(\widehat{A})$ are isomorphic as graded rings.
        \item[(c)] $\mathrm{Gr}(M)$ is a finitely generated graded  $\mathrm{Gr}_\mathfrak{a}(A)$-module if $M$ is a finitely generated $A$-module and $(M_n)$ is a stable $\mathfrak{a}$-filtration of $M$. 
        \item[(d)] In particular, $\mathrm{Gr}_\mathfrak{a}(M)$ is a finitely generated graded $\mathrm{Gr}_\mathfrak{a}(A)$-module if $M$ is a finitely generated $A$-module. [AM, 10.22]
    \end{itemize}
\end{proposition}
\begin{proof}
    (a) Since $A$ is Noetherian, $\mathfrak{a}$ is finitely generated, say by $x_1,\cdots,x_r$. From the discussion above, we have $\mathrm{Gr}(A)=(A/\mathfrak{a})[\overline{x}_1,\cdots,\overline{x}_r]$. As $A/\mathfrak{a}$ is Noetherian, $\mathrm{Gr}(A)$ is also Noetherian by \textbf{Hilbert Basis Thm}.\\
    (b) Recall from \textbf{(296c)} that $\mathfrak{a}^n/\mathfrak{a}^{n+1}\cong\widehat{\mathfrak{a}}^n/\widehat{\mathfrak{a}}^{n+1}$. The result follows from the definition.\\
    (c) Since $(M_n)$ is stable, there exists $k\in\mathbb{N}$ such that $M_{k+r}=\mathfrak{a}^rM_k,\ \forall r\geq0$, hence $\mathrm{Gr}(M)$ is generated by $\bigoplus_{n\leq k}\mathrm{Gr}_n(M)$ as a $\mathrm{Gr}(A)$-module. Furthermore, each $\mathrm{Gr}_n(M)=M_n/M_{n+1}$ is Noetherian and annihilated by $\mathfrak{a}$, and hence a finitely generated $A/\mathfrak{a}$-module, which follows that $\bigoplus_{n\leq k}\mathrm{Gr}_n(M)$ is generated by a finite number of elements (as an $A/\mathfrak{a}$-module). Hence $\mathrm{Gr}(M)$ is a finitely generated $\mathrm{Gr}(A)$-module. Finally, (d) is just a special case from (c).
\end{proof}

\begin{lemma}{306}
    Let $\phi:M\to N$ be an $A$-linear homomorphism of \textbf{filtered modules}, i.e. $\phi(M_n)\subset N_n$, and let $\mathrm{Gr}(\phi):\mathrm{Gr}(M)\to\mathrm{Gr}(N)$ and $\widehat{\phi}:\widehat{M}\to\widehat{N}$ be the induced homomorphisms of the associated graded and completed modules, respectively. Then
    \begin{itemize}
        \item[(a)] $\mathrm{Gr}(\phi)$ is injective. $\Rightarrow\ \widehat{\phi}$ is injective. 
        \item[(b)] $\mathrm{Gr}(\phi)$ is surjective. $\Rightarrow\ \widehat{\phi}$ is surjective. 
    \end{itemize}
    In particular, $\mathrm{Gr}(M)$ and $\mathrm{Gr}(N)$ are isomorphic via such a map $\phi:M\to N$, then the completions are isomorphic. [AM, 10.23]
\end{lemma}
\begin{proof}
    For each $n\in\mathbb{N}$, consider the commutative diagram of exact sequences:
    \begin{center}
        \begin{tikzcd}
            0\rar[] &M_n/M_{n+1}\rar[] \dar["\mathrm{Gr}_n(\phi)"]
            &M/M_{n+1}\rar[] \dar["\phi_{n+1}"]
            &M/M_{n}\rar[] \dar["\phi_{n}"] &0 \\
            0\rar[] &N_n/N_{n+1}\rar[] &N/N_{n+1}\rar[] &N/N_n\rar[] &0. \\
        \end{tikzcd}
    \end{center}
    This gives the exact sequence 
    \begin{align*}
        0\to\mathrm{ker\ Gr}_n(\phi)\to\mathrm{ker}\ \phi_{n+1}\to\mathrm{ker}\ \phi_{n}\to\mathrm{Coker\ Gr}_n(\phi)\to\mathrm{Coker}\ \phi_{n+1}\to\mathrm{Coker}\ \phi_{n}\to0
    \end{align*}
    by \textit{Snake Lemma}. From this we see, by induction on $n$, that\\
    \textsc{case 1.} \textit{$\mathrm{Gr}(\phi)$ is injective}, then each $\mathrm{Gr}_n(\phi)$ is injective. For $n=0$, we have $M=M_0$ and $N=N_0$, so that $\phi_0:M/M=\{0\}\to N/N=\{0\}$ is injective, whence $\phi_1$ is injective by \textit{Short Five Lemma}. Now, proceed by induction to conclude that $\phi_n$ is injective for all $n$. Finally, apply the inverse limit functor to the sequence 
    \begin{align*}
        0\to M/M_n\xhookrightarrow{~\phi_n~} N/N_n\twoheadrightarrow \mathrm{Coker}\ \phi_n\to0.
    \end{align*}
    As this is a left exact functor by \textbf{(235)}, $\widehat{\phi}$ is injective.\\
    \textsc{case 2.} \textit{$\mathrm{Gr}(\phi)$ is surjective}, then each $\mathrm{Gr}_n(\phi)$ is surjective. Similarly, we obtain by \textit{Short Five Lemma} that $\phi_n$ is surjective for all $n$. Finally, apply the inverse limit functor to the sequence 
    \begin{align*}
        0\to \mathrm{ker}\ \phi_n\xhookrightarrow{} M/M_n\xtwoheadrightarrow{\phi_n} N/N_n\to0.
    \end{align*}
    Moreover, the inverse system $\{\mathrm{ker}\ \phi_n\}$ is surjective (from the long exact sequence), hence $\widehat{\phi}$ is sujective by \textbf{(235)}.
\end{proof}

\begin{proposition}{307}
    Let $A$ be a ring, $M$ an $A$-module, and $(M_n)$ an $\mathfrak{a}$-filtration of $M$. Suppose that $A$ is complete in the $\mathfrak{a}$-topology and that $M$ is Hausdorff in its filtration topology, (i.e. $\bigcap_{n\geq0}M_n=0$). If $\mathrm{Gr}(M)$ is a finitely generated $\mathrm{Gr}(A)$-module, then $M$ is a finitely generated $A$-module. [AM, 10.24]
\end{proposition}
\begin{proof}
    Let $\overline{x}_i\coloneqq x_i+M_{n_i+1}\in\mathrm{Gr}_{n_i}(M),\ (x_i\in M_{n_i})$ such that $\{\overline{x}_i\}_{i=1}^\nu$ is a finite set of homogeneous generators of degree $n_i$. Let $F^i$ be the module $A$ with the stable $\mathfrak{a}$-filtration given by $F_k^i=\mathfrak{a}^{k+n_i}$ and put $F\coloneqq\bigoplus_{i=1}^{\nu}F^i\cong A^\nu$, where we endow the filtration $F_n\coloneqq\bigoplus_{i=1}^{\nu}F_n^i=\bigoplus_{i=1}^{\nu}\mathfrak{a}^{n+n_i}$. Now, mapping the generator $a_i=1\in F^i$ to $x_i$ defines a homomorphism $\phi:F\to M$ of filtered modules for $\phi(\mathfrak{a}^{n-n_i})\subset\mathfrak{a}^{n-n_i}M_{n_i}\subset M_n$, and that $\mathrm{Gr}(\phi):\mathrm{Gr}(F)\to\mathrm{Gr}(M)$ is a surjective homomorphism of $\mathrm{Gr}(A)$-module by design. Hence $\widehat{\phi}$ is surjective by \textbf{(306)}. Consider now the diagram
    \begin{center}
        \begin{tikzcd}
            F\rar["\phi"] \dar["\alpha"] &M\dar["\beta"]\\
            \widehat{F}\rar["\widehat{\phi}"] &\widehat{M}.\\
        \end{tikzcd}
    \end{center}
     Since $F\cong A^\nu$ and $\widehat{A}=A$, it follows that $\alpha$ is an isomorphism. As $M$ is Hausdorff, $\beta$ is injective. Moreover, $\widehat{\phi}$ is surjective, which implies that $\phi$ is surjective, i.e. $M$ is finitely generated by $\{x_i\}_{i=1}^\nu$ over $A$.
\end{proof}

\begin{corollary}{308}
    With the hypothesis of \textbf{(307)}, if $\mathrm{Gr}(M)$ is a Noetherian $\mathrm{Gr}(A)$-module, then $M$ is a Noetherian $A$-module. [AM, 10.25]
\end{corollary}
\begin{proof}
    We claim that every $A$-submodule $M'$ of $M$ is finitely generated. Consider the $\mathfrak{a}$-filtration $M'_n\coloneqq M'\cap M_n$ of $M'$, then the embedding $M'_n\xhookrightarrow{} M_n$ gives rise to an injective homomorphism $M'_n/M'_{n+1}\xhookrightarrow{} M_n/M_{n+1}$, and hence an embedding $\mathrm{Gr}(M')\xhookrightarrow{}\mathrm{Gr}(M)$. Since $\mathrm{Gr}(M)$ is Noetherian, $\mathrm{Gr}(M')$ is finitely generated. Moreover, $M'$ is Hausdorff since $\bigcap M'_n\subset\bigcap M_n=0$. Hence $M'$ is finitely generated by \textbf{(307)}.
\end{proof}

\begin{theorem}{309}
    If $A$ is Noetherian, then the $\mathfrak{a}$-completion $\widehat{A}_\mathfrak{a}$ of $A$ is Noetherian. [AM, 10.26]
\end{theorem}
\begin{proof}
    As $A$ is Noetherian, from \textbf{(305)} we know that $\mathrm{Gr}_{\widehat{\mathfrak{a}}}(\widehat{A})=\mathrm{Gr}_{\mathfrak{a}}(A)$ is a Noetherian $\mathrm{Gr}_{\mathfrak{a}}(A)$-module. Now apply \textbf{(308)} to the complete ring $\widehat{A}$ and take $M=\widehat{A}$ (filtered by $(\widehat{\mathfrak{a}}^n)$, and so Hausdorff since $\bigcap\widehat{\mathfrak{a}}^n=0$). Thus, $\widehat{A}$ is Noetherian.
\end{proof}

\begin{corollary}{310}
    If $A$ is Noetherian, then the ring of formal power series $A\llbracket x_1,\cdots,x_n\rrbracket$ is Noetherian. In particular, $k\llbracket x_1,\cdots,x_n\rrbracket$ is Noetherian. [AM, 10.27]
\end{corollary}
\begin{proof}
    Note that $A\llbracket x_1,\cdots,x_n\rrbracket$ is the completion (under the $(x_1,\cdots,x_n)$-adic topology) of $A[x_1,\cdots,x_n]$, which is Noetherian by \textbf{Hilbert Basis Thm}. 
\end{proof}

\newpage


\section{Exercises}
\begin{exe}{1. (Computing Expansions)}
    Compute the $5$-adic expansion of $\frac{1}{3},\ \frac{2}{3}$ and $-\frac{2}{3}$.
\end{exe}
\begin{proof}
    
\end{proof}

\begin{exe}{2. (Units in $\mathbb{Z}_p$)}
    A $p$-adic integer $a\coloneqq a_0+a_1p+a_2p^2+\cdots$ is a unit in $\mathbb{Z}_p$ iff $a_0\neq0$.
\end{exe}
\begin{proof}
    "$\Leftarrow$"\ Suppose $a_0\neq0$. First, we consider the sequence $(s_n)$, where $s_n=\sum_{i=0}^{n-1}a_ip^i$, then $s_n\in\mathbb{Z}\subset\mathbb{Z}_p$, thus it suffices to show that $p\nmid s_n$. But $s_n=a_0+p\cdot b,\ (b\in\mathbb{Z})$, we have $p|s_n\ \Leftrightarrow\ p|a_0\ \Leftrightarrow\ a_0=0$. By assumption, $p\nmid s_n$, so $\nu_p(x_n)=0$, whence $s_n\in\mathbb{Z}_p^\times$ by \textbf{(249)}.\\
    Now, the element $a\in\mathbb{Z}_p$ is given by the Cauchy sequence $(s_n)$. As each $s_n\in\mathbb{Z}_p^\times$, we may consider the sequence $(s_n^{-1})$. Note that $s_ns_{n+1}(s_n^{-1}-s_{n+1}^{-1})=s_ns_{n+1}$ and that $s_ns_{n+1}$ is a unit, it follows that $s_n^{-1}-s_{n+1}^{-1}\in(p^n)$ so that $(s_n^{-1})$ converges with respect to the $(p)$-adic topology. Moreover, we have $a^{-1}\coloneqq\lim_ns_n^{-1}\in\mathbb{Z}_p$, i.e. $a$ is a unit in $\mathbb{Z}_p$.\\
    "$\Rightarrow$"\ Suppose $a$ is a unit in $\mathbb{Z}_p$. We now think of as an inverse system by considering the Cauchy sequence $(s_n\ \mathrm{mod}\ p^n)$. Note that the multiplication of $\mathbb{Z}_p$ is induced by the product ring $\prod_{n\in\mathbb{N}}\mathbb{Z}/p^{n+1}\mathbb{Z}$, thus $s_1=a_0\in\mathbb{Z}/p\mathbb{Z}$ must be a unit. As $\mathbb{Z}/p\mathbb{Z}$ is a field, we obtain the required result $a_0\neq0$.
\end{proof}

\begin{exe}{3. ($p$-Adic Integers \& Formal Power Series)}
    There is a canonical ring isomorphism $\mathbb{Z}_p\cong\mathbb{Z}\llbracket x\rrbracket/(x-p)$.
\end{exe}

\begin{exe}{4. (Solving Quadratic Equations in $\mathbb{Z}_p$)}
    \begin{itemize}
        \item[(a)] The equation $x^2=2$ admits a solution in $\mathbb{Z}_7$.
        \item[(b)] If $b\in\mathbb{Z}_2^\times$ with $b\equiv a^2\ \mathrm{mod}\ (2)\mathbb{Z}_2$, then $f=x^2-b$ has a root in $\mathbb{Z}_2$.
    \end{itemize}
\end{exe}


%	\include{multiToC}
\appendix
\cleardoublepage\part{Appendix}

\chapter{Completions}
\begin{definition}{1. (Topological Group)}
    \newline A \textbf{topological group} consists of a topological space $(G,\mathcal{T})$ and a group $(G,\cdot,1)$, (write $(G,+,0)$ if it is abelian) such that the operation $\cdot:G\to G,\ (x,y)\mapsto x\cdot y$ and the inversion $\iota:G\to G,\ x\mapsto x^{-1}$ are contiuous.
\end{definition}

In this chapter, we only consider the topological abelian group $(G,\mathcal{T},+,0)$:
\begin{nt}{2}
    \begin{itemize}
        \item[(a)] Note that $\iota$ is a homeomorphism, for each neighbourhood $U\in\mathcal{N}_0$, we also have $-U=\iota(U)\in\mathcal{N}_0$. Now, set $V\coloneqq U\cap-U\in\mathcal{N}_0$, then we have $V=-V$.
        \item[(b)] Since $+$ is continuous and sends $(0,0)\mapsto0$, for any neighbourhood $U\in\mathcal{N}_0$, there is a neighbourhood $W\subset G\times G$ of $(0,0)$ such that $+:W\xhookrightarrow{} U$. By the definition of product topology, there are $V_1,V_2\in\mathcal{N}_0$ such that $V_1\times V_2\subset W$.\\ Now, set $V\coloneqq V_1\cap V_2\in\mathcal{N}_0$, then we have $V+V\subset U$.
    \end{itemize}
\end{nt}

\begin{lemma}{3}
    If $\{0\}$ is closed in $G$, then the diagonal $D\coloneqq\{(x,x):x\in G\}$ is closed in $G$. Hence $G$ is Hausdorff.
\end{lemma}
\begin{proof}
    Note that the map 
    \begin{align*}
        \kappa:G\times G\longrightarrow G\times G,\ (x,y)\longmapsto(x,-y)
    \end{align*}
    is continuous since $\mu=(1_G,\ \iota)$. It follows that the map 
    \begin{align*}
        \phi:G\times G\longrightarrow G,\ (x,y)\longmapsto x-y
    \end{align*}
    is continuous since $\phi=+\circ\mu$. Hence by assumption $D=\phi^{-1}(0)$ is closed in $G\times G$. 
\end{proof}
\newpage
\begin{proposition}{4}
    \begin{itemize}
        \item[(a)] If $a$ is a fixed element of $G$, the \textbf{translation} defined by 
        \begin{align*}
            T_a:G\longrightarrow G,\ x\longmapsto x+a
        \end{align*}
        is a homeomorphism.
        \item[(b)] If $U\in\mathcal{N}_0$ is any neighbourhood, then $a+U\in\mathcal{N}_a$; conversely, if $V\in\mathcal{N}_a$, then $V$ is of the form $V\coloneqq a+U$ for some $U\in\mathcal{N}_0$.
    \end{itemize}
    \textbf{Remark.} The topology of $G$ is thus uniquely determined by the neighbourhoods of $0$ in $G$.
\end{proposition}
\begin{proof}
    (a) Note that $T_a=+\circ\kappa$, where $\kappa:G\xhookrightarrow{} G\times G,\ x\mapsto(x,a)$ is continuous. Hence $T_a$ is continuous. Moreover,  $(T_a)^{-1}=T_{-a}$ is also continuous, whence $T_a$ is a homeomorphism.\\
    (b) Since $T_{-a}(a)=0$ and the map $T_{-a}$ is continuous, it follows that $a+U=(T_{-a})^{-1}(U)\in\mathcal{N}_a$. Conversely, since $T_a(0)=a$ and $T_a$ is continuous, $U\coloneqq T_a^{-1}(V)\in\mathcal{N}_0$. Moreover, $T_a$ is a homeomorphism, thus $V=T_a\circ U=a+U$.
\end{proof}

\begin{lemma}{5}
    Let $H\coloneqq\bigcap_{U\in\mathcal{N}_0(G)}U$. Then
    \begin{itemize}
        \item[(a)] $H$ is a subgroup of $G$.
        \item[(b)] $H$ is the closure of $\{0\}$.
        \item[(c)] $G/H$ is Hausdorff.
        \item[(d)] $G$ is Hausdorff $\Longleftrightarrow\ H=0$.\qquad [AM, 10.1]
    \end{itemize}
\end{lemma}
\begin{proof}
    (a) This follows from the continuity of group operations and \textbf{Note 2}. \\ For (b) we have:
    \begin{align*}
        x\in H\ \Leftrightarrow\ x\in U, \forall U\in\mathcal{N}_0(G)\ 
        \Leftrightarrow\ 0\in x-U, \forall U\in\mathcal{N}_0(G)\ \Leftrightarrow\  x\in\overline{\{0\}}.
    \end{align*}
    (c) From (b) the cosets of $H$ are all closed, thus points are closed in $G/H$ and so $G/H$ is Hausdorff. (d) By definition of Hausdorff and construction of $H$, it's clearly that $H=0$; the converse is trivial.
\end{proof}

\begin{definition}{6. (Convergent \& Cauchy Sequences)}
    \begin{itemize}
        \item A sequence $(a_n)$ in $A$ \textbf{converges} converges to $0$ if for all $U\in\mathcal{N}_0$ there exists $N\in\mathbb{N}$ with $a_n\in U$ provided $n\geq N$. In this case we write $a_n\longrightarrow0$.
        \item A \textbf{Cauchy sequence} $(a_n)$ in $A$ is a sequence such that for all $U\in\mathcal{N}_0$ there exists $N\in\mathbb{N}$ with $a_i-a_j\in U$ provided $i,j\geq N$. Moreover, two Cauchy sequences $(a_n)$ and $(b_n)$ are \textbf{equivalent}, write $(a_n)\sim(b_n)$, if $(a_n-b_n)$ converges to $0$. 
    \end{itemize}
\end{definition}
\newpage
\begin{lemma}{7}
    The relation $\sim$ of Cauchy sequences in $G$ is an equivalence relation.
\end{lemma}
\begin{proof}
    \begin{itemize}
        \item $(x_i)\sim(x_i)$ since $x_i-x_i=0\longrightarrow0$.
        \item If $(x_i)\sim(y_i)$, then $x_i-y_i\longrightarrow0$. Since the map $\iota:G\to G,\ x\mapsto-x$ is continuous, it follows that 
        \begin{align*}
            y_i-x_i=\iota(x_i-y_i)\longrightarrow\iota(0)=0,
        \end{align*}
        whence $(y_i)\sim(x_i)$.
        \item Now, suppose $(x_i)\sim(y_i)$ and $(y_i)\sim(z_i)$. For any neighbourhood $U\in\mathcal{N}_0$, we take $V\in\mathcal{N}_0$ such that $V+V\subset U$. By assumption, there exist $N_1,N_2\in\mathbb{N}$ such that 
        \begin{align*}
            x_i-y_i\in V,\ \forall i\geq N_1\qquad \text{and}\qquad y_i-z_i\in V,\ \forall i\geq N_2.
        \end{align*}
        Take $N\coloneqq\mathrm{max}\{N_1,N_2\}$, then for any $i\geq N$, we have
        \begin{align*}
            x_i-z_i=(x_i-y_i)+(y_i-z_i)\in V+V\subset U.
        \end{align*}
    \end{itemize}
\end{proof}

\begin{definition}{8. (Completion)}
    In the following, we assume for simplicity that $0\in G$ has a \textit{countable local base}. We define the \textbf{completion} of $G$ by 
        \begin{align*}
            \widehat{G}\coloneqq\ \text{Cauchy sequences in $G$}/\sim.
        \end{align*}
\end{definition}
\newpage
\begin{proposition}{9}
    Let $(x_i)$ and $(y_i)$ be Cauchy sequences in a topological abelian group $G$. Then
    \begin{itemize}
        \item[(a)] $(x_i+y_i)$ is a Cauchy sequences and its class in $\widehat{G}$ depends only on the classes of $(x_i)$ and $(y_i)$.
        \item[(b)] $(-x_i)$ is a Cauchy sequences and its class in $\widehat{G}$ depends only on the class of $(x_i)$.
        \item[(c)] $(0)$ is a Cauchy sequence.
        \item[(d)] We thus obtain an addition $+$ in $\widehat{G}$ making it an abelian group.
    \end{itemize}
\end{proposition}
\begin{proof}
    (a) For any neighbourhood $U\in\mathcal{N}_0$, we take $V\in\mathcal{N}_0$ such that $V+V\subset U$. Since $(x_i)$ and $(y_i)$ are Cauchy, there exist $N_1,N_2\in\mathbb{N}$ such that
    \begin{align*}
            x_i-x_j\in V,\ \forall i,j\geq N_1\qquad \text{and}\qquad y_i-y_j\in V,\ \forall i\geq N_2.
        \end{align*}
        Take $N\coloneqq\mathrm{max}\{N_1,N_2\}$, then for any $i,j\geq N$, we have
        \begin{align*}
            (x_i+y_i)-(x_j-y_j)=(x_i-x_j)+(y_i-y_j)\in V+V\subset U.
        \end{align*}
        Now, if $\overline{(x'_i)}=\overline{(x_i)}$, then $x_i-x'_i\longrightarrow0$. Note that 
        \begin{align*}
            (x_i+y_i)-(x'_i+y_i)=x_i-x'_i\longrightarrow0, 
        \end{align*}
        whence $(x_i+y_i)\sim(x'_i+y_i)$. Similarly, if $\overline{(y'_i)}=\overline{(y_i)}$, we have $(x'_i+y_i)\sim(x'_i+y'_i)$. By transitivity $(x_i+y_i)$ and $(x'_i+y'_i)$ are equivalent.\\
        (b) For any $U\in\mathcal{N}_0$, since $(x_i)$ is Cauchy, there is $N\in\mathbb{N}$ such that $x_i-x_j\in U$ if $i,j\geq N$. Now, we have
        \begin{align*}
            (-x_i)-(-x_j)=x_j-x_i\in U,\qquad \forall i,j\geq N.
        \end{align*}
        Hence $(-x_i)$ is also Cauchy. If $(x'_i)$ represents the same class as $(x_i)$, so that $x_i-x'_i\longrightarrow0$, then $(-x_i)-(-x'_i)\longrightarrow0$ as well; for any $U\in\mathcal{N}_0$, we may find $V\in\mathcal{N}_0$ with $V=-V\subset U$ and some $N\in\mathbb{N}$ such that $x_i-x'_i\in V=-V,\ \forall i\geq N$, thus 
        \begin{align*}
            (-x_i)-(-x'_i)=-(x_i-x'_i)\in V\subset U,
        \end{align*}
        i.e. $(-x'_i)\sim(-x_i)$.\\
        (c) Clearly, since all differences are zero.\\
        (d) By the above facts on independence of representatives, it will be enough to show that the set of Cauchy sequences in $G$ forms an abelian group. But thus easily follows from the facts that this group is defined as a subgroup of the product group $G^\mathbb{N}$ and that the abelian group identities hold componentwise:
        \begin{align*}
            &\bullet [(x_i)+(y_i)]+(z_i)= (x_i+y_i)+(z_i)=(x_i+y_i+z_i) =(x_i)+(y_i+z_i)=(x_i)+[(y_i)+(z_i)]\\
            &\bullet (x_i)+(0)=(x_i+0)=(x_i)\\
            &\bullet (x_i)+(-x_i)=(x_i-x_i)=(0)\\
            &\bullet (x_i)+(y_i)=(x_i+y_i)=(y_i+x_i)=(y_i)+(x_i).
        \end{align*}
\end{proof}

\begin{proposition}{10}
    There is a natural topology making $\widehat{G}$ a topological group.
\end{proposition}
\begin{proof}
    For each $U\in\mathcal{N}_0$, we define $\widehat{U}\subset\widehat{G}$ to be the set of all elements $\widehat{x}\in\widehat{G}$ such that all representatives $(x_i)$ of $\widehat{x}$ are "eventually" in $U$, i.e.
    \begin{align*}
        \widehat{U}\coloneqq\{\widehat{x}\in\widehat{G}:\ \forall(x_i)\in\widehat{x},\ \exists N\in\mathbb{N}\ \text{with}\ x_i\in U,\ \forall i\geq N\}.
    \end{align*}
    Note that for all $U,V\in\mathcal{N}_0(G)$, we have $\widehat{0}\in\widehat{U\cap V}=\widehat{U}\cap\widehat{V}$. Now we consider the set
    \begin{align*}
        \widehat{\mathcal{N}}\coloneqq\{\widehat{U}:U\in\mathcal{N}_0(G)\}
    \end{align*}
    and take all translates in $\widehat{G}$, these sets together generate a unique topology on $\widehat{G}$ for which $\widehat{\mathcal{N}}$ forms a \textit{local base} of $\widehat{0}$. Now we claim that the addition and inversion are continuous.\\ Suppose $\widehat{x},\widehat{y}\in\widehat{G}$ such that $\widehat{x}+\widehat{y}\in W$ for some open set $W\subset\widehat{G}$. By our definition of topology, $\exists\widehat{U}\in\widehat{\mathcal{N}}$ such that $(\widehat{x}+\widehat{y})+\widehat{U}\subset W$. Moreover, there exists $V\in\mathcal{N}_0(G)$ with $V+V\subset U$. Now, if $(x_i)$ and $(y_i)$ are Cauchy sequences representing elements of $\widehat{V}$, then $\exists N\in\mathbb{N}$ sufficiently large with $x_i,y_i\in V,\ \forall i\geq N$, whence $x_i+y_i\in U$. Hence, we have $\widehat{V}+\widehat{V}\subset\widehat{U}$ and the addition
    \begin{align*}
        \widehat{+}:(\widehat{x}+\widehat{V})\times(\widehat{y}+\widehat{V})\xhookrightarrow{} (\widehat{x}+\widehat{y})+\widehat{U}\subset W
    \end{align*}
    is continuous. Similarly, any open neighbourhood of $-\widehat{x}$ contains a basic open set $-\widehat{x}+\widehat{U}$. Moreover, there exists $V\in\mathcal{N}_0(G)$ with $V=-V\subset U$. Then the "negative" of any sequence eventually in $V$ is also eventually in $V$. Thus $-\widehat{V}=\widehat{V}$ and the inversion
    \begin{align*}
        \widehat{\iota}:-(\widehat{x}+\widehat{V})\xhookrightarrow{=}-\widehat{x}+\widehat{V}\subset-\widehat{x}+\widehat{U}
    \end{align*}
    is continuous.
\end{proof}

\begin{proposition}{11}
    The natural map $\phi:G\to\widehat{G},\ x\mapsto\overline{(x)}$, is a continuous homomorphism of abelian groups. If $\phi$ is an isomorphism, we say that $G$ is \textbf{complete}.
\end{proposition}
\begin{proof}
    Clearly, a constant sequence $(x)$ is a Cauchy sequence.\\
    \textsc{claim.} $\phi$ is a continuous homomorphism.
    \begin{itemize}
        \item $\phi(0)$ is the equivalence class of $(0)$, which is the zero of $\widehat{G}$.
        \item From \textbf{Prop. 9}, we have $(x+y)=(x)+(y)$ and $(-x)=-(x)$. By taking classes we have $\phi(x+y)=\phi(x)+\phi(y)$ and $\phi(-x)-\phi(x)$. 
        \item \textbf{(Continuity).} Let $x\in G$ and consider a basic open set $\phi(x)+\widehat{U}$ of $\widehat{G}$ for some $U\in\mathcal{N}_0(G)$. If $u\in U$, then $\exists V\in\mathcal{N}_0(G)$ such that $u+V\subset U$. Now, if $(y_i)\sim(u)$ in $\widehat{G}$, we have "eventually" $y_i-u\in V$, i.e. eventually $y_i\in u+V\subset U$. It follows that $\phi(u)\in\widehat{U}$ and hence $\phi(x+U)\subset\phi(x)+\widehat{U}$, thus $\phi$ is continuous.
    \end{itemize}
\end{proof}

\begin{remark}{12}
    We have $\mathrm{ker}\phi=\bigcap_{U_\in\mathcal{N}_0(G)}U$, and so by \textbf{Lemma 5}, $\phi$ is injective iff $H$ is Hausdorff. 
\end{remark}
\begin{proof}
    "$\subset$"\ If $x\in\mathrm{ker}\phi$, then $(x)\sim(0)$, i.e. $x=x-0\longrightarrow0$ as the index increases. Now, for any neighbourhood $U\in\mathcal{N}_0(G)$, the sequence $(x)$ is eventually in $U$, but it is constant, whence $x\in U$. Hence $x\in\bigcap_{U\in\mathcal{N}_0(G)}U$. "$\supset$"\ Clearly if $x\in\bigcap_{U\in\mathcal{N}_0(G)}U$, then the constant sequence $(x)$ converges to $0$, whence $(x)\sim(0)$, i.e. $\phi(x)=(x)=(0)$.
\end{proof}

\begin{proposition}{13}
    Let $G,H$ be two abelian topological group.
    \begin{itemize}
        \item[(a)]  If $f:G\to H$ is a continuous homomorphism, then $f$ maps a Cauchy sequence in $G$ into a Cauchy sequence of $H$.
        \item[(b)] From (a), $f$ induces a continuous homomorphism $\widehat{f}:\widehat{G}\to\widehat{H}$.
        \item[(c)] If we have $G\xrightarrow{f} H\xrightarrow{g} K$, then $\widehat{g\circ f}=\widehat{g}\circ\widehat{f}$.
    \end{itemize}
\end{proposition}
\begin{proof}
    (a) Let $(x_i)$ be a Cauchy sequence and let $U\in\mathcal{N}_0(H).\ \because f(0)=0$ and $f$ is continuous. $\therefore\exists V\in\mathcal{N}_0(G)$ such that $f(V)\subset U.\ \because (x_i)$ is Cauchy. $\therefore\exists N\in\mathbb{N}$ such that $x_i-x_j\in V,\ \forall i,j\geq N$. Now, since $f$ is a homomorphism, we have 
    \begin{align*}
        f(x_i)-f(x_j)=f(x_i-x_j)\in f(V)\subset U.
    \end{align*}
    Hence $(f(x_i))$ is a Cauchy sequence in $H$.\\
    (b) We first show that $\widehat{f}$ is well-defined.
    \begin{itemize}
        \item \textbf{(Well-defined).} If $(x_i)\sim(y_i)$ in $\widehat{G}$, then $x_i-y_i\longrightarrow0$. Since $f$ is a continuous homomorphism, we have
        \begin{align*}
            f(x_i)-f(y_i)=f(x_i-y_i)\longrightarrow f(0)=0. 
        \end{align*}
        Thus, $(f(x_i))\sim(f(y_i))$ in $\widehat{H}$, and so  $\widehat{f}(\overline{(x_i)})=\overline{(f(x_i))}=\overline{(f(y_i))}=\widehat{f}(\overline{(y_i)})$.
        \item \textbf{(Homomorphism).} Obviously, $f$ is a homomorphism since operations on Cauchy sequences are defined componentwise and $\widehat{f}$ is induced by applying $f$ componentwise to Cauchy sequences.
        \item \textbf{(Continuity).} Given any basic open set $\widehat{f}(x)+\widehat{V}$ in $\widehat{H}$, where $V\in\mathcal{N}_0(H)$. Since $f$ is continuous and $f(0)=0$, there exists $U\in\mathcal{N}_0(G)$ such that $f(U)\subset V$. Now, for any Cauchy sequence $(u_i)$ eventually in $U$, the image $(f(u_i))$ is a Cauchy sequence eventually in $V$ by (a), hence $\widehat{f}(\widehat{U})\subset\widehat{V}$. It follows that $\widehat{f}(\widehat{x}+\widehat{U})\subset\widehat{f}(\widehat{x})+\widehat{V}$, thus $\widehat{f}$ is continuous.
    \end{itemize}
    (c) For any Cauchy sequence $(x_i)$ in $G$, we have
    \begin{align*}
        \widehat{g\circ f}((x_i)_i)=(g\circ f(x_i))_i=g\circ f((x_i))_i=\widehat{g}((f(x_i))_i)=\widehat{g}\circ\widehat{f}((x_i)_i).
    \end{align*}
\end{proof}
\newpage
From now on, we restrict ourselves to the special kind of topologies on $G$ occurring in commutative algebra:
\begin{definition}{14}
    We assume that $0\in G$ has a fundamental system of neighborhoods (local base) consisting of \textit{subgroups}. Thus we have a sequence of subgroups
    \begin{align*}
        G=G_0\supset G_1\supset\cdots\supset G_n\supset\cdots
    \end{align*}
    and $U\in\mathcal{N}_0(G)$ iff $U\supset G_n$ for some $n\in\mathbb{N}$. Moreover, in such topologies the subgroups $G_n$ of $G$ are both open and closed.
\end{definition}
\begin{proof}
    (i) It is a well-defined topological group in the sense that:
    \begin{itemize}
        \item $g+G_n\subset (g+G_n)\cap(g+G_m)$ for any $m\leq n$.
        \item Clearly, $g\in g+G_n$ for any $n\in\mathbb{N}$.
        \item For any $y\in g+G_n$, we have $y+G_n=g+G_n+G_n=g+G_n$.
        \item $\iota:G\to G,\ x\mapsto-x$ is continuous since $g+G_n\leftrightarrow-g+G_n$.
        \item $+:G\times G\to G,\ (x,y)\mapsto x+y$ is continuous since for any given $(x+y)+G_n\in\mathcal{N}_{x+y}(G)$, we have
        \begin{align*}
            (x+G_n)\times(y+G_n)\ \xhookrightarrow{+}\ (x+y)+G_n.
        \end{align*}
    \end{itemize}
    (ii) If $g\in G_n$, then $g+G_n\in\mathcal{N}_g(G)$. Since $g+G_n\subset G_n$, it follows that $G_n$ is open. Moreover, note that the complement of $G_n$ is indeed the union of nontrivial cosets
    \begin{align*}
        G\setminus G_n=\bigcup_{h\notin G_n}(h+G_n),
    \end{align*}
     which is open since each $h+G_n$ is open. Thus, $G_n$ is also closed in $G$.
\end{proof}

\begin{definition}{15. (Coherent Sequence)}
    Suppose $(x_i)$ is a Cauchy sequence in $G$. Then the image of $x_i$ in $G/G_n$ is eventually constant, say $x_i+G_n\longrightarrow\xi_n$ in $G/G_n$. It we pass from $n+1$ to $n$ it is clear that the projection $G/G_{n+1}\xrightarrow{\theta_{n+1}} G/G_n$ maps $\xi_{n+1}\mapsto\xi_n$.\\ In this case we say that $(x_i)$ induces the \textbf{coherent sequence} $(\xi_n)$.
\end{definition}

\begin{remark}{16}
    \begin{itemize}
        \item It is clear that the equivalent Cauchy sequences define the same coherent sequence.
        \item On the other hand, given any coherent sequence $(\xi_n)$, we can construct a Cauchy sequence $(x_n)$ giving rise to it by taking $x_n\in\xi_n$ so that we have $x_{n+1}-x_n\in G_n$.
        \item $\widehat{G}$ can also be well-defined as the set of coherent sequences $(\xi_n)$ with the obvious group structure.
    \end{itemize}
\end{remark}

\begin{definition}{17. (Inverse System \& Inverse Limit)}
    \newline A set $I$ is \textbf{directed} by $\leq$ if 
    \begin{itemize}
        \item (Reflexive) $i\leq i,\ \forall i\in I$.
        \item (Transitive) If $i\leq j$ and $j\leq k$ in $I$, then $i\leq k$.
        \item For any $i,j\in I$, there exists $k\in I$ such that both $i,j\leq k$.
    \end{itemize}
    Let $\{A_i\}_{i\in I}$ be a family of (abelian) groups together with group homomorphisms $\{\theta_{ij}:A_j\to A_i\}_{i\leq j}$ such that
    \begin{itemize}
        \item $\theta_{ii}=1_{A_i}$ for all $i\in I$.
        \item $\theta_{ik}=\theta_{ij}\circ\theta_{jk}:A_k\to A_i$ whenever $i\leq j\leq k$.
    \end{itemize}
    We also say that $(A_i,\ \theta_{ij})$ (or simply $A_i$ for short) defines an \textbf{inverse system}. The abelian group
    \begin{align*}
        \varprojlim_{i\in I}A_i\coloneqq\{a\in\prod_{i\in I}A_i:a_i=\theta_{ij}(a_j),\ \forall i\leq j\ \text{in}\ I\}
    \end{align*}
    is called the \textbf{inverse limit} of the family.
\end{definition}

\begin{remark}{18}
    The inverse limit comes with natural projections
    \begin{align*}
        \pi_j:\varprojlim A_i\xtwoheadrightarrow{} A_j \qquad\text{satisfying}\qquad \theta_{ij}\circ\pi_j=\pi_i.
    \end{align*}
    This even characterises the inverse limit via the following universal property: If $A$ is an abelian group together with morphisms $\psi_i:A\to A_i$ such that $\theta_{ij}\circ\psi_j=\psi_i$, then there exists a unique morphism 
    \begin{align*}
        \psi:A\to\varprojlim A_i\qquad\text{such that}\qquad \pi_j\circ\psi=\psi_j,\ \forall j\in I.
    \end{align*}
    \begin{center}
        \begin{tikzcd}
            & A \arrow[d, "\psi_j"] \\
            \varprojlim A_i \arrow[ru, leftarrow, dashed, "\exists!\ \psi", red]
            & A_j \arrow[l, leftarrow, "\pi_j"] \\
        \end{tikzcd}
    \end{center}
\end{remark}
\newpage
In this section we always consider $I=\mathbb{N}_{\geq1}$: 
\begin{nt}{19. (Completions as Inverse Limits)}
    \newline Let $G$ be an (abelian) group whose topology is defined by the subgroups
    \begin{align*}
        G=G_0\supset G_1\supset\cdots\supset G_n\supset\cdots.
    \end{align*}
    Let $\pi_{ij}:G/G_j\xtwoheadrightarrow{} G/G_i,\ g+G_j\mapsto g+G_i,\ (i\leq j)$, be the natural projection map. Since the completion $\widehat{G}$ with respect to this topology equals the set of coherent sequences which is just the inverse limit of the system $\{G/G_i\}$, we thus obtain
    \begin{align*}
        \widehat{G}\ \cong\ \varprojlim G/G_i\ =\ \{a\in\prod_{i\in I}G/G_i:a_i=\pi_{ij}(a_j),\ \forall i\leq j\ \text{in}\ I\}.
    \end{align*}
\end{nt}

\begin{definition}{20}
    \begin{itemize}
        \item An inverse system $\{A_i\}$ is called \textbf{surjective} if for each $i,\ \theta_i\coloneqq\theta_{i,i+1}:A_{i+1}\to A_i$ is surjective.
        \item An \textbf{exact sequence} of inverse systems
        \begin{align*}
            0\ \to\ \{A_i\}\ \to\ \{B_i\}\ \to\ \{C_i\}\ \to\ 0
        \end{align*}
        consists of a commutative diagram
        \begin{center}
        \begin{tikzcd}
            0\rar[] &A_{i+1}\rar[] \dar["\theta'_{i+1}"]
            &B_{i+1}\rar[] \dar["\phi_{i+1}"]
            &C_{i+1}\rar[] \dar["\psi_{i+1}"] &0 \\
            0\rar[] &A_i\rar[] &B_i\rar[] &C_i\rar[] &0 \\
        \end{tikzcd}
    \end{center}
    of exact sequences where the maps $\theta_{i+1},\ \phi_{i+1}$ and $\psi_{i+1}$ are provided by the inverse systems $\{A_i\},\ \{B_i\}$ and $\{C_i\}$, respectively.
    \end{itemize}
\end{definition}
\newpage
\begin{proposition}{21}
    If $0\to\{A_i\}\to\{B_i\}\to\{C_i\}\to0$ is an exact sequence of inverse systems, then taking inverse limit is a "left exact functor", i.e.
    \begin{align*}
        0\ \to\ \varprojlim A_i\ \to\ \varprojlim B_i\  \to\ \varprojlim C_i\
    \end{align*}
    is exact. Furthermore, if $\{A_i\}$ is a surjective system, then
    \begin{align*}
        0\ \to\ \varprojlim A_i\ \to\ \varprojlim B_i\ \to\ \varprojlim C_i\ \to\ 0
    \end{align*}
    is exact. [AM, 10.2]
\end{proposition}
\begin{proof}
    For each $i$, we have the commutative diagram of groups
    \begin{center}
        \begin{tikzcd}
            0\rar[] &A_{i+1}\rar["f_{i+1}"] \dar["\theta_{i+1}"]
            &B_{i+1}\rar["g_{i+1}"] \dar["\phi_{i+1}"]
            &C_{i+1}\rar[] \dar["\psi_{i+1}"] &0 \\
            0\rar[] &A_i\rar["f_i"] &B_i\rar["g_i"] &C_i\rar[] &0 \\
        \end{tikzcd}
    \end{center}
    by assumption. Let $A\coloneqq\prod_{i}A_i$ and consider $d^A:A\to A,\ a_i\mapsto a_i-\theta_{i+1}(a_{i+1})$, then, clearly from the above definition, $\mathrm{ker}d^A=\varprojlim A_i$. We define this on $B$ and $C$ similarly.\\ Now, consider the diagram
    \begin{center}
        \begin{tikzcd}
            0\rar[] &A\rar["f"] \dar["d^A"]
            &B\rar["g"] \dar["d^B"]
            &C\rar[] \dar["d^C"] &0 \\
            0\rar[] &A\rar[] &B_i\rar[] &C_i\rar[] &0. \\
        \end{tikzcd}
    \end{center}
    The horizontal rows are exact since $0\to A_i\to B_i\to C_i\to0$ is exact for each $i$ and the diagram is commutative since
    \begin{align*}
        f\circ d^A(a_i) &= f(a_i-\theta_{i+1}(a_{i+1}))=(f_i(a_i-\theta_{i+1}(a_{i+1})))_i=(f_i(a_i-\theta_{i+1}(a_{i+1})))_i \\ &=(f_i(a_i)-\phi_{i+1}\circ f_{i+1}(a_{i+1}))_i=d^B\circ f_i(a_i) 
    \end{align*}
    and similarly, $g\circ d^B=d^C\circ g$. Therefore, by \textbf{Snake Lemma}, we obtain the first result.\\ Now, if $\{A_i\}$ is surjective, we claim that $d^A$ is surjective. Let $a\coloneqq(a_i)_i\in A$. Take $x_1=0$. Since $\theta_2:A_2\to A_1$ is surjective, for $-a_1\in A_1,\ \exists x_2\in A_2$ such that $\theta_2(x_2)=-a_1$. Now, $x_1=0\mapsto x_1-\theta_2(a_2)=a_1$. Inductively for $i\geq2$, since $\theta_{i+1}$ is surjective, for $x_i-a_i\in A_i,\ \exists x_{i+1}\in A_{i+1}$ such that $\theta_{i+1}(x_{i+1})=x_i-a_i$, thus $x_i\mapsto x_i-\theta_{i+1}(x_{i+1})=a_i$. Hence $d^A((x_i)_i)=(a_i)_i$, i.e. $d^A$ is surjective. It follows that $\mathrm{Coker}d^A=0$, thus the required result. 
\end{proof}
\newpage
\begin{lemma}{22}
    Suppose that $0\to M'\xhookrightarrow{i} M\xtwoheadrightarrow{p} M''\to0$ is a short exact sequence as modules. If $N\subset M$ is a submodule, then 
    \begin{align*}
        0\to \tfrac{M'}{M'\cap N}\xhookrightarrow{\iota} \tfrac{M}{N}\xtwoheadrightarrow{\pi} \tfrac{M''}{p(N)}\to 0
    \end{align*}
    is exact as modules.
\end{lemma}
\begin{proof}
    We first proof the well-definedness of
    \begin{align*}
        \iota:x+M'\cap N\mapsto i(x)+N\qquad \text{and}\qquad \pi:m+N\mapsto p(m)+p(N)
    \end{align*}
    \begin{itemize}
        \item If $\overline{x}=x+M'\cap N=y+M'\cap N=\overline{y}$, then $x-y\in M'\cap N\subset N.\ \Rightarrow i(x)-i(y)=i(x-y)\in i(N)\cong N$. It follows that
        \begin{align*}
            \iota(\overline{x})=i(x)+N=i(y)+N=\iota(\overline{y}).
        \end{align*}
        \item If $m+N=n+N$, then $m-n\in N.\ \Rightarrow p(m)-p(n)=p(m-n)\in p(N)$. Thus,
        \begin{align*}
            \pi(m+N)=p(m)+N=p(n)+N=\pi(n+N).
        \end{align*}
    \end{itemize}
    Now, it suffices to show $\iota$ is 1-1, $\pi$ is onto and the exactness at $M/N$:
    \begin{itemize}
        \item \textbf{($\iota$ is 1-1).} Suppose $\iota(\overline{x})\coloneqq\iota(x+M'\cap N)=0+N$, then $i(x)=0$. Since $i$ is 1-1, we have $x=0$. Hence, $\iota$ is 1-1. 
        \item \textbf{($\pi$ is onto).} Let $x+p(N)\in M''/p(N)$ where $x\in M''$. Since $p$ is onto, $\exists m\in M$ with $p(m)=x$. Thus, $\pi(m+N)=p(m)+N=x+N$.
        \item \textbf{(Exactness at $M/N$).} 
        \begin{align*}
            &(\subset):\ m+N\in\mathrm{ker}\pi & &\Rightarrow\ p(m)\in p(N)\ \Rightarrow\ \exists n\in N\ \text{such that}\ p(m)=p(n)\\
            & & &\Rightarrow\ p(m-n)=0.\ \Rightarrow\ m-n\in\mathrm{ker}p=\mathrm{Im}i\\
            & & &\Rightarrow\ \exists x\in M'\ \text{such that}\ i(x)=m-n\\
            & & &\Rightarrow\ m+N=(m-n)+N=\iota(x+M'\cap N)\in\mathrm{Im}\iota.\\
            &(\supset):\ m+N\in\mathrm{Im}\iota & &\Rightarrow\ \exists x\in M'\ \text{such that}\ m+N=\iota(x+M'\cap N)=i(x)+N\\
            & & &\Rightarrow\ \pi(m+N)=\pi(i(x)+N)=p(i(x))+p(N)=0+p(N)\\ 
            & & &\Rightarrow\ m+N\in\mathrm{ker}\pi.
        \end{align*}
    \end{itemize}
\end{proof}
\newpage
\begin{corollary}{23}
    Let $0\to G'\xhookrightarrow{i}  G\xtwoheadrightarrow{p} G''\to0$ be an exact sequence of groups. Let $G$ have the topology defined by a decreasing sequence $\{G_n\}$ of subgroups, and endow $G'$ and $G''$ with the induced topologies given by $\{G'\cap G_n\}$ and $\{p(G_n)\}$, respectively. Then we have the exact sequence of groups
    \begin{align*}
        0\ \to\ \widehat{G}'\ \to\ \widehat{G}\ \to\ \widehat{G}''\ \to\ 0
    \end{align*}
    completed with respect to these topologies. [AM, 10.3]
\end{corollary}
\begin{proof}
    We claim that the inverse systems are exact and that $\{\theta_{n+1}\}$ is surjective. \\ For each $n\in\mathbb{N}_{\geq1}$, consider the diagram:
    \begin{center}
        \begin{tikzcd}
            0\rar[] &\tfrac{G'}{G'\cap G_{n+1}}\rar["\iota_{n+1}"] \dar["\theta_{i+1}"]
            &\tfrac{G}{G_{n+1}}\rar["\pi_{n+1}"] \dar["\phi_{n+1}"]
            &\tfrac{G''}{p(G_{n+1})}\rar[] \dar["\psi_{n+1}"] &0 \\
            0\rar[] &\tfrac{G'}{G'\cap G_n}\rar["\iota_n"] &\tfrac{G}{G_n}\rar["\pi_n"] &\tfrac{G''}{p(G_n)}\rar[] &0 \\
        \end{tikzcd}
    \end{center}
    \begin{itemize}
        \item \textbf{(Commutative).}
        \begin{align*}
            &\phi_{n+1}\circ\iota_{n+1}(x+G'\cap G_{n+1}) & &= \phi_{n+1}(i(x)+G_{n+1})=i(x)+G_n
            =\iota_n(x+G'\cap G_n)\\ & & &=\iota_n\circ\theta_{n+1}(x+G'\cap G_{n+1});\\
            &\psi_{n+1}\circ\pi_{n+1}(g+G_{n+1}) & &=\psi_{n+1}(p(g)+p(G_{n+1}))=p(g)+p(G_n)=\pi_n(g+G_n)\\ & & &= \pi_n\circ\phi_{n+1}(g+G_{n+1}).
        \end{align*}
        \item \textbf{(Exactness).} It follows from the previous lemma for each $n\geq1$.
        \item \textbf{(Surjective).} For $\theta_{n+1}:\frac{G'}{G'\cap G_{n+1}}\to\frac{G'}{G'\cap G_n},\ x+G'\cap G_{n+1}\mapsto x+G'\cap G_n$, we note that:
        \begin{align*}
            x\in G'\cap G_{n+1}\in\mathrm{ker}\theta_{n+1}\ &\Leftrightarrow\ x+G'\cap G_n=0+G'\cap G_n\ \Leftrightarrow\ x\in G'\cap G_n\\ &\Leftrightarrow\ x\in\tfrac{G'\cap G_n}{G'\cap G_{n+1}}.
        \end{align*}
        It follows that $\mathrm{ker}\theta_{n+1}=\frac{G'\cap G_n}{G'\cap G_{n+1}}$. Hence, $\mathrm{Im}\theta_{n+1}\cong\frac{G'/G'\cap G_{n+1}}{\mathrm{ker}\theta_{n+1}}\cong\frac{G'}{G'\cap G_n}$, i.e. $\theta_{n+1}$ is surjective.
    \end{itemize}
    Therefore by \textbf{Prop. 21},  
    \begin{center}
        \begin{tikzcd}
            0\rar[] &\varprojlim\tfrac{G'}{G'\cap G_{n}}\rar[] \dar["\cong"]
            &\varprojlim\tfrac{G}{G_{n}}\rar[] \dar["\cong"]
            &\varprojlim\tfrac{G''}{p(G_n)}\rar[] \dar["\cong"] &0 \\
            0\rar[] &\widehat{G}'\rar[] &\widehat{G}\rar[] &\widehat{G}''\rar[] &0 \\
        \end{tikzcd}
    \end{center}
    is exact.
\end{proof}

\begin{remark}{24}
    \begin{itemize}
        \item[(a)] If we apply the previous corollary with $G'=G_n$ and $G''=G/G_n$, then $G''$ has the discrete topology.
        \item[(b)] We may deduce that $\widehat{G}''=G''$, i.e. $G''$ is complete.
    \end{itemize}
\end{remark}
\begin{proof}
    (a) For any $g+G_n\in G/G_n,\ (g\in G)$. Since $G_n$ is open in $G$ (with the topology), its translation $g+G_n$ is also open in $G$, whence the singleton $\{g+G_n\}$ is open in $G/G_n$. Therefore $G''=G/G_n$ is discrete.\\
    (b) We claim that $\phi:G''\to\widehat{G}''$ is an isomorphism. But from the previous discussion we already know $\phi$ is a surjective group homomorphism. It remains to show that $\phi$ is 1-1.\\ From (a), $G''$ is discrete, thus Hausdorff, and hence by \textbf{Rmk. 12}, $\phi$ is 1-1.
\end{proof}

\begin{corollary}{25}
    (a) $\widehat{G}_n$ is a subgroup of $\widehat{G}$; (b) $\widehat{G}/\widehat{G}_n\cong G/G_n$; (c) $\widehat{\widehat{G}}\cong\widehat{G}$. [AM, 10.4-5]
\end{corollary}
\begin{proof}
    (a)(b) Note that $0\to G_n\xhookrightarrow{} G\twoheadrightarrow G/G_n\to 0$ is exact. Apply \textbf{Cor. 23}, we have
    \begin{align*}
        0\ \to\ \widehat{G}_n\ \xhookrightarrow{}\ \widehat{G}\ \twoheadrightarrow\ \widehat{G/G_n}\ \to\ 0. 
    \end{align*}
    Thus, $\widehat{G}_n\subset\widehat{G}$ is a subgroup and $\widehat{G}/\widehat{G}_n\cong\widehat{G/G_n}\cong G/G_n$.\\
    (c) It follows from the completions as inverse limits and (b) that:
    \begin{align*}
        \widehat{\widehat{G}}\ \cong\ \varprojlim\widehat{G}/\widehat{G}_n\ \cong\ \varprojlim G/G_n\ \cong\ \widehat{G}.
    \end{align*}
\end{proof}

    
\chapter{Sheaves \& Schemes}
    
\section{Sheaves}
\begin{definition}{1. (Presheaf)}
    Let $X$ be a topological space. A \textbf{presheaf} $\mathscr{F}$ of abelian groups on $X$ consists of the data:
    \begin{itemize}
        \item For every open subset $U\subset X$, an abelian group $\mathscr{F}(U)$;
        \item For every inclusion of open sets $V\subset U$ of $X$, a morphism of abelian groups $\rho_{UV}:\mathscr{F}(U)\to\mathscr{F}(V)$,
    \end{itemize}
    subject to the conditions:
    \begin{itemize}
        \item $\mathscr{F}(\varnothing)=0$.
        \item $\rho_{UU}=1_{\mathscr{F}(U)}$.
        \item If $W\subset V\subset U$ are open subsets, then $\rho_{UW}=\rho_{VW}\circ\rho_{UV}$.
    \end{itemize}
\end{definition}

\begin{definition}{2. (Stalk)}
    If $\mathscr{F}$ is a presheaf on $X$ and if $P$ is a point on $X$, we define the \textbf{stalk} $\mathscr{F}_P$ of $\mathscr{F}$ at $P$ to be the direct limit of the groups $\mathscr{F}(U)$ for all open sets $U\ni P$, via the restriction maps $\rho$.
\end{definition}

\begin{remark}{3}
    From the definition, we have the following observation:
    \begin{itemize}
        \item $\mathscr{F}_P=\{(U,s):U\in\mathcal{N}_P,\ s\in\mathscr{F}(U)\}/\sim$.
        \item $(U,s)\sim(V,t).\ \Longleftrightarrow \exists\ W\in\mathcal{N}_P$ such that $W\subset U\cup V$ and $s|_W=t|_W$. Thus we may speak of elements of the stalk $\mathscr{F}_P$ as \textbf{germs} of \textbf{sections} $s$ of $\mathscr{F}$ at the point $P$.
    \end{itemize}
\end{remark}

\begin{definition}{4. (Sheaf)}
    A presheaf $\mathscr{F}$ on a topological space $X$ is \textbf{sheaf} if it satisfies the following supplementary conditions:
    \begin{itemize}
        \item Let $U\subset X$ be an open set and $\{V_i\}$ be an open covering of $U$. If $s\in\mathscr{F}(U)$ is an element such that $s|_{V_i}=0$ for all $i$, then $s=0$.
        \item If we have $s_i\in\mathscr{F}(V_i)$ for each $i$, with the property that $s_i|_{V_i\cap V_j}=s_j|_{V_i\cap V_j}$ for each $i,j$, then there is an element $s\in\mathscr{F}(U)$ such that $s|_{V_i}=s_i$ for each $i$. 
    \end{itemize}
    \textbf{Note.} The first condition implies that $s$ is unique.
\end{definition}

\end{document}